% Generated mechanically from the frozen spaces-morphisms.tex target.
% Do not edit this generated file; edit the bound locale source instead.

% OK, start here.
%

\setcounter{chapter}{66}
\chapter{代数空间的态射}



\phantomsection
\label{spaces-morphisms-section-phantom}


\section{引言}
\label{spaces-morphisms-section-introduction}

\noindent
本章介绍代数空间的若干类态射。
参考文献为 \cite{Kn}。

\medskip\noindent
目标是把《概形》以及《概形的态射》各章中定义的概形态射类型
分别扩展到代数空间范畴。每种情形略有不同，最好分别处理。





\section{约定}
\label{spaces-morphisms-section-conventions}

\noindent
固定假设所有概形都包含于大 fppf 站点 $\Sch_{fppf}$ 中。
并且，所考虑的每个环 $A$ 都满足 $\Spec(A)$ 同构于该大站点的一个对象。

\medskip\noindent
设 $S$ 为概形，设 $X$ 为 $S$ 上的代数空间。
在本章及下一章中，我们将用 $X \times_S X$ 表示 $X$ 与自身的乘积
（在 $S$ 上代数空间范畴中），而不写作 $X \times X$。





\section{可表示态射的性质}
\label{spaces-morphisms-section-representable}

\noindent
设 $S$ 为概形。
设 $f : X \to Y$ 为代数空间的可表示态射。在《空间》第
\ref{spaces-section-representable-properties} 节中，我们定义了当
$f$ 具有性质 $\mathcal{P}$ 的含义，其中 $\mathcal{P}$ 是概形态射的性质，
且满足下列条件：
\begin{enumerate}
\item 在任意基变换下保持，参见《概形》定义
\ref{schemes-definition-preserved-by-base-change}；并且
\item 在底上是 fppf 局部的，参见《下降》定义
\ref{descent-definition-property-morphisms-local}。
\end{enumerate}
也就是说，在此情形下，我们称 $f$ 具有性质 $\mathcal{P}$，当且仅当对每个
概形 $U$ 及任意态射 $U \to Y$，概形态射 $X \times_Y U \to U$
具有性质 $\mathcal{P}$。

\medskip\noindent
根据《空间》第 \ref{spaces-section-lists} 节的列表，这适用于下列性质：
(1)(a) 闭浸入，
(1)(b) 开浸入，
(1)(c) 拟紧浸入，
(2) 拟紧，
(3) 万有闭，
(4) （拟）分离，
(5) 单态射，
(6) 满射，
(7) 万有单射，
(8) 仿射，
(9) 拟仿射，
(10) （局部）有限型，
(11) （局部）拟有限，
(12) （局部）有限表示，
(13) 相对维数 $d$ 的局部有限型，
(14) 万有开，
(15) 平坦，
(16) syntomic，
(17) 光滑，
(18) 非分歧（相应地，G-非分歧），
(19) \'etale，
(20) proper，
(21) 有限或整，
(22) 有限局部自由，
(23) 万有商映射，
(24) 万有同胚，以及
(25) 浸入。

\medskip\noindent
本章将为不一定可表示的代数空间态射重新定义这些概念。
每次这样做时，我们都会确保新定义与旧定义一致，以避免歧义。

\medskip\noindent
注意，只要 $X$ 是概形，上述定义就适用，因为从概形到代数空间的
态射是可表示的。特别地，当 $X$ 与 $Y$ 都是概形时也适用。
在《空间》引理
\ref{spaces-lemma-morphism-schemes-gives-representable-transformation-property}
中我们已经看到，此时各定义相符，不会产生歧义。

\medskip\noindent
此外，在《空间》引理
\ref{spaces-lemma-base-change-representable-transformations-property}
中我们已经看到，如此定义的代数空间可表示态射的性质在任意代数空间
态射的基变换下保持。最后，在《空间》引理
\ref{spaces-lemma-composition-representable-transformations-property} 及
\ref{spaces-lemma-product-representable-transformations-property}
中我们已经看到，若 $\mathcal{P}$ 在复合下保持（对于上述性质中的
(1)(a)、(1)(b)、(1)(c)、(2)--(25)，但不包括 (13)，均如此），
则可表示态射的乘积保持性质 $\mathcal{P}$，并且可表示态射的复合也保持
性质 $\mathcal{P}$。

\medskip\noindent
下文将使用这些事实；每次使用时，我们简称为“本节”的结果。















\section{分离公理}
\label{spaces-morphisms-section-separation-axioms}

\noindent
列出代数空间态射的对角的一些先验性质是有意义的。

\begin{lemma}
\label{spaces-morphisms-lemma-properties-diagonal}
设 $S$ 为包含于 $\Sch_{fppf}$ 中的概形。
设 $f : X \to Y$ 为 $S$ 上代数空间的态射。
设 $\Delta_{X/Y} : X \to X \times_Y X$ 为对角态射。则
\begin{enumerate}
\item $\Delta_{X/Y}$ 可表示；
\item $\Delta_{X/Y}$ 局部有限型；
\item $\Delta_{X/Y}$ 是单态射；
\item $\Delta_{X/Y}$ 是分离的；以及
\item $\Delta_{X/Y}$ 局部拟有限。
\end{enumerate}
\end{lemma}

\begin{proof}
我们将使用如下事实：$\Delta_{X/S}$ 可表示（按代数空间的定义），
且满足性质 (2)--(5)，参见《空间》引理
\ref{spaces-lemma-properties-diagonal}。
注意，我们有如下因子分解
$$
X
\longrightarrow
X \times_Y X
\longrightarrow
X \times_S X
$$
它是对角 $\Delta_{X/S} : X \to X \times_S X$ 的因子分解。由于
$X \times_Y X \to X \times_S X$ 是单态射，且 $\Delta_{X/S}$ 可表示，
形式上可得 $\Delta_{X/Y}$ 可表示。特别地，其余陈述现在都有意义，
参见第 \ref{spaces-morphisms-section-representable} 节。

\medskip\noindent
选取满的 \'etale 态射 $U \to X$，其中 $U$ 为概形。考虑图表
$$
\xymatrix{
R = U \times_X U \ar[r] \ar[d] &
U \times_Y U \ar[d] \ar[r] &
U \times_S U \ar[d] \\
X \ar[r] & X \times_Y X \ar[r] & X \times_S X
}
$$
两个方块都是笛卡尔的，因而外侧矩形也是笛卡尔的。
顶行由概形组成，竖直箭头是满的 \'etale 态射。由《空间》引理
\ref{spaces-lemma-representable-morphisms-spaces-property}，
$\Delta_{X/Y}$ 的性质 (2)--(5) 等价于 $R \to U \times_Y U$ 的相应性质。
在《空间》引理 \ref{spaces-lemma-properties-diagonal} 的证明中，
我们已经看到 $R \to U \times_S U$ 具有性质 (2)--(5)。
态射 $U \times_Y U \to U \times_S U$ 是概形的单态射。
这些事实推出 $R \to U \times_Y U$ 具有性质 (2)--(5)。

\medskip\noindent
具体地，对于 (3)，注意 $R \to U \times_Y U$ 是单态射，因为复合
$R \to U \times_S U$ 是单态射。对于 (2)，注意 $R \to U \times_Y U$
局部有限型，因为复合 $R \to U \times_S U$ 局部有限型（《态射》引理
\ref{morphisms-lemma-permanence-finite-type}）。局部有限型的单态射
局部拟有限，因为其纤维有限（《态射》引理
\ref{morphisms-lemma-finite-fibre}），故得 (5)。单态射是分离的
（《概形》引理 \ref{schemes-lemma-monomorphism-separated}），故得 (4)。
\end{proof}

\begin{definition}
\label{spaces-morphisms-definition-separated}
设 $S$ 为概形。
设 $f : X \to Y$ 为 $S$ 上代数空间的态射。
设 $\Delta_{X/Y} : X \to X \times_Y X$ 为对角态射。
\begin{enumerate}
\item 称 $f$ 为 {\it 分离}，若 $\Delta_{X/Y}$ 是闭浸入。
\item 称 $f$ 为 {\it 局部分离}\footnote{在文献中，该术语通常指拟分离且局部分离的态射。}，
若 $\Delta_{X/Y}$ 是浸入。
\item 称 $f$ 为 {\it 拟分离}，若 $\Delta_{X/Y}$ 拟紧。
\end{enumerate}
\end{definition}

\noindent
由于 $\Delta_{X/Y}$ 可表示，该定义是有意义的，因而我们知道它具有
定义中所述性质之一是什么意思。下文（引理
\ref{spaces-morphisms-lemma-match-separated}）将看到，该定义与已有的概形态射及可表示
态射的定义相符。

\begin{lemma}
\label{spaces-morphisms-lemma-trivial-implications}
设 $S$ 为概形，设 $f : X \to Y$ 为 $S$ 上代数空间的态射。
若 $f$ 分离，则 $f$ 局部分离且 $f$ 拟分离。
\end{lemma}

\begin{proof}
这由一般原理成立，即《空间》引理
\ref{spaces-lemma-representable-transformations-property-implication}；
因为概形的闭浸入是浸入且拟紧。
\end{proof}

\begin{lemma}
\label{spaces-morphisms-lemma-base-change-separated}
定义 \ref{spaces-morphisms-definition-separated} 中列出的所有分离公理在基变换下保持。
\end{lemma}

\begin{proof}
设 $f : X \to Y$ 与 $Y' \to Y$ 为代数空间的态射。
设 $f' : X' \to Y'$ 为 $f$ 沿 $Y' \to Y$ 的基变换。则
$\Delta_{X'/Y'}$ 是 $\Delta_{X/Y}$ 沿态射
$X' \times_{Y'} X' \to X \times_Y X$ 的基变换。由第
\ref{spaces-morphisms-section-representable} 节的结果，定义
\ref{spaces-morphisms-definition-separated} 中对角的每个相关性质在基变换下保持。
因此引理成立。
\end{proof}

\begin{lemma}
\label{spaces-morphisms-lemma-fibre-product-after-map}
\begin{slogan}
代数空间的 ``魔法图'' 的上方箭头具有良好的浸入类性质；
在分离性假设下，这些性质会加强。
\end{slogan}
设 $S$ 为概形。设 $f : X \to Z$、$g : Y \to Z$ 以及 $Z \to T$
为 $S$ 上代数空间的态射。考虑诱导的态射
$i : X \times_Z Y \to X \times_T Y$。则
\begin{enumerate}
\item $i$ 可表示、局有限型、局拟有限、分离且为单态射；
\item 若 $Z \to T$ 局部分离，则 $i$ 为浸入；
\item 若 $Z \to T$ 分离，则 $i$ 为闭浸入；
\item 若 $Z \to T$ 拟分离，则 $i$ 拟紧。
\end{enumerate}
\end{lemma}

\begin{proof}
由一般范畴论，以下图表
$$
\xymatrix{
X \times_Z Y \ar[r]_i \ar[d] & X \times_T Y \ar[d] \\
Z \ar[r]^-{\Delta_{Z/T}} \ar[r] & Z \times_T Z
}
$$
是纤维积图表。因此 $i$ 是对角态射 $\Delta_{Z/T}$ 的基变换。
于是由引理 \ref{spaces-morphisms-lemma-properties-diagonal} 以及
第 \ref{spaces-morphisms-section-representable} 节的内容可得本引理。
\end{proof}

\begin{lemma}
\label{spaces-morphisms-lemma-semi-diagonal}
\begin{slogan}
代数空间态射的图的性质，是目标的分离性质的结果。
\end{slogan}
设 $S$ 为概形，设 $T$ 为 $S$ 上的代数空间。
设 $g : X \to Y$ 为 $T$ 上代数空间的态射。
考虑 $g$ 的图 $i : X \to X \times_T Y$。则
\begin{enumerate}
\item $i$ 可表示、局有限型、局拟有限、分离且为单态射；
\item 若 $Y \to T$ 局部分离，则 $i$ 为浸入；
\item 若 $Y \to T$ 分离，则 $i$ 为闭浸入；
\item 若 $Y \to T$ 拟分离，则 $i$ 拟紧。
\end{enumerate}
\end{lemma}

\begin{proof}
这是引理 \ref{spaces-morphisms-lemma-fibre-product-after-map} 应用于态射
$X = X \times_Y Y \to X \times_T Y$ 的特殊情形。
\end{proof}

\begin{lemma}
\label{spaces-morphisms-lemma-section-immersion}
设 $S$ 为概形。
设 $f : X \to T$ 为 $S$ 上代数空间的态射。
设 $s : T \to X$ 为 $f$ 的截面（即满足公式
$f \circ s = \text{id}_T$）。则
\begin{enumerate}
\item $s$ 可表示、局有限型、局拟有限、分离且为单态射；
\item 若 $f$ 局部分离，则 $s$ 为浸入；
\item 若 $f$ 分离，则 $s$ 为闭浸入；
\item 若 $f$ 拟分离，则 $s$ 拟紧。
\end{enumerate}
\end{lemma}

\begin{proof}
这是引理 \ref{spaces-morphisms-lemma-semi-diagonal} 应用于 $g = s$ 的特殊情形，
此时态射为 $i = s : T \to T \times_T X$。
\end{proof}

\begin{lemma}
\label{spaces-morphisms-lemma-composition-separated}
定义 \ref{spaces-morphisms-definition-separated} 中列出的所有分离公理
在态射复合下保持。
\end{lemma}

\begin{proof}
设 $f : X \to Y$ 与 $g : Y \to Z$ 为满足所考察公理适用条件的
代数空间态射。对角 $\Delta_{X/Z}$ 是如下复合：
$$
X \longrightarrow X \times_Y X \longrightarrow X \times_Z X.
$$
我们的分离公理定义为要求对角具有某个性质 $\mathcal{P}$。
由上面的引理 \ref{spaces-morphisms-lemma-fibre-product-after-map} 可知，第二个箭头
也具有此性质。因此引理成立，因为具有性质 $\mathcal{P}$ 的（可表示）
态射的复合仍是具有性质 $\mathcal{P}$ 的态射，见第
\ref{spaces-morphisms-section-representable} 节。
\end{proof}

\begin{lemma}
\label{spaces-morphisms-lemma-separated-over-separated}
设 $S$ 为概形。
设 $f : X \to Y$ 为 $S$ 上代数空间的态射。
\begin{enumerate}
\item 若 $Y$ 分离且 $f$ 分离，则 $X$ 分离。
\item 若 $Y$ 拟分离且 $f$ 拟分离，则 $X$ 拟分离。
\item 若 $Y$ 局部分离且 $f$ 局部分离，则 $X$ 局部分离。
\item 若 $Y$ 在 $S$ 上分离且 $f$ 分离，则 $X$ 在 $S$ 上分离。
\item 若 $Y$ 在 $S$ 上拟分离且 $f$ 拟分离，则 $X$ 在 $S$ 上拟分离。
\item 若 $Y$ 在 $S$ 上局部分离且 $f$ 局部分离，则 $X$ 在 $S$ 上局部分离。
\end{enumerate}
\end{lemma}

\begin{proof}
第 (4)、(5)、(6) 部分立即由引理
\ref{spaces-morphisms-lemma-composition-separated} 以及《空间》定义
\ref{spaces-definition-separated} 得到。
将 $X$ 和 $Y$ 看作 $\Spec(\mathbf{Z})$ 上的代数空间后，
第 (1)、(2)、(3) 部分归约为第 (4)、(5)、(6) 部分；见
《空间的性质》定义 \ref{spaces-properties-definition-separated}。
\end{proof}

\begin{lemma}
\label{spaces-morphisms-lemma-compose-after-separated}
设 $S$ 为概形。
设 $f : X \to Y$ 与 $g : Y \to Z$ 为 $S$ 上代数空间的态射。
\begin{enumerate}
\item 若 $g \circ f$ 分离，则 $f$ 分离。
\item 若 $g \circ f$ 局部分离，则 $f$ 局部分离。
\item 若 $g \circ f$ 拟分离，则 $f$ 拟分离。
\end{enumerate}
\end{lemma}

\begin{proof}
考虑对角态射 $g \circ f$ 的分解
$$
X \to X \times_Y X \to X \times_Z X
$$
无论哪种情形，最后一个态射都是单态射。因此对于任意概形 $T$
以及态射 $T \to X \times_Y X$，都有等式
$$
X \times_{(X \times_Y X)} T = X \times_{(X \times_Z X)} T.
$$
所以结论显然。
\end{proof}

\begin{lemma}
\label{spaces-morphisms-lemma-separated-implies-morphism-separated}
设 $S$ 为概形，设 $X$ 为 $S$ 上的代数空间。
\begin{enumerate}
\item 若 $X$ 分离，则 $X$ 在 $S$ 上分离。
\item 若 $X$ 局部分离，则 $X$ 在 $S$ 上局部分离。
\item 若 $X$ 拟分离，则 $X$ 在 $S$ 上拟分离。
\end{enumerate}
设 $f : X \to Y$ 为 $S$ 上代数空间的态射。
\begin{enumerate}
\item[(4)] 若 $X$ 在 $S$ 上分离，则 $f$ 分离。
\item[(5)] 若 $X$ 在 $S$ 上局部分离，则 $f$ 局部分离。
\item[(6)] 若 $X$ 在 $S$ 上拟分离，则 $f$ 拟分离。
\end{enumerate}
\end{lemma}

\begin{proof}
第 (4)、(5)、(6) 部分立即由引理
\ref{spaces-morphisms-lemma-compose-after-separated} 以及《空间》定义
\ref{spaces-definition-separated} 得到。
将 $X$ 和 $Y$ 看作 $\Spec(\mathbf{Z})$ 上的代数空间后，
第 (1)、(2)、(3) 部分由第 (4)、(5)、(6) 部分得到；见
《空间的性质》定义 \ref{spaces-properties-definition-separated}。
\end{proof}

\begin{lemma}
\label{spaces-morphisms-lemma-separated-local}
设 $S$ 为概形。
设 $f : X \to Y$ 为 $S$ 上代数空间的态射。
设 $\mathcal{P}$ 为定义 \ref{spaces-morphisms-definition-separated} 中的任一分离公理。
以下命题等价：
\begin{enumerate}
\item $f$ 具有性质 $\mathcal{P}$；
\item 对每个概形 $Z$ 及态射 $Z \to Y$，$f$ 的基变换
$Z \times_Y X \to Z$ 具有性质 $\mathcal{P}$；
\item 对每个仿射概形 $Z$ 及每个态射 $Z \to Y$，$f$ 的基变换
$Z \times_Y X \to Z$ 具有性质 $\mathcal{P}$；
\item 对每个仿射概形 $Z$ 及每个态射 $Z \to Y$，代数空间
$Z \times_Y X$ 具有性质 $\mathcal{P}$（见《空间的性质》定义
\ref{spaces-properties-definition-separated}）；
\item 存在概形 $V$ 及满射 \'etale 态射 $V \to Y$，使得基变换
$V \times_Y X \to V$ 具有性质 $\mathcal{P}$；
\item 存在 Zariski 覆盖 $Y = \bigcup Y_i$，使得每个态射
$f^{-1}(Y_i) \to Y_i$ 都具有性质 $\mathcal{P}$。
\end{enumerate}
\end{lemma}

\begin{proof}
下文将反复使用引理 \ref{spaces-morphisms-lemma-base-change-separated}，不再另行说明。
特别地，显然 (1) 蕴含 (2)，而 (2) 蕴含 (3)。

\medskip\noindent
证明 (3) 与 (4) 等价。注意，若 $Z$ 为仿射概形，则态射
$Z \to \Spec(\mathbf{Z})$ 作为 $\Spec(\mathbf{Z})$ 上代数空间的态射是分离的。
若 $Z \times_Y X \to Z$ 具有性质 $\mathcal{P}$，则由复合（见引理
\ref{spaces-morphisms-lemma-composition-separated}），$Z \times_Y X \to \Spec(\mathbf{Z})$
具有性质 $\mathcal{P}$。因此代数空间 $Z \times_Y X$ 具有性质 $\mathcal{P}$。
反之，若代数空间 $Z \times_Y X$ 具有性质 $\mathcal{P}$，则
$Z \times_Y X \to \Spec(\mathbf{Z})$ 具有性质 $\mathcal{P}$，从而由引理
\ref{spaces-morphisms-lemma-compose-after-separated} 可知 $Z \times_Y X \to Z$ 具有性质
$\mathcal{P}$。

\medskip\noindent
证明 (3) 蕴含 (5)。假设 (3)。设 $V$ 为概形，且
$V \to Y$ 为满射 \'etale 态射。我们需证明
$V \times_Y X \to V$ 具有性质 $\mathcal{P}$。换言之，需证明态射
$$
V \times_Y X \longrightarrow
(V \times_Y X) \times_V (V \times_Y X) = V \times_Y X \times_Y X
$$
具有相应性质（即为闭浸入、浸入或拟紧）。设
$V = \bigcup V_j$ 为 $V$ 的仿射开覆盖。由假设，每个态射
$$
V_j \times_Y X \longrightarrow V_j \times_Y X \times_Y X
$$
都具有相应性质。由于“闭浸入”、“浸入”、“拟紧浸入”或“拟紧”
在目标上是 Zariski 局部的，且 $V_j$ 覆盖 $V$，故得所需结论。

\medskip\noindent
证明 (5) 蕴含 (1)。设 $V \to Y$ 如 (5) 中所述。
则有纤维积图表
$$
\xymatrix{
V \times_Y X \ar[r] \ar[d] &
X \ar[d] \\
V \times_Y X \times_Y X \ar[r] &
X \times_Y X
}
$$
由假设，左侧竖直箭头是闭浸入、浸入、拟紧浸入或拟紧。
由《空间》引理 \ref{spaces-lemma-descent-representable-transformations-property}
可知，右侧竖直箭头也分别是闭浸入、浸入、拟紧浸入或拟紧。

\medskip\noindent
(1) 蕴含 (6) 是显然的，只需取覆盖 $Y = Y$。
假设 $Y = \bigcup Y_i$ 如 (6) 中所述。取概形 $V_i$ 及满射 \'etale
态射 $V_i \to Y_i$。注意，态射 $V_i \times_Y X \to V_i$ 具有
$\mathcal{P}$，因为它们是态射 $f^{-1}(Y_i) \to Y_i$ 的基变换。
令 $V = \coprod V_i$。则 $V \to Y$ 是 (5) 中的态射（细节略）。
因此 (6) 蕴含 (5)，证明完毕。
\end{proof}

\begin{lemma}
\label{spaces-morphisms-lemma-match-separated}
设 $S$ 为概形。
设 $f : X \to Y$ 为 $S$ 上代数空间的可表示态射。
\begin{enumerate}
\item 态射 $f$ 局部分离。
\item 态射 $f$ 按照上述定义 \ref{spaces-morphisms-definition-separated} 的意义
（拟）分离，当且仅当 $f$ 按照第 \ref{spaces-morphisms-section-representable} 节的意义
（拟）分离。
\end{enumerate}
特别地，若 $f : X \to Y$ 为 $S$ 上概形的态射，则 $f$ 按照定义
\ref{spaces-morphisms-definition-separated} 的意义（拟）分离，当且仅当 $f$ 作为概形态射
（拟）分离。
\end{lemma}

\begin{proof}
这是引理 \ref{spaces-morphisms-lemma-separated-local} 中 (1) 与 (2) 的等价性，
结合任意概形态射都局部分离这一事实得到；见《概形》引理
\ref{schemes-lemma-diagonal-immersion}。
\end{proof}











\section{满射态射}
\label{spaces-morphisms-section-surjective}

\noindent
我们已在第 \ref{spaces-morphisms-section-representable} 节定义了代数空间的可表示态射
为满射的含义。

\begin{lemma}
\label{spaces-morphisms-lemma-surjective-representable}
设 $S$ 为概形。设 $f : X \to Y$ 为 $S$ 上代数空间的可表示态射。
则 $f$ 按照第 \ref{spaces-morphisms-section-representable} 节的意义为满射，当且仅当
$|f| : |X| \to |Y|$ 为满射。
\end{lemma}

\begin{proof}
事实上，若 $f : X \to Y$ 可表示，则它为满射，当且仅当对于每个概形
$T$ 及每个态射 $T \to Y$，$f$ 的基变换
$f_T : T \times_Y X \to T$ 是概形的满射态射，换言之，当且仅当
$|f_T|$ 为满射。由《空间的性质》引理
\ref{spaces-properties-lemma-points-cartesian}，映射
$|T \times_Y X| \to |T| \times_{|Y|} |X|$ 总是满射。
因此，若 $|f| : |X| \to |Y|$ 为满射，则
$|f_T| : |T \times_Y X| \to |T|$ 为满射。反之，若对上述每个
$T \to Y$，$|f_T|$ 都为满射，则取 $T$ 为某个域的谱，便得到
$|X| \to |Y|$ 为满射。
\end{proof}

\noindent
这为下述定义扫清了道路。

\begin{definition}
\label{spaces-morphisms-definition-surjective}
设 $S$ 为概形。设 $f : X \to Y$ 为 $S$ 上代数空间的态射。
若关联拓扑空间的映射 $|f| : |X| \to |Y|$ 为满射，
则称 $f$ 为 {\it 满射}。
\end{definition}

\begin{lemma}
\label{spaces-morphisms-lemma-surjective-local}
设 $S$ 为概形。
设 $f : X \to Y$ 为 $S$ 上代数空间的态射。
以下命题等价：
\begin{enumerate}
\item $f$ 为满射；
\item 对每个概形 $Z$ 及任意态射 $Z \to Y$，态射
$Z \times_Y X \to Z$ 为满射；
\item 对每个仿射概形 $Z$ 及任意态射 $Z \to Y$，态射
$Z \times_Y X \to Z$ 为满射；
\item 存在概形 $V$ 及满射 \'etale 态射 $V \to Y$，使得
$V \times_Y X \to V$ 为满射态射；
\item 存在概形 $U$ 及满射 \'etale 态射
$\varphi : U \to X$，使得复合 $f \circ \varphi$ 为满射；
\item 存在交换图表
$$
\xymatrix{
U \ar[d] \ar[r] & V \ar[d] \\
X \ar[r] & Y
}
$$
其中 $U$、$V$ 为概形，竖直箭头为满射 \'etale 态射，
且上方水平箭头为满射；以及
\item 存在 Zariski 覆盖 $Y = \bigcup Y_i$，使得每个态射
$f^{-1}(Y_i) \to Y_i$ 都为满射。
\end{enumerate}
\end{lemma}

\begin{proof}
略。
\end{proof}

\begin{lemma}
\label{spaces-morphisms-lemma-composition-surjective}
满射态射的复合是满射。
\end{lemma}

\begin{proof}
这直接由定义得到。
\end{proof}

\begin{lemma}
\label{spaces-morphisms-lemma-base-change-surjective}
满射态射的基变换是满射。
\end{lemma}

\begin{proof}
直接由《空间的性质》引理
\ref{spaces-properties-lemma-points-cartesian} 得到。
\end{proof}










\section{开态射}
\label{spaces-morphisms-section-open}

\noindent
对于代数空间的可表示态射，我们已在第 \ref{spaces-morphisms-section-representable} 节
定义了“普遍开”的含义。因此，在给出自然定义之前，先检查它在可表示
情形下与该定义一致。

\begin{lemma}
\label{spaces-morphisms-lemma-characterize-representable-universally-open}
设 $S$ 为概形。设 $f : X \to Y$ 为 $S$ 上代数空间的可表示态射。
以下命题等价：
\begin{enumerate}
\item $f$ 按照第 \ref{spaces-morphisms-section-representable} 节的意义是普遍开；
\item 对每个代数空间态射 $Z \to Y$，拓扑空间的映射
$|Z \times_Y X| \to |Z|$ 是开映射。
\end{enumerate}
\end{lemma}

\begin{proof}
假设 (1)，并设 $Z \to Y$ 如 (2) 中所述。取概形 $V$ 及满射 \'etale
态射 $V \to Y$。由假设，概形态射 $V \times_Y X \to V$ 普遍开。
由《空间的性质》第 \ref{spaces-properties-section-points} 节，
在交换图表
$$
\xymatrix{
|V \times_Y X| \ar[r] \ar[d] & |Z \times_Y X| \ar[d] \\
|V| \ar[r] & |Z|
}
$$
水平箭头是开且满射的，此外
$$
|V \times_Y X| \longrightarrow |V| \times_{|Z|} |Z \times_Y X|
$$
是满射。因此，由于左侧竖直箭头是开映射，右侧竖直箭头也是开映射。
这证明了 (2)。蕴含 (2) $\Rightarrow$ (1) 直接由定义得到。
\end{proof}

\noindent
因此可以采用下述自然定义。

\begin{definition}
\label{spaces-morphisms-definition-open}
设 $S$ 为概形。设 $f : X \to Y$ 为 $S$ 上代数空间的态射。
\begin{enumerate}
\item 若拓扑空间的映射 $|f| : |X| \to |Y|$ 为开映射，
则称 $f$ 为 {\it 开}。
\item 若对每个代数空间态射 $Z \to Y$，拓扑空间的态射
$$
|Z \times_Y X| \to |Z|
$$
为开映射，即基变换 $Z \times_Y X \to Z$ 为开态射，
则称 $f$ 为 {\it 普遍开}。
\end{enumerate}
\end{definition}

\noindent
注意，代数空间的 \'etale 态射是普遍开的；见《空间的性质》定义
\ref{spaces-properties-definition-etale} 以及引理
\ref{spaces-properties-lemma-etale-open} 和
\ref{spaces-properties-lemma-base-change-etale}。

\begin{lemma}
\label{spaces-morphisms-lemma-base-change-universally-open}
代数空间的普遍开态射沿任意代数空间态射的基变换仍是普遍开。
\end{lemma}

\begin{proof}
这直接由定义得到。
\end{proof}

\begin{lemma}
\label{spaces-morphisms-lemma-composition-open}
一对（普遍）开代数空间态射的复合是（普遍）开的。
\end{lemma}

\begin{proof}
略。
\end{proof}

\begin{lemma}
\label{spaces-morphisms-lemma-universally-open-local}
设 $S$ 为概形。设 $f : X \to Y$ 为 $S$ 上代数空间的态射。
以下命题等价：
\begin{enumerate}
\item $f$ 普遍开；
\item 对每个概形 $Z$ 及每个态射 $Z \to Y$，投影
$|Z \times_Y X| \to |Z|$ 为开映射；
\item 对每个仿射概形 $Z$ 及每个态射 $Z \to Y$，投影
$|Z \times_Y X| \to |Z|$ 为开映射；以及
\item 存在概形 $V$ 及满射 \'etale 态射 $V \to Y$，使得
$V \times_Y X \to V$ 为代数空间的普遍开态射；以及
\item 存在 Zariski 覆盖 $Y = \bigcup Y_i$，使得每个态射
$f^{-1}(Y_i) \to Y_i$ 都普遍开。
\end{enumerate}
\end{lemma}

\begin{proof}
我们略去 (1) 蕴含 (2) 及 (2) 蕴含 (3) 的证明。

\medskip\noindent
假设 (3)。取满射 \'etale 态射 $V \to Y$。我们将证明
$V \times_Y X \to V$ 是代数空间的普遍开态射。设 $Z \to V$ 为
从代数空间到 $V$ 的态射。取满射 \'etale 态射 $W \to Z$，其中
$W = \coprod W_i$ 是仿射概形的不交并，见《空间的性质》引理
\ref{spaces-properties-lemma-cover-by-union-affines}。于是有如下交换图表
$$
\xymatrix{
\coprod_i |W_i \times_Y X| \ar@{=}[r] \ar[d] &
|W \times_Y X| \ar[r] \ar[d] &
|Z \times_Y X| \ar[d] \ar@{=}[r] &
|Z \times_V (V \times_Y X)| \ar[ld] \\
\coprod |W_i| \ar@{=}[r] &
|W| \ar[r] &
|Z|
}
$$
我们需证明东南箭头为开映射。中间的水平箭头是满射且开映射
（《空间的性质》引理 \ref{spaces-properties-lemma-etale-open}）。
由假设 (3) 及 $W_i$ 为仿射概形这一事实，左侧竖直箭头是开映射。
因此右侧竖直箭头也是开映射。

\medskip\noindent
假设 $V \to Y$ 如 (4) 中所述。我们将证明 $f$ 普遍开。
设 $Z \to Y$ 为代数空间的态射。考虑图表
$$
\xymatrix{
|(V \times_Y Z) \times_V (V \times_Y X)| \ar@{=}[r] \ar[rd] &
|V \times_Y X| \ar[r] \ar[d] &
|Z \times_Y X| \ar[d] \\
 &
|V \times_Y Z| \ar[r] &
|Z|
}
$$
西南箭头由假设为开映射。水平箭头是满射且开映射，因为相应的
代数空间态射是 \'etale（见《空间的性质》引理
\ref{spaces-properties-lemma-etale-open}）。因此右侧竖直箭头为开映射。

\medskip\noindent
显然，取覆盖 $Y = Y$ 即得 (1) 蕴含 (5)。
假设 $Y = \bigcup Y_i$ 如 (5) 中所述。则对于任意 $Z \to Y$，
可得到相应的 Zariski 覆盖 $Z = \bigcup Z_i$，使得 $f$ 到 $Z_i$
的基变换为开映射。由简单的拓扑论证，$Z \times_Y X \to Z$
为开映射。因此 (1) 成立。
\end{proof}

\begin{lemma}
\label{spaces-morphisms-lemma-space-over-field-universally-open}
设 $S$ 为概形。设 $p : X \to \Spec(k)$ 为 $S$ 上代数空间的态射，
其中 $k$ 为域。则 $p : X \to \Spec(k)$ 普遍开。
\end{lemma}

\begin{proof}
取概形 $U$ 及满射 \'etale 态射 $U \to X$。由《态射》引理
\ref{morphisms-lemma-scheme-over-field-universally-open}，复合
$U \to \Spec(k)$（作为概形态射）普遍开。设 $Z \to \Spec(k)$
为概形态射。则
$U \times_{\Spec(k)} Z \to X \times_{\Spec(k)} Z$ 为满射，见引理
\ref{spaces-morphisms-lemma-base-change-surjective}。因此下列映射中的第一个
$$
|U \times_{\Spec(k)} Z| \to |X \times_{\Spec(k)} Z| \to |Z|
$$
是满射。由于复合由上述结果为开映射，故第二个映射也为开映射。
于是由引理 \ref{spaces-morphisms-lemma-universally-open-local}，$p$ 普遍开。
\end{proof}








\section{商映射态射}
\label{spaces-morphisms-section-submersive}

\noindent
对于代数空间的可表示态射，我们已在第 \ref{spaces-morphisms-section-representable} 节
定义了“普遍商映射”的含义。因此，在给出自然定义之前，先检查它在可表示
情形下与该定义一致。

\begin{lemma}
\label{spaces-morphisms-lemma-characterize-representable-universally-submersive}
设 $S$ 为概形。设 $f : X \to Y$ 为 $S$ 上代数空间的可表示态射。
以下命题等价：
\begin{enumerate}
\item $f$ 按照第 \ref{spaces-morphisms-section-representable} 节的意义是普遍商映射；
\item 对每个代数空间态射 $Z \to Y$，拓扑空间的映射
$|Z \times_Y X| \to |Z|$ 为商映射。
\end{enumerate}
\end{lemma}

\begin{proof}
假设 (1)，并设 $Z \to Y$ 如 (2) 中所述。取概形 $V$ 及满射 \'etale
态射 $V \to Y$。由假设，概形态射 $V \times_Y X \to V$ 是普遍商映射。
由《空间的性质》第 \ref{spaces-properties-section-points} 节，在交换图表
$$
\xymatrix{
|V \times_Y X| \ar[r] \ar[d] & |Z \times_Y X| \ar[d] \\
|V| \ar[r] & |Z|
}
$$
水平箭头是开且满射的，此外
$$
|V \times_Y X| \longrightarrow |V| \times_{|Z|} |Z \times_Y X|
$$
为满射。因此由于左侧竖直箭头为商映射，右侧竖直箭头也为商映射。
这证明了 (2)。蕴含 (2) $\Rightarrow$ (1) 直接由定义得到。
\end{proof}

\noindent
因此可以采用下述自然定义。

\begin{definition}
\label{spaces-morphisms-definition-submersive}
设 $S$ 为概形。
设 $f : X \to Y$ 为 $S$ 上代数空间的态射。
\begin{enumerate}
\item 若连续映射 $|X| \to |Y|$ 为商映射，见《拓扑》定义
\ref{topology-definition-submersive}，则称 $f$ 为 {\it 商映射}\footnote{这与
微分流形的浸没概念有很大不同。}。
\item 若对于每个代数空间态射 $Y' \to Y$，基变换
$Y' \times_Y X \to Y'$ 为商映射，则称 $f$ 为 {\it 普遍商映射}。
\end{enumerate}
\end{definition}

\noindent
注意，商映射态射特别是满射。

\begin{lemma}
\label{spaces-morphisms-lemma-base-change-universally-submersive}
代数空间的普遍商映射态射沿任意代数空间态射的基变换仍是普遍商映射。
\end{lemma}

\begin{proof}
这直接由定义得到。
\end{proof}

\begin{lemma}
\label{spaces-morphisms-lemma-composition-universally-submersive}
一对（普遍）商映射代数空间态射的复合是（普遍）商映射。
\end{lemma}

\begin{proof}
略。
\end{proof}












\section{拟紧态射}
\label{spaces-morphisms-section-quasi-compact}

\noindent
由第 \ref{spaces-morphisms-section-representable} 节，我们知道代数空间的可表示态射
为拟紧的含义。为给出一般代数空间态射的定义，先作如下观察。

\begin{lemma}
\label{spaces-morphisms-lemma-characterize-representable-quasi-compact}
设 $S$ 为概形。
设 $f : X \to Y$ 为 $S$ 上代数空间的可表示态射。
以下命题等价：
\begin{enumerate}
\item $f$ 按照第 \ref{spaces-morphisms-section-representable} 节的意义为拟紧；
\item 对每个拟紧代数空间 $Z$ 及任意态射 $Z \to Y$，代数空间
$Z \times_Y X$ 拟紧。
\end{enumerate}
\end{lemma}

\begin{proof}
假设 (1)，并设 $Z \to Y$ 为代数空间态射，其中 $Z$ 拟紧。由《空间的性质》
定义 \ref{spaces-properties-definition-quasi-compact}，存在拟紧概形 $U$
及满射 \'etale 态射 $U \to Z$。由于 $f$ 可表示且拟紧，按定义
$U \times_Y X$ 是概形，且 $U \times_Y X \to U$ 拟紧。因此
$U \times_Y X$ 是拟紧概形。态射 $U \times_Y X \to Z \times_Y X$
是 \'etale 且满射的（因为它是可表示的满射 \'etale 态射 $U \to Z$
的基变换，见第 \ref{spaces-morphisms-section-representable} 节）。所以按定义
$Z \times_Y X$ 拟紧。

\medskip\noindent
假设 (2)。设 $Z \to Y$ 为态射，其中 $Z$ 为概形。我们需证明
$p : Z \times_Y X \to Z$ 拟紧。设 $U \subset Z$ 为仿射开集。
则 $p^{-1}(U) = U \times_Y Z$，且由假设 (2)，概形
$U \times_Y Z$ 拟紧。因此 $p$ 拟紧，见《概形》第
\ref{schemes-section-quasi-compact} 节。
\end{proof}

\noindent
这促成下述定义。

\begin{definition}
\label{spaces-morphisms-definition-quasi-compact}
设 $S$ 为概形。
设 $f : X \to Y$ 为 $S$ 上代数空间的态射。
若对每个拟紧代数空间 $Z$ 及态射 $Z \to Y$，纤维积
$Z \times_Y X$ 拟紧，则称 $f$ 为 {\it 拟紧}。
\end{definition}

\noindent
由上面的引理 \ref{spaces-morphisms-lemma-characterize-representable-quasi-compact}，
该定义与代数空间可表示态射的既有概念一致。

\begin{lemma}
\label{spaces-morphisms-lemma-quasi-compact-is-quasi-compact}
设 $S$ 为概形。若 $f : X \to Y$ 为 $S$ 上代数空间的拟紧态射，则其
底层拓扑空间映射 $|f| : |X| \to |Y|$ 拟紧。
\end{lemma}

\begin{proof}
设 $V \subset |Y|$ 为拟紧开集。由《空间的性质》引理
\ref{spaces-properties-lemma-open-subspaces}，存在开子空间
$Y' \subset Y$，使 $V = |Y'|$。由《空间的性质》引理
\ref{spaces-properties-lemma-quasi-compact-space}，$Y'$ 是拟紧代数空间，
因而由定义 \ref{spaces-morphisms-definition-quasi-compact}，$X' = Y' \times_Y X$ 是拟紧
代数空间。另一方面，$X' \subset X$ 是开子空间（《空间》引理
\ref{spaces-lemma-base-change-immersions}），且由《空间的性质》引理
\ref{spaces-properties-lemma-points-cartesian}，
$|X'| = |f|^{-1}(|X'|) = |f|^{-1}(V)$。再次应用《空间的性质》引理
\ref{spaces-properties-lemma-quasi-compact-space}，即得 $|X'|$ 是
$|X|$ 的拟紧开集。
\end{proof}

\begin{lemma}
\label{spaces-morphisms-lemma-base-change-quasi-compact}
代数空间的拟紧态射沿任意代数空间态射的基变换仍是拟紧。
\end{lemma}

\begin{proof}
略。提示：纤维积的传递性。
\end{proof}

\begin{lemma}
\label{spaces-morphisms-lemma-composition-quasi-compact}
一对拟紧代数空间态射的复合是拟紧。
\end{lemma}

\begin{proof}
略。提示：纤维积的传递性。
\end{proof}

\begin{lemma}
\label{spaces-morphisms-lemma-surjection-from-quasi-compact}
\begin{slogan}
拟紧代数空间在满射态射下的像是拟紧的。
\end{slogan}
设 $S$ 为概形。
\begin{enumerate}
\item 若 $X \to Y$ 是 $S$ 上代数空间的满射态射，且 $X$ 拟紧，
则 $Y$ 拟紧。
\item 若
$$
\xymatrix{
X \ar[rr]_f \ar[rd]_p & &
Y \ar[dl]^q \\
& Z
}
$$
是 $S$ 上代数空间态射的交换图表，且 $f$ 满射、$p$ 拟紧，则 $q$ 拟紧。
\end{enumerate}
\end{lemma}

\begin{proof}
假设 $X$ 拟紧且 $X \to Y$ 满射。由定义 \ref{spaces-morphisms-definition-surjective}，
$|X| \to |Y|$ 满射。因此由《空间的性质》引理
\ref{spaces-properties-lemma-quasi-compact-space} 以及拟紧空间在连续
映射下的像仍拟紧这一拓扑事实，$Y$ 拟紧；见《拓扑》引理
\ref{topology-lemma-image-quasi-compact}。
设 $f,p,q$ 如 (2) 中所述。设 $T \to Z$ 为态射，且其源是拟紧代数空间。
由假设，$T \times_Z X$ 拟紧。由引理 \ref{spaces-morphisms-lemma-base-change-surjective}，
态射 $T \times_Z X \to T \times_Z Y$ 满射。因此由第 (1) 部分，
$T \times_Z Y$ 也拟紧。故 $q$ 拟紧。
\end{proof}

\begin{lemma}
\label{spaces-morphisms-lemma-descent-quasi-compact}
设 $S$ 为概形。
设 $f : X \to Y$ 为 $S$ 上代数空间的态射。
设 $g : Y' \to Y$ 为代数空间的普遍开满射，且其基变换
$f' : X' \to Y'$ 拟紧。则 $f$ 拟紧。
\end{lemma}

\begin{proof}
设 $Z \to Y$ 为代数空间态射，且 $Z$ 拟紧。由于 $g$ 普遍开且满射，
$Y' \times_Y Z \to Z$ 开且满射。由于 $|Y' \times_Y Z|$ 的每一点
都有拟紧开邻域的基本系（见《空间的性质》引理
\ref{spaces-properties-lemma-space-locally-quasi-compact}），可取
$W \subset |Y' \times_Y Z|$ 中的拟紧开集，使其映到 $Z$ 满射。记
$f'' : W \times_Y X \to W$ 为 $f'$ 沿 $W \to Y'$ 的基变换。
由假设，$W \times_Y X$ 拟紧。由于 $W \to Z$ 满射，
$W \times_Y X \to Z \times_Y X$ 满射。因此由引理
\ref{spaces-morphisms-lemma-surjection-from-quasi-compact}，$Z \times_Y X$ 拟紧。
故 $f$ 拟紧。
\end{proof}

\begin{lemma}
\label{spaces-morphisms-lemma-quasi-compact-local}
\begin{slogan}
代数空间的拟紧态射在拉回下保持，并且在目标上是局部性质。
\end{slogan}
设 $S$ 为概形。
设 $f : X \to Y$ 为 $S$ 上代数空间的态射。
以下命题等价：
\begin{enumerate}
\item $f$ 拟紧；
\item 对每个概形 $Z$ 及任意态射 $Z \to Y$，代数空间态射
$Z \times_Y X \to Z$ 拟紧；
\item 对每个仿射概形 $Z$ 及任意态射 $Z \to Y$，代数空间
$Z \times_Y X$ 拟紧；
\item 存在概形 $V$ 及满射 \'etale 态射 $V \to Y$，使得
$V \times_Y X \to V$ 为拟紧代数空间态射；以及
\item 存在代数空间的满射 \'etale 态射 $Y' \to Y$，使得
$Y' \times_Y X \to Y'$ 为拟紧代数空间态射；以及
\item 存在 Zariski 覆盖 $Y = \bigcup Y_i$，使得每个态射
$f^{-1}(Y_i) \to Y_i$ 拟紧。
\end{enumerate}
\end{lemma}

\begin{proof}
下文将使用引理 \ref{spaces-morphisms-lemma-base-change-quasi-compact}，不再另行说明。
显然 (1) 蕴含 (2)，(2) 蕴含 (3)。假设 (3)。设 $Z$ 为 $S$ 上的拟紧代数空间，
且 $Z \to Y$ 为态射。由《空间的性质》引理
\ref{spaces-properties-lemma-quasi-compact-affine-cover}，存在仿射概形 $U$
及满射 \'etale 态射 $U \to Z$。于是 $U \times_Y X \to Z \times_Y X$
是代数空间的满射态射，见引理 \ref{spaces-morphisms-lemma-base-change-surjective}。
由假设，$|U \times_Y X|$ 拟紧。它满射到 $|Z \times_Y X|$，故由《拓扑》引理
\ref{topology-lemma-image-quasi-compact}，$|Z \times_Y X|$ 拟紧。
这证明了 (3) 蕴含 (1)。

\medskip\noindent
蕴含 (1) $\Rightarrow$ (4) 及 (4) $\Rightarrow$ (5) 是显然的。
蕴含 (5) $\Rightarrow$ (1) 由引理 \ref{spaces-morphisms-lemma-descent-quasi-compact} 以及
代数空间的 \'etale 态射普遍开这一事实得到（见定义
\ref{spaces-morphisms-definition-open} 后的讨论）。

\medskip\noindent
显然，取覆盖 $Y = Y$ 即得 (1) 蕴含 (6)。假设 $Y = \bigcup Y_i$
如 (6) 中所述。设 $Z$ 为仿射概形，且 $Z \to Y$ 为态射。
则存在有限标准仿射覆盖 $Z = Z_1 \cup \ldots \cup Z_n$，使得每个
$Z_j \to Y$ 对某个 $i_j$ 经过 $Y_{i_j}$ 分解。因此代数空间
$$
Z_j \times_Y X = Z_j \times_{Y_{i_j}} f^{-1}(Y_{i_j})
$$
拟紧。由于
$Z \times_Y X = \bigcup_{j = 1, \ldots, n} Z_j \times_Y X$
是 Zariski 覆盖，故
$|Z \times_Y X| = \bigcup_{j = 1, \ldots, n} |Z_j \times_Y X|$
（见《空间的性质》引理 \ref{spaces-properties-lemma-open-subspaces}）
是拟紧空间的有限并，从而拟紧。因此 (6) 蕴含 (3)。
\end{proof}

\noindent
下述引理（以及下一个引理）特别保证：只要 $X$ 是拟紧代数空间，
态射 $X \to \Spec(A)$ 就是拟紧的。

\begin{lemma}
\label{spaces-morphisms-lemma-quasi-compact-permanence}
设 $S$ 为概形。
设 $f : X \to Y$ 与 $g : Y \to Z$ 为 $S$ 上代数空间的态射。
若 $g \circ f$ 拟紧且 $g$ 拟分离，则 $f$ 拟紧。
\end{lemma}

\begin{proof}
这是因为 $f$ 等于复合
$(1, f) : X \to X \times_Z Y \to Y$。第一个映射是拟紧的，
由引理 \ref{spaces-morphisms-lemma-section-immersion}，因为它是拟分离态射
$X \times_Z Y \to X$ 的截面（该态射是 $g$ 的基变换，见引理
\ref{spaces-morphisms-lemma-base-change-separated}）。第二个映射是拟紧的，因为它是
$g \circ f$ 的基变换，见引理 \ref{spaces-morphisms-lemma-base-change-quasi-compact}。
拟紧态射的复合是拟紧的，见引理 \ref{spaces-morphisms-lemma-composition-quasi-compact}。
\end{proof}

\begin{lemma}
\label{spaces-morphisms-lemma-quasi-compact-quasi-separated-permanence}
设 $f : X \to Y$ 为概形 $S$ 上代数空间的态射。
\begin{enumerate}
\item 若 $X$ 拟紧且 $Y$ 拟分离，则 $f$ 拟紧。
\item 若 $X$ 拟紧且拟分离，且 $Y$ 拟分离，则 $f$ 拟紧且拟分离。
\item 拟紧且拟分离代数空间的纤维积拟紧且拟分离。
\end{enumerate}
\end{lemma}

\begin{proof}
第 (1) 部分由引理 \ref{spaces-morphisms-lemma-quasi-compact-permanence}（取
$Z = S = \Spec(\mathbf{Z})$）得到。第 (2) 部分由 (1) 及引理
\ref{spaces-morphisms-lemma-compose-after-separated} 得到。对于 (3)，设 $X \to Y$ 与
$Z \to Y$ 为拟紧且拟分离代数空间的态射。由 (2) 及引理
\ref{spaces-morphisms-lemma-base-change-quasi-compact} 和
\ref{spaces-morphisms-lemma-base-change-separated}，作为 $X \to Y$ 的基变换，
$X \times_Y Z \to Z$ 拟紧且拟分离。因此 $X \times_Y Z$ 作为
$Z$ 上拟紧且拟分离的代数空间本身拟紧且拟分离；见引理
\ref{spaces-morphisms-lemma-separated-over-separated} 和
\ref{spaces-morphisms-lemma-composition-quasi-compact}。
\end{proof}


\section{普遍闭态射}
\label{spaces-morphisms-section-universally-closed}

\noindent
对于代数空间的可表示态射，我们已在第 \ref{spaces-morphisms-section-representable} 节
定义了“普遍闭”的含义。因此，在给出自然定义之前，先检查它在可表示
情形下与该定义一致。

\begin{lemma}
\label{spaces-morphisms-lemma-characterize-representable-universally-closed}
设 $S$ 为概形。设 $f : X \to Y$ 为 $S$ 上代数空间的可表示态射。
以下命题等价：
\begin{enumerate}
\item $f$ 按照第 \ref{spaces-morphisms-section-representable} 节的意义普遍闭；
\item 对每个代数空间态射 $Z \to Y$，拓扑空间的映射
$|Z \times_Y X| \to |Z|$ 为闭映射。
\end{enumerate}
\end{lemma}

\begin{proof}
假设 (1)，并设 $Z \to Y$ 如 (2) 中所述。取概形 $V$ 及满射 \'etale
态射 $V \to Y$。由假设，概形态射 $V \times_Y X \to V$ 普遍闭。
由《空间的性质》第 \ref{spaces-properties-section-points} 节，在交换图表
$$
\xymatrix{
|V \times_Y X| \ar[r] \ar[d] & |Z \times_Y X| \ar[d] \\
|V| \ar[r] & |Z|
}
$$
水平箭头是开且满射的，此外
$$
|V \times_Y X| \longrightarrow |V| \times_{|Z|} |Z \times_Y X|
$$
为满射。因此，由于左侧竖直箭头是闭映射，右侧竖直箭头也是闭映射。
这证明了 (2)。蕴含 (2) $\Rightarrow$ (1) 直接由定义得到。
\end{proof}

\noindent
因此可以采用下述自然定义。

\begin{definition}
\label{spaces-morphisms-definition-closed}
设 $S$ 为概形。设 $f : X \to Y$ 为 $S$ 上代数空间的态射。
\begin{enumerate}
\item 若拓扑空间映射 $|X| \to |Y|$ 为闭映射，则称 $f$ 为 {\it 闭}。
\item 若对每个代数空间态射 $Z \to Y$，拓扑空间的态射
$$
|Z \times_Y X| \to |Z|
$$
为闭映射，即基变换 $Z \times_Y X \to Z$ 为闭态射，则称 $f$ 为
{\it 普遍闭}。
\end{enumerate}
\end{definition}

\begin{lemma}
\label{spaces-morphisms-lemma-base-change-universally-closed}
代数空间的普遍闭态射沿任意代数空间态射的基变换仍是普遍闭。
\end{lemma}

\begin{proof}
这直接由定义得到。
\end{proof}

\begin{lemma}
\label{spaces-morphisms-lemma-composition-universally-closed}
一对（普遍）闭代数空间态射的复合是（普遍）闭的。
\end{lemma}

\begin{proof}
略。
\end{proof}

\begin{lemma}
\label{spaces-morphisms-lemma-universally-closed-local}
设 $S$ 为概形。设 $f : X \to Y$ 为 $S$ 上代数空间的态射。
以下命题等价：
\begin{enumerate}
\item $f$ 普遍闭；
\item 对每个概形 $Z$ 及每个态射 $Z \to Y$，投影
$|Z \times_Y X| \to |Z|$ 为闭映射；
\item 对每个仿射概形 $Z$ 及每个态射 $Z \to Y$，投影
$|Z \times_Y X| \to |Z|$ 为闭映射；
\item 存在概形 $V$ 及满射 \'etale 态射 $V \to Y$，使得
$V \times_Y X \to V$ 为代数空间的普遍闭态射；以及
\item 存在 Zariski 覆盖 $Y = \bigcup Y_i$，使得每个态射
$f^{-1}(Y_i) \to Y_i$ 都普遍闭。
\end{enumerate}
\end{lemma}

\begin{proof}
我们略去 (1) 蕴含 (2) 及 (2) 蕴含 (3) 的证明。

\medskip\noindent
假设 (3)。取满射 \'etale 态射 $V \to Y$。我们将证明
$V \times_Y X \to V$ 是代数空间的普遍闭态射。设 $Z \to V$ 为
从代数空间到 $V$ 的态射。取满射 \'etale 态射 $W \to Z$，其中
$W = \coprod W_i$ 是仿射概形的不交并，见《空间的性质》引理
\ref{spaces-properties-lemma-cover-by-union-affines}。于是有如下交换图表
$$
\xymatrix{
\coprod_i |W_i \times_Y X| \ar@{=}[r] \ar[d] &
|W \times_Y X| \ar[r] \ar[d] &
|Z \times_Y X| \ar[d] \ar@{=}[r] &
|Z \times_V (V \times_Y X)| \ar[ld] \\
\coprod |W_i| \ar@{=}[r] &
|W| \ar[r] &
|Z|
}
$$
我们需证明东南箭头为闭映射。中间的水平箭头是满射且开映射
（《空间的性质》引理 \ref{spaces-properties-lemma-etale-open}）。
由假设 (3) 及 $W_i$ 为仿射概形这一事实，左侧竖直箭头为闭映射。
因此右侧竖直箭头也为闭映射。

\medskip\noindent
假设 (4)。我们将证明 $f$ 普遍闭。设 $Z \to Y$ 为代数空间的态射。
考虑图表
$$
\xymatrix{
|(V \times_Y Z) \times_V (V \times_Y X)| \ar@{=}[r] \ar[rd] &
|V \times_Y X| \ar[r] \ar[d] &
|Z \times_Y X| \ar[d] \\
 &
|V \times_Y Z| \ar[r] &
|Z|
}
$$
西南箭头由假设为闭映射。水平箭头是满射且开映射，因为相应的代数空间
态射是 \'etale（见《空间的性质》引理
\ref{spaces-properties-lemma-etale-open}）。因此右侧竖直箭头为闭映射。

\medskip\noindent
显然，取覆盖 $Y = Y$ 即得 (1) 蕴含 (5)。假设 $Y = \bigcup Y_i$
如 (5) 中所述。则对于任意 $Z \to Y$，可得到相应的 Zariski 覆盖
$Z = \bigcup Z_i$，使得 $f$ 到 $Z_i$ 的基变换为闭映射。由简单的
拓扑论证，$Z \times_Y X \to Z$ 为闭映射。因此 (1) 成立。
\end{proof}

\begin{example}
\label{spaces-morphisms-example-strange-universally-closed}
普遍闭态射的奇异例子。
设 $\mathbf{Q} \subset k$ 为特征零的域。令
$X = \mathbf{A}^1_k/\mathbf{Z}$，如《空间》例
\ref{spaces-example-affine-line-translation}。我们断言结构态射
$p : X \to \Spec(k)$ 普遍闭。
事实上，若 $Z/k$ 为概形，且 $T \subset |X \times_k Z|$ 为闭集，
则 $T$ 对应于 $T' \subset |\mathbf{A}^1 \times Z|$ 中的
$\mathbf{Z}$-不变闭集。容易看出，这意味着 $T'$ 是 $Z$ 的某个子集
$T''$ 的逆像。由《态射》引理
\ref{morphisms-lemma-fpqc-quotient-topology}，$T'' \subset Z$ 为闭集。
显然 $T''$ 是 $T$ 的像。因此由引理 \ref{spaces-morphisms-lemma-universally-closed-local}，
$p$ 普遍闭。
\end{example}

\begin{lemma}
\label{spaces-morphisms-lemma-universally-closed-quasi-compact}
设 $S$ 为概形。
$S$ 上代数空间的普遍闭态射是拟紧的。
\end{lemma}

\begin{proof}
该证明重复概形情形的证明，见《态射》引理
\ref{morphisms-lemma-universally-closed-quasi-compact}。
设 $f : X \to Y$ 为 $S$ 上代数空间的态射，并假设 $f$ 不拟紧。
我们的目标是证明 $f$ 不普遍闭。由引理 \ref{spaces-morphisms-lemma-quasi-compact-local}，
存在仿射概形 $Z$ 及态射 $Z \to Y$，使得 $Z \times_Y X \to Z$
不拟紧。为达到目标，只需证明 $Z \times_Y X \to Z$ 不普遍闭，
因此可假设 $Y = \Spec(B)$，其中 $B$ 为某个环。

\medskip\noindent
写成 $X = \bigcup_{i \in I} X_i$，其中 $X_i$ 是 $X$ 的拟紧开子空间。
例如，取满射 \'etale 态射 $U \to X$，其中 $U$ 为概形；取仿射开覆盖
$U = \bigcup U_i$，并令 $X_i \subset X$ 为 $U_i$ 的像。稍后将使用
$X_i \to Y$ 拟紧这一事实，见引理 \ref{spaces-morphisms-lemma-quasi-compact-permanence}。
令 $T = \Spec(B[a_i ; i \in I])$，令 $T_i = D(a_i) \subset T$。
令 $Z \subset T \times_Y X$ 为约化闭子空间，其底层点集为
$|T \times_Y Z| \setminus \bigcup_{i \in I} |T_i \times_Y X_i|$，见
《空间的性质》引理 \ref{spaces-properties-lemma-reduced-closed-subspace}。
（注意，$T_i \times_Y X_i$ 是 $T \times_Y X$ 的开子空间，因为
$T_i \to T$ 与 $X_i \to X$ 是开浸入；见《空间》引理
\ref{spaces-lemma-base-change-immersions} 和
\ref{spaces-lemma-composition-immersions}。）图表如下：
$$
\xymatrix{
Z \ar[r] \ar[rd] &
T \times_Y X \ar[d]^{f_T} \ar[r]_q &
X \ar[d]^f \\
& T \ar[r]^p & Y
}
$$
只需证明像 $f_T(|Z|)$ 在 $|T|$ 中不是闭集。

\medskip\noindent
我们断言，存在点 $y \in Y$，使得 $Y$ 中不存在 $y$ 的仿射开邻域
$V$ 满足 $X_V$ 拟紧。否则，我们可以用有限多个这样的 $V$ 覆盖 $Y$；
对每个 $V$，由引理 \ref{spaces-morphisms-lemma-quasi-compact-permanence}，态射
$Y_V \to V$ 拟紧，继而由引理 \ref{spaces-morphisms-lemma-quasi-compact-local} 得到
$f$ 拟紧，矛盾。固定断言中的一个 $y \in Y$。

\medskip\noindent
令 $t \in T$ 为位于 $y$ 上方且满足 $\kappa(t) = \kappa(y)$ 的点，
并使得对所有 $i$，$a_i = 1$ 于 $\kappa(t)$ 中。若存在 $z \in |Z|$
满足 $f_T(z) = t$，则对某个 $i$ 有 $q(t) \in X_i$。于是按 $Z$ 的构造，
$f_T(z) \not \in T_i$，这与按构造 $t \in T_i$ 矛盾。因此
$t \in |T| \setminus f_T(|Z|)$。

\medskip\noindent
假设 $f_T(|Z|)$ 在 $|T|$ 中闭。则存在元素
$g \in B[a_i; i \in I]$，使 $f_T(|Z|) \subset V(g)$ 但 $t \not \in V(g)$。
于是 $g$ 在 $\kappa(t)$ 中的像非零。特别地，$g$ 的某个系数在
$\kappa(y)$ 中的像非零。因此，该系数在 $y$ 的某个仿射开邻域 $V$ 上可逆。
令 $J$ 为满足 $j \in I$ 且变量 $a_j$ 出现在 $g$ 中的有限集合。
由于 $X_V$ 不拟紧而每个 $X_{i, V}$ 拟紧，我们可以取点
$x \in |X_V| \setminus \bigcup_{j \in J} |X_{j, V}|$。
换言之，$x \in |X| \setminus \bigcup_{j \in J} |X_j|$，且 $x$ 位于某个
$v \in V$ 上方。由于 $g$ 有一个在 $V$ 上可逆的系数，我们可以找到位于
$v$ 上方的点 $t' \in T$，使 $t' \not \in V(g)$ 且对所有 $i \notin J$
都有 $t' \in V(a_i)$。这是因为
$V(a_i; i \in I \setminus J) = \Spec(B[a_j; j\in J])$，
且该概形中位于 $v$ 上方的点集与
$\Spec(\kappa(v)[a_j; j \in J])$ 双射，而按构造 $g$ 限制为该多项式环
中的非零元素。换言之，对于每个 $i \not \in J$，$t' \not \in T_i$。
由《空间的性质》引理 \ref{spaces-properties-lemma-points-cartesian}，
可找到 $X \times_Y T$ 中的点 $z$，其映射到 $x \in X$ 以及 $t' \in T$。
由于当 $j \in J$ 时 $x \not \in |X_j|$，且当
$i \in I \setminus J$ 时 $t' \not \in T_i$，可知 $z \in |Z|$。
另一方面，$f_T(z) = t' \not \in V(g)$，这与 $f_T(Z) \subset V(g)$ 矛盾。
因此假设“$f_T(|Z|)$ 闭”错误，故如所需，$f_T$ 非闭。
\end{proof}

\noindent
分离代数空间在满射普遍闭态射下的目标仍然分离。

\begin{lemma}
\label{spaces-morphisms-lemma-image-universally-closed-separated}
设 $S$ 为概形，设 $B$ 为 $S$ 上的代数空间。
设 $f : X \to Y$ 为 $B$ 上代数空间的满射普遍闭态射。
\begin{enumerate}
\item 若 $X$ 拟分离，则 $Y$ 拟分离。
\item 若 $X$ 分离，则 $Y$ 分离。
\item 若 $X$ 在 $B$ 上拟分离，则 $Y$ 在 $B$ 上拟分离。
\item 若 $X$ 在 $B$ 上分离，则 $Y$ 在 $B$ 上分离。
\end{enumerate}
\end{lemma}

\begin{proof}
第 (1)、(2) 部分是 $S = B = \Spec(\mathbf{Z})$ 时第 (3)、(4) 部分的
结果（见《空间的性质》定义 \ref{spaces-properties-definition-separated}）。
考虑交换图表
$$
\xymatrix{
X \ar[d] \ar[rr]_{\Delta_{X/B}} & & X \times_B X \ar[d] \\
Y \ar[rr]^{\Delta_{Y/B}} & & Y \times_B Y
}
$$
左侧竖直箭头为满射（即普遍满射）。右侧竖直箭头是普遍闭的，
因为它是普遍闭态射
$X \times_B X \to X \times_B Y \to Y \times_B Y$ 的复合。因此它也拟紧，
见引理 \ref{spaces-morphisms-lemma-universally-closed-quasi-compact}。

\medskip\noindent
假设 $X$ 在 $B$ 上拟分离，即 $\Delta_{X/B}$ 拟紧。若 $Z$ 拟紧且
$Z \to Y \times_B Y$ 为态射，则
$Z \times_{Y \times_B Y} X \to Z \times_{Y \times_B Y} Y$
满射，并且由上述说明 $Z \times_{Y \times_B Y} X$ 拟紧。
因此 $\Delta_{Y/B}$ 拟紧，即 $Y$ 在 $B$ 上拟分离。

\medskip\noindent
假设 $X$ 在 $B$ 上分离，即 $\Delta_{X/B}$ 为闭浸入。若 $Z$ 为仿射概形，
且 $Z \to Y \times_B Y$ 为态射，则
$Z \times_{Y \times_B Y} X \to Z \times_{Y \times_B Y} Y$
满射，并且由上述说明 $Z \times_{Y \times_B Y} X \to Z$ 普遍闭。
因此 $\Delta_{Y/B}$ 普遍闭。由此 $\Delta_{Y/B}$ 可表示、局有限型且为
单态射（见引理 \ref{spaces-morphisms-lemma-properties-diagonal}），并且普遍闭，
所以是闭浸入；见《\'Etale 态射》引理
\ref{etale-lemma-characterize-closed-immersion}（也见抽象原理
《空间》引理 \ref{spaces-lemma-representable-transformations-property-implication}）。
于是 $Y$ 在 $B$ 上分离。
\end{proof}













\section{单态射}
\label{spaces-morphisms-section-monomorphisms}

\noindent
按照第 \ref{spaces-morphisms-section-representable} 节，代数空间的可表示态射
$X \to Y$ 是单态射，是指对于每个概形 $Z$ 及态射 $Z \to Y$，
态射 $Z \times_Y X \to Z$ 可由概形的单态射表示。
这恰好意味着 $Z \times_Y X \to Z$ 是 $(\Sch/S)_{fppf}$ 上层的
单射。由于这应对所有 $Z$ 及所有 $Z \to Y$ 成立，因而等价于
$X \to Y$ 是 $(\Sch/S)_{fppf}$ 上层的单射。因此，对（可能不可表示的）
代数空间态射定义单态射如下：\footnote{我们不知道代数空间的任意单态射
是否可表示。讨论见《空间态射的更多内容》第
\ref{spaces-more-morphisms-section-monomorphisms} 节。}

\begin{definition}
\label{spaces-morphisms-definition-monomorphism}
设 $S$ 为概形。
若 $S$ 上的代数空间态射是层的单射，即 $(\Sch/S)_{fppf}$ 上层范畴中的
单态射，则称其为 {\it 单态射}。
\end{definition}

\noindent
下述引理表明，这也等价于它是 $S$ 上代数空间范畴中的单态射。

\begin{lemma}
\label{spaces-morphisms-lemma-monomorphism}
设 $S$ 为概形。
设 $j : X \to Y$ 为 $S$ 上代数空间的态射。
以下命题等价：
\begin{enumerate}
\item $j$ 是单态射（如定义 \ref{spaces-morphisms-definition-monomorphism}）；
\item $j$ 是 $S$ 上代数空间范畴中的单态射；
\item 对角态射 $\Delta_{X/Y} : X \to X \times_Y X$ 是同构。
\end{enumerate}
\end{lemma}

\begin{proof}
注意，$X \times_Y X$ 既是 $(\Sch/S)_{fppf}$ 上层范畴中的纤维积，
也是 $S$ 上代数空间范畴中的纤维积，见《空间》引理
\ref{spaces-lemma-fibre-product-spaces}。(1) 与 (3) 的等价性是任意
网站上层单射的一般刻画。(2) 与 (3) 的等价性是具有纤维积的任意范畴中
对单态射的刻画。
\end{proof}

\begin{lemma}
\label{spaces-morphisms-lemma-monomorphism-separated}
代数空间的单态射是分离的。
\end{lemma}

\begin{proof}
这是因为同构是闭浸入，且可应用上面的引理
\ref{spaces-morphisms-lemma-monomorphism}。
\end{proof}

\begin{lemma}
\label{spaces-morphisms-lemma-composition-monomorphism}
单态射的复合是单态射。
\end{lemma}

\begin{proof}
这是因为层单射的复合仍为单射。
\end{proof}

\begin{lemma}
\label{spaces-morphisms-lemma-base-change-monomorphism}
单态射的基变换是单态射。
\end{lemma}

\begin{proof}
这是层范畴中纤维积的一般事实。
\end{proof}

\begin{lemma}
\label{spaces-morphisms-lemma-monomorphism-local}
设 $S$ 为概形。
设 $f : X \to Y$ 为 $S$ 上代数空间的态射。
以下命题等价：
\begin{enumerate}
\item $f$ 是单态射；
\item 对每个概形 $Z$ 及态射 $Z \to Y$，$f$ 的基变换
$Z \times_Y X \to Z$ 是单态射；
\item 对每个仿射概形 $Z$ 及每个态射 $Z \to Y$，$f$ 的基变换
$Z \times_Y X \to Z$ 是单态射；
\item 存在概形 $V$ 及满射 \'etale 态射 $V \to Y$，使得基变换
$V \times_Y X \to V$ 是单态射；以及
\item 存在 Zariski 覆盖 $Y = \bigcup Y_i$，使得每个态射
$f^{-1}(Y_i) \to Y_i$ 是单态射。
\end{enumerate}
\end{lemma}

\begin{proof}
下文将不再另行说明单态射的基变换仍是单态射这一事实，见引理
\ref{spaces-morphisms-lemma-base-change-monomorphism}。特别地，显然
(1) $\Rightarrow$ (2) $\Rightarrow$ (3) $\Rightarrow$ (4)，其中 (4) 中
$V$ 取为与 $Y$ \'etale 的仿射概形的不交并（见《空间的性质》引理
\ref{spaces-properties-lemma-cover-by-union-affines}）。设 $V$ 为概形，
且 $V \to Y$ 为满射 \'etale 态射。若 $V \times_Y X \to V$ 是单态射，
则 $X \to Y$ 也是单态射。事实上，给定任意层的笛卡尔图表
$$
\vcenter{
\xymatrix{
\mathcal{F} \ar[r]_a \ar[d]_b & \mathcal{G} \ar[d]^c \\
\mathcal{H} \ar[r]^d & \mathcal{I}
}
}
\quad
\quad
\mathcal{F} = \mathcal{H} \times_\mathcal{I} \mathcal{G}
$$
若 $c$ 是层的满射且 $a$ 为单射，则 $d$ 也是单射。因此 (4) 蕴含 (1)。
(5) 与 (1) 的等价性证明略。
\end{proof}

\begin{lemma}
\label{spaces-morphisms-lemma-immersions-monomorphisms}
代数空间的浸入是单态射。
特别地，任意浸入都是分离的。
\end{lemma}

\begin{proof}
设 $f : X \to Y$ 为代数空间的浸入。对于任意态射 $Z \to Y$，其中 $Z$
可表示，基变换 $Z \times_Y X \to Z$ 是概形的浸入，因而为单态射，
见《概形》引理 \ref{schemes-lemma-immersions-monomorphisms}。
因此 $f$ 可表示且为单态射。
\end{proof}

\noindent
我们将在《良好空间》引理
\ref{decent-spaces-lemma-monomorphism-toward-disjoint-union-dim-0-rings}
中加强下述引理。

\begin{lemma}
\label{spaces-morphisms-lemma-monomorphism-toward-field}
设 $S$ 为概形。设 $k$ 为域，且 $Z \to \Spec(k)$ 为 $S$ 上代数空间的
单态射。则或者 $Z = \emptyset$，或者 $Z = \Spec(k)$。
\end{lemma}

\begin{proof}
由引理 \ref{spaces-morphisms-lemma-monomorphism-separated} 和
\ref{spaces-morphisms-lemma-separated-over-separated}，可知 $Z$ 是分离代数空间。因此存在
概形形式的开稠密子空间 $Z' \subset Z$，见《空间的性质》命题
\ref{spaces-properties-proposition-locally-quasi-separated-open-dense-scheme}。
由《概形》引理 \ref{schemes-lemma-mono-towards-spec-field}，或者
$Z' = \emptyset$，或者 $Z' \cong \Spec(k)$。在第一种情形，得到
$Z = \emptyset$；在第二种情形，由于 $Z \to \Spec(k)$ 是在 $Z'$ 上为
同构的单态射，得到 $Z' = Z = \Spec(k)$。
\end{proof}

\begin{lemma}
\label{spaces-morphisms-lemma-monomorphism-injective-points}
设 $S$ 为概形。若 $X \to Y$ 为 $S$ 上代数空间的单态射，则
$|X| \to |Y|$ 为单射。
\end{lemma}

\begin{proof}
直接由定义得到。
\end{proof}















\section{拟凝聚层的推前}
\label{spaces-morphisms-section-pushforward}

\noindent
我们先证明一个简单引理，将代数空间态射的模层推前与概形态射的模层
推前联系起来。

\begin{lemma}
\label{spaces-morphisms-lemma-compute-pushforward}
设 $S$ 为概形。
设 $f : X \to Y$ 为 $S$ 上代数空间的态射。
设 $U \to X$ 为从概形到 $X$ 的满射 \'etale 态射。
令 $R = U \times_X U$，并按通常记号令 $t, s : R \to U$ 为投影态射。
记 $a : U \to Y$ 与 $b : R \to Y$ 为诱导态射。对于
$\textit{Mod}(\mathcal{O}_X)$ 中任意对象 $\mathcal{F}$，存在正合序列
$$
0 \to f_*\mathcal{F} \to a_*(\mathcal{F}|_U) \to b_*(\mathcal{F}|_R)
$$
其中第二个箭头为差 $t^* - s^*$。
\end{lemma}

\begin{proof}
我们也把 $\mathcal{F}$ 记作其延拓到 $X_{spaces, \etale}$ 上的模层，见
《空间的性质》注记 \ref{spaces-properties-remark-explain-equivalence}。
设 $V \to Y$ 为 $Y_\etale$ 中的对象。则 $V \times_Y X$ 是
$X_{spaces, \etale}$ 中的对象，且按定义
$f_*\mathcal{F}(V) = \mathcal{F}(V \times_Y X)$。由于 $U \to X$
满射 \'etale，故 $\{V \times_Y U \to V \times_Y X\}$ 是覆盖。
此外有
$(V \times_Y U) \times_X (V \times_Y U) = V \times_Y R$.
因此由 $\mathcal{F}$ 在 $X_{spaces, \etale}$ 上的层条件，有短正合序列
$$
0 \to \mathcal{F}(V \times_Y X)
\to \mathcal{F}(V \times_Y U) \to \mathcal{F}(V \times_Y R)
$$
其中第二个箭头是沿 $t$ 或 $s$ 限制所得的差。该正合序列关于 $V$
函子性，故得到本引理。
\end{proof}

\noindent
设 $S$ 为概形。设 $f : X \to Y$ 为 $S$ 上可表示代数空间 $X$ 与 $Y$
之间的拟紧且拟分离态射。由《下降》命题
\ref{descent-proposition-equivalence-quasi-coherent-functorial}，函子
$f_* : \QCoh(\mathcal{O}_X) \to \QCoh(\mathcal{O}_Y)$
在将 $X$ 与 $Y$ 看作概形时与通常函子一致。

\medskip\noindent
更一般地，设 $f : X \to Y$ 为 $S$ 上代数空间的可表示、拟紧且拟分离态射。
设 $V$ 为概形，且 $V \to Y$ 为满射 \'etale 态射。令
$U = V \times_Y X$，令 $f' : U \to V$ 为 $f$ 的基变换。则对于任意
拟凝聚 $\mathcal{O}_X$-模 $\mathcal{F}$，有
\begin{equation}
\label{spaces-morphisms-equation-representable-pushforward}
f'_*(\mathcal{F}|_U) = (f_*\mathcal{F})|_V,
\end{equation}
见《空间的性质》引理
\ref{spaces-properties-lemma-pushforward-etale-base-change-modules}。
又由于 $f' : U \to V$ 是概形的拟紧且拟分离态射，由上一段的注记，
我们可以将 $\mathcal{F}|_U$ 看作概形 $U$ 上的拟凝聚层、将 $f'$ 看作
概形态射来计算 $f'_*(\mathcal{F}|_U)$。下文经常使用这一点，不再说明。

\medskip\noindent
下一层一般性是考虑任意拟紧且拟分离的代数空间态射。

\begin{lemma}
\label{spaces-morphisms-lemma-pushforward}
设 $S$ 为概形。
设 $f : X \to Y$ 为 $S$ 上代数空间的态射。
若 $f$ 拟紧且拟分离，则 $f_*$ 将拟凝聚 $\mathcal{O}_X$-模变为
拟凝聚 $\mathcal{O}_Y$-模。
\end{lemma}

\begin{proof}
设 $\mathcal{F}$ 为 $X$ 上的拟凝聚层。需证明 $f_*\mathcal{F}$ 是 $Y$ 上的
拟凝聚层。只需证明对于任意仿射概形 $V$ 及 \'etale 态射 $V \to Y$，
$f_*\mathcal{F}$ 到 $V$ 的限制是拟凝聚的；见《空间的性质》引理
\ref{spaces-properties-lemma-characterize-quasi-coherent}。
令 $f' : V \times_Y X \to V$ 为 $f$ 沿 $V \to Y$ 的基变换。注意由引理
\ref{spaces-morphisms-lemma-base-change-quasi-compact} 和
\ref{spaces-morphisms-lemma-base-change-separated}，$f'$ 也拟紧且拟分离。
由 (\ref{spaces-morphisms-equation-representable-pushforward})，$f_*\mathcal{F}$ 到 $V$ 的
限制是 $\mathcal{F}$ 到 $V \times_Y X$ 的限制的 $f'_*$。因此可将 $f$
替换为 $f'$，并假设 $Y$ 是仿射概形。

\medskip\noindent
假设 $Y$ 为仿射概形。由于 $f$ 拟紧，$X$ 拟紧。因此可取仿射概形 $U$
及满射 \'etale 态射 $U \to X$，见《空间的性质》引理
\ref{spaces-properties-lemma-quasi-compact-affine-cover}。
由引理 \ref{spaces-morphisms-lemma-compute-pushforward}，有正合序列
$$
0 \to f_*\mathcal{F} \to a_*(\mathcal{F}|_U) \to b_*(\mathcal{F}|_R).
$$
其中 $R = U \times_X U$。由于 $X \to Y$ 拟分离，
$R \to U \times_Y U$ 是拟紧单态射。这意味着 $R$ 是拟紧分离概形
（此时 $U$ 与 $Y$ 均为仿射）。因此 $a : U \to Y$ 与 $b : R \to Y$
是概形的拟紧且拟分离态射。于是由《下降》命题
\ref{descent-proposition-equivalence-quasi-coherent-functorial}，
层 $a_*(\mathcal{F}|_U)$ 与 $b_*(\mathcal{F}|_R)$ 拟凝聚（也见本引理前的
讨论）。这说明 $f_*\mathcal{F}$ 是拟凝聚模的核，因而自身拟凝聚；见
《空间的性质》引理 \ref{spaces-properties-lemma-properties-quasi-coherent}。
\end{proof}

\noindent
高阶推前见《空间的上同调》第
\ref{spaces-cohomology-section-higher-direct-image} 节。









\section{浸入}
\label{spaces-morphisms-section-immersions}

\noindent
代数空间的开浸入、闭浸入及局部闭浸入已在《空间》第
\ref{spaces-section-Zariski} 节定义。也就是说，若代数空间态射可表示，
且按照第 \ref{spaces-morphisms-section-representable} 节的意义为闭浸入（分别为开浸入、浸入），
则称其为 {\it 闭浸入}（分别为 {\it 开浸入}、{\it 浸入}）。

\medskip\noindent
特别地，这些类型的态射在 $S$ 上代数空间范畴中的基变换与态射复合下保持；
见《空间》引理 \ref{spaces-lemma-composition-immersions} 和
\ref{spaces-lemma-base-change-immersions}。

\begin{lemma}
\label{spaces-morphisms-lemma-closed-immersion-local}
设 $S$ 为概形。设 $f : X \to Y$ 为 $S$ 上代数空间的态射。
以下命题等价：
\begin{enumerate}
\item $f$ 是闭浸入（分别为开浸入、浸入）；
\item 对每个概形 $Z$ 及任意态射 $Z \to Y$，态射
$Z \times_Y X \to Z$ 是闭浸入（分别为开浸入、浸入）；
\item 对每个仿射概形 $Z$ 及任意态射 $Z \to Y$，态射
$Z \times_Y X \to Z$ 是闭浸入（分别为开浸入、浸入）；
\item 存在概形 $V$ 及满射 \'etale 态射 $V \to Y$，使得
$V \times_Y X \to V$ 是闭浸入（分别为开浸入、浸入）；以及
\item 存在 Zariski 覆盖 $Y = \bigcup Y_i$，使得每个态射
$f^{-1}(Y_i) \to Y_i$ 是闭浸入（分别为开浸入、浸入）。
\end{enumerate}
\end{lemma}

\begin{proof}
由于闭浸入（分别为开浸入、浸入）的基变换仍为同类态射，显然
(1) 蕴含 (2)，(2) 蕴含 (3)。又 (3) 蕴含 (4)，因为可取 $V$ 为仿射概形
的不交并；见《空间的性质》引理
\ref{spaces-properties-lemma-cover-by-union-affines}。

\medskip\noindent
假设 $V \to Y$ 如 (4) 中所述。令 $\mathcal{P}$ 为概形态射的性质
“闭浸入（分别为开浸入、浸入）”。注意，性质 $\mathcal{P}$ 在任意基变换下
保持，且在基上是 fppf 局部的（见第 \ref{spaces-morphisms-section-representable} 节）。
此外，$\mathcal{P}$ 类型的态射是分离且局拟有限的（在三种情形中均如此；见
《概形》引理 \ref{schemes-lemma-immersions-monomorphisms} 以及
《态射》引理 \ref{morphisms-lemma-immersion-locally-quasi-finite}）。
因此由《态射的更多内容》引理
\ref{more-morphisms-lemma-separated-locally-quasi-finite-morphisms-fppf-descend}，
$\mathcal{P}$ 类型态射满足 fppf 覆盖下的下降。于是可应用《空间》引理
\ref{spaces-lemma-morphism-sheaves-with-P-effective-descent-etale}，得到
$X \to Y$ 可表示且具有性质 $\mathcal{P}$，即 (1) 成立。

\medskip\noindent
由 $\mathcal{P}$ 在目标上是 Zariski 局部性质，(1) 与 (5) 等价（上面已知
$\mathcal{P}$ 实际上在目标上是 fppf 局部的）。
\end{proof}

\begin{lemma}
\label{spaces-morphisms-lemma-immersion-permanence}
设 $S$ 为概形。
设 $Z \to Y \to X$ 为 $S$ 上代数空间的态射。
\begin{enumerate}
\item 若 $Z \to X$ 可表示、局有限型、局拟有限、分离且为单态射，
则 $Z \to Y$ 也可表示、局有限型、局拟有限、分离且为单态射。
\item 若 $Z \to X$ 为浸入且 $Y \to X$ 局部分离，则 $Z \to Y$ 为浸入。
\item 若 $Z \to X$ 为闭浸入且 $Y \to X$ 分离，则 $Z \to Y$ 为闭浸入。
\end{enumerate}
\end{lemma}

\begin{proof}
每种情形都考虑交换图表
$$
\xymatrix{
Z \ar[r] \ar[rd] & Y \times_X Z \ar[r] \ar[d] & Z \ar[d] \\
& Y \ar[r] & X
}
$$
其中上方水平箭头的复合为恒等。证明 (1)。第一个水平箭头是
$Y \times_X Z \to Z$ 的截面，故由引理 \ref{spaces-morphisms-lemma-section-immersion}，
它可表示、局有限型、局拟有限、分离且为单态射。
态射 $Y \times_X Z \to Y$ 是 $Z \to X$ 的基变换，因而也可表示、
局有限型、局拟有限、分离且为单态射（这些概形态射性质在基变换下保持；见
《空间》注记 \ref{spaces-remark-list-properties-stable-base-change}）。
因此复合也具有同样性质（这些概形态射性质在复合下保持；见《空间》注记
\ref{spaces-remark-list-properties-stable-composition}）。这证明 (1)。
其余结论完全同理。
\end{proof}

\begin{lemma}
\label{spaces-morphisms-lemma-immersion-when-closed}
设 $S$ 为概形。设 $i : Z \to X$ 为 $S$ 上代数空间的浸入。则
$|i| : |Z| \to |X|$ 是到局部闭子集上的同胚，且 $i$ 为闭浸入，当且仅当
像 $|i|(|Z|) \subset |X|$ 是闭集。
\end{lemma}

\begin{proof}
第一条断言是《空间的性质》引理
\ref{spaces-properties-lemma-subspace-induced-topology}。设 $U$ 为概形，
且 $U \to X$ 为满射 \'etale 态射。由假设，$T = U \times_X Z$ 为概形，
且 $j : T \to U$ 为概形的浸入。由引理 \ref{spaces-morphisms-lemma-closed-immersion-local}，
$i$ 为闭浸入，当且仅当 $j$ 为闭浸入。由《概形》引理
\ref{schemes-lemma-immersion-when-closed}，这又当且仅当 $j(T)$ 在 $U$
中为闭集。然而子集 $j(T) \subset U$ 是 $|i|(|Z|) \subset |X|$ 的逆像，
见《空间的性质》引理 \ref{spaces-properties-lemma-points-cartesian}。
证明完毕。
\end{proof}

\begin{remark}
\label{spaces-morphisms-remark-immersion}
设 $S$ 为概形。设 $i : Z \to X$ 为 $S$ 上代数空间的浸入。由于 $i$ 是
单态射，可以将 $|Z|$ 看作 $|X|$ 的子集；本注记中均如此处理。
令 $\partial |Z|$ 为 $|Z|$ 在拓扑空间 $|X|$ 中的边界。公式为
$$
\partial |Z| = \overline{|Z|} \setminus |Z|.
$$
令 $\partial Z$ 为 $X$ 的约化闭子空间，满足
$|\partial Z| = \partial |Z|$，其由约化诱导闭子空间结构得到；见《空间的性质》
定义 \ref{spaces-properties-definition-reduced-induced-space}。由构造，
$|Z|$ 在 $|X| \setminus |\partial Z| = |X \setminus \partial Z|$ 中为闭集。
因此任何代数空间浸入都可分解为先闭浸入后开浸入（但一般不能反向分解；见
《态射》例 \ref{morphisms-example-thibaut}）。
\end{remark}

\begin{remark}
\label{spaces-morphisms-remark-space-structure-locally-closed-subset}
设 $S$ 为概形，设 $X$ 为 $S$ 上的代数空间。设 $T \subset |X|$ 为局部闭子集。
令 $\partial T$ 为 $T$ 在拓扑空间 $|X|$ 中的边界。公式为
$$
\partial T = \overline{T} \setminus T.
$$
令 $U \subset X$ 为 $X$ 的开子空间，满足 $|U| = |X| \setminus \partial T$，见《空间的性质》
引理 \ref{spaces-properties-lemma-open-subspaces}。令 $Z$ 为 $U$ 的约化闭子空间，
满足 $|Z| = T$，由约化诱导闭子空间结构得到，见《空间的性质》定义
\ref{spaces-properties-definition-reduced-induced-space}。由构造，
$Z \to U$ 是代数空间的闭浸入，而 $U \to X$ 是开浸入，故
$Z \to X$ 是 $S$ 上代数空间的浸入（见《空间》引理
\ref{spaces-lemma-composition-immersions}）。注意 $Z$ 是约化代数空间，且作为
$|X|$ 的子集有 $|Z| = T$。有时称 $Z$ 为 $T$ 上的 {\it 约化诱导子空间结构}。
\end{remark}

\begin{lemma}
\label{spaces-morphisms-lemma-factor-the-other-way}
设 $S$ 为概形。设 $Z \to X$ 为 $S$ 上代数空间的浸入，并假设 $Z \to X$
拟紧。存在分解 $Z \to \overline{Z} \to X$，其中
$Z \to \overline{Z}$ 为开浸入，$\overline{Z} \to X$ 为闭浸入。
\end{lemma}

\begin{proof}
设 $U$ 为概形，且 $U \to X$ 为满射 \'etale 态射。按通常记号，令
$R = U \times_X U$，投影为 $s, t : R \to U$。令 $T = Z \times_U X$。
令 $\overline{T} \subset U$ 为 $T \to U$ 的概形论像。注意，由于拟紧态射的
概形论像与平坦基变换交换，
$s^{-1}\overline{T} = t^{-1}\overline{T}$；见《态射》引理
\ref{morphisms-lemma-flat-base-change-scheme-theoretic-image}。
因此得到闭子空间 $\overline{Z} \subset X$，其到 $U$ 的拉回为 $\overline{T}$；
见《空间的性质》引理 \ref{spaces-properties-lemma-subspaces-presentation}。
由《态射》引理 \ref{morphisms-lemma-quasi-compact-immersion}，
$T \to \overline{T}$ 为开浸入。因此 $Z \to \overline{Z}$ 为开浸入，结论成立。
\end{proof}


















\section{闭浸入}
\label{spaces-morphisms-section-closed-immersions}

\noindent
本节阐明前面关于代数空间浸入所得的一些结果。见《空间》第
\ref{spaces-section-Zariski} 节及本章第 \ref{spaces-morphisms-section-immersions} 节。
本节是代数空间版本的《态射》第
\ref{morphisms-section-closed-immersions} 节。

\begin{lemma}
\label{spaces-morphisms-lemma-closed-immersion-ideals}
设 $S$ 为概形。
设 $X$ 为 $S$ 上的代数空间。对于每个闭浸入 $i : Z \to X$，层
$i_*\mathcal{O}_Z$ 是拟凝聚 $\mathcal{O}_X$-模，映射
$i^\sharp : \mathcal{O}_X \to i_*\mathcal{O}_Z$ 为满射，且其核是拟凝聚理想层。
对应规则
$Z \mapsto \Ker(\mathcal{O}_X \to i_*\mathcal{O}_Z)$
定义一个反包含双射
$$
\begin{matrix}
\text{闭子空间}\\
Z \subset X
\end{matrix}
\longrightarrow
\begin{matrix}
\text{拟凝聚层}\\
\text{理想 }\mathcal{I} \subset \mathcal{O}_X
\end{matrix}
$$
此外，给定对应于拟凝聚理想层 $\mathcal{I} \subset \mathcal{O}_X$ 的闭子概形
$Z$，代数空间态射 $h : Y \to X$ 可经由 $Z$ 分解，当且仅当映射
$h^*\mathcal{I} \to h^*\mathcal{O}_X = \mathcal{O}_Y$ 为零。
\end{lemma}

\begin{proof}
设 $U \to X$ 为满射 \'etale 态射，且源为概形。考虑图表
$$
\xymatrix{
U \times_X Z \ar[r] \ar[d]_{i'} & Z \ar[d]^i \\
U \ar[r] & X
}
$$
由《态射》引理 \ref{spaces-morphisms-lemma-closed-immersion-local}
可知 $i$ 为闭浸入
当且仅当 $i'$ 为闭浸入。由
《空间的性质》
引理 \ref{spaces-properties-lemma-pushforward-etale-base-change-modules}
表明 $i'_*\mathcal{O}_{U \times_X Z}$ 是 $i_*\mathcal{O}_Z$ 到 $U$ 的限制。
因此关于 $\mathcal{O}_X \to i_*\mathcal{O}_Z$ 的断言，等价于关于
$\mathcal{O}_U \to i'_*\mathcal{O}_{U \times_X Z}$ 的相应断言。
又由于 $i'$ 是概形的闭浸入，这些结果由《态射》引理
\ref{morphisms-lemma-closed-immersion} 得到。

\medskip\noindent
 现在证明：给定拟凝聚理想层 $\mathcal{I} \subset \mathcal{O}_X$，公式
$$
Z(T) = \{h : T \to X \mid h^*\mathcal{I} \to \mathcal{O}_T
\text{ 为零}\}
$$
定义 $X$ 的一个闭子空间。显然它是 $X$ 的子函子。
为证明 $Z \to X$ 可由闭浸入表示，设
$\varphi : U \to X$ 是从概形到 $X$ 的态射。那么
$Z \times_X U$ 由 $U$ 的相应子函子表示，该子函子对应于理想层
$\Im(\varphi^*\mathcal{I} \to \mathcal{O}_U)$。由《空间的性质》
引理 \ref{spaces-properties-lemma-pullback-quasi-coherent}，
$\mathcal{O}_U$-模 $\varphi^*\mathcal{I}$ 在 $U$ 上是拟凝聚的，
从而 $\Im(\varphi^*\mathcal{I} \to \mathcal{O}_U)$
是 $U$ 上的拟凝聚理想层。由《概形》引理
\ref{schemes-lemma-characterize-closed-subspace}，可知
$Z \times_X U$ 由与 $\Im(\varphi^*\mathcal{I} \to \mathcal{O}_U)$
关联的 $U$ 的闭子概形表示。因此 $Z$ 是 $X$ 的闭子空间。

\medskip\noindent
在上述 $Z$ 的公式中，输入 $T$ 是概形，这是因为代数空间是
$(\Sch/S)_{fppf}$ 上的层。我们省略验证当 $T$ 为代数空间时
同一公式仍然成立的过程。
\end{proof}

\begin{definition}
\label{spaces-morphisms-definition-inverse-image-closed-subspace}
设 $S$ 为概形。设 $f : Y \to X$ 是 $S$ 上代数空间的态射。
设 $Z \subset X$ 为闭子空间。闭子空间 $Z$ 的
{\it 逆像 $f^{-1}(Z)$} 定义为 $Y$ 的闭子空间 $Z \times_X Y$。
\end{definition}

\noindent
由引理 \ref{spaces-morphisms-lemma-closed-immersion-local}，此定义是有意义的。
若 $\mathcal{I} \subset \mathcal{O}_X$ 是通过引理
\ref{spaces-morphisms-lemma-closed-immersion-ideals} 与 $Z$ 对应的拟凝聚理想层，则
$f^{-1}\mathcal{I}\mathcal{O}_Y = \Im(f^*\mathcal{I} \to \mathcal{O}_Y)$
是与 $f^{-1}(Z)$ 对应的理想层。

\begin{lemma}
\label{spaces-morphisms-lemma-closed-immersion-quasi-compact}
代数空间的闭浸入是拟紧的。
\end{lemma}

\begin{proof}
这由一般原理及《概形》引理
\ref{schemes-lemma-closed-immersion-quasi-compact} 得到；见《空间》引理
\ref{spaces-lemma-representable-transformations-property-implication}。
\end{proof}

\begin{lemma}
\label{spaces-morphisms-lemma-closed-immersion-separated}
代数空间的闭浸入是分离的。
\end{lemma}

\begin{proof}
这由一般原理及《概形》引理
\ref{schemes-lemma-immersions-monomorphisms} 得到；见《空间》引理
\ref{spaces-lemma-representable-transformations-property-implication}。
\end{proof}

\begin{lemma}
\label{spaces-morphisms-lemma-closed-immersion-push-pull}
设 $S$ 为概形。设 $i : Z \to X$ 是 $S$ 上代数空间的闭浸入。
\begin{enumerate}
\item 函子
$$
i_{small, *} :
\Sh(Z_\etale)
\longrightarrow
\Sh(X_\etale)
$$
是全忠实的，其本质像由 $X_\etale$ 上满足下述条件的集合层
$\mathcal{F}$ 组成：其到 $X \setminus Z$ 的限制同构于 $*$；以及
\item 函子
$$
i_{small, *} :
\textit{Ab}(Z_\etale)
\longrightarrow
\textit{Ab}(X_\etale)
$$
是全忠实的，其本质像由 $X_\etale$ 上支集包含于 $|Z|$ 的阿贝尔层组成。
\end{enumerate}
在两种情形中，$i_{small}^{-1}$ 都是函子 $i_{small, *}$ 的左逆。
\end{lemma}

\begin{proof}
设 $U$ 为概形，且 $U \to X$ 为满射 \'etale 态射。令
$V = Z \times_X U$。则 $V$ 为概形，且 $i' : V \to U$ 是概形的闭浸入。由
《空间的性质》引理 \ref{spaces-properties-lemma-pushforward-etale-base-change}，
对 $Z$ 上的任意层 $\mathcal{G}$，有
$$
(i_{small}^{-1}i_{small, *}\mathcal{G})|_V =
(i')_{small}^{-1}i'_{small, *}(\mathcal{G}|_V)
$$
由《空间的上同调》命题
\ref{etale-cohomology-proposition-closed-immersion-pushforward}，映射
$(i')_{small}^{-1}i'_{small, *}(\mathcal{G}|_V) \to \mathcal{G}|_V$
是同构。由于 $V \to Z$ 是满射 \'etale 态射，这意味着
$i_{small}^{-1}i_{small, *}\mathcal{G} \to \mathcal{G}$ 是同构。
这显然推出 $i_{small, *}$ 是全忠实的；见《网站》引理
\ref{sites-lemma-exactness-properties}。
为证明关于本质像的断言，考虑 $X_\etale$ 上的集合层 $\mathcal{F}$，
其到 $X \setminus Z$ 的限制同构于 $*$。如《空间的上同调》命题
\ref{etale-cohomology-proposition-closed-immersion-pushforward} 的证明中那样，
我们考虑伴随映射
$$
\mathcal{F} \longrightarrow i_{small, *}i_{small}^{-1}\mathcal{F}.
$$
如第一部分，我们由概形情形的相应结果看到该映射到
$U$ 的限制是同构。由于 $U$ 是 $X$ 的 \'etale 覆盖，故它是同构。
\end{proof}

\begin{lemma}
\label{spaces-morphisms-lemma-stalk-push-closed}
设 $S$ 为概形。设 $i : Z \to X$ 是 $S$ 上代数空间的闭浸入。
设 $\overline{z}$ 为 $Z$ 的几何点，其在 $X$ 中的像为 $\overline{x}$。则对
$Z_\etale$ 上的任意层 $\mathcal{F}$，有
$(i_{small, *}\mathcal{F})_{\overline{z}} = \mathcal{F}_{\overline{x}}$。
\end{lemma}

\begin{proof}
取 $\overline{x}$ 的一个 \'etale 邻域 $(U, \overline{u})$。
那么茎 $(i_{small, *}\mathcal{F})_{\overline{z}}$ 是
$i_{small, *}\mathcal{F}|_U$ 在 $\overline{u}$ 处的茎。由《空间的性质》引理
\ref{spaces-properties-lemma-pushforward-etale-base-change}，可将 $X$ 替换为 $U$，
并将 $Z$ 替换为 $Z \times_X U$。此时 $Z \to X$ 是概形的闭浸入，结论就是
《空间的上同调》引理
\ref{etale-cohomology-lemma-stalk-pushforward-closed-immersion}。
\end{proof}

\noindent
下述引理在拓扑斯的闭浸入情形下更一般地成立（未来补充引用）。

\begin{lemma}
\label{spaces-morphisms-lemma-closed-immersion-rings}
设 $S$ 为概形。设 $i : Z \to X$ 是 $S$ 上代数空间的闭浸入。设
$\mathcal{A}$ 为 $X_\etale$ 上的环层，$\mathcal{B}$ 为 $Z_\etale$ 上的环层。
设 $\varphi : \mathcal{A} \to i_{small, *}\mathcal{B}$ 为环层同态，
从而得到环化拓扑斯的态射
$$
f : (\Sh(Z_\etale), \mathcal{B}) \longrightarrow (\Sh(X_\etale), \mathcal{A}).
$$
对于 $\mathcal{A}$-模层 $\mathcal{F}$ 和 $\mathcal{B}$-模层 $\mathcal{G}$，典范映射
$$
\mathcal{F} \otimes_\mathcal{A} f_*\mathcal{G}
\longrightarrow
f_*(f^*\mathcal{F} \otimes_\mathcal{B} \mathcal{G}).
$$
是同构。
\end{lemma}

\begin{proof}
该映射是下述映射的伴随映射：
$$
f^*\mathcal{F} \otimes_\mathcal{B}
f^* f_*\mathcal{G} =
f^*(\mathcal{F} \otimes_\mathcal{A} f_*\mathcal{G})
\longrightarrow
f^*\mathcal{F} \otimes_\mathcal{B} \mathcal{G}
$$
它来自 $\text{id} : f^*\mathcal{F} \to f^*\mathcal{F}$
以及伴随映射 $f^* f_*\mathcal{G} \to \mathcal{G}$。
为证明该映射是同构，可以在茎上检验（《空间的性质》定理
\ref{spaces-properties-theorem-exactness-stalks}）。
设 $\overline{z} : \Spec(k) \to Z$ 为几何点，其像为
$\overline{x} = i \circ \overline{z} : \Spec(k) \to X$。
逐点计算这些映射在茎上的作用可知，需要证明
$$
\mathcal{F}_{\overline{x}}
\otimes_{\mathcal{A}_{\overline{x}}}
\mathcal{G}_{\overline{z}} =
(\mathcal{F}_{\overline{x}}
\otimes_{\mathcal{A}_{\overline{x}}}
\mathcal{B}_{\overline{z}}) \otimes_{\mathcal{B}_{\overline{z}}}
\mathcal{G}_{\overline{z}}
$$
这成立。这里使用了张量积与取茎可交换、拉回下茎的行为（《空间的性质》引理
\ref{spaces-properties-lemma-stalk-pullback}），以及沿闭浸入推前时茎的行为
（引理 \ref{spaces-morphisms-lemma-stalk-push-closed}）。
\end{proof}






\section{闭浸入与拟凝聚层}
\label{spaces-morphisms-section-closed-immersions-quasi-coherent}

\noindent
本节是《态射》第 \ref{morphisms-section-closed-immersions-quasi-coherent} 节的代数空间版本。

\begin{lemma}
\label{spaces-morphisms-lemma-i-star-equivalence}
设 $S$ 为概形。设 $i : Z \to X$ 是 $S$ 上代数空间的闭浸入。
设 $\mathcal{I} \subset \mathcal{O}_X$ 为定义 $Z$ 的拟凝聚理想层。
\begin{enumerate}
\item 对任意 $\mathcal{O}_X$-模 $\mathcal{F}$，伴随映射
$\mathcal{F} \to i_*i^*\mathcal{F}$ 诱导同构
$\mathcal{F}/\mathcal{I}\mathcal{F} \cong i_*i^*\mathcal{F}$。
\item 函子 $i^*$ 是 $i_*$ 的左逆，即对任意
$\mathcal{O}_Z$-模 $\mathcal{G}$，伴随映射
$i^*i_*\mathcal{G} \to \mathcal{G}$ 是同构。
\item 函子
$$
i_* :
\QCoh(\mathcal{O}_Z)
\longrightarrow
\QCoh(\mathcal{O}_X)
$$
是正合且全忠实的，其本质像由满足 $\mathcal{I}\mathcal{F} = 0$ 的
拟凝聚 $\mathcal{O}_X$-模 $\mathcal{F}$ 组成。
\end{enumerate}
\end{lemma}

\begin{proof}
在本证明中我们只处理小 \'etale 位点上的层，并使用 $i_*, i^{-1}, \ldots$
表示阿贝尔群层的推前与拉回，而不使用 $i_{small, *}, i_{small}^{-1}$。

\medskip\noindent
设 $\mathcal{F}$ 为 $\mathcal{O}_X$-模。将引理
\ref{spaces-morphisms-lemma-closed-immersion-rings} 应用于
$\mathcal{A} = \mathcal{O}_X$ 及
$\mathcal{G} = \mathcal{B} = \mathcal{O}_Z$，可知
$i_*i^*\mathcal{F} = \mathcal{F} \otimes_{\mathcal{O}_X} \mathcal{O}_Z$.
由引理 \ref{spaces-morphisms-lemma-closed-immersion-ideals}，有短正合列
$$
0 \to \mathcal{I} \to \mathcal{O}_X \to i_*\mathcal{O}_Z \to 0
$$
由张量积的性质，
$\mathcal{F} \otimes_{\mathcal{O}_X} i_*\mathcal{O}_Z
= \mathcal{F}/\mathcal{I}\mathcal{F}$。这证明了 (1)（但我们省略验证该映射
由伴随映射诱导）。

\medskip\noindent
设 $\mathcal{G}$ 为任意 $\mathcal{O}_Z$-模。由引理
\ref{spaces-morphisms-lemma-closed-immersion-push-pull}，有 $i^{-1}i_*\mathcal{G} = \mathcal{G}$。
因此为证明 (2)，需证明典范映射
$\mathcal{G} \otimes_{i^{-1}\mathcal{O}_X} \mathcal{O}_Z \to \mathcal{G}$
是同构。若能证明 $i^{-1}\mathcal{O}_X \to \mathcal{O}_Z$ 为满射，这由张量积的一般性质得到。
由引理 \ref{spaces-morphisms-lemma-closed-immersion-push-pull}，只需证明
$i_*i^{-1}\mathcal{O}_X \to i_*\mathcal{O}_Z$
是满射。由于满射映射 $\mathcal{O}_X \to i_*\mathcal{O}_Z$
经由此映射因子化，故 (2) 成立。

\medskip\noindent
最后证明引理最重要的部分，即 (3)。闭浸入是拟紧且分离的，见引理
\ref{spaces-morphisms-lemma-closed-immersion-quasi-compact} 和
\ref{spaces-morphisms-lemma-closed-immersion-separated}。因此引理
\ref{spaces-morphisms-lemma-pushforward} 适用，$Z$ 上拟凝聚层的推前确实是 $X$ 上的拟凝聚层。
于是得到函子
$i^{QCoh}_* : \QCoh(\mathcal{O}_Z)
\to \QCoh(\mathcal{O}_X)$.
由 (2) 可知 $i^{QCoh}_*$ 是全忠实的，因为它具有左逆 $i^*$。

\medskip\noindent
现在转而描述函子 $i_*$ 的本质像。对任意 $\mathcal{O}_Z$-模，显然有
$\mathcal{I}(i_*\mathcal{G}) = 0$，因为 $\mathcal{I}$ 是映射
$\mathcal{O}_X \to i_*\mathcal{O}_Z$ 的核，而我们用该映射在 $i_*\mathcal{G}$ 上定义
$\mathcal{O}_X$-模结构。反之，设 $\mathcal{F}$ 是满足
$\mathcal{I}\mathcal{F} = 0$ 的任意拟凝聚 $\mathcal{O}_X$-模。则
$\mathcal{F}$ 是 $i_*\mathcal{O}_Z$-模，因为 $i_*\mathcal{O}_Z = \mathcal{O}_X/\mathcal{I}$。
特别地，其支集包含于 $|Z|$。应用引理
\ref{spaces-morphisms-lemma-closed-immersion-push-pull} 可知存在某个 $\mathcal{O}_Z$-模 $\mathcal{G}$ 使
$\mathcal{F} \cong i_*\mathcal{G}$。剩下的唯一细节是说明 $\mathcal{G}$ 为拟凝聚。
这是因为由 (2) 和《空间的性质》引理
\ref{spaces-properties-lemma-pullback-quasi-coherent}，$\mathcal{G} \cong i^*\mathcal{F}$。
\end{proof}

\noindent
设 $i : Z \to X$ 为代数空间的闭浸入。由上述引理，滥用记号时我们常把
$i_*\mathcal{F}$ 在 $X$ 上的层也记为 $\mathcal{F}$。

\begin{lemma}
\label{spaces-morphisms-lemma-largest-quasi-coherent-subsheaf}
设 $S$ 为概形，设 $X$ 为 $S$ 上的代数空间。
设 $\mathcal{F}$ 为拟凝聚 $\mathcal{O}_X$-模，设 $\mathcal{G} \subset \mathcal{F}$ 为
$\mathcal{O}_X$-子模。存在唯一的拟凝聚 $\mathcal{O}_X$-子模
$\mathcal{G}' \subset \mathcal{G}$，具有如下性质：对每个拟凝聚
$\mathcal{O}_X$-模 $\mathcal{H}$，映射
$$
\Hom_{\mathcal{O}_X}(\mathcal{H}, \mathcal{G}')
\longrightarrow
\Hom_{\mathcal{O}_X}(\mathcal{H}, \mathcal{G})
$$
是双射。特别地，$\mathcal{G}'$ 是包含于 $\mathcal{G}$ 的 $\mathcal{F}$ 的最大拟凝聚
$\mathcal{O}_X$-子模。
\end{lemma}

\begin{proof}
令 $\mathcal{G}_a$（$a \in A$）为包含于 $\mathcal{G}$ 的拟凝聚
$\mathcal{O}_X$-子模的集合。令映射
$$
\bigoplus\nolimits_{a \in A} \mathcal{G}_a \longrightarrow \mathcal{F}
$$
的像 $\mathcal{G}'$。它是拟凝聚的，因为 $X$ 上拟凝聚层之间的映射的像仍为拟凝聚层，
且拟凝聚层的直和是拟凝聚的；见《空间的性质》引理
\ref{spaces-properties-lemma-properties-quasi-coherent}。
模 $\mathcal{G}'$ 包含于 $\mathcal{G}$，故它是包含于 $\mathcal{G}$ 的最大拟凝聚
$\mathcal{O}_X$-模。

\medskip\noindent
为证明该公式，设 $\mathcal{H}$ 为拟凝聚 $\mathcal{O}_X$-模，且
$\alpha : \mathcal{H} \to \mathcal{G}$ 为 $\mathcal{O}_X$-模映射。复合映射
$\mathcal{H} \to \mathcal{G} \to \mathcal{F}$ 的像是拟凝聚的，因为它是拟凝聚层映射的像。
因此该像包含于 $\mathcal{G}'$，从而 $\alpha$ 经由 $\mathcal{G}'$ 分解，如所需。
\end{proof}

\begin{lemma}
\label{spaces-morphisms-lemma-i-upper-shriek}
设 $S$ 为概形。设 $i : Z \to X$ 是 $S$ 上代数空间的闭浸入。
存在函子\footnote{这很可能是非标准记号。}
$i^! : \QCoh(\mathcal{O}_X) \to \QCoh(\mathcal{O}_Z)$
它是 $i_*$ 的右伴随。（比较《模》引理
\ref{modules-lemma-i-star-right-adjoint}。）
\end{lemma}

\begin{proof}
给定拟凝聚 $\mathcal{O}_X$-模 $\mathcal{G}$，考虑 $\mathcal{G}$ 的子层
$\mathcal{H}_Z(\mathcal{G})$，它由被 $\mathcal{I}$ 湮灭的局部截面组成。由引理
\ref{spaces-morphisms-lemma-largest-quasi-coherent-subsheaf}，存在典范的最大拟凝聚
$\mathcal{O}_X$-子模 $\mathcal{H}_Z(\mathcal{G})'$。按构造有
$$
\Hom_{\mathcal{O}_X}(i_*\mathcal{F}, \mathcal{H}_Z(\mathcal{G})')
=
\Hom_{\mathcal{O}_X}(i_*\mathcal{F}, \mathcal{G})
$$
对任意拟凝聚 $\mathcal{O}_Z$-模 $\mathcal{F}$ 成立。
因此可令 $i^!\mathcal{G} = i^*(\mathcal{H}_Z(\mathcal{G})')$。细节略。
\end{proof}

\noindent
利用拟凝聚理想层与闭子空间之间的 $1$ 对 $1$ 对应（见引理
\ref{spaces-morphisms-lemma-closed-immersion-ideals}），可以定义闭子概形的概形论交与概形论并。

\begin{definition}
\label{spaces-morphisms-definition-scheme-theoretic-intersection-union}
设 $S$ 为概形，设 $X$ 为 $S$ 上的代数空间。设 $Z, Y \subset X$ 为闭子空间，
分别对应于拟凝聚理想层 $\mathcal{I}, \mathcal{J} \subset \mathcal{O}_X$。
$Z$ 与 $Y$ 的 {\it 概形论交} 是由 $\mathcal{I} + \mathcal{J}$
定义的 $X$ 的闭子空间。$Z$ 与 $Y$ 的 {\it 概形论并} 是由
$\mathcal{I} \cap \mathcal{J}$ 定义的 $X$ 的闭子空间。
\end{definition}

\noindent
显然，概形论交的构造与 'etale 局部化相容，概形论并亦然。

\begin{lemma}
\label{spaces-morphisms-lemma-scheme-theoretic-intersection}
设 $S$ 为概形，设 $X$ 为 $S$ 上的代数空间。设 $Z, Y \subset X$ 为闭子空间。
令 $Z \cap Y$ 为 $Z$ 与 $Y$ 的概形论交。则 $Z \cap Y \to Z$ 与
$Z \cap Y \to Y$ 是闭浸入，并且
$$
\xymatrix{
Z \cap Y \ar[r] \ar[d] & Z \ar[d] \\
Y \ar[r] & X
}
$$
是 $S$ 上代数空间的笛卡尔图，即 $Z \cap Y = Z \times_X Y$。
\end{lemma}

\begin{proof}
由引理 \ref{spaces-morphisms-lemma-closed-immersion-ideals}，态射 $Z \cap Y \to Z$ 与
$Z \cap Y \to Y$ 是闭浸入。由于概形论交的构造与 'etale 局部化相容，
由概形情形可知该图为笛卡尔图。见《态射》引理
\ref{morphisms-lemma-scheme-theoretic-intersection}。
\end{proof}

\begin{lemma}
\label{spaces-morphisms-lemma-scheme-theoretic-union}
设 $S$ 为概形，设 $X$ 为 $S$ 上的代数空间。设 $Y, Z \subset X$ 为闭子空间。
令 $Y \cup Z$ 为 $Y$ 与 $Z$ 的概形论并，令 $Y \cap Z$ 为 $Y$ 与 $Z$ 的概形论交。
则 $Y \to Y \cup Z$ 与 $Z \to Y \cup Z$ 是闭浸入，且存在短正合列
$$
0 \to \mathcal{O}_{Y \cup Z} \to \mathcal{O}_Y \times \mathcal{O}_Z
\to \mathcal{O}_{Y \cap Z} \to 0
$$
的 $\mathcal{O}_Z$-模，并且图
$$
\xymatrix{
Y \cap Z \ar[r] \ar[d] & Y \ar[d] \\
Z \ar[r] & Y \cup Z
}
$$
在 $S$ 上代数空间范畴中为余笛卡尔图，即
$Y \cup Z = Y \amalg_{Y \cap Z} Z$。
\end{lemma}

\begin{proof}
由引理 \ref{spaces-morphisms-lemma-closed-immersion-ideals}，态射 $Y \to Y \cup Z$ 与
$Z \to Y \cup Z$ 是闭浸入。在短正合列中，我们利用引理
\ref{spaces-morphisms-lemma-i-star-equivalence} 的等价性，把 $X$ 的闭子空间上的拟凝聚模视为 $X$ 上的拟凝聚模。
对列中的第一个映射，使用典范映射
$\mathcal{O}_{Y \cup Z} \to \mathcal{O}_Y$ 和
$\mathcal{O}_{Y \cup Z} \to \mathcal{O}_Z$；
对第二个映射，使用典范映射
$\mathcal{O}_Y \to \mathcal{O}_{Y \cap Z}$ 以及
典范映射的负映射
$\mathcal{O}_Z \to \mathcal{O}_{Y \cap Z}$。为检验正合性，可以在 \'etale 局部工作，并由概形情形推出正合性
（《态射》引理 \ref{morphisms-lemma-scheme-theoretic-union}）。

\medskip\noindent
为证明该图为余笛卡尔图，设给定 $S$ 上的代数空间 $T$ 以及态射
$f : Y \to T$、$g : Z \to T$，它们在 $Y \cap Z \to T$ 上相同。目标是证明存在唯一的
态射 $h : Y \cup Z \to T$，其与 $f$、$g$ 相容。
构造 $h$ 时，可以在 $Y \cup Z$ 上作 \'etale 局部工作（因为作为代数空间，$Y \cup Z$ 是 \'etale 层）。
因此可假设 $X$ 为概形。此时由《态射》引理
\ref{morphisms-lemma-scheme-theoretic-union}，$Y \cup Z$ 是概形范畴中沿 $Y \cap Z$
将 $Y$ 与 $Z$ 粘合所得的推出。取概形 $T'$ 及满射 \'etale 态射 $T' \to T$。
令 $Y' = T' \times_{T, f} Y$，$Z' = T' \times_{T, g} Z$。
则 $Y'$ 与 $Z'$ 为概形，并有概形的典范同构
$\varphi : Y' \times_Y (Y \cap Z) \to Z' \times_Z (Y \cap Z)$。由《态射补篇》引理
\ref{more-morphisms-lemma-pushout-along-closed-immersions}
推出 $W' = Y' \amalg_{Y' \times_Y (Y \cap Z), \varphi} Z'$ 在概形范畴中存在。
由《态射补篇》引理
\ref{more-morphisms-lemma-pushout-along-closed-immersions-properties-above}，
态射 $W' \to Y \cup Z$ 是 \'etale 的。该态射是满射，因为 $Y' \to Y$ 与 $Z' \to Z$ 是满射。
态射 $f' : Y' \to T'$ 与 $g' : Z' \to T'$ 粘合为唯一的概形态射
$h' : W' \to T'$。由唯一性，复合 $W' \to T' \to T$ 下降为所需的态射
$h : Y \cup Z \to T$。细节略。
\end{proof}






\section{模的支集}
\label{spaces-morphisms-section-support}

\noindent
本节汇集代数空间上拟凝聚模支集的一些基本结果。设 $X$ 为代数
空间。$X_\etale$ 上阿贝尔层的支集已在《空间的性质》第
\ref{spaces-properties-section-support} 节中定义。我们对模的支集使用同一定义。
下述引理说明，这与概形上拟凝聚模所定义的概念一致。

\begin{lemma}
\label{spaces-morphisms-lemma-support-covering}
设 $S$ 为概形，设 $X$ 为 $S$ 上的代数空间。设 $\mathcal{F}$ 为拟凝聚
$\mathcal{O}_X$-模。设 $U$ 为概形，且 $\varphi : U \to X$ 为 \'etale 态射。则
$$
\text{Supp}(\varphi^*\mathcal{F}) = |\varphi|^{-1}(\text{Supp}(\mathcal{F}))
$$
其中左边是 $\varphi^*\mathcal{F}$ 作为概形 $U$ 上拟凝聚模时的支集。
\end{lemma}

\begin{proof}
 设 $u\in U$ 为普通点，设 $\overline{x}$ 为位于 $u$ 上方的几何点。由《空间的性质》引理
\ref{spaces-properties-lemma-stalk-quasi-coherent}，有
$(\varphi^*\mathcal{F})_u \otimes_{\mathcal{O}_{U, u}}
\mathcal{O}_{X, \overline{x}} = \mathcal{F}_{\overline{x}}$.
由于由《空间的性质》引理
\ref{spaces-properties-lemma-describe-etale-local-ring}，$\mathcal{O}_{U, u} \to \mathcal{O}_{X, \overline{x}}$
是严格 Hensel 化，故它是忠实平坦的（见《代数补篇》引理
\ref{more-algebra-lemma-dumb-properties-henselization}）。因此
$(\varphi^*\mathcal{F})_u = 0$ 当且仅当 $\mathcal{F}_{\overline{x}} = 0$。引理得证。
\end{proof}

\noindent
对有限型拟凝聚模，支集是闭的，可在纤维上检验，并与基变换相容。

\begin{lemma}
\label{spaces-morphisms-lemma-support-finite-type}
设 $S$ 为概形，设 $X$ 为 $S$ 上的代数空间。设 $\mathcal{F}$ 为有限型拟凝聚
$\mathcal{O}_X$-模。则
\begin{enumerate}
\item $\mathcal{F}$ 的支集是闭集。
\item 对于位于 $x \in |X|$ 上方的几何点 $\overline{x}$，有
$$
x \in \text{Supp}(\mathcal{F})
\Leftrightarrow
\mathcal{F}_{\overline{x}} \not = 0
\Leftrightarrow
\mathcal{F}_{\overline{x}} \otimes_{\mathcal{O}_{X, \overline{x}}}
\kappa(\overline{x}) \not = 0.
$$
\item 对任意代数空间态射 $f : Y \to X$，拉回 $f^*\mathcal{F}$ 也是有限型的，且
$\text{Supp}(f^*\mathcal{F}) = f^{-1}(\text{Supp}(\mathcal{F}))$。
\end{enumerate}
\end{lemma}

\begin{proof}
取概形 $U$ 及满射 \'etale 态射 $\varphi : U \to X$。由引理
\ref{spaces-morphisms-lemma-support-covering}，$\mathcal{F}$ 的支集的逆像是 $\varphi^*\mathcal{F}$ 的支集，
后者由《态射》引理 \ref{morphisms-lemma-support-finite-type} 为闭集。
因此 (1) 由 $|X|$ 上拓扑的定义得到。

\medskip\noindent
(2) 中第一个等价关系就是支集的定义。第二个等价关系由 Nakayama 引理得到，见
《代数》引理 \ref{algebra-lemma-NAK}。

\medskip\noindent
设 $f : Y \to X$ 如 (3) 所述。由《空间的性质》第
\ref{spaces-properties-section-properties-modules} 节，$f^*\mathcal{F}$ 是有限型的。
为证明最后一个断言，设 $\overline{y}$ 为 $Y$ 的几何点，映到 $X$ 上的几何点 $\overline{x}$。回忆
$$
(f^*\mathcal{F})_{\overline{y}} =
\mathcal{F}_{\overline{x}} \otimes_{\mathcal{O}_{X, \overline{x}}}
\mathcal{O}_{Y, \overline{y}},
$$
见《空间的性质》引理
\ref{spaces-properties-lemma-stalk-pullback-quasi-coherent}。
因此 $(f^*\mathcal{F})_{\overline{y}} \otimes \kappa(\overline{y})$ 非零当且仅当
$\mathcal{F}_{\overline{x}} \otimes \kappa(\overline{x})$ 非零。
由 (2)，这意味着 $x \in \text{Supp}(\mathcal{F})$ 当且仅当
$y \in \text{Supp}(f^*\mathcal{F})$，这正是断言 (3)。
\end{proof}

\noindent
下一个任务是说明，有限型拟凝聚模的概形论支集（见《态射》定义
\ref{morphisms-definition-scheme-theoretic-support}）
也可以对代数空间上的有限型拟凝聚模定义。

\begin{lemma}
\label{spaces-morphisms-lemma-scheme-theoretic-support}
设 $S$ 为概形，设 $X$ 为 $S$ 上的代数空间。设 $\mathcal{F}$ 为有限型拟凝聚
$\mathcal{O}_X$-模。存在最小闭子空间 $i : Z \to X$，使得存在拟凝聚
$\mathcal{O}_Z$-模 $\mathcal{G}$ 满足 $i_*\mathcal{G} \cong \mathcal{F}$。此外：
\begin{enumerate}
\item 若 $U$ 为概形且 $\varphi : U \to X$ 为 \'etale 态射，则 $Z \times_X U$ 是
$\varphi^*\mathcal{F}$ 的概形论支集。
\item 拟凝聚层 $\mathcal{G}$ 在唯一同构意义下唯一。
\item 拟凝聚层 $\mathcal{G}$ 是有限型的。
\item $\mathcal{G}$ 与 $\mathcal{F}$ 的支集均为 $|Z|$。
\end{enumerate}
\end{lemma}

\begin{proof}
取概形 $U$ 及满射 \'etale 态射 $\varphi : U \to X$。令 $R = U \times_X U$，投影为
$s, t : R \to U$。令 $i' : Z' \to U$ 为 $\varphi^*\mathcal{F}$ 的概形论支集，
并令 $\mathcal{G}'$ 为（在唯一同构意义下）唯一的有限型拟凝聚
$\mathcal{O}_{Z'}$-模，满足 $i'_*\mathcal{G}' = \varphi^*\mathcal{F}$；见《态射》引理
\ref{morphisms-lemma-scheme-theoretic-support}。
由于 $s^*\varphi^*\mathcal{F} = t^*\varphi^*\mathcal{F}$，可知
由《态射》引理 \ref{morphisms-lemma-flat-pullback-support}，
$R' = s^{-1}Z' = t^{-1}Z'$ 是 $R$ 的闭子概形。于是可应用《空间的性质》引理
\ref{spaces-properties-lemma-subspaces-presentation}，得到闭子空间 $i : Z \to X$，
其到 $U$ 的拉回为 $Z'$。记投影为 $s', t' : R' \to Z'$，并记给定的闭浸入为
$j' : R' \to R$，则
$$
j'_* (s')^*\mathcal{G}' = s^* i'_*\mathcal{G}' =
s^*\varphi^*\mathcal{F} = t^*\varphi^*\mathcal{F} =
t^*i'_*\mathcal{G}' = j'_*(t')^*\mathcal{G}'
$$
（第一个和最后一个等式由《概形的上同调》引理
\ref{coherent-lemma-flat-base-change-cohomology} 给出）。因此将《态射》引理
\ref{morphisms-lemma-flat-pullback-support} 应用于 $R' \to R$ 的唯一性，给出同构
$\alpha : (t')^*\mathcal{G}' \to (s')^*\mathcal{G}'$
它经由 $j'_*$ 与典范同构
$t^*\varphi^*\mathcal{F} = s^*\varphi^*\mathcal{F}$ 相容。显然 $\alpha$ 满足上链条件，
因此可应用《空间的性质》命题
\ref{spaces-properties-proposition-quasi-coherent}
得到 $Z$ 上的拟凝聚模 $\mathcal{G}$，其到 $Z'$ 的限制为与 $\alpha$ 相容的 $\mathcal{G}'$。
再次使用上述命题的等价性（这次对 $X$ 使用），可知 $i_*\mathcal{G} \cong \mathcal{F}$。

\medskip\noindent
这就证明了存在性。将它与概形的结果比较，并使用引理
\ref{spaces-morphisms-lemma-support-covering}，即可得到引理的其他性质。详细证明略。
\end{proof}

\begin{definition}
\label{spaces-morphisms-definition-scheme-theoretic-support}
设 $S$ 为概形，设 $X$ 为 $S$ 上的代数空间。设 $\mathcal{F}$ 为有限型拟凝聚
$\mathcal{O}_X$-模。$\mathcal{F}$ 的 {\it 概形论支集} 是引理
\ref{spaces-morphisms-lemma-scheme-theoretic-support} 中构造的闭子空间 $Z \subset X$。
\end{definition}

\noindent
在此情形中，我们常通过引理 \ref{spaces-morphisms-lemma-i-star-equivalence} 的范畴等价，
把 $\mathcal{F}$ 视为 $Z$ 上的有限型拟凝聚层。






\section{概形论像}
\label{spaces-morphisms-section-scheme-theoretic-image}

\noindent
注意：本节的部分内容非常一般，其行为可能与预期不同。

\begin{lemma}
\label{spaces-morphisms-lemma-scheme-theoretic-image}
\begin{slogan}
代数空间态射的概形论像存在。
\end{slogan}
设 $S$ 为概形。设 $f : X \to Y$ 是 $S$ 上代数空间的态射。存在闭子空间
$Z \subset Y$，使得 $f$ 经由 $Z$ 分解，并且对任意其他闭子空间
$Z' \subset Y$，使 $f$ 经由 $Z'$ 分解，都有 $Z \subset Z'$。
\end{lemma}

\begin{proof}
令 $\mathcal{I} = \Ker(\mathcal{O}_Y \to f_*\mathcal{O}_X)$。若 $\mathcal{I}$ 是拟凝聚的，
则取由 $\mathcal{I}$ 定义的闭子概形 $Z$ 即可，见引理
\ref{spaces-morphisms-lemma-closed-immersion-ideals}。一般情形下，需要证明存在包含于
$\mathcal{I}$ 的最大拟凝聚理想层 $\mathcal{I}'$。这由引理
\ref{spaces-morphisms-lemma-largest-quasi-coherent-subsheaf} 得到。
\end{proof}

\noindent
假设上述引理 \ref{spaces-morphisms-lemma-scheme-theoretic-image} 的情形中 $X$ 和 $Y$ 可表示。
引理中得到的闭子空间 $Z \subset Y$ 与《态射》引理
\ref{morphisms-lemma-scheme-theoretic-image} 中得到的闭子概形 $Z \subset Y$ 一致。
原因是，闭子空间（或子概形）与 $X_\etale$ 和 $X$ 上的拟凝聚理想层成反包含对应。
由于 $X_\etale$ 与 $X$ 上拟凝聚模的范畴相同（《空间的性质》第
\ref{spaces-properties-section-quasi-coherent} 节），结论得证。因此下述定义与早先对概形态射的定义一致。

\begin{definition}
\label{spaces-morphisms-definition-scheme-theoretic-image}
设 $S$ 为概形。设 $f : X \to Y$ 是 $S$ 上代数空间的态射。$f$ 的
{\it 概形论像}是使 $f$ 经由其分解的最小闭子空间 $Z \subset Y$；见上面的引理
\ref{spaces-morphisms-lemma-scheme-theoretic-image}。
\end{definition}

\noindent
我们通常只用 $f : X \to Z$ 表示 $f$ 的分解。若态射 $f$ 不是拟紧的，
则一般而言，概形论像的构造与限制到 $Y$ 的开子空间不相容。

\begin{lemma}
\label{spaces-morphisms-lemma-quasi-compact-scheme-theoretic-image}
设 $S$ 为概形。设 $f : X \to Y$ 是 $S$ 上代数空间的态射。
设 $Z \subset Y$ 为 $f$ 的概形论像。若 $f$ 拟紧，则
\begin{enumerate}
\item 理想层
$\mathcal{I} = \Ker(\mathcal{O}_Y \to f_*\mathcal{O}_X)$
是拟凝聚的；
\item 概形论像 $Z$ 是与 $\mathcal{I}$ 对应的闭子空间；
\item 对任意 \'etale 态射 $V \to Y$，$X \times_Y V \to V$ 的概形论像等于
$Z \times_Y V$；以及
\item 像 $|f|(|X|) \subset |Z|$ 是 $|Z|$ 的稠密子集。
\end{enumerate}
\end{lemma}

\begin{proof}
由于 $\mathcal{I}$ 的构造与 \'etale 局部化相容，为证明 (3)，只需证明 (1) 和 (2)。
若 (1) 成立，则我们已在引理 \ref{spaces-morphisms-lemma-scheme-theoretic-image} 的证明中证明 (2)。
现在证明 $\mathcal{I}$ 为拟凝聚。拟凝聚性是 \'etale 局部性质，故可假设 $Y$ 为仿射概形。
由于 $f$ 拟紧，可取仿射概形 $U$ 及满射 \'etale 态射 $U \to X$。
记 $f'$ 为复合 $U \to X \to Y$。则 $f_*\mathcal{O}_X$ 是
$f'_*\mathcal{O}_U$ 的子层，从而 $\mathcal{I} = \Ker(\mathcal{O}_Y \to \mathcal{O}_{X'})$。
由引理 \ref{spaces-morphisms-lemma-pushforward}，层 $f'_*\mathcal{O}_U$ 在 $Y$ 上拟凝聚。
因此，作为凝聚模间映射的核，$\mathcal{I}$ 是拟凝聚的。
最后，(4) 由 (1)、(2) 和 (3) 得到，因为在 $|Y|$ 中任何不包含于
$|f|(|X|)$ 的闭包的点上，理想 $\mathcal{I}$ 都是单位理想。
\end{proof}

\begin{lemma}
\label{spaces-morphisms-lemma-scheme-theoretic-image-reduced}
设 $S$ 为概形。设 $f : X \to Y$ 是 $S$ 上代数空间的态射。假设 $X$ 约化。则
\begin{enumerate}
\item $f$ 的概形论像 $Z$ 是 $\overline{|f|(|X|)}$ 上的约化诱导代数空间结构；以及
\item 对任意 \'etale 态射 $V \to Y$，$X \times_Y V \to V$ 的概形论像等于
$Z \times_Y V$。
\end{enumerate}
\end{lemma}

\begin{proof}
(1) 成立，因为 $\overline{|f|(|X|)}$ 上的约化诱导代数空间结构是 $Y$ 中
使 $f$ 经由其分解的最小闭子空间；见《空间的性质》引理
\ref{spaces-properties-lemma-map-into-reduction}。
(2) 由 (1)、$|V| \to |Y|$ 为开映射，以及约化性在 \'etale 局部化下保持不变得到。
\end{proof}

\begin{lemma}
\label{spaces-morphisms-lemma-reach-points-scheme-theoretic-image}
设 $S$ 为概形。设 $f : X \to Y$ 是 $S$ 上代数空间的拟紧态射。
设 $Z$ 为 $f$ 的概形论像。设 $z \in |Z|$。存在分式域为 $K$ 的赋值环 $A$，
以及交换图
$$
\xymatrix{
\Spec(K) \ar[rr] \ar[d] & & X \ar[d] \ar[ld] \\
\Spec(A) \ar[r] & Z \ar[r] & Y
}
$$
使得 $\Spec(A)$ 的闭点映到 $z$。
\end{lemma}

\begin{proof}
取含点 $z' \in V$ 的仿射概形 $V$，以及把 $z'$ 映到 $z$ 的 \'etale 态射
$V \to Y$。令 $Z' \subset V$ 为 $X \times_Y V \to V$ 的概形论像。
由引理 \ref{spaces-morphisms-lemma-quasi-compact-scheme-theoretic-image}，有
$Z' = Z \times_Y V$，故 $z' \in Z'$。由于 $f$ 拟紧且 $V$ 仿射，
$X \times_Y V$ 拟紧。因此存在仿射概形 $W$ 及满射 \'etale 态射
$W \to X \times_Y V$。此时 $Z' \subset V$ 也是 $W \to V$ 的概形论像。
由《态射》引理 \ref{morphisms-lemma-reach-points-scheme-theoretic-image}，可取图
$$
\xymatrix{
\Spec(K) \ar[r] \ar[d] &
W \ar[r] \ar[d] &
X \times_Y V \ar[d] \ar[r] &
X \ar[d] \\
\Spec(A) \ar[r] &
Z' \ar[r] &
V \ar[r] &
Y
}
$$
使得 $\Spec(A)$ 的闭点映到 $z'$。再与 $Z' \to Z$ 及
$W \to X \times_Y V \to X$ 复合，即得所求。
\end{proof}

\begin{lemma}
\label{spaces-morphisms-lemma-factor-factor}
设 $S$ 为概形。设
$$
\xymatrix{
X_1 \ar[d] \ar[r]_{f_1} & Y_1 \ar[d] \\
X_2 \ar[r]^{f_2} & Y_2
}
$$
为 $S$ 上代数空间的交换图。令 $Z_i \subset Y_i$（$i = 1, 2$）为 $f_i$ 的概形论像。
则态射 $Y_1 \to Y_2$ 诱导态射 $Z_1 \to Z_2$，并得到交换图
$$
\xymatrix{
X_1 \ar[r] \ar[d] & Z_1 \ar[d] \ar[r] & Y_1 \ar[d] \\
X_2 \ar[r] & Z_2 \ar[r] & Y_2
}
$$
\end{lemma}

\begin{proof}
$Z_2$ 在 $Y_1$ 中的概形论逆像是 $Y_1$ 的闭子空间，且 $f_1$ 经由它分解。
因此 $Z_1$ 包含于该逆像中。引理得证。
\end{proof}

\begin{lemma}
\label{spaces-morphisms-lemma-scheme-theoretic-image-of-partial-section}
设 $S$ 为概形。设 $f : X \to Y$ 是 $S$ 上代数空间的分离态射。
设 $V \subset Y$ 为开子空间，且 $V \to Y$ 拟紧。设 $s : V \to X$ 为满足
$f \circ s = \text{id}_V$ 的态射。令 $Y'$ 为 $s$ 的概形论像。
则 $Y' \to Y$ 在 $V$ 上是同构。
\end{lemma}

\begin{proof}
由引理 \ref{spaces-morphisms-lemma-quasi-compact-permanence}，态射 $s : V \to X$ 拟紧。
因此由引理 \ref{spaces-morphisms-lemma-quasi-compact-scheme-theoretic-image}，$s$ 的概形论像 $Y'$ 的构造
与限制到开子空间相容。特别地，$Y' \cap f^{-1}(V)$ 是分离态射
$f^{-1}(V) \to V$ 的一个截面的概形论像。分离态射的截面是闭浸入
（引理 \ref{spaces-morphisms-lemma-section-immersion}），故
$Y' \cap f^{-1}(V) \to V$ 是同构，如所需。
\end{proof}








\section{概形论闭包与稠密性}
\label{spaces-morphisms-section-scheme-theoretic-closure}

\noindent
本节是《态射》第 \ref{morphisms-section-scheme-theoretic-closure} 节的代数空间版本。

\begin{lemma}
\label{spaces-morphisms-lemma-scheme-theoretically-dense-representable}
设 $S$ 为概形。设 $W \subset S$ 为概形论稠密的开子概形
（《态射》定义 \ref{morphisms-definition-scheme-theoretically-dense}）。
设 $f : X \to S$ 为概形态射，且平坦、局部有限表示并局部拟有限。
则 $f^{-1}(W)$ 在 $X$ 中概形论稠密。
\end{lemma}

\begin{proof}
我们将使用《态射》引理
\ref{morphisms-lemma-characterize-scheme-theoretically-dense} 的刻画。
设 $V \subset X$ 为开集，且 $g \in \Gamma(V, \mathcal{O}_V)$ 是在
$f^{-1}(W) \cap V$ 上限制为零的函数。需证明 $g = 0$。假设 $g \not = 0$ 以导出矛盾。
由《态射补篇》引理 \ref{more-morphisms-lemma-go-down-with-annihilators}，
可缩小 $V$，并找到开集 $U \subset S$ 使其落入交换图
$$
\xymatrix{
V \ar[r] \ar[d]_\pi & X \ar[d]^f \\
U \ar[r] & S,
}
$$
一个拟凝聚子层 $\mathcal{F} \subset \mathcal{O}_U$、整数 $r > 0$，以及单射
$\mathcal{O}_U$-模映射 $\mathcal{F}^{\oplus r} \to \pi_*\mathcal{O}_V$，
其像包含 $g|_V$。设 $(g_1, \ldots, g_r) \in \Gamma(U, \mathcal{F}^{\oplus r})$ 映到 $g$。
由于 $g|_{f^{-1}W \cap V} = 0$，可知 $g_i|_{W \cap U} = 0$。
又因 $\mathcal{F} \subset \mathcal{O}_U$ 且 $W$ 在 $S$ 中概形论稠密，故 $g_i = 0$。
这推出 $g = 0$，与假设矛盾。
\end{proof}

\begin{lemma}
\label{spaces-morphisms-lemma-scheme-theoretically-dense}
设 $S$ 为概形。设 $X$ 为 $S$ 上的代数空间。设 $U \subset X$ 为开子空间。
下列条件等价：
\begin{enumerate}
\item 对每个代数空间的 \'etale 态射 $\varphi : V \to X$，$\varphi^{-1}(U)$ 在 $V$ 中的
概形论闭包等于 $V$；
\item 存在概形 $V$ 及满射 \'etale 态射 $\varphi : V \to X$，使得 $\varphi^{-1}(U)$ 在 $V$ 中的
概形论闭包等于 $V$；
\end{enumerate}
\end{lemma}

\begin{proof}
注意，若 $V \to V'$ 是在 $X$ 上 \'etale 的代数空间态射，且 $Z \subset V$、
相应地 $Z' \subset V'$ 分别是 $U \times_X V$、$U \times_X V'$ 在 $V$、$V'$ 中的概形论闭包，
则 $Z$ 映到 $Z'$。因此若 $V \to V'$ 满射且 \'etale，则 $Z = V$ 推出 $Z' = V'$。
此外，注意 \'etale 态射是平坦、局部有限表示且局部拟有限的
（见《态射》第 \ref{morphisms-section-etale} 节）。因此由引理
\ref{spaces-morphisms-lemma-scheme-theoretically-dense-representable}，若 $V$ 与 $V'$ 为概形，
则 $Z' = V'$ 推出 $Z = V$。利用每个代数空间都有概形的 \'etale 覆盖作形式论证，
可知 (1) 与 (2) 等价。
\end{proof}

\noindent
由引理 \ref{spaces-morphisms-lemma-scheme-theoretically-dense}，下述定义与概形情形的定义相容。

\begin{definition}
\label{spaces-morphisms-definition-scheme-theoretically-dense}
设 $S$ 为概形。设 $X$ 为 $S$ 上的代数空间。设 $U \subset X$ 为开子空间。
\begin{enumerate}
\item 态射 $U \to X$ 的概形论像称为 $U$ 在 $X$ 中的 {\it 概形论闭包}。
\item 若引理 \ref{spaces-morphisms-lemma-scheme-theoretically-dense} 的等价条件成立，则称 $U$ 在 $X$ 中
{\it 概形论稠密}。
\end{enumerate}
\end{definition}

\noindent
按此定义，$U$ 在 $X$ 中概形论稠密并不等价于 $U$ 的概形论闭包为 $X$。
这略显不优雅，但在适当的有限性条件下，我们将看到它确实成立。

\begin{lemma}
\label{spaces-morphisms-lemma-scheme-theoretically-dense-quasi-compact}
设 $S$ 为概形，设 $X$ 为 $S$ 上的代数空间，设 $U \subset X$ 为开子空间。
若 $U \to X$ 拟紧，则 $U$ 在 $X$ 中概形论稠密，当且仅当 $U$ 在 $X$ 中的
概形论闭包为 $X$。
\end{lemma}

\begin{proof}
由引理 \ref{spaces-morphisms-lemma-quasi-compact-scheme-theoretic-image} 的 (3) 得到。
\end{proof}

\begin{lemma}
\label{spaces-morphisms-lemma-characterize-scheme-theoretically-dense}
设 $S$ 为概形。设 $j : U \to X$ 是 $S$ 上代数空间的开浸入。
则 $U$ 在 $X$ 中概形论稠密，当且仅当
$\mathcal{O}_X \to j_*\mathcal{O}_U$ 是单射。
\end{lemma}

\begin{proof}
若 $\mathcal{O}_X \to j_*\mathcal{O}_U$ 为单射，
则将其限制到任意在 $X$ 上 \'etale 的代数空间 $V$ 仍为单射。
因此 $U \times_X V$ 在 $V$ 中的概形论闭包等于 $V$，见引理
\ref{spaces-morphisms-lemma-scheme-theoretic-image} 的证明。
反之，假设 $U \times_X V$ 的概形论闭包等于 $V$，对所有在 $X$ 上 \'etale 的 $V$ 都成立。
若 $\mathcal{O}_X \to j_*\mathcal{O}_U$ 不是单射，则可找到一个仿射概形（记为 $V = \Spec(A)$），
它在 $X$ 上 \'etale，并有非零元素 $f \in A$，使得 $f$ 在
$\Gamma(V \times_X U, \mathcal{O})$ 中映为零。此时 $V \times_X U$ 在 $V$ 中的概形论闭包
显然包含于 $\Spec(A/(f))$，矛盾。
\end{proof}

\begin{lemma}
\label{spaces-morphisms-lemma-intersection-scheme-theoretically-dense}
设 $S$ 为概形，设 $X$ 为 $S$ 上的代数空间。若 $U$、$V$ 是 $X$ 中概形论稠密的
开子空间，则 $U \cap V$ 亦然。
\end{lemma}

\begin{proof}
设 $W \to X$ 为任意 \'etale 态射。考虑映射
$\mathcal{O}(W) \to \mathcal{O}(W \times_X V)
\to \mathcal{O}(W \times_X (V \cap U))$.
由引理 \ref{spaces-morphisms-lemma-characterize-scheme-theoretically-dense}，两个映射都是单射，
故复合映射为单射。再由引理 \ref{spaces-morphisms-lemma-characterize-scheme-theoretically-dense}，
$U \cap V$ 在 $X$ 中概形论稠密。
\end{proof}

\begin{lemma}
\label{spaces-morphisms-lemma-quasi-compact-immersion}
设 $S$ 为概形。设 $h : Z \to X$ 是 $S$ 上代数空间的浸入。假设 $Z \to X$ 拟紧或 $Z$ 约化。
设 $\overline{Z} \subset X$ 为 $h$ 的概形论像。则态射 $Z \to \overline{Z}$ 是开浸入，
并将 $Z$ 识别为 $\overline{Z}$ 中概形论稠密的开子空间。此外，$Z$ 在 $\overline{Z}$ 中拓扑稠密。
\end{lemma}

\begin{proof}
在两种情形中，$\overline{Z}$ 的构造都与 \'etale 局部化相容，见引理
\ref{spaces-morphisms-lemma-quasi-compact-scheme-theoretic-image} 和
\ref{spaces-morphisms-lemma-scheme-theoretic-image-reduced}。因此由概形情形可得本引理，见《态射》引理
\ref{morphisms-lemma-quasi-compact-immersion}。
\end{proof}

\begin{lemma}
\label{spaces-morphisms-lemma-equality-of-morphisms}
设 $S$ 为概形，设 $B$ 为 $S$ 上的代数空间。设 $f, g : X \to Y$ 为 $B$ 上代数空间的态射。
设 $U \subset X$ 为开子空间且 $f|_U = g|_U$。若 $U$ 在 $X$ 中的概形论闭包为 $X$，
且 $Y \to B$ 分离，则 $f = g$。
\end{lemma}

\begin{proof}
由于 $Y \to B$ 分离，纤维积
$Y \times_{\Delta, Y \times_B Y, (f, g)} X$ 是闭子空间 $Z \subset X$。
由于 $f|_U = g|_U$，有 $U \subset Z$。因假设 $U$ 在 $X$ 中概形论稠密，故 $Z = X$。
\end{proof}






\section{主导态射}
\label{spaces-morphisms-section-dominant}

\noindent
我们复制概形主导态射的定义，以得到代数空间主导态射的概念。注意，除非态射是拟紧的且代数空间满足某些分离性公理，该定义的性质不太好。

\begin{definition}
\label{spaces-morphisms-definition-dominant}
设 $S$ 为概形。若 $S$ 上代数空间的态射 $f : X \to Y$ 所诱导的
$|f| : |X| \to |Y|$ 的像在 $|Y|$ 中稠密，则称为{\it 主导}。
\end{definition}






\section{普遍单射态射}
\label{spaces-morphisms-section-universally-injective}

\noindent
在第 \ref{spaces-morphisms-section-representable} 节中，我们已经定义了代数空间的可表示态射普遍单射的含义。对 $S$ 上的域 $K$（回忆，这意味着给定结构态射 $\Spec(K) \to S$）以及 $S$ 上的代数空间 $X$，我们写作
$X(K) = \Mor_S(\Spec(K), X)$。我们先将可表示态射的条件转换为关于点函子的条件。

\begin{lemma}
\label{spaces-morphisms-lemma-universally-injective-representable}
设 $S$ 为概形。设 $f : X \to Y$ 为 $S$ 上代数空间的可表示态射。则 $f$ 普遍单射
（按第 \ref{spaces-morphisms-section-representable} 节的含义）当且仅当对所有域 $K$，映射
$X(K) \to Y(K)$ 是单射。
\end{lemma}

\begin{proof}
以下不再另行说明，我们将使用《态射》引理
\ref{morphisms-lemma-universally-injective}。
假设 $f$ 普遍单射。则对任意域 $K$ 及任意态射 $\Spec(K) \to Y$，概形态射
$\Spec(K) \times_Y X \to \Spec(K)$ 普遍单射。因此该态射至多有一个截面
$\Spec(K) \times_Y X \to \Spec(K)$，故 $X(K) \to Y(K)$ 为单射。反之，假设对每个域 $K$，映射 $X(K) \to Y(K)$ 为单射。
设 $T \to Y$ 为从概形到 $Y$ 的态射，考虑基变换 $f_T : T \times_Y X \to T$。
对任意域 $K$，有
$$
(T \times_Y X)(K) = T(K) \times_{Y(K)} X(K)
$$
按纤维积的定义，故 $X(K) \to Y(K)$ 的单射性保证
$(T \times_Y X)(K) \to T(K)$ 单射，这意味着 $f_T$ 普遍单射，如所需。
\end{proof}

\noindent
接下来，把域值点之间的变换为单射这一性质转化为更具几何性的条件。

\begin{lemma}
\label{spaces-morphisms-lemma-universally-injective}
设 $S$ 为概形。
设 $f : X \to Y$ 为 $S$ 上代数空间的态射。
以下各条件等价：
\begin{enumerate}
\item 对每个 $S$ 上的域 $K$，映射 $X(K) \to Y(K)$ 是单射；
\item 对每个 $S$ 上代数空间的态射 $Y' \to Y$，诱导映射
$|Y' \times_Y X| \to |Y'|$ 是单射；
\item 对角态射 $X \to X \times_Y X$ 是满射。
\end{enumerate}
\end{lemma}

\begin{proof}
假设 (1)。设 $g : Y' \to Y$ 为代数空间的态射，并记
$f' : Y' \times_Y X \to Y'$ 为 $f$ 的基变换。
设 $K_i$（$i = 1, 2$）为域，并设
$\varphi_i : \Spec(K_i) \to Y' \times_Y X$ 为态射，
使得 $f' \circ \varphi_1$ 与 $f' \circ \varphi_2$ 在 $|Y'|$ 中定义同一点。
按定义，这意味着存在域 $\Omega$ 以及嵌入
$\alpha_i : K_i \subset \Omega$，使得两个态射
$f' \circ \varphi_i \circ \alpha_i : \Spec(\Omega) \to Y'$ 相等。
相应的交换图如下：
$$
\xymatrix{
\Spec(\Omega) \ar@/_5ex/[ddrr] \ar[rd]^{\alpha_1} \ar[r]_{\alpha_2} &
\Spec(K_2) \ar[rd]^{\varphi_2} \\
& \Spec(K_1) \ar[r]^{\varphi_1} &
Y' \times_Y X  \ar[d]^{f'} \ar[r]^{g'} &
 X \ar[d]^f \\
& & Y' \ar[r]^g & Y.
}
$$
特别地，复合态射 $g \circ f' \circ \varphi_i \circ \alpha_i$
相等。由假设 (1)，这意味着态射
$g' \circ \varphi_i \circ \alpha_i$ 相等，其中 $g' : Y' \times_Y X \to X$
是投影。由纤维积的泛性质可知，态射
$\varphi_i \circ \alpha_i : \Spec(\Omega) \to Y' \times_Y X$ 相等。
换言之，$\varphi_1$ 与 $\varphi_2$ 在 $Y' \times_Y X$ 中定义同一点。
因此 (2) 成立。

\medskip\noindent
假设 (2)。设 $K$ 为 $S$ 上的域，设 $a, b : \Spec(K) \to X$
为满足 $f \circ a = f \circ b$ 的两个态射。记
$c : \Spec(K) \to Y$ 为公共值。由假设，
$|\Spec(K) \times_{c, Y} X| \to |\Spec(K)|$ 是单射。
这意味着存在域 $\Omega$ 以及嵌入
$\alpha_i : K \to \Omega$，使得
$$
\xymatrix{
\Spec(\Omega) \ar[r]_{\alpha_1} \ar[d]_{\alpha_2} &
\Spec(K) \ar[d]^a \\
\Spec(K) \ar[r]^-b &
\Spec(K) \times_{c, Y} X
}
$$
是交换的。与到 $\Spec(K)$ 的投影复合可得
$\alpha_1 = \alpha_2$。记公共值为 $\alpha$。
于是 $\{\alpha : \Spec(\Omega) \to \Spec(K)\}$
是 $\Spec(K)$ 的 fpqc 覆盖，并且两个态射
$a, b$ 在该覆盖的各成员上相等。由《代数空间的性质》命题
\ref{spaces-properties-proposition-sheaf-fpqc}
可知 $a = b$。因此 (1) 成立。

\medskip\noindent
假设 (3)。设 $x, x' \in |X|$ 为满足
$f(x) = f(x')$（在 $|Y|$ 中）的两个点。由《代数空间的性质》引理
\ref{spaces-properties-lemma-points-cartesian}
可知存在 $x'' \in |X \times_Y X|$，其两个投影
分别为 $x$ 与 $x'$。由假设以及《代数空间的性质》
引理 \ref{spaces-properties-lemma-characterize-surjective}
可知存在 $x''' \in |X|$，使得 $\Delta_{X/Y}(x''') = x''$。
因此 $x = x'$。换言之，$f$ 是单射。
由于条件 (3) 在基变换下保持不变，故 $f$ 满足 (2)。

\medskip\noindent
假设 (2)。特别地，$|X \times_Y X| \to |X|$ 是单射，
这立即推出
$|\Delta_{X/Y}| : |X| \to |X \times_Y X|$ 是满射；进而由《代数空间的性质》
引理 \ref{spaces-properties-lemma-characterize-surjective}
可知 $\Delta_{X/Y}$ 是满射。
\end{proof}

\noindent
根据上面的两个引理，以下定义与已经定义的代数空间可表示态射的普遍单射概念不冲突。

\begin{definition}
\label{spaces-morphisms-definition-universally-injective}
设 $S$ 为概形。设 $f : X \to Y$ 为 $S$ 上代数空间的态射。若对每个态射
$Y' \to Y$，诱导映射
$|Y' \times_Y X| \to |Y'|$ 都是单射，则称 $f$ 为{\it 普遍单射}。
\end{definition}

\noindent
为明确起见，这意味着引理 \ref{spaces-morphisms-lemma-universally-injective} 中的任一（或全部）等价条件均成立。

\begin{remark}
\label{spaces-morphisms-remark-universally-injective-not-separated}
概形的普遍单射态射是分离的，见《态射》引理
\ref{morphisms-lemma-universally-injective-separated}。
但代数空间的态射不一定如此。具体地，代数空间
$X = \mathbf{A}^1_k/\{x \sim -x \mid x \not = 0\}$
是在《代数空间》例 \ref{spaces-example-affine-line-involution} 中构造的，
它带有一个态射 $X \to \mathbf{A}^1_k$，将坐标为 $x$ 的点映到坐标为 $x^2$ 的点。
该态射在 $0$ 之外是同构，并且 $X$ 中位于 $0$ 上方的点唯一。由于 $X$ 不是分离的，
这是一个不分离的代数空间普遍单射态射。
\end{remark}

\begin{lemma}
\label{spaces-morphisms-lemma-base-change-universally-injective}
普遍单射态射的基变换仍是普遍单射。
\end{lemma}

\begin{proof}
省略。提示：这是形式性的。
\end{proof}

\begin{lemma}
\label{spaces-morphisms-lemma-universally-injective-local}
设 $S$ 为概形。
设 $f : X \to Y$ 为 $S$ 上代数空间的态射。
以下各条件等价：
\begin{enumerate}
\item $f$ 是普遍单射；
\item 对每个概形 $Z$ 及任意态射 $Z \to Y$，态射
$Z \times_Y X \to Z$ 是普遍单射；
\item 对每个仿射概形 $Z$ 及任意态射
$Z \to Y$，态射 $Z \times_Y X \to Z$ 是普遍单射；
\item 存在概形 $Z$ 及满射态射
$Z \to Y$，使得 $Z \times_Y X \to Z$ 是普遍单射；
\item 存在 Zariski 覆盖 $Y = \bigcup Y_i$，使得
每个态射 $f^{-1}(Y_i) \to Y_i$ 都是普遍单射。
\end{enumerate}
\end{lemma}

\begin{proof}
在本证明中，以下不再另行说明，我们使用普遍单射在基变换下保持这一事实
（引理 \ref{spaces-morphisms-lemma-base-change-universally-injective}）。
显然有 (1) $\Rightarrow$ (2) $\Rightarrow$ (3)
$\Rightarrow$ (4)。

\medskip\noindent
设 $g : Z \to Y$ 如 (4) 所述。设 $y : \Spec(K) \to Y$ 为
从域的谱到 $Y$ 的态射。由假设，可找到扩域
$\alpha : K \subset K'$ 以及态射
$z : \Spec(K') \to Z$，使得 $y \circ \alpha = g \circ z$
（这里对记号作了显然的滥用）。由假设，
态射 $Z \times_Y X \to Z$ 是普遍单射，因此
$g \circ z : \Spec(K') \to Y$ 到 $X$ 的态射提升至多有一个。
由于 $\{\alpha : \Spec(K') \to \Spec(K)\}$ 是
fpqc 覆盖，这意味着 $y : \Spec(K) \to Y$ 到 $X$ 的态射提升至多有一个，见
《代数空间的性质》命题
\ref{spaces-properties-proposition-sheaf-fpqc}。
因此 (1) 成立。

\medskip\noindent
我们省略 (5) 与 (1) 等价性的验证。
\end{proof}

\begin{lemma}
\label{spaces-morphisms-lemma-composition-universally-injective}
普遍单射态射的复合仍是普遍单射。
\end{lemma}

\begin{proof}
省略。
\end{proof}











\section{仿射态射}
\label{spaces-morphisms-section-affine}

\noindent
我们已经在第 \ref{spaces-morphisms-section-representable} 节定义了代数空间的可表示态射
为仿射的含义。

\begin{lemma}
\label{spaces-morphisms-lemma-affine-representable}
设 $S$ 为概形。设 $f : X \to Y$ 为 $S$ 上代数空间的可表示态射。则
$f$ 是仿射的（即第 \ref{spaces-morphisms-section-representable} 节的含义）
当且仅当对所有仿射概形 $Z$ 及态射 $Z \to Y$，概形 $X \times_Y Z$ 是仿射的。
\end{lemma}

\begin{proof}
这直接由概形仿射态射的定义得到（《态射》定义
\ref{morphisms-definition-affine}）。
\end{proof}

\noindent
这为下面的定义铺平了道路。

\begin{definition}
\label{spaces-morphisms-definition-affine}
设 $S$ 为概形。
设 $f : X \to Y$ 为 $S$ 上代数空间的态射。
若对每个仿射概形 $Z$ 及态射 $Z \to Y$，代数空间 $X \times_Y Z$
都可由仿射概形表示，则称 $f$ 为{\it 仿射}。
\end{definition}

\begin{lemma}
\label{spaces-morphisms-lemma-affine-local}
设 $S$ 为概形。
设 $f : X \to Y$ 为 $S$ 上代数空间的态射。
以下各条件等价：
\begin{enumerate}
\item $f$ 可表示且仿射；
\item $f$ 仿射；
\item 对每个仿射概形 $V$ 及 'etale 态射 $V \to Y$，
概形 $X \times_Y V$ 是仿射的；
\item 存在概形 $V$ 及满射 'etale 态射
$V \to Y$，使得 $V \times_Y X \to V$ 是仿射的；
\item 存在 Zariski 覆盖 $Y = \bigcup Y_i$，使得
每个态射 $f^{-1}(Y_i) \to Y_i$ 都是仿射的。
\end{enumerate}
\end{lemma}

\begin{proof}
显然 (1) 推出 (2)，(2) 推出 (3)；取 $V$ 为在 $Y$ 上 'etale 的仿射概形的不交并，
由此 (3) 推出 (4)，见《代数空间的性质》引理
\ref{spaces-properties-lemma-cover-by-union-affines}。
设 $V \to Y$ 如 (4) 所述。对 $V$ 的每个仿射开子概形 $W$，有
$W \times_Y X$ 是 $V \times_Y X$ 的仿射开子概形。因此由《代数空间的性质》引理
\ref{spaces-properties-lemma-subscheme}
可知 $V \times_Y X$ 是概形。此外，态射
$V \times_Y X \to V$ 是仿射的。于是可以应用《代数空间》引理
\ref{spaces-lemma-morphism-sheaves-with-P-effective-descent-etale}，
因为仿射态射类满足所有所需性质（见《态射》引理
\ref{morphisms-lemma-base-change-affine} 以及《下降》引理
\ref{descent-lemma-descending-property-affine} 与
\ref{descent-lemma-affine}）。应用该引理的结论是
$f$ 可表示且仿射，即 (1) 成立。

\medskip\noindent
由于仿射性在目标上是 Zariski 局部的，(1) 与 (5) 等价（上面的参考还表明，
仿射性事实上在目标上是 fpqc 局部的）。
\end{proof}

\begin{lemma}
\label{spaces-morphisms-lemma-composition-affine}
仿射态射的复合是仿射的。
\end{lemma}

\begin{proof}
省略。提示：纤维积的传递性。
\end{proof}

\begin{lemma}
\label{spaces-morphisms-lemma-base-change-affine}
仿射态射的基变换是仿射的。
\end{lemma}

\begin{proof}
省略。提示：纤维积的传递性。
\end{proof}

\begin{lemma}
\label{spaces-morphisms-lemma-closed-immersion-affine}
闭浸入是仿射的。
\end{lemma}

\begin{proof}
立即由概形态射的相应命题得到，见《态射》引理
\ref{morphisms-lemma-closed-immersion-affine}。
\end{proof}

\begin{lemma}
\label{spaces-morphisms-lemma-affine-equivalence-algebras}
设 $S$ 为概形。设 $X$ 为 $S$ 上的代数空间。
存在范畴的反等价
$$
\begin{matrix}
\text{代数空间} \\
\text{仿射于 }X
\end{matrix}
\longleftrightarrow
\begin{matrix}
\text{拟凝聚层} \\
\text{作为 }\mathcal{O}_X\text{-代数}
\end{matrix}
$$
它将 $f : Y \to X$ 对应到层 $f_*\mathcal{O}_Y$。
此外，该等价与任意基变换相容。
\end{lemma}

\begin{proof}
本引理是《态射》引理
\ref{morphisms-lemma-affine-equivalence-algebras}.
设 $\mathcal{A}$ 为拟凝聚的 $\mathcal{O}_X$-代数层。
我们将构造代数空间的仿射态射
 $\pi : Y = \underline{\Spec}_X(\mathcal{A}) \to X$，满足
$\pi_*\mathcal{O}_Y \cong \mathcal{A}$。为此，取一个概形
$U$ 以及满射 \'etale 态射 $\varphi : U \to X$。照例，记
$R = U \times_X U$，其投影为 $s, t : R \to U$。记
$\psi : R \to X$ 为复合 $\psi = \varphi \circ s = \varphi \circ t$。
由上述引理，存在概形的仿射态射
$\pi_0 : V \to U$ 与 $\pi_1 : W \to R$，满足
$\pi_{0, *}\mathcal{O}_V \cong \varphi^*\mathcal{A}$ 以及
$\pi_{1, *}\mathcal{O}_W \cong \psi^*\mathcal{A}$。
由于该构造与基变换相容，存在态射 $s', t' : W \to V$，使得下列图
$$
\vcenter{
\xymatrix{
W \ar[r]_{s'} \ar[d] & V \ar[d] \\
R \ar[r]^s & U
}
}
\quad\text{且}\quad
\vcenter{
\xymatrix{
W \ar[r]_{t'} \ar[d] & V \ar[d] \\
R \ar[r]^t & U
}
}
$$
这些图是笛卡尔的。因此 $s', t'$ 是 \'etale 态射。由上述事实形式地可得
$(t', s') : W \to V \times_S V$ 是单态射。我们省略验证
$W \to V \times_S V$ 是等价关系（提示：考虑将 $\mathcal{A}$
拉回到 $U \times_X U \times_X U = R \times_{s, U, t} R$）。
商层 $Y = V/W$ 是代数空间，见《代数空间》定理
\ref{spaces-theorem-presentation}。
由《群胚》引理 \ref{groupoids-lemma-criterion-fibre-product}
可知 $Y \times_X U \cong V$。因此由引理
\ref{spaces-morphisms-lemma-affine-local}，$Y \to X$ 是仿射的。最后，下面的同构
$$
(Y \times_X U \to U)_*\mathcal{O}_{Y \times_X U} =
\pi_{0, *}\mathcal{O}_V \cong \varphi^*\mathcal{A}
$$
与粘合同构相容，因此由《代数空间的性质》命题
\ref{spaces-properties-proposition-quasi-coherent}
有 $(Y \to X)_*\mathcal{O}_Y \cong \mathcal{A}$。
我们省略验证该构造与基变换相容。
\end{proof}

\begin{definition}
\label{spaces-morphisms-definition-relative-spec}
设 $S$ 为概形。设 $X$ 为 $S$ 上的代数空间。
设 $\mathcal{A}$ 为拟凝聚的
$\mathcal{O}_X$-代数层。$\mathcal{A}$ 在 $X$ 上的{\it 相对谱}，
或简称 $\mathcal{A}$ 在 $X$ 上的{\it 谱}，是仿射态射
$\underline{\Spec}(\mathcal{A}) \to X$，它在引理
\ref{spaces-morphisms-lemma-affine-equivalence-algebras} 的范畴等价下对应于 $\mathcal{A}$。
\end{definition}

\noindent
构造相对谱与任意基变换相容。

\begin{remark}
\label{spaces-morphisms-remark-factorization-quasi-compact-quasi-separated}
设 $S$ 为概形。设 $f : Y \to X$ 为 $S$ 上代数空间的拟紧且
拟分离态射。则 $f$ 有规范分解
$$
Y \longrightarrow \underline{\Spec}_X(f_*\mathcal{O}_Y) \longrightarrow X
$$
这是有意义的，因为由引理 \ref{spaces-morphisms-lemma-pushforward}，
$f_*\mathcal{O}_Y$ 是拟凝聚的。态射
$Y \to \underline{\Spec}_X(f_*\mathcal{O}_Y)$ 来自规范的
$\mathcal{O}_Y$-代数态射
$f^*f_*\mathcal{O}_Y \to \mathcal{O}_Y$，该态射又对应于一个在 $Y$ 上的规范态射
$Y \to Y \times_X \underline{\Spec}_X(f_*\mathcal{O}_Y)$（见
引理 \ref{spaces-morphisms-lemma-affine-equivalence-algebras}），从而得到上面的
$f$ 的分解。
\end{remark}

\begin{lemma}
\label{spaces-morphisms-lemma-affine-equivalence-modules}
设 $S$ 为概形。
设 $f : Y \to X$ 为 $S$ 上代数空间的仿射态射。
设 $\mathcal{A} = f_*\mathcal{O}_Y$。
函子 $\mathcal{F} \mapsto f_*\mathcal{F}$ 诱导出范畴等价
$$
\left\{
\begin{matrix}
\text{拟凝聚}\\
\mathcal{O}_Y\text{-模范畴}
\end{matrix}
\right\}
\longrightarrow
\left\{
\begin{matrix}
\text{拟凝聚}\\
\mathcal{A}\text{-模范畴}
\end{matrix}
\right\}
$$
此外，一个 $\mathcal{A}$-模作为 $\mathcal{O}_X$-模是拟凝聚的，当且仅当
它作为 $\mathcal{A}$-模是拟凝聚的。
\end{lemma}

\begin{proof}
省略。
\end{proof}

\begin{lemma}
\label{spaces-morphisms-lemma-affine-permanence}
设 $S$ 为概形。设 $B$ 为 $S$ 上的代数空间。
假设 $g : X \to Y$ 为 $B$ 上代数空间的态射。
\begin{enumerate}
\item 若 $X$ 在 $B$ 上仿射且 $\Delta : Y \to Y \times_B Y$ 仿射，
则 $g$ 仿射。
\item 若 $X$ 在 $B$ 上仿射且 $Y$ 在 $B$ 上分离，
则 $g$ 仿射。
\item 从仿射概形到对角态射在 $\mathbf{Z}$ 上仿射的代数空间的态射
（如《代数空间的性质》定义 \ref{spaces-properties-definition-separated}）是仿射的。
\item 从仿射概形到分离代数空间的态射是仿射的。
\end{enumerate}
\end{lemma}

\begin{proof}
证明 (1)。由引理 \ref{spaces-morphisms-lemma-base-change-affine}，基变换 $X \times_B Y \to Y$ 是仿射的。
态射 $(1, g) : X \to X \times_B Y$ 是 $Y \to Y \times_B Y$ 沿
$X \times_B Y \to Y \times_B Y$ 的基变换。因此由引理
\ref{spaces-morphisms-lemma-base-change-affine} 它是仿射的。
仿射态射的复合是仿射的（见引理 \ref{spaces-morphisms-lemma-composition-affine}），故 (1) 成立。
由于闭浸入是仿射的（见引理 \ref{spaces-morphisms-lemma-closed-immersion-affine}），且 $Y/B$ 分离
意味着 $\Delta$ 是闭浸入，(2) 由 (1) 得到。(3) 与 (4) 分别是 (1) 与 (2) 的特例。
\end{proof}

\begin{lemma}
\label{spaces-morphisms-lemma-Artinian-affine}
设 $S$ 为概形。设 $X$ 为 $S$ 上拟分离的代数空间。
设 $A$ 为 Artin 环。任意态射 $\Spec(A) \to X$ 都是仿射的。
\end{lemma}

\begin{proof}
设 $U \to X$ 为 'etale 态射且 $U$ 仿射。要证明本引理，需证明
$\Spec(A) \times_X U$ 仿射，见引理 \ref{spaces-morphisms-lemma-affine-local}。
由于 $X$ 拟分离，概形 $\Spec(A) \times_X U$ 拟紧。此外，
投影态射 $\Spec(A) \times_X U \to \Spec(A)$ 是 'etale 的。
因此该态射的纤维是有限离散的，而且 $\Spec(A)$ 上的拓扑是离散的。
于是 $\Spec(A) \times_X U$ 是一个概形，其底层拓扑空间是有限离散集。
由《概形》引理 \ref{schemes-lemma-scheme-finite-discrete-affine} 即得结论。
\end{proof}












\section{准仿射态射}
\label{spaces-morphisms-section-quasi-affine}

\noindent
我们已经在第 \ref{spaces-morphisms-section-representable} 节定义了代数空间的可表示态射
为准仿射的含义。

\begin{lemma}
\label{spaces-morphisms-lemma-quasi-affine-representable}
设 $S$ 为概形。设 $f : X \to Y$ 为 $S$ 上代数空间的可表示
态射。则 $f$ 是准仿射的（即第 \ref{spaces-morphisms-section-representable} 节的含义）
当且仅当对所有仿射概形 $Z$ 及态射 $Z \to Y$，概形 $X \times_Y Z$ 是准仿射的。
\end{lemma}

\begin{proof}
这直接由概形准仿射态射的定义得到
（《态射》定义 \ref{morphisms-definition-quasi-affine}）。
\end{proof}

\noindent
这为下面的定义铺平了道路。

\begin{definition}
\label{spaces-morphisms-definition-quasi-affine}
设 $S$ 为概形。
设 $f : X \to Y$ 为 $S$ 上代数空间的态射。
若对每个仿射概形 $Z$ 及态射 $Z \to Y$，代数空间 $X \times_Y Z$
都可由准仿射概形表示，则称 $f$ 为{\it 准仿射}。
\end{definition}

\begin{lemma}
\label{spaces-morphisms-lemma-quasi-affine-local}
设 $S$ 为概形。
设 $f : X \to Y$ 为 $S$ 上代数空间的态射。
以下各条件等价：
\begin{enumerate}
\item $f$ 可表示且准仿射；
\item $f$ 准仿射；
\item 存在概形 $V$ 及满射 \'etale 态射
$V \to Y$，使得 $V \times_Y X \to V$ 是准仿射的；
\item 存在 Zariski 覆盖 $Y = \bigcup Y_i$，使得
每个态射 $f^{-1}(Y_i) \to Y_i$ 都是准仿射的。
\end{enumerate}
\end{lemma}

\begin{proof}
显然 (1) 推出 (2)；取 $V$ 为在 $Y$ 上 'etale 的仿射概形的不交并，
由《代数空间的性质》引理
\ref{spaces-properties-lemma-cover-by-union-affines}，(2) 推出 (3)。
设 $V \to Y$ 如 (3) 所述。对 $V$ 的每个仿射开子概形 $W$，有
$W \times_Y X$ 是 $V \times_Y X$ 的准仿射开子概形。因此由《代数空间的性质》引理
\ref{spaces-properties-lemma-subscheme}
可知 $V \times_Y X$ 是概形。此外态射
$V \times_Y X \to V$ 是准仿射的。于是可以应用《代数空间》引理
\ref{spaces-lemma-morphism-sheaves-with-P-effective-descent-etale}，
因为准仿射态射类满足所有所需性质（见《态射》引理
\ref{morphisms-lemma-base-change-quasi-affine} 以及《下降》引理
\ref{descent-lemma-descending-property-quasi-affine} 与
\ref{descent-lemma-quasi-affine}）。应用该引理的结论是
$f$ 可表示且准仿射，即 (1) 成立。

\medskip\noindent
由于准仿射性在目标上是 Zariski 局部的，(1) 与 (4) 等价（上面的参考还表明，
准仿射性事实上在目标上是 fpqc 局部的）。
\end{proof}

\begin{lemma}
\label{spaces-morphisms-lemma-composition-quasi-affine}
准仿射态射的复合是准仿射的。
\end{lemma}

\begin{proof}
省略。
\end{proof}

\begin{lemma}
\label{spaces-morphisms-lemma-base-change-quasi-affine}
准仿射态射的基变换是准仿射的。
\end{lemma}

\begin{proof}
省略。
\end{proof}

\begin{lemma}
\label{spaces-morphisms-lemma-characterize-quasi-affine}
设 $S$ 为概形。
代数空间的拟紧且拟分离态射
$f : Y \to X$ 是准仿射的，当且仅当其规范分解
$Y \to \underline{\Spec}_X(f_*\mathcal{O}_Y)$
（注记 \ref{spaces-morphisms-remark-factorization-quasi-compact-quasi-separated}）
是开浸入。
\end{lemma}

\begin{proof}
设 $U \to X$ 为满射态射，其中 $U$ 是概形。由于可以在沿 $U$ 的基变换后检验
$f$ 是否准仿射（引理 \ref{spaces-morphisms-lemma-quasi-affine-local}），且
$f_*\mathcal{O}_Y|_V$ 等于 $(Y \times_X U \to U)_*\mathcal{O}_{Y \times_X U}$
（《代数空间的性质》引理
\ref{spaces-properties-lemma-pushforward-etale-base-change-modules}），并且
相对谱的构造与基变换相容（引理 \ref{spaces-morphisms-lemma-affine-equivalence-algebras}），
可知断言归结为 $X$ 为概形的情形。
若 $X$ 为概形且 $f$ 准仿射，或
$Y \to \underline{\Spec}_X(f_*\mathcal{O}_Y)$ 为开浸入，
则 $Y$ 也是概形。因此问题归结为《态射》引理
\ref{morphisms-lemma-characterize-quasi-affine}。
\end{proof}







\section{源与目标上 \'etale 局部的态射类型}
\label{spaces-morphisms-section-local-source-target}

\noindent
给定一种概形态射的性质，它在源与目标上是{\it \'etale 局部的}（见
《下降》定义 \ref{descent-definition-local-source-target}），
我们可以用它定义代数空间态射的相应性质，即施加下面引理中的任一等价条件。

\begin{lemma}
\label{spaces-morphisms-lemma-local-source-target}
设 $\mathcal{P}$ 为概形态射的一种在源与目标上 \'etale 局部的性质。
设 $S$ 为概形。
设 $f : X \to Y$ 为 $S$ 上代数空间的态射。
考虑交换图
$$
\xymatrix{
U \ar[d]_a \ar[r]_h & V \ar[d]^b \\
X \ar[r]^f & Y
}
$$
其中 $U$ 与 $V$ 为概形，竖直箭头为 \'etale 态射。
以下条件等价：
\begin{enumerate}
\item 对任意上述图，态射 $h$ 具有性质
$\mathcal{P}$；
\item 对某个上述图（其中 $a : U \to X$ 为满射），
态射 $h$ 具有性质 $\mathcal{P}$。
\end{enumerate}
若 $X$ 与 $Y$ 可表示，这还等价于 $f$（作为概形态射）具有性质 $\mathcal{P}$。
若 $\mathcal{P}$ 还在任意基变换下保持，且在基上 fppf 局部，
那么对于可表示态射 $f$，这也等价于 $f$ 按第 \ref{spaces-morphisms-section-representable} 节的含义
具有性质 $\mathcal{P}$。
\end{lemma}

\begin{proof}
我们证明 (1) 与 (2) 的等价性。
(1) $\Rightarrow$ (2) 立即成立（注意《代数空间》引理
\ref{spaces-lemma-lift-morphism-presentations}）。
设
$$
\xymatrix{
U \ar[d] \ar[r]_h & V \ar[d] \\
X \ar[r]^f & Y
}
\quad\quad
\xymatrix{
U' \ar[d] \ar[r]_{h'} & V' \ar[d] \\
X \ar[r]^f & Y
}
$$
为引理中的两个图。假设 $U \to X$ 满射且 $h$ 具有性质 $\mathcal{P}$。
为证明 (2) 推出 (1)，需证明 $h'$ 具有 $\mathcal{P}$。为此考虑图
$$
\xymatrix{
U \ar[d]_h &
U \times_X U' \ar[l] \ar[d]^{(h, h')} \ar[r] &
U' \ar[d]^{h'} \\
V &
V \times_Y V' \ar[l] \ar[r] &
V'
}
$$
由《下降》引理 \ref{descent-lemma-local-source-target-characterize}，
$h$ 具有 $\mathcal{P}$ 推出 $(h, h')$ 具有 $\mathcal{P}$；又由于
$U \times_X U' \to U'$ 是满射，由同一引理可知 $h'$ 具有 $\mathcal{P}$。

\medskip\noindent
若 $X$ 与 $Y$ 可表示，则可应用《下降》引理
\ref{descent-lemma-local-source-target-characterize}，从而 (1) 与 (2) 等价于 $f$
具有 $\mathcal{P}$。

\medskip\noindent
最后，设 $f$ 可表示，且 $U,V,a,b,h$ 如引理 (2) 所述，并假设
$\mathcal{P}$ 在任意基变换下保持。我们需证明对任意概形
$Z$ 及态射 $Z \to X$，基变换 $Z \times_Y X \to Z$
具有性质 $\mathcal{P}$。考虑图
$$
\xymatrix{
Z \times_Y U \ar[d] \ar[r] &
Z \times_Y V \ar[d] \\
Z \times_Y X \ar[r] &
Z
}
$$
注意上方水平箭头是 $h$ 的基变换，故具有性质 $\mathcal{P}$。
左侧竖直箭头是满射 \'etale，右侧竖直箭头是 \'etale。因此再次应用《下降》引理
\ref{descent-lemma-local-source-target-characterize}，可知 $Z \times_Y X \to Z$
具有性质 $\mathcal{P}$。
\end{proof}

\begin{definition}
\label{spaces-morphisms-definition-P}
设 $S$ 为概形。
设 $\mathcal{P}$ 为概形态射的一种在源与目标上 \'etale 局部的性质。
若引理 \ref{spaces-morphisms-lemma-local-source-target} 的等价条件成立，
则称 $S$ 上代数空间的态射 $f : X \to Y${\it 具有性质 $\mathcal{P}$}。
\end{definition}

\noindent
下面是几个显然的注记。

\begin{remark}
\label{spaces-morphisms-remark-composition-P}
设 $S$ 为概形。设 $\mathcal{P}$ 为概形态射的一种在源与目标上 \'etale 局部的性质，
并且假设 $\mathcal{P}$ 在复合下稳定。则具有性质 $\mathcal{P}$ 的代数空间态射类
在复合下稳定。
\end{remark}

\begin{remark}
\label{spaces-morphisms-remark-base-change-P}
设 $S$ 为概形。设 $\mathcal{P}$ 为概形态射的一种在源与目标上 \'etale 局部的性质，
并且假设 $\mathcal{P}$ 在基变换下稳定。则具有性质 $\mathcal{P}$ 的代数空间态射类
在基变换下稳定。
\end{remark}

\noindent
给定一种概形芽态射的性质，它在源与目标上是{\it \'etale 局部的}（见
《下降》定义 \ref{descent-definition-local-source-target-at-point}），
我们可以用它定义代数空间在一点处的态射性质，即施加下面引理中的任一等价条件。

\begin{lemma}
\label{spaces-morphisms-lemma-local-source-target-at-point}
设 $\mathcal{Q}$ 为概形芽态射的一种在源与目标上 \'etale 局部的性质。
设 $S$ 为概形。
设 $f : X \to Y$ 为 $S$ 上代数空间的态射。
设 $x \in |X|$ 为 $X$ 的一点。
考虑图
$$
\xymatrix{
U \ar[d]_a \ar[r]_h & V \ar[d]^b \\
X \ar[r]^f & Y
}
\quad\quad
\xymatrix{
u \ar[d] \ar[r] & v \ar[d] \\
x \ar[r] & y
}
$$
其中 $U$ 与 $V$ 为概形，$a,b$ 为 \'etale 态射，且 $u,v,x,y$
分别是相应空间的点。
以下条件等价：
\begin{enumerate}
\item 对任意上述图，有 $\mathcal{Q}((U, u) \to (V, v))$；
\item 对某个上述图，有 $\mathcal{Q}((U, u) \to (V, v))$。
\end{enumerate}
若 $X$ 与 $Y$ 可表示，这还等价于 $\mathcal{Q}((X, x) \to (Y, y))$。
\end{lemma}

\begin{proof}
省略。提示：与引理 \ref{spaces-morphisms-lemma-local-source-target} 的证明非常相似。
\end{proof}

\begin{definition}
\label{spaces-morphisms-definition-P-at-point}
设 $\mathcal{Q}$ 为概形芽态射的一种在源与目标上 \'etale 局部的性质。
设 $S$ 为概形。给定 $S$ 上代数空间的态射 $f : X \to Y$ 及一点
$x \in |X|$，若引理 \ref{spaces-morphisms-lemma-local-source-target-at-point} 的等价条件成立，
则称 $f${\it 在 $x$ 处具有性质 $\mathcal{Q}$}。
\end{definition}

\noindent
下面的引理不应被盲目地用来把态射的性质转化为一点处态射的性质。
例如，若 $\mathcal{P}$ 是平坦性，则下面引理中的性质
$Q$ 意味着“$f$ 在 $x$ 的某个开邻域内平坦”，这不同于“$f$ 在 $x$ 处平坦”。

\begin{lemma}
\label{spaces-morphisms-lemma-local-source-target-global-implies-local}
设 $\mathcal{P}$ 为概形态射的一种在源与目标上 \'etale 局部的性质。
考虑《下降》引理
\ref{descent-lemma-local-source-target-global-implies-local} 中与 $\mathcal{P}$
关联的芽态射性质 $\mathcal{Q}$。则
\begin{enumerate}
\item $\mathcal{Q}$ 在源与目标上 \'etale 局部。
\item 给定代数空间态射 $f : X \to Y$ 及 $x \in |X|$，以下条件等价：
\begin{enumerate}
\item $f$ 在 $x$ 处具有 $\mathcal{Q}$；
\item 存在 $x$ 的开邻域 $X' \subset X$，使得 $X' \to Y$ 具有 $\mathcal{P}$。
\end{enumerate}
\item 给定代数空间态射 $f : X \to Y$，以下条件等价：
\begin{enumerate}
\item $f$ 具有 $\mathcal{P}$；
\item 对每个 $x \in |X|$，态射 $f$ 在 $x$ 处具有 $\mathcal{Q}$。
\end{enumerate}
\end{enumerate}
\end{lemma}

\begin{proof}
关于 (1) 见《下降》引理
\ref{descent-lemma-local-source-target-global-implies-local}。(1)(a) $\Rightarrow$ (2)(b)
是因为对引理 \ref{spaces-morphisms-lemma-local-source-target-at-point} 中的图取
$X' = a(U) \subset X$。 (2)(b) $\Rightarrow$ (1)(a) 显然。
(3)(a) 与 (3)(b) 的等价性来自概形态射的相应结果，见《下降》引理
\ref{descent-lemma-local-source-target-local-implies-global}。
\end{proof}

\begin{remark}
\label{spaces-morphisms-remark-when-apply}
我们将把上面的引理 \ref{spaces-morphisms-lemma-local-source-target-global-implies-local}
应用于《下降》注记 \ref{descent-remark-list-local-source-target} 列出的所有情形，
但不包括“平坦”。在每种情形中，我们定义：若 $f$ 在 $x$ 的某个邻域内具有
$\mathcal{P}$，则称 $f$ 在 $x$ 处具有性质 $\mathcal{P}$。
\end{remark}




\section{有限型态射}
\label{spaces-morphisms-section-finite-type}

\noindent
概形态射的“局部有限型”性质在源与目标上 \'etale 局部，见
《下降》注记 \ref{descent-remark-list-local-source-target}。
它还在基变换下保持，并在目标上 fpqc 局部，见《态射》引理
\ref{morphisms-lemma-base-change-finite-type} 以及《下降》引理
\ref{descent-lemma-descending-property-locally-finite-type}。
因此由上面的引理 \ref{spaces-morphisms-lemma-local-source-target}，我们可以如下定义代数空间态射
为局部有限型；当态射可表示时，这与第 \ref{spaces-morphisms-section-representable} 节已有的定义一致。

\begin{definition}
\label{spaces-morphisms-definition-locally-finite-type}
设 $S$ 为概形。
设 $f : X \to Y$ 为 $S$ 上代数空间的态射。
\begin{enumerate}
\item 若引理 \ref{spaces-morphisms-lemma-local-source-target} 的等价条件在
$\mathcal{P} = \text{局部有限型}$ 下成立，则称 $f$
{\it 局部有限型}。
\item 设 $x \in |X|$。若存在 $x$ 的开邻域 $X' \subset X$，使得
$f|_{X'} : X' \to Y$ 是局部有限型，则称 $f${\it 在 $x$ 处有限型}。
\item 若 $f$ 局部有限型且拟紧，则称{\it 为有限型}。
\end{enumerate}
\end{definition}

\noindent
考虑《代数空间》例 \ref{spaces-example-affine-line-translation} 中的代数空间
$\mathbf{A}^1_k/\mathbf{Z}$。
态射 $\mathbf{A}^1_k/\mathbf{Z} \to \Spec(k)$ 是有限型的。

\begin{lemma}
\label{spaces-morphisms-lemma-composition-finite-type}
有限型态射的复合是有限型的。
局部有限型也有同样结论。
\end{lemma}

\begin{proof}
见注记 \ref{spaces-morphisms-remark-composition-P} 及《态射》引理
\ref{morphisms-lemma-composition-finite-type}。
\end{proof}

\begin{lemma}
\label{spaces-morphisms-lemma-base-change-finite-type}
有限型态射的基变换是有限型的。
局部有限型也有同样结论。
\end{lemma}

\begin{proof}
见注记 \ref{spaces-morphisms-remark-base-change-P} 及《态射》引理
\ref{morphisms-lemma-base-change-finite-type}。
\end{proof}

\begin{lemma}
\label{spaces-morphisms-lemma-finite-type-local}
设 $S$ 为概形。
设 $f : X \to Y$ 为 $S$ 上代数空间的态射。
以下各条件等价：
\begin{enumerate}
\item $f$ 局部有限型；
\item 对每个 $x \in |X|$，态射 $f$ 在 $x$ 处有限型；
\item 对每个概形 $Z$ 及任意态射 $Z \to Y$，态射
$Z \times_Y X \to Z$ 局部有限型；
\item 对每个仿射概形 $Z$ 及任意态射
$Z \to Y$，态射 $Z \times_Y X \to Z$ 局部有限型；
\item 存在概形 $V$ 及满射 \'etale 态射
$V \to Y$，使得 $V \times_Y X \to V$ 局部有限型；
\item 存在概形 $U$ 及满射 \'etale 态射
$\varphi : U \to X$，使得复合 $f \circ \varphi$
局部有限型；
\item 对每个交换图
$$
\xymatrix{
U \ar[d] \ar[r] & V \ar[d] \\
X \ar[r] & Y
}
$$
其中 $U$ 与 $V$ 为概形且竖直箭头为 \'etale，
上方水平箭头局部有限型；
\item 存在交换图
$$
\xymatrix{
U \ar[d] \ar[r] & V \ar[d] \\
X \ar[r] & Y
}
$$
其中 $U$ 与 $V$ 为概形，竖直箭头为 \'etale，$U \to X$
为满射，且上方水平箭头局部有限型；
\item 存在 Zariski 覆盖 $Y = \bigcup_{i \in I} Y_i$，
以及 $f^{-1}(Y_i) = \bigcup X_{ij}$，使得
每个态射 $X_{ij} \to Y_i$ 都局部有限型。
\end{enumerate}
\end{lemma}

\begin{proof}
条件 (2)、(3)、(4)、(5)、(6)、(7) 与 (9) 均以直接方式推出条件 (8)。例如，若 (5) 成立，
则可取概形 $V$ 为仿射概形的不交并，并取满射态射 $V \to Y$
（见《代数空间的性质》引理
\ref{spaces-properties-lemma-cover-by-union-affines}）。
由 (5)，$V \times_Y X \to V$ 局部有限型。
取概形 $U$ 及满射 \'etale 态射
$U \to V \times_Y X$。由引理 \ref{spaces-morphisms-lemma-composition-finite-type}，
$U \to V$ 局部有限型。因此 (8) 成立。

\medskip\noindent
条件 (1)、(7) 与 (8) 按定义等价。

\medskip\noindent
最后，我们证明 (1) 推出 (2)、(3)、(4)、(5)、(6) 与 (9)。(2) 立即成立。
(3)、(4)、(5) 与 (9) 来自局部有限型在基变换下保持这一事实，见
引理 \ref{spaces-morphisms-lemma-base-change-finite-type}。
对于 (6)，取 $U = X$ 即可。
\end{proof}

\begin{lemma}
\label{spaces-morphisms-lemma-locally-finite-type-locally-noetherian}
设 $S$ 为概形。
设 $f : X \to Y$ 为 $S$ 上代数空间的态射。
若 $f$ 局部有限型且 $Y$ 局部 Noether，则 $X$ 局部 Noether。
\end{lemma}

\begin{proof}
设有
$$
\xymatrix{
U \ar[d] \ar[r] & V \ar[d] \\
X \ar[r] & Y
}
$$
为交换图，其中 $U$ 与 $V$ 为概形且竖直箭头为满射 \'etale。若 $f$ 局部有限型，则
$U \to V$ 局部有限型。若 $Y$ 局部 Noether，则 $V$ 局部 Noether。由《态射》引理
\ref{morphisms-lemma-finite-type-noetherian}
可知 $U$ 局部 Noether，这意味着 $X$ 局部 Noether。
\end{proof}

\begin{lemma}
\label{spaces-morphisms-lemma-permanence-finite-type}
设 $S$ 为概形。
设 $f : X \to Y$、$g : Y \to Z$ 为 $S$ 上代数空间的态射。
若 $g \circ f : X \to Z$ 局部有限型，则 $f : X \to Y$
局部有限型。
\end{lemma}

\begin{proof}
可取交换图
$$
\xymatrix{
U \ar[r] \ar[d] & V \ar[r] \ar[d] & W \ar[d] \\
X \ar[r] & Y \ar[r] & Z
}
$$
其中 $U$、$V$、$W$ 为概形，竖直箭头为满射 \'etale，见
《代数空间》引理 \ref{spaces-lemma-lift-morphism-presentations}。
此时可应用引理 \ref{spaces-morphisms-lemma-finite-type-local}
以及《态射》引理 \ref{morphisms-lemma-permanence-finite-type} 得出结论。
\end{proof}

\begin{lemma}
\label{spaces-morphisms-lemma-immersion-locally-finite-type}
浸入是局部有限型的。
\end{lemma}

\begin{proof}
由一般原则《代数空间》引理
\ref{spaces-lemma-representable-transformations-property-implication}
以及《态射》引理 \ref{morphisms-lemma-immersion-locally-finite-type} 得到。
\end{proof}







\section{点与几何点}
\label{spaces-morphisms-section-points-fields}

\noindent
本节对点与几何点作一些注记（见《代数空间的性质》定义
\ref{spaces-properties-definition-geometric-point}）。
理解 $X$ 的几何点的一种方式是：
考虑 $S$ 的几何点 $\overline{s} : \Spec(k) \to S$
以及将 $\overline{s}$ 提升为到 $X$ 的态射 $\overline{x}$。
相应图为
$$
\xymatrix{
\Spec(k) \ar[r]_-{\overline{x}} \ar[rd]_{\overline{s}} & X \ar[d] \\
& S.
}
$$
我们常说“设 $k$ 为 $S$ 上的代数闭域”，表示 $\Spec(k)$ 带有态射
$\Spec(k) \to S$。此时记
$$
X(k) = \Mor_S(\Spec(k), X) =
\{\overline{x} \in X\text{ 位于 }\overline{s}\}
$$
这是 $X$ 的 $k$-值点集。在这种情况下，映射
$X(k) \to |X|$ 的像包含于子集 $|X_s| \subset |X|$ 中。
这里 $X_s = \Spec(\kappa(s)) \times_S X$，其中 $s \in S$
是与 $\overline{s}$ 对应的点。由于
$\Spec(\kappa(s)) \to S$ 是单态射，其基变换
$X_s \to X$ 是单态射，且 $|X_s|$ 确实是 $|X|$ 的子集。

\begin{lemma}
\label{spaces-morphisms-lemma-locally-finite-type-surjective-geometric-points}
设 $S$ 为概形。设 $f : X \to Y$ 为代数空间的态射。
假设 $f$ 在 $S$ 上局部有限型。以下条件等价：
\begin{enumerate}
\item $f$ 是满射；
\item 对每个 $S$ 上的代数闭域 $k$，诱导映射
$X(k) \to Y(k)$ 是满射。
\end{enumerate}
\end{lemma}

\begin{proof}
取交换图
$$
\xymatrix{
U \ar[d] \ar[r] & V \ar[d] \\
X \ar[r] & Y
}
$$
其中 $U$、$V$ 为 $S$ 上的概形，竖直箭头为满射且 \'etale，见
《代数空间》引理 \ref{spaces-lemma-lift-morphism-presentations}。
由于 $f$ 局部有限型，$U \to V$ 局部有限型。

\medskip\noindent
假设 (1)，并设 $\overline{y} \in Y(k)$。由引理
\ref{spaces-morphisms-lemma-composition-surjective} 与 \ref{spaces-morphisms-lemma-composition-finite-type}，
$U \to Y$ 是满射且局部有限型。
令 $Z = U \times_{Y, \overline{y}} \Spec(k)$。这是一个概形。
由引理 \ref{spaces-morphisms-lemma-base-change-surjective} 与
\ref{spaces-morphisms-lemma-base-change-finite-type}，投影 $Z \to \Spec(k)$
是满射且局部有限型。
由《簇》引理 \ref{varieties-lemma-locally-finite-type-Jacobson}，
$Z$ 有一个 $k$-值点 $\overline{z}$。点 $\overline{z}$ 的像
$\overline{x} \in X(k)$ 映到 $\overline{y}$，即得所需结论。

\medskip\noindent
假设 (2)。由《代数空间的性质》引理
\ref{spaces-properties-lemma-characterize-surjective}，只需证明 $|X| \to |Y|$ 满射。
设 $y \in |Y|$。取 $u \in U$，其映到 $y$。
取 $k \supset \kappa(u)$ 为代数闭包。记相应的点为
$\overline{u} \in U(k)$，其像为 $\overline{y} \in Y(k)$。
由假设，存在映到 $\overline{y}$ 的 $\overline{x} \in X(k)$。
显然，$\overline{x}$ 的像 $x \in |X|$ 映到 $y$。
\end{proof}

\noindent
为表述下一个引理，引入如下记号。对概形 $T$，记
$$
\lambda(T) = \sup\{\aleph_0, |\kappa(t)| ; t \in T\}.
$$
换言之，$\lambda(T)$ 是界住 $T$ 的所有剩余域基数的最小无限基数。
注意，若 $R$ 是环，则 $R$ 的任意剩余域 $\kappa(\mathfrak p)$
的基数都被 $R$ 的基数所界住（细节略）。
这意味着 $\lambda(T) \leq \text{size}(T)$，其中
$\text{size}(T)$ 是《集合》一节 \ref{sets-section-categories-schemes} 中引入的概形 $T$ 的大小。
若 $L/K$ 是有限生成域扩张，则
$|K| \leq |L| \leq \max\{\aleph_0, |K|\}$。因此，若 $T' \to T$
是局部有限型的概形态射，则 $\lambda(T') \leq \lambda(T)$；若 $T' \to T$ 还
是满射，则等号成立。接下来，设 $S$ 为概形，$X$ 为 $S$ 上的代数空间。此时定义
$$
\lambda(X) := \lambda(U)
$$
其中 $U$ 是任意一个在 $S$ 上、带有指向 $X$ 的满射 \'etale 态射的概形。
该定义与 $U$ 的选择无关，因为给定两个这样的概形 $U$ 与 $U'$，纤维积
$U \times_X U'$ 是概形，并且有分别指向 $U$ 与 $U'$ 的满射 \'etale 态射；
由上面的讨论，$\lambda(U) = \lambda(U \times_X U') =
\lambda(U')$。

\begin{lemma}
\label{spaces-morphisms-lemma-large-enough}
设 $S$ 为概形。设 $X$、$Y$ 为 $S$ 上的代数空间。
\begin{enumerate}
\item 当 $k$ 遍历所有 $S$ 上的代数闭域时，几何点
$\overline{y} \in Y(k)$ 的集合覆盖 $|Y|$。
\item 当 $k$ 遍历所有满足
$|k| \geq \lambda(Y)$ 且 $|k| > \lambda(X)$ 的 $S$ 上代数闭域时，几何点
$\overline{y} \in Y(k)$ 的集合覆盖 $|Y|$。
\item 对任意几何点
$\overline{s} : \Spec(k) \to S$，若
$k$ 的基数 $> \lambda(X)$，则映射
$$
X(k) \longrightarrow |X_s|
$$
是满射。
\item 设 $X \to Y$ 为 $S$ 上代数空间的态射。
对任意几何点 $\overline{s} : \Spec(k) \to S$，若
$k$ 的基数 $> \lambda(X)$，则映射
$$
X(k) \longrightarrow |X| \times_{|Y|} Y(k)
$$
是满射。
\item 设 $X \to Y$ 为 $S$ 上代数空间的态射。
以下条件等价：
\begin{enumerate}
\item 映射 $X \to Y$ 是满射；
\item 对所有满足 $|k| > \lambda(X)$ 且 $|k| \geq \lambda(Y)$ 的
$S$ 上代数闭域 $k$，映射 $X(k) \to Y(k)$ 是满射。
\end{enumerate}
\end{enumerate}
\end{lemma}

\begin{proof}
为证明 (1)，取满射 \'etale 态射 $V \to Y$，其中 $V$ 为概形。
对每个 $v \in V$，取代数闭包 $\kappa(v) \subset k_v$。考虑态射
$\overline{x} : \Spec(k_v) \to V \to Y$。
由 $|Y|$ 的构造，这些态射覆盖 $|Y|$。

\medskip\noindent
为证明 (2)，使用以下两个事实（证明略）：
(\romannumeral1) 若 $K$ 为域且 $\overline{K}$ 为其代数闭包，则
$|\overline{K}| \leq \max\{\aleph_0, |K|\}$。
(\romannumeral2) 对任意代数闭域 $k$ 及基数 $\aleph$，若
$\aleph \geq |k|$，则存在代数闭域扩张 $k'/k$，满足 $|k'| = \aleph$。
现在令 $\aleph = \max\{\lambda(X), \lambda(Y)\}^+$。
这里 $\lambda^+ > \lambda$ 表示下一个更大的基数，见《集合》一节
\ref{sets-section-cardinals}。
由 (\romannumeral1)，证明第一段中构造的域 $k_u$ 的基数都被
$\lambda(X)$ 所界住。因此由 (\romannumeral2)，可找到扩张
$k_u \subset k'_u$，使得 $|k'_u| = \aleph$。态射
$\overline{x}' : \Spec(k'_u) \to X$ 覆盖 $|X|$，如所需。
要真正完成 (2) 的证明，还需说明概形 $\Spec(k'_u)$
（同构于）$\Sch_{fppf}$ 的对象，因为我们的约定是所有概形都是
$\Sch_{fppf}$ 的对象；对集合论问题不感兴趣的读者可跳过本段其余部分。
由构造，存在 $\Sch_{fppf}$ 的对象 $T$，使得
$\lambda(X)$ 与 $\lambda(Y)$ 被 $\text{size}(T)$ 所界住。
按照我们在《拓扑》定义 \ref{topologies-definition-big-fppf-site} 中
将 $\Sch_{fppf}$ 构造为《集合》引理 \ref{sets-lemma-construct-category}
中构造的范畴 $\Sch_\alpha$ 的方式，大小满足 $\leq \text{size}(T)^+$ 的任意概形
都同构于 $\Sch_{fppf}$ 的对象。见《集合》方程
(\ref{sets-equation-bound}) 中函数 $Bound$ 的表达式。
由于 $\aleph \leq \text{size}(T)^+$，结论成立。

\medskip\noindent
第 (3) 部分中的记号 $X_s$ 表示纤维积
$\Spec(\kappa(s)) \times_S X$，其中 $s \in S$
是与 $\overline{s}$ 对应的点。因此 (2) 由 (4) 在
$Y = \Spec(\kappa(s))$ 时推出。

\medskip\noindent
证明 (4)。设 $X \to Y$ 为 $S$ 上代数空间的态射。
设 $k$ 为 $S$ 上基数满足
$> \lambda(X)$ 的代数闭域。设 $\overline{y} \in Y(k)$ 与 $x \in |X|$
映到 $|Y|$ 中的同一点 $y$。
我们要找一个 $\overline{x} \in X(k)$，使其映到 $x$ 与 $\overline{y}$。
取交换图
$$
\xymatrix{
U \ar[d] \ar[r] & V \ar[d] \\
X \ar[r] & Y
}
$$
其中 $U$、$V$ 为 $S$ 上的概形，竖直箭头为满射且 \'etale，见
《代数空间》引理 \ref{spaces-lemma-lift-morphism-presentations}。
取 $u \in |U|$ 映到 $x$，并记其像为 $v \in |V|$。
我们把 $u = \Spec(\kappa(u))$ 与
$v = \Spec(\kappa(v))$ 看作概形。
注意 $V \times_Y \Spec(k)$ 是在 $k$ 上 \'etale 的概形。
因此它是 $k$ 的有限可分扩张的谱的不交并，见
《态射》引理 \ref{morphisms-lemma-etale-over-field}。
由于 $v$ 映到 $y$，$v \times_Y \Spec(k)$ 是非空概形。
由于 $v \to V$ 是单态射，$v \times_Y \Spec(k) \to V \times_Y \Spec(k)$
也是单态射。因此由《概形》引理
\ref{schemes-lemma-mono-towards-spec-field}，$v \times_Y \Spec(k)$ 是
$k$ 的有限可分扩张的谱的不交并。
于是态射 $v \times_Y \Spec(k) \to \Spec(k)$ 有截面；也就是说，存在态射
$\overline{v} : \Spec(k) \to V$，它位于 $v$ 之上并位于
$\overline{y}$ 之上。最后考虑概形
$$
u \times_{V, \overline{v}} \Spec(k)
=
\Spec(\kappa(u) \otimes_{\kappa(v)} k)
$$
其中 $\kappa(v) \to k$ 是定义态射 $\overline{v}$ 的域映射。
由假设，$k$ 的基数大于 $\kappa(u)$ 的基数，因此可应用
《代数》引理 \ref{algebra-lemma-base-change-Jacobson}，
得到任意极大理想
$\mathfrak m \subset \kappa(u) \otimes_{\kappa(v)} k$
的剩余域都代数于 $k$，因而等于 $k$。
这样的极大理想给出态射
$\overline{u} : \Spec(k) \to U$，它位于 $u$ 之上并映到
$\overline{v}$。复合 $\Spec(k) \to U \to X$
就是所需的几何点 $\overline{x} \in X(k)$。
(4) 的证明完成。

\medskip\noindent
(5) 是 (2) 与 (4) 以及《代数空间的性质》引理
\ref{spaces-properties-lemma-characterize-surjective} 的形式推论。
\end{proof}








\section{有限型点}
\label{spaces-morphisms-section-points-finite-type}

\noindent
设 $S$ 为概形。设 $X$ 为 $S$ 上的代数空间。
有限型点 $x \in |X|$ 是可由局部有限型态射
$\Spec(k) \to X$ 表示的点。
对代数空间与代数栈而言，有限型点是闭点的适当替代。例如，总有“足够多”的有限型点。

\begin{lemma}
\label{spaces-morphisms-lemma-point-finite-type}
设 $S$ 为概形。设 $X$ 为 $S$ 上的代数空间。
设 $x \in |X|$。以下条件等价：
\begin{enumerate}
\item 存在局部有限型态射 $\Spec(k) \to X$，它表示 $x$。
\item 存在概形 $U$、闭点 $u \in U$ 以及 \'etale 态射
$\varphi : U \to X$，使得 $\varphi(u) = x$。
\end{enumerate}
\end{lemma}

\begin{proof}
设 $u \in U$ 与 $U \to X$ 如 (2) 所述。则
$\Spec(\kappa(u)) \to U$ 是有限型的，而 $U \to X$
可表示且局部有限型（由一般原则《代数空间》引理
\ref{spaces-lemma-representable-transformations-property-implication}
以及《态射》引理 \ref{morphisms-lemma-etale-locally-finite-presentation} 与
\ref{morphisms-lemma-finite-presentation-finite-type}）。因此由引理
\ref{spaces-morphisms-lemma-composition-finite-type}，(1) 成立。

\medskip\noindent
反之，假设 $\Spec(k) \to X$ 局部有限型并表示 $x$。
设 $U \to X$ 为满射 \'etale 态射，其中 $U$ 为概形。由假设，
$U \times_X \Spec(k) \to U$ 局部有限型。取
$U \times_X \Spec(k)$ 的有限型点 $v$（至少存在一个，见《态射》引理
\ref{morphisms-lemma-identify-finite-type-points}）。
由《态射》引理 \ref{morphisms-lemma-finite-type-points-morphism}，
$v$ 的像 $u \in U$ 是 $U$ 的有限型点。
因此由《态射》引理 \ref{morphisms-lemma-identify-finite-type-points}，
缩小 $U$ 后可假设 $u$ 是 $U$ 的闭点，即 (2) 成立。
\end{proof}

\begin{definition}
\label{spaces-morphisms-definition-finite-type-point}
设 $S$ 为概形。设 $X$ 为 $S$ 上的代数空间。
若引理 \ref{spaces-morphisms-lemma-point-finite-type} 的等价条件成立，
则称点 $x \in |X|$ 为{\it 有限型点}\footnote{这略微滥用了术语；更准确的说法或许是“局部有限型点”。}。
以 $X_{\text{ft-pts}}$ 表示 $X$ 的有限型点集。
\end{definition}

\noindent
有限型点集可描述如下。

\begin{lemma}
\label{spaces-morphisms-lemma-identify-finite-type-points}
设 $S$ 为概形。设 $X$ 为 $S$ 上的代数空间。则
$$
X_{\text{ft-pts}} =
\bigcup\nolimits_{\varphi : U \to X\text{ \'etale }} |\varphi|(U_0)
$$
其中 $U_0$ 是 $U$ 的闭点集。
这里可令 $U$ 遍历所有在 $X$ 上 \'etale 的概形，或所有在 $X$ 上 \'etale 的仿射概形。
\end{lemma}

\begin{proof}
立即由引理 \ref{spaces-morphisms-lemma-point-finite-type} 得到。
\end{proof}

\begin{lemma}
\label{spaces-morphisms-lemma-finite-type-points-morphism}
设 $S$ 为概形。
设 $f : X \to Y$ 为 $S$ 上代数空间的态射。
若 $f$ 局部有限型，则
$f(X_{\text{ft-pts}}) \subset Y_{\text{ft-pts}}$.
\end{lemma}

\begin{proof}
取 $x \in X_{\text{ft-pts}}$。以局部有限型态射
$x : \Spec(k) \to X$ 表示 $x$。由引理
\ref{spaces-morphisms-lemma-composition-finite-type}，$f \circ x$ 局部有限型。
因此 $f(x) \in Y_{\text{ft-pts}}$。
\end{proof}

\begin{lemma}
\label{spaces-morphisms-lemma-finite-type-points-surjective-morphism}
设 $S$ 为概形。
设 $f : X \to Y$ 为 $S$ 上代数空间的态射。
若 $f$ 局部有限型且满射，则
$f(X_{\text{ft-pts}}) = Y_{\text{ft-pts}}$.
\end{lemma}

\begin{proof}
由引理 \ref{spaces-morphisms-lemma-finite-type-points-morphism}，有
$f(X_{\text{ft-pts}}) \subset Y_{\text{ft-pts}}$。
设 $y \in |Y|$ 为有限型点。以局部有限型态射
$\Spec(k) \to Y$ 表示 $y$。
由于 $f$ 满射，代数空间
$X_k = \Spec(k) \times_Y X$ 非空，因而由引理
\ref{spaces-morphisms-lemma-identify-finite-type-points}，它有一个有限型点 $x \in |X_k|$。
态射 $X_k \to X$ 是 $\Spec(k) \to Y$ 的基变换，故局部有限型
（引理 \ref{spaces-morphisms-lemma-base-change-finite-type}）。
因此由引理 \ref{spaces-morphisms-lemma-finite-type-points-morphism}，$x$ 在 $X$ 中的像是有限型点，
而按构造它映到 $y$。
\end{proof}

\begin{lemma}
\label{spaces-morphisms-lemma-enough-finite-type-points}
设 $S$ 为概形。设 $X$ 为 $S$ 上的代数空间。
对任意局部闭子集 $T \subset |X|$，有
$$
T \not = \emptyset
\Rightarrow
T \cap X_{\text{ft-pts}} \not = \emptyset.
$$
特别地，对任意闭子集 $T \subset |X|$，
$T \cap X_{\text{ft-pts}}$ 在 $T$ 中稠密。
\end{lemma}

\begin{proof}
设 $i : Z \to X$ 为 $T$ 上诱导的既约子空间结构，见注记
\ref{spaces-morphisms-remark-space-structure-locally-closed-subset}。
任意浸入都局部有限型，见引理
\ref{spaces-morphisms-lemma-immersion-locally-finite-type}。
因此由引理 \ref{spaces-morphisms-lemma-finite-type-points-morphism}，有
$Z_{\text{ft-pts}} \subset X_{\text{ft-pts}} \cap T$。
最后，任意非空仿射概形 $U$ 若有指向 $Z$ 的 \'etale 态射，则至少有一个闭点。
因此由引理 \ref{spaces-morphisms-lemma-identify-finite-type-points}，$Z$ 至少有一个有限型点。
引理得证。
\end{proof}

\noindent
下面给出代数空间有限型点的另一个更具技术性的刻画。

\begin{lemma}
\label{spaces-morphisms-lemma-point-finite-type-monomorphism}
设 $S$ 为概形。设 $X$ 为 $S$ 上的代数空间。
设 $x \in |X|$。以下条件等价：
\begin{enumerate}
\item $x$ 是有限型点；
\item 存在代数空间 $Z$，其底层拓扑空间 $|Z|$ 为单点集，并存在局部有限型态射
$f : Z \to X$，使得 $\{x\} = |f|(|Z|)$；
\item 存在代数空间 $Z$ 与态射 $f : Z \to X$，满足以下性质：
\begin{enumerate}
\item 存在满射 \'etale 态射
$z : \Spec(k) \to Z$，其中 $k$ 为域；
\item $f$ 局部有限型；
\item $f$ 是单态射；
\item $x = f(z)$。
\end{enumerate}
\end{enumerate}
\end{lemma}

\begin{proof}
假设 $x$ 为有限型点。取仿射概形 $U$、闭点 $u \in U$ 以及 \'etale 态射
$\varphi : U \to X$，满足 $\varphi(u) = x$，见引理
\ref{spaces-morphisms-lemma-identify-finite-type-points}。
照例令 $u = \Spec(\kappa(u))$。投影态射
$u \times_X u \to u$ 是复合
$$
u \times_X u \to u \times_X U \to u \times_X X = u
$$
其中第一箭头是闭浸入（$u \to U$ 的基变换），第二箭头是 \'etale 态射
（\'etale 态射 $U \to X$ 的基变换）。因此 $u \times_X U$ 是
$k$ 的有限可分扩张的谱的不交并（见《态射》引理
\ref{morphisms-lemma-etale-over-field}），
从而闭子概形 $u \times_X u$ 是 $k$ 的有限可分扩张的不交并；也就是说，
$u \times_X u \to u$ 是 \'etale 态射。由《代数空间》定理
\ref{spaces-theorem-presentation}，$Z = u/u \times_X u$ 是代数空间。
按构造，图
$$
\xymatrix{
u \ar[d] \ar[r] & U \ar[d] \\
Z \ar[r] & X
}
$$
交换且竖直箭头为 \'etale。因此 $Z \to X$ 局部有限型（见引理
\ref{spaces-morphisms-lemma-finite-type-local}）。
按构造，态射 $Z \to X$ 是单态射，且 $z$ 的像为 $x$。因此 (3) 成立。

\medskip\noindent
显然 (3) 推出 (2)。
若 (2) 成立，则由引理 \ref{spaces-morphisms-lemma-finite-type-points-morphism}，$x$ 是 $X$ 的有限型点
（同时由引理 \ref{spaces-morphisms-lemma-enough-finite-type-points} 可知
$Z_{\text{ft-pts}}$ 非空，即 $Z$ 的唯一一点是 $Z$ 的有限型点）。
\end{proof}





\section{Nagata 空间}
\label{spaces-morphisms-section-nagata}

\noindent
Nagata 代数空间的定义见《空间的性质》第
\ref{spaces-properties-section-types-properties} 节。

\begin{lemma}
\label{spaces-morphisms-lemma-finite-type-nagata}
设 $S$ 为概形。
设 $f : X \to Y$ 为 $S$ 上代数空间的态射。
若 $Y$ 是 Nagata 的且 $f$ 局部有限型，则 $X$ 是 Nagata 的。
\end{lemma}

\begin{proof}
取概形 $V$ 及满 \'etale 态射 $V \to Y$。
取概形 $U$ 及满 \'etale 态射 $U \to X \times_Y V$。
若 $Y$ 是 Nagata 的，则 $V$ 是 Nagata 概形。
若 $X \to Y$ 局部有限型，则 $U \to V$ 局部有限型。故由
《态射》引理 \ref{morphisms-lemma-finite-type-nagata}，
$V$ 是 Nagata 概形。于是按定义 $X$ 是 Nagata 的。
\end{proof}

\begin{lemma}
\label{spaces-morphisms-lemma-ubiquity-nagata}
下列类型的代数空间是 Nagata 的。
\begin{enumerate}
\item Nagata 概形上任意局部有限型的代数空间。
\item 域上任意局部有限型的代数空间。
\item Noether 完备局部环上任意局部有限型的代数空间。
\item $\mathbf{Z}$ 上任意局部有限型的代数空间。
\item 特征零 Dedekind 环上任意局部有限型的代数空间。
\item 等等。
\end{enumerate}
\end{lemma}

\begin{proof}
第一条由引理 \ref{spaces-morphisms-lemma-finite-type-nagata} 成立。
因此其余各条也成立，见《态射》引理
\ref{morphisms-lemma-ubiquity-nagata}。
\end{proof}








\section{拟有限态射}
\label{spaces-morphisms-section-quasi-finite}

\noindent
概形态射的“局部拟有限”性质在源与靶上都是 \'etale 局部的，见
《下降》评注 \ref{descent-remark-list-local-source-target}。
此性质也在基变换下稳定，并且在靶上是 fpqc 局部的，见
《态射》引理 \ref{morphisms-lemma-base-change-quasi-finite} 以及
《下降》引理 \ref{descent-lemma-descending-property-quasi-finite}。
因此，由上面的引理 \ref{spaces-morphisms-lemma-local-source-target}，
我们可以如下定义代数空间的态射局部拟有限的含义；当该态射可表时，
这与第 \ref{spaces-morphisms-section-representable} 节中已有的概念一致。

\begin{definition}
\label{spaces-morphisms-definition-locally-quasi-finite}
设 $S$ 为概形。
设 $f : X \to Y$ 为 $S$ 上代数空间的态射。
\begin{enumerate}
\item 若引理 \ref{spaces-morphisms-lemma-local-source-target} 中取
$\mathcal{P} = \text{局部拟有限}$.
时其中的等价条件成立，则称 $f$ {\it 局部拟有限}。
\item 设 $x \in |X|$。若存在 $x$ 的开邻域 $X' \subset X$，使得
$f|_{X'} : X' \to Y$ 局部拟有限，则称 $f$ {\it 在 $x$ 处拟有限}。
\item 若代数空间的态射 $f : X \to Y$ 局部拟有限且拟紧，
则称其{\it 拟有限}。
\end{enumerate}
\end{definition}

\noindent
由《态射》引理
\ref{morphisms-lemma-quasi-finite-locally-quasi-compact}，
最后一条与概形态射的拟有限概念相容。

\begin{lemma}
\label{spaces-morphisms-lemma-base-change-quasi-finite-locus}
设 $S$ 为概形。设 $f : X \to Y$ 与 $g : Y' \to Y$ 为
$S$ 上代数空间的态射。将 $f$ 沿 $g$ 的基变换记为
$f' : X' \to Y'$，并将投影记为 $g' : X' \to X$。
假设 $f$ 局部有限型。分别以 $W \subset |X|$ 与 $W' \subset |X'|$
表示 $f$ 与 $f'$ 拟有限的点集。
\begin{enumerate}
\item $W \subset |X|$ 与 $W' \subset |X'|$ 都是开的；
\item $W' = (g')^{-1}(W)$，即形成 $f$ 拟有限的点集与基变换可交换；
\item 局部拟有限态射的基变换局部拟有限；
\item 拟有限态射的基变换拟有限。
\end{enumerate}
\end{lemma}

\begin{proof}
取概形 $V$ 及满 \'etale 态射 $V \to Y$。
取概形 $U$ 及满 \'etale 态射 $U \to V \times_Y X$。
取概形 $V'$ 及满 \'etale 态射 $V' \to Y' \times_Y V$。
令 $U' = V' \times_V U$，于是 $U' \to X'$ 也是满 \'etale 态射。图示为
$$
\vcenter{
\xymatrix{
U' \ar[d] \ar[r] & U \ar[d] \\
V' \ar[r] & V
}
}
\quad\text{位于下图之上}\quad
\vcenter{
\xymatrix{
X' \ar[d] \ar[r] & X \ar[d] \\
Y' \ar[r] & Y
}
}
$$
取 $u \in |U|$，其像为 $x \in |X|$。
“局部拟有限”性质在源与靶上都是 \'etale 局部的，见
《下降》评注 \ref{descent-remark-list-local-source-target}。
因此可应用引理 \ref{spaces-morphisms-lemma-local-source-target-at-point} 与
\ref{spaces-morphisms-lemma-local-source-target-global-implies-local}，从而
$f : X \to Y$ 在 $x$ 处拟有限，当且仅当 $U \to V$ 在 $u$ 处拟有限。
对 $f' : X' \to Y'$ 及态射 $U' \to V'$ 也同样成立。
故 (1)、(2)、(3) 归结为《态射》引理
\ref{morphisms-lemma-base-change-quasi-finite} 与
\ref{morphisms-lemma-quasi-finite-points-open}。
(4) 由 (3) 和引理 \ref{spaces-morphisms-lemma-base-change-quasi-compact} 得出。
\end{proof}

\begin{lemma}
\label{spaces-morphisms-lemma-composition-quasi-finite}
拟有限态射的复合拟有限。
对于局部拟有限态射也有同样结论。
\end{lemma}

\begin{proof}
见评注 \ref{spaces-morphisms-remark-composition-P} 和
《态射》引理 \ref{morphisms-lemma-composition-quasi-finite}。
\end{proof}

\begin{lemma}
\label{spaces-morphisms-lemma-base-change-quasi-finite}
拟有限态射的基变换拟有限。
对于局部拟有限态射也有同样结论。
\end{lemma}

\begin{proof}
这是引理 \ref{spaces-morphisms-lemma-base-change-quasi-finite-locus} 的直接推论。
\end{proof}

\noindent
下述引理把局部拟有限态射刻画为局部有限型且具有“离散纤维”的态射。
不过，这不同于要求 $|X| \to |Y|$ 具有离散纤维，参见
《例》第 \ref{examples-section-specializations-fibre-etale} 节中的讨论。

\begin{lemma}
\label{spaces-morphisms-lemma-locally-quasi-finite}
设 $S$ 为概形。设 $f : X \to Y$ 为代数空间的态射。
假设 $f$ 局部有限型。下列条件等价：
\begin{enumerate}
\item $f$ 局部拟有限；
\item 对每个态射 $\Spec(k) \to Y$，其中 $k$ 为域，空间 $|X_k|$
都是离散的。这里 $X_k = \Spec(k) \times_Y X$。
\end{enumerate}
\end{lemma}

\begin{proof}
假设 $f$ 局部拟有限。令 $\Spec(k) \to Y$ 如 (2) 所述。
取满 \'etale 态射 $U \to X$，其中 $U$ 为概形。由《空间的性质》引理
\ref{spaces-properties-lemma-base-change-etale}，
$U_k = \Spec(k) \times_Y U \to X_k$ 是代数空间的 \'etale 态射。
由引理 \ref{spaces-morphisms-lemma-base-change-quasi-finite}，
$X_k \to \Spec(k)$ 局部拟有限。按定义，这意味着
$U_k \to \Spec(k)$ 局部拟有限。故由《态射》引理
\ref{morphisms-lemma-locally-quasi-finite-fibres}，$|U_k|$ 是离散的。
由于 $|U_k| \to |X_k|$ 满且开，遂得 $|X_k|$ 离散。

\medskip\noindent
反之，假设 (2)。取满 \'etale 态射 $V \to Y$，其中 $V$ 为概形。
再取满 \'etale 态射 $U \to V \times_Y X$，其中 $U$ 为概形。
注意，由于 $f$ 局部有限型，$U \to V$ 也局部有限型。图示为
$$
\xymatrix{
U \ar[r] \ar[rd] & X \times_Y V \ar[d] \ar[r] & V \ar[d] \\
& X \ar[r] & Y
}
$$
若 $f$ 不局部拟有限，则 $U \to V$ 不局部拟有限。因此存在专门化
$u \leadsto u'$，其中某些 $u, u' \in U$ 位于同一点 $v \in V$ 上方，见
《态射》引理 \ref{morphisms-lemma-quasi-finite-at-point-characterize}。
我们断言 $u, u'$ 在 $X_v = \Spec(\kappa(v)) \times_Y X$ 中的像不同；
这将按所需方式与 $|X_v|$ 离散的假设矛盾。令
$d = \text{trdeg}_{\kappa(v)}(\kappa(u))$ 以及
$d' = \text{trdeg}_{\kappa(v)}(\kappa(u'))$。
由《态射》引理 \ref{morphisms-lemma-dimension-fibre-specialization}，
可见 $d > d'$。注意 $U_v$（即 $U \to V$ 在 $v$ 上的纤维）是
$U$ 与 $X_v$ 在 $X \times_Y V$ 上的纤维积，故 $U_v \to X_v$ 是
\'etale 的（它是 \'etale 态射 $U \to X \times_Y V$ 的基变换）。
若 $u, u' \in U_v$ 映到 $|X_v|$ 的同一元素，则由《空间的性质》引理
\ref{spaces-properties-lemma-points-cartesian}，存在一点
$r \in R_v = U_v \times_{X_v} U_v$，满足 $t(r) = u$ 且 $s(r) = u'$。
注意 $s, t : R_v \to U_v$ 是 $\kappa(v)$ 上概形的 \'etale 态射，故
$\kappa(u) \subset \kappa(r) \supset \kappa(u')$ 是 $\kappa(v)$ 上域的
有限可分扩张（见《态射》引理 \ref{morphisms-lemma-etale-over-field}）。
因此这些超越次数相等。这一矛盾完成证明。
\end{proof}

\begin{lemma}
\label{spaces-morphisms-lemma-quasi-finite-local}
设 $S$ 为概形。
设 $f : X \to Y$ 为 $S$ 上代数空间的态射。
下列条件等价：
\begin{enumerate}
\item $f$ 局部拟有限；
\item 对每个 $x \in |X|$，态射 $f$ 在 $x$ 处拟有限；
\item 对每个概形 $Z$ 及任意态射 $Z \to Y$，态射
$Z \times_Y X \to Z$ 局部拟有限；
\item 对每个仿射概形 $Z$ 及任意态射 $Z \to Y$，态射
$Z \times_Y X \to Z$ 局部拟有限；
\item 存在概形 $V$ 及满 \'etale 态射 $V \to Y$，使得
$V \times_Y X \to V$ 局部拟有限；
\item 存在概形 $U$ 及满 \'etale 态射 $\varphi : U \to X$，使得复合
$f \circ \varphi$ 局部拟有限；
\item 对每个交换图
$$
\xymatrix{
U \ar[d] \ar[r] & V \ar[d] \\
X \ar[r] & Y
}
$$
其中 $U$、$V$ 为概形，竖直箭头为 \'etale 态射，顶部水平箭头局部拟有限；
\item 存在交换图
$$
\xymatrix{
U \ar[d] \ar[r] & V \ar[d] \\
X \ar[r] & Y
}
$$
其中 $U$、$V$ 为概形，竖直箭头为 \'etale 态射，且 $U \to X$ 满，
使得顶部水平箭头局部拟有限；
\item 存在 Zariski 覆盖 $Y = \bigcup_{i \in I} Y_i$ 以及
$f^{-1}(Y_i) = \bigcup X_{ij}$，使得每个态射
$X_{ij} \to Y_i$ 局部拟有限。
\end{enumerate}
\end{lemma}

\begin{proof}
略。
\end{proof}

\begin{lemma}
\label{spaces-morphisms-lemma-immersion-quasi-finite}
浸入局部拟有限。
\end{lemma}

\begin{proof}
略。
\end{proof}

\begin{lemma}
\label{spaces-morphisms-lemma-permanence-quasi-finite}
设 $S$ 为概形。
设 $X \to Y \to Z$ 为 $S$ 上代数空间的态射。
若 $X \to Z$ 局部拟有限，则 $X \to Y$ 局部拟有限。
\end{lemma}

\begin{proof}
取交换图
$$
\xymatrix{
U \ar[d] \ar[r] & V \ar[d] \ar[r] & W \ar[d] \\
X \ar[r] & Y \ar[r] & Z
}
$$
使竖直箭头均为满 \'etale 态射。（见《空间》引理
\ref{spaces-lemma-lift-morphism-presentations}。）
对顶行应用《态射》引理
\ref{morphisms-lemma-permanence-quasi-finite}。
\end{proof}

\begin{lemma}
\label{spaces-morphisms-lemma-quasi-finite-at-a-finite-number-of-points}
设 $S$ 为概形。设 $f : X \to Y$ 为 $S$ 上代数空间的有限型态射。
设 $y \in |Y|$。在 $y$ 上方且使 $f$ 拟有限的 $|X|$ 中点至多有限个。
\end{lemma}

\begin{proof}
取域 $k$ 及由 $y$ 所确定的等价类中的态射 $\Spec(k) \to Y$。
纤维 $X_k = \Spec(k) \times_Y X$ 是域上的有限型代数空间，尤其拟紧。
映射 $|X_k| \to |X|$ 满射到 $|X| \to |Y|$ 在 $y$ 上的纤维
（《空间的性质》引理 \ref{spaces-properties-lemma-points-cartesian}）。
而且，由引理 \ref{spaces-morphisms-lemma-base-change-quasi-finite-locus}，
$X_k \to \Spec(k)$ 拟有限的点集映满到 $y$ 上方使 $f$ 拟有限的点集。
取仿射概形 $U$ 及满 \'etale 态射 $U \to X_k$
（《空间的性质》引理
\ref{spaces-properties-lemma-quasi-compact-affine-cover}）。
于是 $U \to \Spec(k)$ 是有限型态射，且使该态射拟有限的点至多有限个，
见《态射》引理
\ref{morphisms-lemma-quasi-finite-at-a-finite-number-of-points}。
由于 $X_k \to \Spec(k)$ 在一点 $x'$ 处拟有限，当且仅当该点是 $U$
中某个使 $U \to \Spec(k)$ 拟有限之点的像，结论随即成立。
\end{proof}

\begin{lemma}
\label{spaces-morphisms-lemma-monomorphism-loc-finite-type-loc-quasi-finite}
设 $S$ 为概形。
设 $f : X \to Y$ 为 $S$ 上代数空间的态射。
若 $f$ 局部有限型且为单态射，则 $f$ 分离并且局部拟有限。
\end{lemma}

\begin{proof}
单态射是分离的，见引理 \ref{spaces-morphisms-lemma-monomorphism-separated}。
由引理 \ref{spaces-morphisms-lemma-quasi-finite-local}，只需在沿 $Z \to Y$ 作基变换后
证明此引理，其中 $Z$ 仿射。因此可以假设 $Y$ 是仿射概形。
取仿射概形 $U$ 及 \'etale 态射 $U \to X$。由于 $X \to Y$ 局部有限型，
仿射概形的态射 $U \to Y$ 有限型。由于 $X \to Y$ 是单态射，有
$U \times_X U = U \times_Y U$。特别地，映射 $U \times_Y U \to U$
是 \'etale 的。设 $y \in Y$。则或者 $U_y$ 为空，或者对位于 $y$
上方的某个 $u \in U$，
$\Spec(\kappa(u)) \times_{\Spec(\kappa(y))} U_y$
同构于 $U \times_Y U \to U$ 在 $u$ 上的纤维。这说明 $U \to Y$ 的
纤维是有限离散集（因为 $U \times_Y U \to U$ 是仿射概形的 \'etale 态射，
见《态射》引理 \ref{morphisms-lemma-etale-over-field}）。
因此 $U \to Y$ 拟有限，见《态射》引理
\ref{morphisms-lemma-quasi-finite-at-point-characterize}。
由于 $U \to X$ 是任意一个以仿射概形 $U$ 为源的 \'etale 态射，
这说明 $X \to Y$ 局部拟有限。
\end{proof}













\section{有限表示态射}
\label{spaces-morphisms-section-finite-presentation}

\noindent
概形态射的“局部有限表示”性质在源与靶上都是 \'etale 局部的，见
《下降》评注 \ref{descent-remark-list-local-source-target}。
此性质也在基变换下稳定，并且在靶上是 fpqc 局部的，见
《态射》引理 \ref{morphisms-lemma-base-change-finite-presentation} 以及
《下降》引理
\ref{descent-lemma-descending-property-locally-finite-presentation}。
因此，由上面的引理 \ref{spaces-morphisms-lemma-local-source-target}，
我们可以如下定义代数空间的态射局部有限表示的含义；当该态射可表时，
这与第 \ref{spaces-morphisms-section-representable} 节中已有的概念一致。

\begin{definition}
\label{spaces-morphisms-definition-locally-finite-presentation}
设 $S$ 为概形。
设 $f : X \to Y$ 为 $S$ 上代数空间的态射。
\begin{enumerate}
\item 若引理 \ref{spaces-morphisms-lemma-local-source-target} 中取
$\mathcal{P} =$“局部有限表示”时其中的等价条件成立，则称 $f$
{\it 局部有限表示}。
\item 设 $x \in |X|$。若存在 $x$ 的开邻域 $X' \subset X$，使得
$f|_{X'} : X' \to Y$ 局部有限表示，则称 $f$ {\it 在 $x$ 处有限表示}
\footnote{使用“在 $x$ 处局部有限表示”似乎很别扭；不过，现用术语也可能
引起误解，因为“在 $x$ 处有限表示”{\bf 并不}意味着存在开邻域
$X' \subset X$，使得 $f|_{X'}$ 有限表示。}。
\item 若代数空间的态射 $f : X \to Y$ 局部有限表示、拟紧且拟分离，
则称其{\it 有限表示}。
\end{enumerate}
\end{definition}

\noindent
注意，有限表示态射{\bf 并不}只是局部有限表示的拟紧态射。

\begin{lemma}
\label{spaces-morphisms-lemma-composition-finite-presentation}
有限表示态射的复合有限表示。
对于局部有限表示态射也有同样结论。
\end{lemma}

\begin{proof}
见评注 \ref{spaces-morphisms-remark-composition-P} 和《态射》引理
\ref{morphisms-lemma-composition-finite-presentation}。
还要使用拟紧态射与拟分离态射的相应结论
（引理 \ref{spaces-morphisms-lemma-composition-quasi-compact} 与
\ref{spaces-morphisms-lemma-composition-separated}）。
\end{proof}

\begin{lemma}
\label{spaces-morphisms-lemma-base-change-finite-presentation}
有限表示态射的基变换有限表示。
对于局部有限表示态射也有同样结论。
\end{lemma}

\begin{proof}
见评注 \ref{spaces-morphisms-remark-base-change-P} 和《态射》引理
\ref{morphisms-lemma-base-change-finite-presentation}。
还要使用拟紧态射与拟分离态射的相应结论
（引理 \ref{spaces-morphisms-lemma-base-change-quasi-compact} 与
\ref{spaces-morphisms-lemma-base-change-separated}）。
\end{proof}

\begin{lemma}
\label{spaces-morphisms-lemma-finite-presentation-local}
设 $S$ 为概形。
设 $f : X \to Y$ 为 $S$ 上代数空间的态射。
下列条件等价：
\begin{enumerate}
\item $f$ 局部有限表示；
\item 对每个 $x \in |X|$，态射 $f$ 在 $x$ 处有限表示；
\item 对每个概形 $Z$ 及任意态射 $Z \to Y$，态射
$Z \times_Y X \to Z$ 局部有限表示；
\item 对每个仿射概形 $Z$ 及任意态射 $Z \to Y$，态射
$Z \times_Y X \to Z$ 局部有限表示；
\item 存在概形 $V$ 及满 \'etale 态射 $V \to Y$，使得
$V \times_Y X \to V$ 局部有限表示；
\item 存在概形 $U$ 及满 \'etale 态射 $\varphi : U \to X$，使得复合
$f \circ \varphi$ 局部有限表示；
\item 对每个交换图
$$
\xymatrix{
U \ar[d] \ar[r] & V \ar[d] \\
X \ar[r] & Y
}
$$
其中 $U$、$V$ 为概形，竖直箭头为 \'etale 态射，顶部水平箭头局部有限表示；
\item 存在交换图
$$
\xymatrix{
U \ar[d] \ar[r] & V \ar[d] \\
X \ar[r] & Y
}
$$
其中 $U$、$V$ 为概形，竖直箭头为 \'etale 态射，且 $U \to X$ 满，
使得顶部水平箭头局部有限表示；
\item 存在 Zariski 覆盖 $Y = \bigcup_{i \in I} Y_i$ 以及
$f^{-1}(Y_i) = \bigcup X_{ij}$，使得每个态射
$X_{ij} \to Y_i$ 局部有限表示。
\end{enumerate}
\end{lemma}

\begin{proof}
略。
\end{proof}

\begin{lemma}
\label{spaces-morphisms-lemma-finite-presentation-finite-type}
局部有限表示的态射局部有限型。
有限表示态射有限型。
\end{lemma}

\begin{proof}
设 $f : X \to Y$ 为局部有限表示的代数空间态射。这意味着存在
引理 \ref{spaces-morphisms-lemma-local-source-target} 中的图，其中 $h$ 局部有限表示，
且竖直箭头 $a$ 满。由《态射》引理
\ref{morphisms-lemma-finite-presentation-finite-type}，$h$ 局部有限型。
故按定义 $X \to Y$ 局部有限型。若 $f$ 有限表示，则它拟紧，
从而 $f$ 有限型。
\end{proof}

\begin{lemma}
\label{spaces-morphisms-lemma-finite-presentation-noetherian}
设 $S$ 为概形。
设 $f : X \to Y$ 为 $S$ 上代数空间的态射。
若 $f$ 有限表示且 $Y$ 是 Noether 的，则 $X$ 是 Noether 的。
\end{lemma}

\begin{proof}
假设 $f$ 有限表示且 $Y$ 是 Noether 的。由引理
\ref{spaces-morphisms-lemma-finite-presentation-finite-type} 与
\ref{spaces-morphisms-lemma-locally-finite-type-locally-noetherian}，可见 $X$ 局部 Noether。
由于 $f$ 拟紧且 $Y$ 拟紧，可见 $X$ 拟紧。
由于 $f$ 有限表示，它拟分离（见定义
\ref{spaces-morphisms-definition-locally-finite-presentation}）；又由于 $Y$ 是 Noether 的，
它拟分离（见《空间的性质》定义
\ref{spaces-properties-definition-noetherian}）。
故由引理 \ref{spaces-morphisms-lemma-separated-over-separated}，$X$ 拟分离。
于是我们已验证《空间的性质》定义
\ref{spaces-properties-definition-noetherian} 中的全部三个条件，证明完成。
\end{proof}

\begin{lemma}
\label{spaces-morphisms-lemma-noetherian-finite-type-finite-presentation}
设 $S$ 为概形。
设 $f : X \to Y$ 为 $S$ 上代数空间的态射。
\begin{enumerate}
\item 若 $Y$ 局部 Noether 且 $f$ 局部有限型，则 $f$ 局部有限表示。
\item 若 $Y$ 局部 Noether，且 $f$ 有限型并拟分离，则 $f$ 有限表示。
\end{enumerate}
\end{lemma}

\begin{proof}
假设 $f : X \to Y$ 局部有限型且 $Y$ 局部 Noether。
这意味着存在引理 \ref{spaces-morphisms-lemma-local-source-target} 中的图，
其中 $h$ 局部有限型，且竖直箭头 $a$ 满。由《态射》引理
\ref{morphisms-lemma-noetherian-finite-type-finite-presentation}，
$h$ 局部有限表示。故按定义 $X \to Y$ 局部有限表示。这证明 (1)。
若 $f$ 有限型且拟分离，则它也拟紧且拟分离，(2) 立即得出。
\end{proof}

\begin{lemma}
\label{spaces-morphisms-lemma-finite-presentation-quasi-compact-quasi-separated}
设 $S$ 为概形。设 $Y$ 为 $S$ 上拟紧且拟分离的代数空间。
若 $X$ 在 $Y$ 上有限表示，则 $X$ 拟紧且拟分离。
\end{lemma}

\begin{proof}
略。
\end{proof}

\begin{lemma}
\label{spaces-morphisms-lemma-finite-presentation-permanence}
设 $S$ 为概形。
设 $f : X \to Y$ 与 $Y \to Z$ 为 $S$ 上代数空间的态射。
若 $X$ 在 $Z$ 上局部有限表示，且 $Y$ 在 $Z$ 上局部有限型，
则 $f$ 局部有限表示。
\end{lemma}

\begin{proof}
取概形 $W$ 及满 \'etale 态射 $W \to Z$。再取概形 $V$ 及满 \'etale 态射
$V \to W \times_Z Y$。最后取概形 $U$ 及满 \'etale 态射
$U \to V \times_Y X$。按定义，$U$ 在 $W$ 上局部有限表示，而 $V$
在 $W$ 上局部有限型。由《态射》引理
\ref{morphisms-lemma-finite-presentation-permanence}，态射 $U \to V$
局部有限表示。因此 $f$ 局部有限表示。
\end{proof}

\begin{lemma}
\label{spaces-morphisms-lemma-diagonal-morphism-finite-type}
设 $S$ 为概形。设 $f : X \to Y$ 为 $S$ 上代数空间的态射，其对角态射为
$\Delta : X \to X \times_Y X$。若 $f$ 局部有限型，则 $\Delta$ 局部有限表示。
若 $f$ 拟分离且局部有限型，则 $\Delta$ 有限表示。
\end{lemma}

\begin{proof}
注意，$\Delta$ 是 $X$ 上的态射（通过第二投影 $X \times_Y X \to X$）。
假设 $f$ 局部有限型。注意，$X$ 在 $X$ 上有限表示，而 $X \times_Y X$
在 $X$ 上有限型（由引理 \ref{spaces-morphisms-lemma-base-change-finite-type}）。
因此第一条由引理 \ref{spaces-morphisms-lemma-finite-presentation-permanence} 成立。
第二条由第一条、定义以及对角态射分离这一事实
（引理 \ref{spaces-morphisms-lemma-properties-diagonal}）得出。
\end{proof}

\begin{lemma}
\label{spaces-morphisms-lemma-open-immersion-locally-finite-presentation}
代数空间的开浸入局部有限表示。
\end{lemma}

\begin{proof}
开浸入按定义可表，因此可使用《空间》引理
\ref{spaces-lemma-representable-transformations-property-implication}
的一般原理以及《态射》引理
\ref{morphisms-lemma-open-immersion-locally-finite-presentation}。
\end{proof}

\begin{lemma}
\label{spaces-morphisms-lemma-closed-immersion-finite-presentation}
闭浸入 $i : Z \to X$ 有限表示，当且仅当相伴的拟凝聚理想层
$\mathcal{I} = \Ker(\mathcal{O}_X \to i_*\mathcal{O}_Z)$
有限型（作为 $\mathcal{O}_X$-模）。
\end{lemma}

\begin{proof}
设 $U$ 为概形，并设 $U \to X$ 为满 \'etale 态射。由引理
\ref{spaces-morphisms-lemma-finite-presentation-local}，可见
$i' : Z \times_X U \to U$ 有限表示，当且仅当 $i$ 有限表示。
由《空间的性质》第
\ref{spaces-properties-section-properties-modules} 节，可见 $\mathcal{I}$
有限型，当且仅当
$\mathcal{I}|_U = \Ker(\mathcal{O}_U \to i'_*\mathcal{O}_{Z \times_X U})$
有限型。因此结论归结为概形情形，见《态射》引理
\ref{morphisms-lemma-closed-immersion-finite-presentation}。
\end{proof}







\section{可构造集}
\label{spaces-morphisms-section-constructible}

\noindent
本节接续《空间的性质》第
\ref{spaces-properties-section-constructible} 节。

\begin{lemma}
\label{spaces-morphisms-lemma-inverse-image-constructible}
设 $S$ 为概形。
设 $f : X \to Y$ 为 $S$ 上代数空间的态射。
设 $E \subset |Y|$ 为子集。若 $E$ 在 $Y$ 中 \'etale 局部可构造，
则 $f^{-1}(E)$ 在 $X$ 中 \'etale 局部可构造。
\end{lemma}

\begin{proof}
取概形 $V$ 及满 \'etale 态射 $\varphi : V \to Y$。
取概形 $U$ 及满 \'etale 态射 $U \to V \times_Y X$。
于是 $U \to X$ 满且 \'etale，而 $f^{-1}(E)$ 在 $U$ 中的原像是
$\varphi^{-1}(E)$ 沿 $U \to V$ 的原像。因此引理归结为态射 $U \to V$
的概形情形（《态射》引理
\ref{morphisms-lemma-inverse-image-constructible}）以及定义
（《空间的性质》定义
\ref{spaces-properties-definition-locally-constructible}）。
\end{proof}

\begin{theorem}[Chevalley 定理]
\label{spaces-morphisms-theorem-chevalley}
设 $S$ 为概形。
设 $f : X \to Y$ 为 $S$ 上代数空间的态射。
假设 $f$ 拟紧且局部有限表示。则 $|X|$ 的每个 \'etale 局部可构造子集的像
都是 $|Y|$ 的 \'etale 局部可构造子集。
\end{theorem}

\begin{proof}
设 $E \subset |X|$ 为 \'etale 局部可构造子集。
设 $V \to Y$ 为 \'etale 态射，且 $V$ 仿射。只需证明 $f(E)$ 在 $V$
中的原像可构造，见《空间的性质》定义
\ref{spaces-properties-definition-locally-constructible}。
由于 $f$ 拟紧，$V \times_Y X$ 是拟紧代数空间。取仿射概形 $U$ 及满
\'etale 态射 $U \to V \times_Y X$（《空间的性质》引理
\ref{spaces-properties-lemma-quasi-compact-affine-cover}）。
由《空间的性质》引理 \ref{spaces-properties-lemma-points-cartesian}，
$f(E)$ 在 $V$ 中的原像，是 $E$ 在 $U$ 中的原像沿 $U \to V$ 的像。
因此结论归结为概形情形，见《态射》引理
\ref{morphisms-lemma-chevalley}。
\end{proof}







\section{平坦态射}
\label{spaces-morphisms-section-flat}

\noindent
概形态射的“平坦”性质在源与靶上都是 \'etale 局部的，见
《下降》评注 \ref{descent-remark-list-local-source-target}。
此性质也在基变换下稳定，并且在靶上是 fpqc 局部的，见
《态射》引理 \ref{morphisms-lemma-base-change-flat} 与
《下降》引理 \ref{descent-lemma-descending-property-flat}。
因此，由上面的引理 \ref{spaces-morphisms-lemma-local-source-target}，
我们可以如下定义代数空间的平坦态射；当该态射可表时，
这与第 \ref{spaces-morphisms-section-representable} 节中已有的概念一致。

\begin{definition}
\label{spaces-morphisms-definition-flat}
设 $S$ 为概形。
设 $f : X \to Y$ 为 $S$ 上代数空间的态射。
\begin{enumerate}
\item 若引理 \ref{spaces-morphisms-lemma-local-source-target} 中取
$\mathcal{P} =$“平坦”时其中的等价条件成立，则称 $f$ {\it 平坦}。
\item 设 $x \in |X|$。若引理 \ref{spaces-morphisms-lemma-local-source-target-at-point}
中取 $\mathcal{Q} =$“所诱导的局部环映射平坦”时其中的等价条件成立，
则称 $f$ {\it 在 $x$ 处平坦}。
\end{enumerate}
注意，由《下降》引理 \ref{descent-lemma-flat-at-point}，第二条定义良好。
\end{definition}

\noindent
先作一个简单核对。

\begin{lemma}
\label{spaces-morphisms-lemma-flat-is-flat-at-all-points}
设 $S$ 为概形。设 $f : X \to Y$ 为 $S$ 上代数空间的态射。
则 $f$ 平坦，当且仅当 $f$ 在 $|X|$ 的每一点处平坦。
\end{lemma}

\begin{proof}
取交换图
$$
\xymatrix{
U \ar[d]_a \ar[r]_h & V \ar[d]^b \\
X \ar[r]^f & Y
}
$$
其中 $U$ 与 $V$ 为概形，竖直箭头为 \'etale 态射，且 $a$ 满。
按定义，$f$ 平坦，当且仅当 $h$ 平坦（定义 \ref{spaces-morphisms-definition-P}）。
按定义，$f$ 在 $x \in |X|$ 处平坦，当且仅当 $h$ 在映到 $x$ 的某个
（等价地，任意一个）$u \in U$ 处平坦（定义
\ref{spaces-morphisms-definition-P-at-point}）。因此本引理由如下事实得出：
概形态射平坦，当且仅当它在源的每一点处平坦
（《态射》定义 \ref{morphisms-definition-flat}）。
\end{proof}

\begin{lemma}
\label{spaces-morphisms-lemma-composition-flat}
平坦态射的复合平坦。
\end{lemma}

\begin{proof}
见评注 \ref{spaces-morphisms-remark-composition-P} 和
《态射》引理 \ref{morphisms-lemma-composition-flat}。
\end{proof}

\begin{lemma}
\label{spaces-morphisms-lemma-base-change-flat}
平坦态射的基变换平坦。
\end{lemma}

\begin{proof}
见评注 \ref{spaces-morphisms-remark-base-change-P} 和
《态射》引理 \ref{morphisms-lemma-base-change-flat}。
\end{proof}

\begin{lemma}
\label{spaces-morphisms-lemma-flat-local}
设 $S$ 为概形。
设 $f : X \to Y$ 为 $S$ 上代数空间的态射。
下列条件等价：
\begin{enumerate}
\item $f$ 平坦；
\item 对每个 $x \in |X|$，态射 $f$ 在 $x$ 处平坦；
\item 对每个概形 $Z$ 及任意态射 $Z \to Y$，态射
$Z \times_Y X \to Z$ 平坦；
\item 对每个仿射概形 $Z$ 及任意态射 $Z \to Y$，态射
$Z \times_Y X \to Z$ 平坦；
\item 存在概形 $V$ 及满 \'etale 态射 $V \to Y$，使得
$V \times_Y X \to V$ 平坦；
\item 存在概形 $U$ 及满 \'etale 态射 $\varphi : U \to X$，使得复合
$f \circ \varphi$ 平坦；
\item 对每个交换图
$$
\xymatrix{
U \ar[d] \ar[r] & V \ar[d] \\
X \ar[r] & Y
}
$$
其中 $U$、$V$ 为概形，竖直箭头为 \'etale 态射，顶部水平箭头平坦；
\item 存在交换图
$$
\xymatrix{
U \ar[d] \ar[r] & V \ar[d] \\
X \ar[r] & Y
}
$$
其中 $U$、$V$ 为概形，竖直箭头为 \'etale 态射，且 $U \to X$ 满，
使得顶部水平箭头平坦；
\item 存在 Zariski 覆盖 $Y = \bigcup Y_i$ 以及
$f^{-1}(Y_i) = \bigcup X_{ij}$，使得每个态射 $X_{ij} \to Y_i$ 平坦。
\end{enumerate}
\end{lemma}

\begin{proof}
略。
\end{proof}

\begin{lemma}
\label{spaces-morphisms-lemma-fppf-open}
局部有限表示的平坦态射是普遍开的。
\end{lemma}

\begin{proof}
设 $f : X \to Y$ 为 $S$ 上代数空间的局部有限表示平坦态射。取图
$$
\xymatrix{
U \ar[r]_\alpha \ar[d] & V \ar[d] \\
X \ar[r] & Y
}
$$
其中 $U$ 与 $V$ 为概形，竖直箭头满且 \'etale，见《空间》引理
\ref{spaces-lemma-lift-morphism-presentations}。
由引理 \ref{spaces-morphisms-lemma-flat-local} 与 \ref{spaces-morphisms-lemma-finite-presentation-local}，
态射 $\alpha$ 平坦且局部有限表示。因此由《态射》引理
\ref{morphisms-lemma-fppf-open}，可见 $\alpha$ 普遍开。
故由引理 \ref{spaces-morphisms-lemma-universally-open-local}，$X \to Y$ 普遍开。
\end{proof}

\begin{lemma}
\label{spaces-morphisms-lemma-fpqc-quotient-topology}
设 $S$ 为概形。
设 $f : X \to Y$ 为 $S$ 上代数空间的平坦、拟紧满态射。
子集 $T \subset |Y|$ 是开的（相应地，闭的），当且仅当
$f^{-1}(|T|)$ 在 $|X|$ 中是开的（相应地，闭的）。
换言之，$f$ 是商映射；事实上，它是普遍商映射。
\end{lemma}

\begin{proof}
取仿射概形 $V_i$ 及 \'etale 态射 $V_i \to Y$，使得
$V = \coprod V_i \to Y$ 满，见《空间的性质》引理
\ref{spaces-properties-lemma-cover-by-union-affines}。
对每个 $i$，代数空间 $V_i \times_Y X$ 拟紧。因此可取仿射概形 $U_i$
及满 \'etale 态射 $U_i \to V_i \times_Y X$，见《空间的性质》引理
\ref{spaces-properties-lemma-quasi-compact-affine-cover}。
于是复合 $U_i \to V_i \times_Y X \to V_i$ 是仿射概形之间的满平坦态射。
当然，$U = \coprod U_i \to X$ 满且 \'etale，并且
$U \to V \times_Y X$ 满。此外，态射 $U \to V$ 是各态射
$U_i \to V_i$ 的不交并。故 $U \to V$ 满、拟紧且平坦。考虑图
$$
\xymatrix{
U \ar[r] \ar[d] & X \ar[d] \\
V \ar[r] & Y
}
$$
按 $|Y|$ 上拓扑的定义，集合 $T$ 是闭的（相应地，开的），当且仅当
$g^{-1}(T) \subset |V|$ 是闭的（相应地，开的）。对 $f^{-1}(T)$ 及其在
$|U|$ 中的原像也同样成立。由于 $U \to V$ 拟紧、满且平坦，
结论由《态射》引理 \ref{morphisms-lemma-fpqc-quotient-topology} 得出。
\end{proof}

\begin{lemma}
\label{spaces-morphisms-lemma-flat-at-point-etale-local-rings}
设 $S$ 为概形。
设 $f : X \to Y$ 为 $S$ 上代数空间的态射。
设 $\overline{x}$ 为 $X$ 的一个几何点，位于点 $x \in |X|$ 上方。
令 $\overline{y} = f \circ \overline{x}$。下列条件等价：
\begin{enumerate}
\item $f$ 在 $x$ 处平坦；
\item \'etale 局部环之间的映射
$\mathcal{O}_{Y, \overline{y}} \to \mathcal{O}_{X, \overline{x}}$
平坦。
\end{enumerate}
\end{lemma}

\begin{proof}
取交换图
$$
\xymatrix{
U \ar[d]_a \ar[r]_h & V \ar[d]^b \\
X \ar[r]^f & Y
}
$$
其中 $U$ 与 $V$ 为概形，$a, b$ 为 \'etale 态射，且 $u \in U$ 映到 $x$。
可取位于 $u$ 上方的几何点 $\overline{u} : \Spec(k) \to U$，使得
$\overline{x} = a \circ \overline{u}$，见《空间的性质》引理
\ref{spaces-properties-lemma-geometric-lift-to-usual}。
令 $\overline{v} = h \circ \overline{u}$，其像为 $v \in V$。我们知道
$$
\mathcal{O}_{X, \overline{x}} = \mathcal{O}_{U, u}^{sh}
\quad\text{且}\quad
\mathcal{O}_{Y, \overline{y}} = \mathcal{O}_{V, v}^{sh}
$$
见《空间的性质》引理
\ref{spaces-properties-lemma-describe-etale-local-ring}。于是得到局部环的交换图
$$
\xymatrix{
\mathcal{O}_{U, u} \ar[r] &
\mathcal{O}_{X, \overline{x}} \\
\mathcal{O}_{V, v} \ar[u] \ar[r] &
\mathcal{O}_{Y, \overline{y}} \ar[u]
}
$$
其中水平箭头平坦。我们要证明左侧竖直箭头平坦，当且仅当右侧竖直箭头平坦。
《代数》引理 \ref{algebra-lemma-flatness-descends-more-general} 告诉我们，
$\mathcal{O}_{U, u}$ 在 $\mathcal{O}_{V, v}$ 上平坦，当且仅当
$\mathcal{O}_{X, \overline{x}}$ 在 $\mathcal{O}_{V, v}$ 上平坦。
因此结论由《平坦性的进一步内容》引理
\ref{flat-lemma-flat-up-down-henselization} 得出。
\end{proof}

\begin{lemma}
\label{spaces-morphisms-lemma-flat-morphism-sites}
设 $S$ 为概形。
设 $f : X \to Y$ 为 $S$ 上代数空间的态射。
则 $f$ 平坦，当且仅当与 $f$ 相伴的站点态射
$
(f_{small}, f^\sharp) :
(X_\etale, \mathcal{O}_X)
\to
(Y_\etale, \mathcal{O}_Y)
$
平坦。
\end{lemma}

\begin{proof}
$(f_{small}, f^\sharp)$ 的平坦性由 $\mathcal{O}_X$ 作为
$f_{small}^{-1}\mathcal{O}_Y$-模的平坦性定义。这可以在茎上检验，见
《站点上的模》引理 \ref{sites-modules-lemma-check-flat-stalks} 以及
《空间的性质》定理 \ref{spaces-properties-theorem-exactness-stalks}。
而我们已经看到，$f$ 的平坦性可以在茎上检验，见引理
\ref{spaces-morphisms-lemma-flat-at-point-etale-local-rings}。
\end{proof}

\begin{lemma}
\label{spaces-morphisms-lemma-flat-pullback-support}
设 $S$ 为概形。设 $f : Y \to X$ 为 $S$ 上代数空间的态射。
设 $\mathcal{F}$ 为有限型拟凝聚 $\mathcal{O}_X$-模，其概形论支集为
$Z \subset X$。若 $f$ 平坦，则 $f^{-1}(Z)$ 是 $f^*\mathcal{F}$ 的概形论支集。
\end{lemma}

\begin{proof}
利用引理 \ref{spaces-morphisms-lemma-scheme-theoretic-support} 给出的概形论支集刻画，
以及引理 \ref{spaces-morphisms-lemma-flat-local} 中用 \'etale 覆盖给出的平坦态射刻画，
可将问题归结为概形情形，即《态射》引理
\ref{morphisms-lemma-flat-pullback-support}。
\end{proof}

\begin{lemma}
\label{spaces-morphisms-lemma-flat-morphism-scheme-theoretically-dense-open}
设 $S$ 为概形。
设 $f : X \to Y$ 为 $S$ 上代数空间的平坦态射。
设 $V \to Y$ 为拟紧开浸入。若 $V$ 在 $Y$ 中概形论稠密，
则 $f^{-1}V$ 在 $X$ 中概形论稠密。
\end{lemma}

\begin{proof}
利用引理 \ref{spaces-morphisms-lemma-scheme-theoretically-dense} 中对概形论稠密开集的刻画，
以及引理 \ref{spaces-morphisms-lemma-flat-local} 中用 \'etale 覆盖给出的平坦态射刻画，
可将问题归结为概形情形，即《态射》引理
\ref{morphisms-lemma-flat-morphism-scheme-theoretically-dense-open}。
\end{proof}

\begin{lemma}
\label{spaces-morphisms-lemma-flat-base-change-scheme-theoretic-image}
设 $S$ 为概形。设 $f : X \to Y$ 为 $S$ 上代数空间的平坦态射。
设 $g : V \to Y$ 为代数空间的拟紧态射。设 $Z \subset Y$ 为 $g$
的概形论像，并设 $Z' \subset X$ 为基变换 $V \times_Y X \to X$
的概形论像。则 $Z' = f^{-1}Z$。
\end{lemma}

\begin{proof}
取满 \'etale 态射 $Y' \to Y$，使得 $Y'$ 是仿射概形的不交并
（《空间的性质》引理
\ref{spaces-properties-lemma-cover-by-union-affines}）。
取满 \'etale 态射 $X' \to X \times_Y Y'$，使得 $X'$ 是仿射概形的不交并。
由引理 \ref{spaces-morphisms-lemma-flat-local}，态射 $X' \to Y'$ 平坦。令
$V' = V \times_Y Y'$。由引理
\ref{spaces-morphisms-lemma-quasi-compact-scheme-theoretic-image}，$Z$ 在 $Y'$ 中的原像是
$V' \to Y'$ 的概形论像，而 $Z'$ 在 $X'$ 中的原像是
$V' \times_{Y'} X' \to X'$ 的概形论像。由于 $X' \to X$ 满且 \'etale，
只需对态射 $X' \to Y'$ 与 $V' \to Y'$ 证明结论。因此可以假设
$X$ 与 $Y$ 是仿射概形。此时 $V$ 是拟紧代数空间。取仿射概形 $W$
及满 \'etale 态射 $W \to V$（《空间的性质》引理
\ref{spaces-properties-lemma-quasi-compact-affine-cover}）。显然，
$V \to Y$ 的概形论像与 $W \to Y$ 的概形论像一致；对于
$V \times_Y X \to Y$ 与 $W \times_Y X \to X$ 也同样成立。
因此问题归结为概形情形，即《态射》引理
\ref{morphisms-lemma-flat-base-change-scheme-theoretic-image}。
\end{proof}





\section{平坦模}
\label{spaces-morphisms-section-flat-modules}

\noindent
本节定义模在一点处平坦的含义。为此要用到代数空间的小 \'etale 站点
$X_\etale$ 上层的茎这一概念，见《空间的性质》定义
\ref{spaces-properties-definition-stalk}。

\begin{lemma}
\label{spaces-morphisms-lemma-flat-at-point}
设 $S$ 为概形。设 $f : X \to Y$ 为 $S$ 上代数空间的态射。
设 $\mathcal{F}$ 为 $X$ 上的拟凝聚层。设 $x \in |X|$。下列条件等价：
\begin{enumerate}
\item 对某个交换图
$$
\xymatrix{
U \ar[d]_a \ar[r]_h & V \ar[d]^b \\
X \ar[r]^f & Y
}
$$
其中 $U$ 与 $V$ 为概形，$a, b$ 为 \'etale 态射，且 $u \in U$ 映到 $x$，
模 $a^*\mathcal{F}$ 在 $u$ 处相对于 $V$ 平坦；
\item 茎 $\mathcal{F}_{\overline{x}}$ 在 \'etale 局部环
$\mathcal{O}_{Y, \overline{y}}$ 上平坦，其中 $\overline{x}$ 是位于 $x$
上方的任意几何点，且 $\overline{y} = f \circ \overline{x}$。
\end{enumerate}
\end{lemma}

\begin{proof}
% P08-E239: frozen source says "geometric proof" rather than "geometric point".
在此证明中，我们固定 $x$ 上方的几何证明
$\overline{x} : \Spec(k) \to X$，并将其在 $Y$ 中的像记为
$\overline{y} = f \circ \overline{x}$。给定 (1) 中的图，可以找到位于 $u$
上方的几何点 $\overline{u} : \Spec(k) \to U$，使得
$\overline{x} = a \circ \overline{u}$，见《空间的性质》引理
\ref{spaces-properties-lemma-geometric-lift-to-usual}。
令 $\overline{v} = h \circ \overline{u}$，其像为 $v \in V$。我们知道
$$
\mathcal{O}_{X, \overline{x}} = \mathcal{O}_{U, u}^{sh}
\quad\text{且}\quad
\mathcal{O}_{Y, \overline{y}} = \mathcal{O}_{V, v}^{sh}
$$
见《空间的性质》引理
\ref{spaces-properties-lemma-describe-etale-local-ring}。于是得到交换图
$$
\xymatrix{
\mathcal{O}_{U, u} \ar[r] &
\mathcal{O}_{X, \overline{x}} \\
\mathcal{O}_{V, v} \ar[u] \ar[r] &
\mathcal{O}_{Y, \overline{y}} \ar[u]
}
$$
，其中各对象为局部环。最后，由《空间的性质》引理
\ref{spaces-properties-lemma-stalk-quasi-coherent}，有
$$
\mathcal{F}_{\overline{x}} =
(a^*\mathcal{F})_u \otimes_{\mathcal{O}_{U, u}}
\mathcal{O}_{X, \overline{x}}
$$
。因此《代数》引理 \ref{algebra-lemma-flatness-descends-more-general}
告诉我们，$(a^*\mathcal{F})_u$ 在 $\mathcal{O}_{V, v}$ 上平坦，当且仅当
$\mathcal{F}_{\overline{x}}$ 在 $\mathcal{O}_{V, v}$ 上平坦。
故结论由《平坦性的进一步内容》引理
\ref{flat-lemma-flat-up-down-henselization} 得出。
\end{proof}

\begin{definition}
\label{spaces-morphisms-definition-flat-module}
设 $S$ 为概形。
设 $f : X \to Y$ 为 $S$ 上代数空间的态射。
设 $\mathcal{F}$ 为 $X$ 上的拟凝聚层。
\begin{enumerate}
\item 设 $x \in |X|$。若引理 \ref{spaces-morphisms-lemma-flat-at-point} 中的等价条件成立，
则称 $\mathcal{F}$ {\it 在 $x$ 处相对于 $Y$ 平坦}。
\item 若 $\mathcal{F}$ 在每个 $x \in |X|$ 处相对于 $Y$ 平坦，
则称 $\mathcal{F}$ {\it 在 $Y$ 上平坦}。
\end{enumerate}
\end{definition}

\noindent
作出定义后，我们得到必需的基变换引理。该引理说明，
形成拟凝聚层的平坦点集与平坦基变换可交换。

\begin{lemma}
\label{spaces-morphisms-lemma-base-change-module-flat}
设 $S$ 为概形。设
$$
\xymatrix{
X' \ar[d]_{f'} \ar[r]_{g'} & X \ar[d]^f \\
Y' \ar[r]^g & Y
}
$$
为 $S$ 上代数空间的 Cartesian 图。设 $x' \in |X'|$ 的像为 $x \in |X|$。
设 $\mathcal{F}$ 为 $X$ 上的拟凝聚层，并记
$\mathcal{F}' = (g')^*\mathcal{F}$。
\begin{enumerate}
\item 若 $\mathcal{F}$ 在 $x$ 处相对于 $Y$ 平坦，则 $\mathcal{F}'$
在 $x'$ 处相对于 $Y'$ 平坦。
\item 若 $g$ 在 $f'(x')$ 处平坦，且 $\mathcal{F}'$ 在 $x'$ 处相对于
$Y'$ 平坦，则 $\mathcal{F}$ 在 $x$ 处相对于 $Y$ 平坦。
\end{enumerate}
特别地，若 $\mathcal{F}$ 在 $Y$ 上平坦，则 $\mathcal{F}'$ 在 $Y'$ 上平坦。
\end{lemma}

\begin{proof}
取概形 $V$ 及满 \'etale 态射 $V \to Y$。
取概形 $U$ 及满 \'etale 态射 $U \to V \times_Y X$。
取概形 $V'$ 及满 \'etale 态射 $V' \to V \times_Y Y'$。
于是 $U' = V' \times_V U$ 是概形，并带有满 \'etale 态射
$U' = V' \times_V U \to Y' \times_Y X = X'$。取 $u' \in U'$ 映到
$x' \in |X'|$。于是可以用 $\mathcal{F}'|_{U'}$ 在 $u'$ 处相对于 $V'$
的平坦性，检验 $\mathcal{F}'$ 在 $x'$ 处相对于 $Y'$ 的平坦性。
故本引理由《态射的进一步内容》引理
\ref{more-morphisms-lemma-flat-locus-base-change} 得出。
\end{proof}

\noindent
下述引理用模来讨论平坦态射的“复合”。它还说明平坦性满足一种自上而下的下降。

\begin{lemma}
\label{spaces-morphisms-lemma-composition-module-flat}
设 $S$ 为概形。设 $X \to Y \to Z$ 为 $S$ 上代数空间的态射。
设 $\mathcal{F}$ 为 $X$ 上的拟凝聚层。设 $x \in |X|$ 的像为 $y \in |Y|$。
\begin{enumerate}
\item 若 $\mathcal{F}$ 在 $x$ 处相对于 $Y$ 平坦，且 $Y$ 在 $y$ 处相对于
$Z$ 平坦，则 $\mathcal{F}$ 在 $x$ 处相对于 $Z$ 平坦。
\item 设 $x : \Spec(K) \to X$ 为 $x$ 的一个代表。若
\begin{enumerate}
\item $\mathcal{F}$ 在 $x$ 处相对于 $Y$ 平坦；
\item $x^*\mathcal{F} \not = 0$；
\item $\mathcal{F}$ 在 $x$ 处相对于 $Z$ 平坦；
\end{enumerate}
则 $Y$ 在 $y$ 处相对于 $Z$ 平坦。
\item 设 $\overline{x}$ 为 $X$ 的一个几何点，位于 $x$ 上方，
并在 $Y$ 中具有像 $\overline{y}$。若 $\mathcal{F}_{\overline{x}}$ 是
忠实平坦 $\mathcal{O}_{Y, \overline{y}}$-模，且 $\mathcal{F}$ 在 $x$
处相对于 $Z$ 平坦，则 $Y$ 在 $y$ 处相对于 $Z$ 平坦。
\end{enumerate}
\end{lemma}

\begin{proof}
取 (3) 中的 $\overline{x}$ 与 $\overline{y}$，并以 $\overline{z}$ 表示所诱导的
$Z$ 的几何点。利用引理 \ref{spaces-morphisms-lemma-flat-at-point} 与
\ref{spaces-morphisms-lemma-flat-at-point-etale-local-rings} 中的平坦性刻画，本引理归结为
关于局部环映射
$\mathcal{O}_{Z, \overline{z}} \to \mathcal{O}_{Y, \overline{y}}$
和模 $\mathcal{F}_{\overline{x}}$ 的纯代数问题。(1) 由《代数》引理
\ref{algebra-lemma-composition-flat} 得出。注意条件 (2)(b) 保证
$\mathcal{F}_{\overline{x}}/
\mathfrak m_{\overline{y}} \mathcal{F}_{\overline{x}}$
非零。因此 (2)(a) $+$ (2)(b) 蕴含 $\mathcal{F}_{\overline{x}}$ 是忠实平坦
$\mathcal{O}_{Y, \overline{y}}$-模，见《代数》引理
\ref{algebra-lemma-ff}。故 (2) 是 (3) 的特例。最后，(3) 由《代数》引理
\ref{algebra-lemma-flat-permanence} 得出。
\end{proof}

\noindent
有时基变换发生在“上层”。下面给出精确表述。

\begin{lemma}
\label{spaces-morphisms-lemma-flat-permanence}
设 $S$ 为概形。设 $f : X \to Y$、$g : Y \to Z$ 为 $S$ 上代数空间的态射。
设 $\mathcal{G}$ 为 $Y$ 上的拟凝聚层。设 $x \in |X|$ 的像为 $y \in |Y|$。
若 $f$ 在 $x$ 处平坦，则
$$
\mathcal{G}\text{ 相对于 }Z\text{ 在此点平坦：}y
\Leftrightarrow
f^*\mathcal{G}\text{ 相对于 }Z\text{ 在此点平坦：}x.
$$
特别地，若 $f$ 满且平坦，则 $\mathcal{G}$ 在 $Z$ 上平坦，
当且仅当 $f^*\mathcal{G}$ 在 $Z$ 上平坦。
\end{lemma}

\begin{proof}
取 $X$ 的几何点 $\overline{x}$，并分别以 $\overline{y}$、$\overline{z}$
表示它在 $Y$、$Z$ 中的像。利用引理 \ref{spaces-morphisms-lemma-flat-at-point} 与
\ref{spaces-morphisms-lemma-flat-at-point-etale-local-rings} 中的平坦性刻画，
以及《空间的性质》引理
\ref{spaces-properties-lemma-stalk-pullback-quasi-coherent} 对
$f^*\mathcal{G}$ 在 $\overline{x}$ 处之茎的描述，本引理归结为关于局部环映射
$\mathcal{O}_{Z, \overline{z}} \to \mathcal{O}_{Y, \overline{y}}
\to \mathcal{O}_{X, \overline{x}}$
和模 $\mathcal{G}_{\overline{y}}$ 的纯代数问题。该代数命题即《代数》引理
\ref{algebra-lemma-flatness-descends-more-general}。
\end{proof}

\begin{lemma}
\label{spaces-morphisms-lemma-pf-flat-module-open}
设 $S$ 为概形。
设 $f : X \to Y$ 为 $S$ 上代数空间的态射。
设 $\mathcal{F}$ 为拟凝聚 $\mathcal{O}_X$-模。假设 $f$ 局部有限表示，
$\mathcal{F}$ 有限型，$X = \text{Supp}(\mathcal{F})$，且
$\mathcal{F}$ 在 $Y$ 上平坦。则 $f$ 普遍开。
\end{lemma}

\begin{proof}
取满 \'etale 态射 $\varphi : V \to Y$，其中 $V$ 为概形。
取满 \'etale 态射 $U \to V \times_Y X$，其中 $U$ 为概形。
于是只需对 $U \to V$ 及拟凝聚 $\mathcal{O}_V$-模
$\varphi^*\mathcal{F}$ 证明引理。因此本引理归结为概形情形，见
《态射》引理 \ref{morphisms-lemma-pf-flat-module-open}。
\end{proof}






\section{泛平坦性}
\label{spaces-morphisms-section-generic-flatness}

\noindent
本节是《态射》第 \ref{morphisms-section-generic-flatness} 节的对应版本。

\begin{proposition}
\label{spaces-morphisms-proposition-generic-flatness-reduced}
设 $S$ 为概形。
设 $f : X \to Y$ 为 $S$ 上代数空间的态射。
设 $\mathcal{F}$ 为拟凝聚 $\mathcal{O}_X$-模层。假设
\begin{enumerate}
\item $Y$ 约化；
\item $f$ 有限型；
\item $\mathcal{F}$ 是有限型 $\mathcal{O}_X$-模。
\end{enumerate}
则存在稠密开子空间 $W \subset Y$，使得 $f$ 的基变换 $X_W \to W$
平坦、局部有限表示且拟紧，并使得 $\mathcal{F}|_{X_W}$ 在 $W$ 上平坦，
且在 $\mathcal{O}_{X_W}$ 上有限表示。
\end{proposition}

\begin{proof}
设 $V$ 为概形，并设 $V \to Y$ 为满 \'etale 态射。令
$X_V = V \times_Y X$，并以 $\mathcal{F}_V$ 表示 $\mathcal{F}$ 在 $X_V$
上的限制。假设结论对态射 $X_V \to V$ 与层 $\mathcal{F}_V$ 成立。
于是存在开子概形 $V' \subset V$，使得 $X_{V'} \to V'$ 平坦且有限表示，
并且 $\mathcal{F}_{V'}$ 是在 $V'$ 上平坦的有限表示
$\mathcal{O}_{X_{V'}}$-模。令 $W \subset Y$ 为 \'etale 态射 $V' \to Y$
的像，见《空间的性质》引理
\ref{spaces-properties-lemma-etale-image-open}。于是 $V' \to W$ 是满
\'etale 态射，故由引理 \ref{spaces-morphisms-lemma-finite-presentation-local}、
\ref{spaces-morphisms-lemma-flat-local} 与 \ref{spaces-morphisms-lemma-quasi-compact-local}，可见
$X_W \to W$ 平坦、局部有限表示且拟紧。由《空间的性质》第
\ref{spaces-properties-section-properties-modules} 节中的讨论，
可见 $\mathcal{F}_W$ 作为 $\mathcal{O}_{X_W}$-模有限表示；
由引理 \ref{spaces-morphisms-lemma-base-change-module-flat}，可见 $\mathcal{F}_W$
在 $W$ 上平坦。这一论证把命题归结为 $Y$ 是概形的情形。

\medskip\noindent
假设可以在 $Y$ 为仿射概形时证明命题。设 $f : X \to Y$ 为 $S$ 上
代数空间的有限型态射，其中 $Y$ 为概形；并设 $\mathcal{F}$ 为有限型拟凝聚
$\mathcal{O}_X$-模。取仿射开覆盖 $Y = \bigcup V_j$。按假设，可以找到
稠密开集 $W_j \subset V_j$，使得 $X_{W_j} \to W_j$ 平坦、局部有限表示且拟紧，
并使得 $\mathcal{F}|_{X_{W_j}}$ 在 $W_j$ 上平坦，且作为
$\mathcal{O}_{X_{W_j}}$-模有限表示。此时只需取 $W = \bigcup W_j$，
结论即得。因此命题归结为 $Y$ 是仿射概形的情形。

\medskip\noindent
设 $Y$ 为 $S$ 上的仿射概形，设 $f : X \to Y$ 为 $S$ 上代数空间的有限型态射，
并设 $\mathcal{F}$ 为有限型拟凝聚 $\mathcal{O}_X$-模。由于 $f$ 有限型，
它拟紧，故 $X$ 拟紧。因此可取仿射概形 $U$ 及满 \'etale 态射 $U \to X$，
见《空间的性质》引理
\ref{spaces-properties-lemma-quasi-compact-affine-cover}。注意，$U \to Y$
有限型（在此情形下，这正是 $f$ 有限型的含义）。因此可应用《态射》命题
\ref{morphisms-proposition-generic-flatness-reduced}，从而存在稠密开集
$W \subset Y$，使得 $U_W \to W$ 平坦且有限表示，并使得
$\mathcal{F}|_{U_W}$ 在 $W$ 上平坦，且作为 $\mathcal{O}_{U_W}$-模有限表示。
按我们的定义，这意味着 $f$ 的基变换 $X_W \to W$ 平坦、局部有限表示且拟紧，
并且 $\mathcal{F}|_{X_W}$ 在 $W$ 上平坦，且在 $\mathcal{O}_{X_W}$ 上有限表示。
\end{proof}

\noindent
不能把上面引理的结论加强为要求 $X_W \to W$ 有限表示，因为
$\mathbf{A}^1_{\mathbf{Q}}/\mathbf{Z} \to \Spec(\mathbf{Q})$ 给出了反例。
问题在于对角态射 $\Delta_{X/Y}$ 可能不拟紧，即 $f$ 可能不拟分离。
显然，这也是唯一的问题。

\begin{proposition}
\label{spaces-morphisms-proposition-generic-flatness-reduced-quasi-separated}
设 $S$ 为概形。
设 $f : X \to Y$ 为 $S$ 上代数空间的态射。
设 $\mathcal{F}$ 为拟凝聚 $\mathcal{O}_X$-模层。假设
\begin{enumerate}
\item $Y$ 约化；
\item $f$ 拟分离；
\item $f$ 有限型；
\item $\mathcal{F}$ 是有限型 $\mathcal{O}_X$-模。
\end{enumerate}
则存在稠密开子空间 $W \subset Y$，使得 $f$ 的基变换 $X_W \to W$
平坦且有限表示，并使得 $\mathcal{F}|_{X_W}$ 在 $W$ 上平坦，
且在 $\mathcal{O}_{X_W}$ 上有限表示。
\end{proposition}

\begin{proof}
这立即由命题 \ref{spaces-morphisms-proposition-generic-flatness-reduced} 以及下述事实得出：
“有限表示” $=$ “局部有限表示” $+$ “拟紧” $+$ “拟分离”。
\end{proof}















\section{相对维数}
\label{spaces-morphisms-section-relative-dimension}

\noindent
本节定义代数空间态射在一点处的相对维数，以及若干密切相关的性质。

\begin{definition}
\label{spaces-morphisms-definition-dimension-fibre}
设 $S$ 为概形。
设 $f : X \to Y$ 为 $S$ 上代数空间的态射。
设 $x \in |X|$。设 $d, r \in \{0, 1, 2, \ldots, \infty\}$。
\begin{enumerate}
\item 若引理 \ref{spaces-morphisms-lemma-local-source-target-at-point} 中的等价条件对
《下降》引理 \ref{descent-lemma-dimension-local-ring-fibre} 所描述的性质
$\mathcal{P}_d$ 成立，则称{\it $f$ 在 $x$ 处之纤维的局部环维数}为 $d$。
\item 若引理 \ref{spaces-morphisms-lemma-local-source-target-at-point} 中的等价条件对
《下降》引理 \ref{descent-lemma-transcendence-degree-at-point} 所描述的性质
$\mathcal{P}_r$ 成立，则称{\it $x/f(x)$ 的超越次数}为 $r$。
\item 若引理 \ref{spaces-morphisms-lemma-local-source-target-at-point} 中的等价条件对
《下降》引理 \ref{descent-lemma-dimension-at-point} 所描述的性质
$\mathcal{P}_d$ 成立，则称{\it $f$ 在 $x$ 处具有相对维数 $d$}。
\end{enumerate}
\end{definition}

\noindent
把其含义明确写出。具体地，取引理
\ref{spaces-morphisms-lemma-local-source-target-at-point} 中的图
$$
\xymatrix{
U \ar[d]_a \ar[r]_h & V \ar[d]^b \\
X \ar[r]^f & Y
}
\quad\quad
\xymatrix{
u \ar[d] \ar[r] & v \ar[d] \\
x \ar[r] & y
}
$$
。则有
$$
\begin{matrix}
\text{相对维数：态射 }f\text{ 于点 }x & = &
\dim_u (U_v) \\
\text{纤维局部环的维数：态射 }f\text{ 于点 }x & = &
\dim(\mathcal{O}_{U_v, u})\\
\text{超越次数：}x/f(x) & = &
\text{trdeg}_{\kappa(v)}(\kappa(u))
\end{matrix}
$$
注意，若 $Y = \Spec(k)$ 是域的谱，则 $X/Y$ 在 $x$ 处的相对维数等于
$\dim_x(X)$；$x/f(x)$ 的超越次数就是在 $k$ 上的超越次数；而 $f$ 在 $x$
处之纤维的局部环维数就是 $x$ 处局部环的维数。换言之，在这种情形下，
相对概念变成绝对概念。

\begin{definition}
\label{spaces-morphisms-definition-relative-dimension}
设 $S$ 为概形。
设 $f : X \to Y$ 为 $S$ 上代数空间的态射。
设 $d \in \{0, 1, 2, \ldots\}$。
\begin{enumerate}
\item 若 $f$ 在所有 $x \in |X|$ 处都具有相对维数 $\leq d$，
则称 $f$ {\it 具有相对维数 $\leq d$}。
\item 若 $f$ 在所有 $x \in |X|$ 处都具有相对维数 $d$，
则称 $f$ {\it 具有相对维数 $d$}。
\end{enumerate}
\end{definition}

\noindent
粗略地说，具有{\it 等于} $d$ 的相对维数意味着所有非空纤维都是维数为
$d$ 的等维空间。

\begin{lemma}
\label{spaces-morphisms-lemma-compare-tr-deg}
设 $S$ 为概形。设 $X \to Y \to Z$ 为 $S$ 上代数空间的态射。
设 $x \in |X|$ 的像分别为 $y \in |Y|$、$z \in |Z|$。
假设 $X \to Y$ 局部拟有限，且 $Y \to Z$ 局部有限型。
则 $x/z$ 的超越次数等于 $y/z$ 的超越次数。
\end{lemma}

\begin{proof}
可以选取交换图
$$
\xymatrix{
U \ar[d] \ar[r] & V \ar[d] \ar[r] & W \ar[d] \\
X \ar[r] & Y \ar[r] & Z
}
\quad\quad
\xymatrix{
u \ar[d] \ar[r] & v \ar[d] \ar[r] & w \ar[d] \\
x \ar[r] & y \ar[r] & z
}
$$
其中 $U, V, W$ 为概形，竖直箭头为 \'etale 态射。按定义，态射
$U \to V$ 局部拟有限，这蕴含 $\kappa(v) \subset \kappa(u)$ 有限，见
《态射》引理 \ref{morphisms-lemma-residue-field-quasi-finite}。
故结论显然。
\end{proof}

\begin{lemma}
\label{spaces-morphisms-lemma-jacobson-finite-type-points}
设 $S$ 为概形。设 $f : X \to Y$ 为 $S$ 上代数空间的态射。
若 $f$ 局部有限型，$Y$ 是 Jacobson 的（《空间的性质》评注
\ref{spaces-properties-remark-list-properties-local-etale-topology}），
且 $x \in |X|$ 是 $X$ 的有限型点，则 $x/f(x)$ 的超越次数为 $0$。
\end{lemma}

\begin{proof}
取概形 $V$ 及满 \'etale 态射 $V \to Y$。
取概形 $U$ 及满 \'etale 态射 $U \to X \times_Y V$。
由引理 \ref{spaces-morphisms-lemma-finite-type-points-surjective-morphism}，
可以找到映到 $x$ 的有限型点 $u \in U$。缩小 $U$ 后，可以假设
$u \in U$ 是闭点（《态射》引理
\ref{morphisms-lemma-identify-finite-type-points}）。
设 $v \in V$ 为 $u$ 的像。由《态射》引理
\ref{morphisms-lemma-jacobson-finite-type-points}，扩张
$\kappa(u)/\kappa(v)$ 有限。这就完成证明。
\end{proof}

\begin{lemma}
\label{spaces-morphisms-lemma-rel-dimension-dimension}
设 $S$ 为概形。设 $f : X \to Y$ 为 $S$ 上局部 Noether 代数空间的态射，
且它平坦、局部有限型并具有相对维数 $d$。对 $|X|$ 中每个点 $x$，
若其在 $|Y|$ 中的像为 $y$，则 $\dim_x(X) = \dim_y(Y) + d$。
\end{lemma}

\begin{proof}
由代数空间在一点处之维数的定义（《空间的性质》定义
\ref{spaces-properties-definition-dimension-at-point}）以及具有相对维数 $d$
的定义，问题归结为概形的相应命题（《态射》引理
\ref{morphisms-lemma-rel-dimension-dimension}）。
\end{proof}







\section{态射与纤维维数}
\label{spaces-morphisms-section-dimension-fibres}

\noindent
本节是《态射》第 \ref{morphisms-section-dimension-fibres} 节的对应版本。
由于我们没有代数空间在点处的局部环，本节的表述稍显别扭。

\begin{lemma}
\label{spaces-morphisms-lemma-dimension-fibre-at-a-point}
设 $S$ 为概形。
设 $f : X \to Y$ 为 $S$ 上代数空间的态射。
设 $x \in |X|$。假设 $f$ 局部有限型。则有
$$
\begin{matrix}
\text{相对维数：态射 }f\text{ 于点 }x \\
= \\
\text{纤维局部环的维数：态射 }f\text{ 于点 }x \\
+ \\
\text{超越次数：}x/f(x)
\end{matrix}
$$
，其中记号如定义 \ref{spaces-morphisms-definition-dimension-fibre} 所示。
\end{lemma}

\begin{proof}
这立即由《态射》引理
\ref{morphisms-lemma-dimension-fibre-at-a-point} 得出：将其应用于
$h : U \to V$ 与 $u \in U$，其中这些对象如引理
\ref{spaces-morphisms-lemma-local-source-target-at-point} 所示。
\end{proof}

\begin{lemma}
\label{spaces-morphisms-lemma-dimension-fibre-at-a-point-additive}
设 $S$ 为概形。
设 $f : X \to Y$ 与 $g : Y \to Z$ 为 $S$ 上代数空间的态射。
设 $x \in |X|$，并令 $y = f(x)$。假设 $f$ 与 $g$ 局部有限型。则
\begin{enumerate}
\item
$$
\begin{matrix}
\text{相对维数：态射 }g \circ f\text{ 于点 }x \\
\leq \\
\text{相对维数：态射 }f\text{ 于点 }x \\
+ \\
\text{相对维数：态射 }g\text{ 于点 }y
\end{matrix}
$$
\item 若存在来自某个域之谱的态射 $\Spec(k) \to Z$，它属于
$g(f(x)) = g(y)$ 的等价类，且态射 $X_k \to Y_k$ 在 $x$ 处平坦，
则 (1) 中取等号；例如，当 $f$ 在 $x$ 处平坦时即如此；
\item
$$
\begin{matrix}
\text{超越次数：}x/g(f(x)) \\
= \\
\text{超越次数：}x/f(x) \\
+ \\
\text{超越次数：}f(x)/g(f(x))
\end{matrix}
$$
\end{enumerate}
\end{lemma}

\begin{proof}
取图
$$
\xymatrix{
U \ar[d] \ar[r] & V \ar[d] \ar[r] & W \ar[d] \\
X \ar[r] & Y \ar[r] & Z
}
$$
其中 $U, V, W$ 为概形，竖直箭头满且 \'etale。（见《空间》引理
\ref{spaces-lemma-lift-morphism-presentations}。）取映到 $x$ 的 $u \in U$。
分别以 $v, w$ 表示 $u$ 在 $V, W$ 中的像。对顶行及点 $u, v, w$
应用《态射》引理
\ref{morphisms-lemma-dimension-fibre-at-a-point-additive}。细节从略。
\end{proof}

\begin{lemma}
\label{spaces-morphisms-lemma-dimension-fibre-after-base-change}
设 $S$ 为概形。设
$$
\xymatrix{
X' \ar[r]_{g'} \ar[d]_{f'} & X \ar[d]^f \\
Y' \ar[r]^g & Y
}
$$
为 $S$ 上代数空间的纤维积图。设 $x' \in |X'|$。令 $x = g'(x')$。
假设 $f$ 局部有限型。则
\begin{enumerate}
\item
$$
\begin{matrix}
\text{相对维数：态射 }f\text{ 于点 }x \\
= \\
\text{相对维数：态射 }f'\text{ 于点 }x'
\end{matrix}
$$
\item 有
$$
\begin{matrix}
\text{纤维局部环的维数：态射 }f'\text{ 于点 }x' \\
- \\
\text{纤维局部环的维数：态射 }f\text{ 于点 }x \\
= \\
\text{超越次数：}x/f(x) \\
- \\
\text{超越次数：}x'/f'(x')
\end{matrix}
$$
且共同的值为 $\geq 0$；
\item 给定 $x$ 以及映到同一 $y \in |Y|$ 的 $y' \in |Y'|$，
存在 $x'$ 的一个选取，使得 (2) 中的整数为 $0$。
\end{enumerate}
\end{lemma}

\begin{proof}
取满 \'etale 态射 $V \to Y$，其中 $V$ 为概形。
取满 \'etale 态射 $U \to V \times_Y X$，其中 $U$ 为概形。
取满 \'etale 态射 $V' \to V \times_Y Y'$，其中 $V'$ 为概形。
令 $U' = V' \times_V U$。于是所诱导的态射 $U' \to X'$ 也满且
\'etale（论证从略）。取映到 $x'$ 的 $u' \in U'$。此时，对涉及
$U', U, V', V$ 的概形图及点 $u'$ 应用《态射》引理
\ref{morphisms-lemma-dimension-fibre-after-base-change}，即可得到 (1) 与 (2)。
为证明 (3)，先取映到 $y$ 的 $v \in V$。再利用《空间的性质》引理
\ref{spaces-properties-lemma-points-cartesian}，可以选取映到 $y'$ 与 $v$
的 $v' \in V'$，以及映到 $x$ 与 $v$ 的 $u \in U$。最后，由《态射》引理
\ref{morphisms-lemma-dimension-fibre-after-base-change}，可以选取映到
$v'$ 与 $u$ 的 $u' \in U'$，使该整数为零。于是取 $u'$ 的像
$x' \in |X'|$ 即可。
\end{proof}

\begin{lemma}
\label{spaces-morphisms-lemma-openness-bounded-dimension-fibres}
设 $S$ 为概形。
设 $f : X \to Y$ 为 $S$ 上代数空间的态射。
设 $n \geq 0$。假设 $f$ 局部有限型。集合
$$
W_n = \{x \in |X|
\text{ 满足相对维数：态射 }f\text{ 于点 } x \leq n\}
$$
在 $|X|$ 中是开的。
\end{lemma}

\begin{proof}
取图
$$
\xymatrix{
U \ar[r]_h \ar[d]_a & V \ar[d] \\
X \ar[r] & Y
}
$$
其中 $U$ 与 $V$ 为概形，竖直箭头满且 \'etale，见《空间》引理
\ref{spaces-lemma-lift-morphism-presentations}。由《态射》引理
\ref{morphisms-lemma-openness-bounded-dimension-fibres}，$h$ 具有相对维数
$\leq n$ 的点集 $U_n$ 在 $U$ 中是开的。由代数空间态射在点处之相对维数的
定义，可见 $U_n = a^{-1}(W_n)$。本引理由 $|X|$ 上拓扑的定义得出。
\end{proof}

\begin{lemma}
\label{spaces-morphisms-lemma-openness-bounded-dimension-fibres-finite-presentation}
设 $S$ 为概形。
设 $f : X \to Y$ 为 $S$ 上代数空间的态射。
设 $n \geq 0$。假设 $f$ 局部有限表示。开集
$$
W_n = \{x \in |X|
\text{ 满足相对维数：态射 }f\text{ 于点 } x \leq n\}
$$
（引理 \ref{spaces-morphisms-lemma-openness-bounded-dimension-fibres}）在 $|X|$ 中逆紧。
（见《拓扑》定义 \ref{topology-definition-quasi-compact}。）
\end{lemma}

\begin{proof}
取图
$$
\xymatrix{
U \ar[r]_h \ar[d]_a & V \ar[d] \\
X \ar[r] & Y
}
$$
其中 $U$ 与 $V$ 为概形，竖直箭头满且 \'etale，见《空间》引理
\ref{spaces-lemma-lift-morphism-presentations}。在引理
\ref{spaces-morphisms-lemma-openness-bounded-dimension-fibres} 的证明中，我们已看到
$a^{-1}(W_n) = U_n$ 是态射 $h$ 的相应点集。由《态射》引理
\ref{morphisms-lemma-openness-bounded-dimension-fibres-finite-presentation}，
可见 $U_n$ 在 $U$ 中逆紧。本引理由 $|X|$ 上拓扑的定义得出；
可与《空间的性质》引理
\ref{spaces-properties-lemma-space-locally-quasi-compact} 及其证明比较。
\end{proof}

\begin{lemma}
\label{spaces-morphisms-lemma-locally-quasi-finite-rel-dimension-0}
设 $S$ 为概形。
设 $f : X \to Y$ 为 $S$ 上代数空间的态射。
假设 $f$ 局部有限型。则 $f$ 局部拟有限，当且仅当 $f$ 在每个
$x \in |X|$ 处具有相对维数 $0$。
\end{lemma}

\begin{proof}
取图
$$
\xymatrix{
U \ar[r]_h \ar[d]_a & V \ar[d] \\
X \ar[r] & Y
}
$$
其中 $U$ 与 $V$ 为概形，竖直箭头满且 \'etale，见《空间》引理
\ref{spaces-lemma-lift-morphism-presentations}。定义蕴含：$h$ 局部拟有限，
当且仅当 $f$ 局部拟有限；而 $f$ 在所有 $x \in |X|$ 处具有相对维数 $0$，
当且仅当 $h$ 在所有 $u \in U$ 处具有相对维数 $0$。
因此结论由 $h$ 的相应结果得出，即《态射》引理
\ref{morphisms-lemma-locally-quasi-finite-rel-dimension-0}。
\end{proof}

\begin{lemma}
\label{spaces-morphisms-lemma-locally-finite-type-quasi-finite-part}
设 $S$ 为概形。
设 $f : X \to Y$ 为 $S$ 上代数空间的态射。
假设 $f$ 局部有限型。则存在典范开子空间 $X' \subset X$，使得
$f|_{X'} : X' \to Y$ 局部拟有限，并且对任意 $x \in |X|$，
若 $x \not \in |X'|$，则 $f$ 在该点的相对维数为 $\geq 1$。
形成 $X'$ 与任意基变换可交换。
\end{lemma}

\begin{proof}
合用引理 \ref{spaces-morphisms-lemma-openness-bounded-dimension-fibres}、
\ref{spaces-morphisms-lemma-locally-quasi-finite-rel-dimension-0} 与
\ref{spaces-morphisms-lemma-dimension-fibre-after-base-change}。
\end{proof}

\begin{lemma}
\label{spaces-morphisms-lemma-quasi-finite-at-point}
设 $S$ 为概形。考虑 Cartesian 图
$$
\xymatrix{
X \ar[d] & F \ar[l]^p \ar[d] \\
Y & \Spec(k) \ar[l]
}
$$
其中 $X \to Y$ 是 $S$ 上代数空间的局部有限型态射，而 $k$ 是 $S$ 上的域。
设 $z \in |F|$ 满足 $\dim_z(F) = 0$。则以包含 $p(z)$ 的开子空间替换 $X$ 后，
态射
$$
X \longrightarrow Y
$$
局部拟有限。
\end{lemma}

\begin{proof}
令 $X' \subset X$ 为引理 \ref{spaces-morphisms-lemma-locally-finite-type-quasi-finite-part}
给出的、使 $f$ 局部拟有限的开子空间。由于形成 $X'$ 与任意基变换可交换，
可见 $z \in X' \times_Y \Spec(k)$。故引理显然。
\end{proof}





\section{维数公式}
\label{spaces-morphisms-section-dimension-formula}

\noindent
维数公式的对应版本（《态射》引理
\ref{morphisms-lemma-dimension-formula}）表述起来有些棘手，
因为我们必须定义整代数空间（稍后会这样做），还要定义普遍链状代数空间。
不过，下面的版本很直接。

\begin{lemma}
\label{spaces-morphisms-lemma-dimension-formula-general}
设 $S$ 为概形。设 $f : X \to Y$ 为 $S$ 上代数空间的态射。
假设 $Y$ 局部 Noether，且 $f$ 局部有限型。设 $x \in |X|$ 的像为
$y \in |Y|$。则有
\begin{align*}
& \text{局部环维数：空间 }X\text{ 于点 }x \leq \\
& \text{局部环维数：空间 }Y\text{ 于点 }y + E - \\
& \text{超越次数：}x/y
\end{align*}
这里 $E$ 是 $\xi/f(\xi)$ 的超越次数之最大值，其中 $\xi \in |X|$
遍历专门化到 $x$ 且使 $X$ 的局部环维数为 $0$ 的点。
\end{lemma}

\begin{proof}
取仿射概形 $V$、\'etale 态射 $V \to Y$，以及映到 $y$ 的点 $v \in V$。
取仿射概形 $U$、\'etale 态射 $U \to X \times_Y V$，以及点 $u \in U$，
它在 $V$ 中映到 $v$，在 $X$ 中映到 $x$。展开定义
\ref{spaces-morphisms-definition-dimension-fibre} 与《空间的性质》定义
\ref{spaces-properties-definition-dimension-local-ring}，我们要证明
$$
\dim(\mathcal{O}_{U, u}) \leq
\dim(\mathcal{O}_{V, v}) + E - \text{trdeg}_{\kappa(v)}(\kappa(u))
$$
。设 $\xi_U \in U$ 为 $U$ 的某个包含 $u$ 的不可约分支之泛点。
则 $\xi_U$ 映到一点 $\xi \in |X|$，它位于定义量 $E$ 时所用的列表中；
事实上，$E$ 的定义中使用的每个 $\xi$ 都以此方式出现（略去一个小细节）。
特别地，这样的 $\xi$ 只有有限个，因而可以取最大值（即它确实是最大值，
而非上确界）。$\xi$ 在 $f(\xi)$ 上的超越次数是
$\text{trdeg}_{\kappa(\xi_V)}(\kappa(\xi_U))$，其中 $\xi_V \in V$
是 $\xi_U$ 的像。因此本引理由《态射》引理
\ref{morphisms-lemma-dimension-formula-general} 得出。
\end{proof}

\begin{lemma}
\label{spaces-morphisms-lemma-alteration-dimension-general}
设 $S$ 为概形。
设 $f : X \to Y$ 为 $S$ 上代数空间的态射。假设 $Y$ 局部 Noether，
且 $f$ 局部有限型。则
$$
\dim(X) \leq \dim(Y) + E
$$
，其中 $E$ 是 $\xi/f(\xi)$ 的超越次数之上确界，而 $\xi$ 遍历使 $X$
的局部环维数为 $0$ 的点。
\end{lemma}

\begin{proof}
这是引理 \ref{spaces-morphisms-lemma-dimension-formula-general} 与《空间的性质》引理
\ref{spaces-properties-lemma-dimension} 的直接推论。
\end{proof}







\section{Syntomic 态射}
\label{spaces-morphisms-section-syntomic}

\noindent
概形态射的“syntomic”性质在源与靶上都是 \'etale 局部的，见
《下降》评注 \ref{descent-remark-list-local-source-target}。
此性质也在基变换下稳定，并且在靶上是 fpqc 局部的，见《态射》引理
\ref{morphisms-lemma-base-change-syntomic} 与《下降》引理
\ref{descent-lemma-descending-property-syntomic}。因此，由上面的引理
\ref{spaces-morphisms-lemma-local-source-target}，我们可以如下定义代数空间的 syntomic 态射；
当该态射可表时，这与第 \ref{spaces-morphisms-section-representable} 节中已有的概念一致。

\begin{definition}
\label{spaces-morphisms-definition-syntomic}
设 $S$ 为概形。
设 $f : X \to Y$ 为 $S$ 上代数空间的态射。
\begin{enumerate}
\item 若引理 \ref{spaces-morphisms-lemma-local-source-target} 中取
$\mathcal{P} =$“syntomic”时其中的等价条件成立，则称 $f$ {\it syntomic}。
\item 设 $x \in |X|$。若存在 $x$ 的开邻域 $X' \subset X$，使得
$f|_{X'} : X' \to Y$ syntomic，则称 $f$ {\it 在 $x$ 处 syntomic}。
\end{enumerate}
\end{definition}

\begin{lemma}
\label{spaces-morphisms-lemma-composition-syntomic}
Syntomic 态射的复合 syntomic。
\end{lemma}

\begin{proof}
见评注 \ref{spaces-morphisms-remark-composition-P} 和《态射》引理
\ref{morphisms-lemma-composition-syntomic}。
\end{proof}

\begin{lemma}
\label{spaces-morphisms-lemma-base-change-syntomic}
Syntomic 态射的基变换 syntomic。
\end{lemma}

\begin{proof}
见评注 \ref{spaces-morphisms-remark-base-change-P} 和《态射》引理
\ref{morphisms-lemma-base-change-syntomic}。
\end{proof}

\begin{lemma}
\label{spaces-morphisms-lemma-syntomic-local}
设 $S$ 为概形。
设 $f : X \to Y$ 为 $S$ 上代数空间的态射。
下列条件等价：
\begin{enumerate}
\item $f$ syntomic；
\item 对每个 $x \in |X|$，态射 $f$ 在 $x$ 处 syntomic；
\item 对每个概形 $Z$ 及任意态射 $Z \to Y$，态射
$Z \times_Y X \to Z$ syntomic；
\item 对每个仿射概形 $Z$ 及任意态射 $Z \to Y$，态射
$Z \times_Y X \to Z$ syntomic；
\item 存在概形 $V$ 及满 \'etale 态射 $V \to Y$，使得
$V \times_Y X \to V$ 为 syntomic 态射；
\item 存在概形 $U$ 及满 \'etale 态射 $\varphi : U \to X$，使得复合
$f \circ \varphi$ syntomic；
\item 对每个交换图
$$
\xymatrix{
U \ar[d] \ar[r] & V \ar[d] \\
X \ar[r] & Y
}
$$
其中 $U$、$V$ 为概形，竖直箭头为 \'etale 态射，顶部水平箭头 syntomic；
\item 存在交换图
$$
\xymatrix{
U \ar[d] \ar[r] & V \ar[d] \\
X \ar[r] & Y
}
$$
其中 $U$、$V$ 为概形，竖直箭头为 \'etale 态射，且 $U \to X$ 满，
使得顶部水平箭头 syntomic；
\item 存在 Zariski 覆盖 $Y = \bigcup_{i \in I} Y_i$ 以及
$f^{-1}(Y_i) = \bigcup X_{ij}$，使得每个态射
$X_{ij} \to Y_i$ syntomic。
\end{enumerate}
\end{lemma}

\begin{proof}
略。
\end{proof}

\begin{lemma}
\label{spaces-morphisms-lemma-syntomic-locally-finite-presentation}
Syntomic 态射局部有限表示。
\end{lemma}

\begin{proof}
这立即由概形情形得出（《态射》引理
\ref{morphisms-lemma-syntomic-locally-finite-presentation}）。
\end{proof}

\begin{lemma}
\label{spaces-morphisms-lemma-syntomic-flat}
Syntomic 态射平坦。
\end{lemma}

\begin{proof}
这立即由概形情形得出（《态射》引理
\ref{morphisms-lemma-syntomic-flat}）。
\end{proof}

\begin{lemma}
\label{spaces-morphisms-lemma-syntomic-open}
Syntomic 态射普遍开。
\end{lemma}

\begin{proof}
合用引理 \ref{spaces-morphisms-lemma-syntomic-locally-finite-presentation}、
\ref{spaces-morphisms-lemma-syntomic-flat} 与 \ref{spaces-morphisms-lemma-fppf-open}。
\end{proof}





\section{光滑态射}
\label{spaces-morphisms-section-smooth}

\noindent
概形态射的“光滑”性质在源与靶上都是 \'etale 局部的，见
《下降》评注 \ref{descent-remark-list-local-source-target}。
此性质也在基变换下稳定，并且在靶上是 fpqc 局部的，见《态射》引理
\ref{morphisms-lemma-base-change-smooth} 与《下降》引理
\ref{descent-lemma-descending-property-smooth}。因此，由上面的引理
\ref{spaces-morphisms-lemma-local-source-target}，我们可以如下定义代数空间的光滑态射；
当该态射可表时，这与第 \ref{spaces-morphisms-section-representable} 节中已有的概念一致。

\begin{definition}
\label{spaces-morphisms-definition-smooth}
设 $S$ 为概形。
设 $f : X \to Y$ 为 $S$ 上代数空间的态射。
\begin{enumerate}
\item 若引理 \ref{spaces-morphisms-lemma-local-source-target} 中取
$\mathcal{P} =$“光滑”时其中的等价条件成立，则称 $f$ {\it 光滑}。
\item 设 $x \in |X|$。若存在 $x$ 的开邻域 $X' \subset X$，使得
$f|_{X'} : X' \to Y$ 光滑，则称 $f$ {\it 在 $x$ 处光滑}。
\end{enumerate}
\end{definition}

\begin{lemma}
\label{spaces-morphisms-lemma-composition-smooth}
光滑态射的复合光滑。
\end{lemma}

\begin{proof}
见评注 \ref{spaces-morphisms-remark-composition-P} 和《态射》引理
\ref{morphisms-lemma-composition-smooth}。
\end{proof}

\begin{lemma}
\label{spaces-morphisms-lemma-base-change-smooth}
光滑态射的基变换光滑。
\end{lemma}

\begin{proof}
见评注 \ref{spaces-morphisms-remark-base-change-P} 和《态射》引理
\ref{morphisms-lemma-base-change-smooth}。
\end{proof}

\begin{lemma}
\label{spaces-morphisms-lemma-smooth-local}
设 $S$ 为概形。
设 $f : X \to Y$ 为 $S$ 上代数空间的态射。
下列条件等价：
\begin{enumerate}
\item $f$ 光滑；
\item 对每个 $x \in |X|$，态射 $f$ 在 $x$ 处光滑；
\item 对每个概形 $Z$ 及任意态射 $Z \to Y$，态射
$Z \times_Y X \to Z$ 光滑；
\item 对每个仿射概形 $Z$ 及任意态射 $Z \to Y$，态射
$Z \times_Y X \to Z$ 光滑；
\item 存在概形 $V$ 及满 \'etale 态射 $V \to Y$，使得
$V \times_Y X \to V$ 为光滑态射；
\item 存在概形 $U$ 及满 \'etale 态射 $\varphi : U \to X$，使得复合
$f \circ \varphi$ 光滑；
\item 对每个交换图
$$
\xymatrix{
U \ar[d] \ar[r] & V \ar[d] \\
X \ar[r] & Y
}
$$
其中 $U$、$V$ 为概形，竖直箭头为 \'etale 态射，顶部水平箭头光滑；
\item 存在交换图
$$
\xymatrix{
U \ar[d] \ar[r] & V \ar[d] \\
X \ar[r] & Y
}
$$
其中 $U$、$V$ 为概形，竖直箭头为 \'etale 态射，且 $U \to X$ 满，
使得顶部水平箭头光滑；
\item 存在 Zariski 覆盖 $Y = \bigcup_{i \in I} Y_i$ 以及
$f^{-1}(Y_i) = \bigcup X_{ij}$，使得每个态射 $X_{ij} \to Y_i$ 光滑。
\end{enumerate}
\end{lemma}

\begin{proof}
略。
\end{proof}

\begin{lemma}
\label{spaces-morphisms-lemma-smooth-locally-finite-presentation}
代数空间的光滑态射局部有限表示。
\end{lemma}

\begin{proof}
设 $X \to Y$ 为代数空间的光滑态射。按定义，这意味着存在引理
\ref{spaces-morphisms-lemma-local-source-target} 中的图，其中 $h$ 光滑，且竖直箭头 $a$ 满。
由《态射》引理 \ref{morphisms-lemma-smooth-locally-finite-presentation}，
$h$ 局部有限表示。因此按定义 $X \to Y$ 局部有限表示。
\end{proof}

\begin{lemma}
\label{spaces-morphisms-lemma-smooth-locally-finite-type}
代数空间的光滑态射局部有限型。
\end{lemma}

\begin{proof}
合用引理 \ref{spaces-morphisms-lemma-smooth-locally-finite-presentation} 与
\ref{spaces-morphisms-lemma-finite-presentation-finite-type}。
\end{proof}

\begin{lemma}
\label{spaces-morphisms-lemma-smooth-flat}
代数空间的光滑态射平坦。
\end{lemma}

\begin{proof}
设 $X \to Y$ 为代数空间的光滑态射。按定义，这意味着存在引理
\ref{spaces-morphisms-lemma-local-source-target} 中的图，其中 $h$ 光滑，且竖直箭头 $a$ 满。
由《态射》引理 \ref{morphisms-lemma-smooth-locally-finite-presentation}，
$h$ 平坦。因此按定义 $X \to Y$ 平坦。
\end{proof}

\begin{lemma}
\label{spaces-morphisms-lemma-smooth-syntomic}
代数空间的光滑态射 syntomic。
\end{lemma}

\begin{proof}
设 $X \to Y$ 为代数空间的光滑态射。按定义，这意味着存在引理
\ref{spaces-morphisms-lemma-local-source-target} 中的图，其中 $h$ 光滑，且竖直箭头 $a$ 满。
由《态射》引理 \ref{morphisms-lemma-smooth-syntomic}，$h$ syntomic。
因此按定义 $X \to Y$ syntomic。
\end{proof}

\begin{lemma}
\label{spaces-morphisms-lemma-where-smooth}
设 $S$ 为概形。设 $f : X \to Y$ 为 $S$ 上代数空间的态射。
存在最大的开子空间 $U \subset X$，使得 $f|_U : U \to Y$ 光滑。
此外，形成此开子空间与下列基变换可交换：
\begin{enumerate}
\item 平坦且局部有限表示的态射；
\item 当 $f$ 局部有限表示时的平坦态射。
\end{enumerate}
\end{lemma}

\begin{proof}
$U$ 的存在性由如下事实得出：光滑性质在源上是 Zariski 局部的
（甚至是 \'etale 局部的），见引理 \ref{spaces-morphisms-lemma-smooth-local}。
而且，该引理使我们可以把性质 (1)、(2) 转化为概形态射的情形。
概形情形即《态射》引理
\ref{morphisms-lemma-set-points-where-fibres-smooth}。略去若干细节。
\end{proof}

\begin{lemma}
\label{spaces-morphisms-lemma-smoothness-dimension-spaces}
设 $X$ 与 $Y$ 为概形 $S$ 上局部 Noether 的代数空间，并设
$f : X \to Y$ 为光滑态射。对每个像为 $y \in |Y|$ 的点 $x \in |X|$，
$$
\dim_x(X) = \dim_y(Y) + \dim_x(X_y)
$$
，其中 $\dim_x(X_y)$ 是定义 \ref{spaces-morphisms-definition-dimension-fibre} 中
$f$ 在 $x$ 处的相对维数。
\end{lemma}

\begin{proof}
由代数空间在一点处之维数的定义（《空间的性质》定义
\ref{spaces-properties-definition-dimension-at-point}），问题归结为概形的相应命题
（《态射》引理 \ref{morphisms-lemma-smoothness-dimension}）。
\end{proof}



\section{非分歧态射}
\label{spaces-morphisms-section-unramified}

\noindent
概形态射的“非分歧”性质（相应地，“G-非分歧”性质）在源与靶上都是
\'etale 局部的，见《下降》评注
\ref{descent-remark-list-local-source-target}。此性质也在基变换下稳定，
并且在靶上是 fpqc 局部的，见《态射》引理
\ref{morphisms-lemma-base-change-unramified} 与《下降》引理
\ref{descent-lemma-descending-property-unramified}。因此，由上面的引理
\ref{spaces-morphisms-lemma-local-source-target}，我们可以如下定义代数空间的非分歧态射
（相应地，G-非分歧态射）；当该态射可表时，这与第
\ref{spaces-morphisms-section-representable} 节中已有的概念一致。

\begin{definition}
\label{spaces-morphisms-definition-unramified}
设 $S$ 为概形。
设 $f : X \to Y$ 为 $S$ 上代数空间的态射。
\begin{enumerate}
\item 若引理 \ref{spaces-morphisms-lemma-local-source-target} 中取
$\mathcal{P} = \text{非分歧}$ 时其中的等价条件成立，则称 $f$ {\it 非分歧}。
\item 设 $x \in |X|$。若存在 $x$ 的开邻域 $X' \subset X$，使得
$f|_{X'} : X' \to Y$ 非分歧，则称 $f$ {\it 在 $x$ 处非分歧}。
\item 若引理 \ref{spaces-morphisms-lemma-local-source-target} 中取
$\mathcal{P} = \text{G-非分歧}$ 时其中的等价条件成立，则称 $f$
{\it G-非分歧}。
\item 设 $x \in |X|$。若存在 $x$ 的开邻域 $X' \subset X$，使得
$f|_{X'} : X' \to Y$ G-非分歧，则称 $f$ {\it 在 $x$ 处 G-非分歧}。
\end{enumerate}
\end{definition}

\noindent
由于下一个引理，从此以后我们只发展非分歧态射的理论；需要使用
G-非分歧态射时，我们将直接称其为“局部有限表示的非分歧态射”。

\begin{lemma}
\label{spaces-morphisms-lemma-unramified-G-unramified}
设 $S$ 为概形，$f : X \to Y$ 为 $S$ 上代数空间的态射。
则 $f$ G-非分歧，当且仅当 $f$ 非分歧且局部有限表示。
\end{lemma}

\begin{proof}
考虑引理 \ref{spaces-morphisms-lemma-local-source-target} 中的任意图。所要说的正是：态射
$h$ G-非分歧，当且仅当它非分歧且局部有限表示。这由《态射》定义
\ref{morphisms-definition-unramified} 立即可见。
\end{proof}

\begin{lemma}
\label{spaces-morphisms-lemma-composition-unramified}
非分歧态射的复合非分歧。
\end{lemma}

\begin{proof}
见评注 \ref{spaces-morphisms-remark-composition-P} 和《态射》引理
\ref{morphisms-lemma-composition-unramified}。
\end{proof}

\begin{lemma}
\label{spaces-morphisms-lemma-base-change-unramified}
非分歧态射的基变换非分歧。
\end{lemma}

\begin{proof}
见评注 \ref{spaces-morphisms-remark-base-change-P} 和《态射》引理
\ref{morphisms-lemma-base-change-unramified}。
\end{proof}

\begin{lemma}
\label{spaces-morphisms-lemma-unramified-local}
设 $S$ 为概形。
设 $f : X \to Y$ 为 $S$ 上代数空间的态射。
下列条件等价：
\begin{enumerate}
\item $f$ 非分歧；
\item 对每个 $x \in |X|$，态射 $f$ 在 $x$ 处非分歧；
\item 对每个概形 $Z$ 及任意态射 $Z \to Y$，态射
$Z \times_Y X \to Z$ 非分歧；
\item 对每个仿射概形 $Z$ 及任意态射 $Z \to Y$，态射
$Z \times_Y X \to Z$ 非分歧；
\item 存在概形 $V$ 及满 \'etale 态射 $V \to Y$，使得
$V \times_Y X \to V$ 为非分歧态射；
\item 存在概形 $U$ 及满 \'etale 态射 $\varphi : U \to X$，使得复合
$f \circ \varphi$ 非分歧；
\item 对每个交换图
$$
\xymatrix{
U \ar[d] \ar[r] & V \ar[d] \\
X \ar[r] & Y
}
$$
其中 $U$、$V$ 为概形，竖直箭头为 \'etale 态射，顶部水平箭头非分歧；
\item 存在交换图
$$
\xymatrix{
U \ar[d] \ar[r] & V \ar[d] \\
X \ar[r] & Y
}
$$
其中 $U$、$V$ 为概形，竖直箭头为 \'etale 态射，且 $U \to X$ 满，
使得顶部水平箭头非分歧；
\item 存在 Zariski 覆盖 $Y = \bigcup_{i \in I} Y_i$ 以及
$f^{-1}(Y_i) = \bigcup X_{ij}$，使得每个态射
$X_{ij} \to Y_i$ 非分歧。
\end{enumerate}
\end{lemma}

\begin{proof}
略。
\end{proof}


\begin{lemma}
\label{spaces-morphisms-lemma-unramified-locally-finite-type}
代数空间的非分歧态射局部有限型。
\end{lemma}

\begin{proof}
借助引理 \ref{spaces-morphisms-lemma-local-source-target} 中的图，这转化为《态射》引理
\ref{morphisms-lemma-unramified-locally-finite-type}。
\end{proof}

\begin{lemma}
\label{spaces-morphisms-lemma-unramified-quasi-finite}
若 $f$ 在 $x$ 处非分歧，则 $f$ 在 $x$ 处拟有限。
特别地，非分歧态射局部拟有限。
\end{lemma}

\begin{proof}
借助引理 \ref{spaces-morphisms-lemma-local-source-target} 中的图，这转化为《态射》引理
\ref{morphisms-lemma-unramified-quasi-finite}。
\end{proof}

\begin{lemma}
\label{spaces-morphisms-lemma-immersion-unramified}
代数空间的浸入非分歧。
\end{lemma}

\begin{proof}
设 $i : X \to Y$ 为代数空间的浸入。取概形 $V$ 及满 \'etale 态射
$V \to Y$。则 $V \times_Y X \to V$ 为概形的浸入，因而非分歧
（见《态射》引理 \ref{morphisms-lemma-open-immersion-unramified} 与
\ref{morphisms-lemma-closed-immersion-unramified}）。故按定义，$i$ 非分歧。
\end{proof}

\begin{lemma}
\label{spaces-morphisms-lemma-diagonal-unramified-morphism}
设 $S$ 为概形。
设 $f : X \to Y$ 为 $S$ 上代数空间的态射。
\begin{enumerate}
\item 若 $f$ 非分歧，则对角态射
$\Delta_{X/Y} : X \to X \times_Y X$ 为开浸入。
\item 若 $f$ 局部有限型且 $\Delta_{X/Y}$ 为开浸入，则 $f$ 非分歧。
\end{enumerate}
\end{lemma}

\begin{proof}
无论如何，$\Delta_{X/Y}$ 都是可表单态射，见引理
\ref{spaces-morphisms-lemma-properties-diagonal}。取概形 $V$ 及满 \'etale 态射 $V \to Y$。
再取概形 $U$ 及满 \'etale 态射 $U \to X \times_Y V$。考虑交换图
$$
\xymatrix{
U \ar[d] \ar[rr]_-{\Delta_{U/V}} & &
U \times_V U \ar[d] \ar[r] &
V \ar[d]^{\Delta_{V/Y}} \\
X \ar[rr]^-{\Delta_{X/Y}} & &
X \times_Y X \ar[r] &
V \times_Y V
}
$$
其右方形为笛卡尔方块。左侧竖直箭头满且 \'etale。右侧竖直箭头
作为 $Y$ 上 \'etale 概形之间的态射是 \'etale 的，见《空间的性质》引理
\ref{spaces-properties-lemma-etale-permanence}。因此中间的竖直箭头也
\'etale（但未必满）。

\medskip\noindent
设 $f$ 非分歧。则 $U \to V$ 非分歧，故由《态射》引理
\ref{morphisms-lemma-diagonal-unramified-morphism}，$\Delta_{U/V}$ 为开浸入。
考察上图左方形，由《空间的性质》引理
\ref{spaces-properties-lemma-etale-local} 可知 $\Delta_{X/Y}$ 是 \'etale 态射。
因此 $\Delta_{X/Y}$ 是可表的 \'etale 单态射；由《\'Etale 态射》定理
\ref{etale-theorem-etale-radicial-open}，它是开浸入。（关于从概形语言到
函子语言的转译，另见《空间》引理
\ref{spaces-lemma-representable-transformations-property-implication}。）

\medskip\noindent
设 $f$ 局部有限型，且 $\Delta_{X/Y}$ 为开浸入。由代数空间的态射局部
有限型之定义，$U \to V$ 也局部有限型。考察上面的图，由《空间的性质》
引理 \ref{spaces-properties-lemma-etale-permanence}，可知 $\Delta_{U/V}$
作为 $X \times_Y X$ 上 \'etale 概形之间的态射是 \'etale 的。另一方面，
$\Delta_{U/V}$ 是概形之间态射的对角态射，因而无论如何都是浸入，见
《概形》引理 \ref{schemes-lemma-diagonal-immersion}。故它是开浸入，进而由
《态射》引理 \ref{morphisms-lemma-diagonal-unramified-morphism}，
$U \to V$ 非分歧。这又按定义说明 $f$ 非分歧。
\end{proof}

\begin{lemma}
\label{spaces-morphisms-lemma-where-unramified}
设 $S$ 为概形。考虑交换图
$$
\xymatrix{
X \ar[rr]_f \ar[rd]_p & & Y \ar[ld]^q \\
& Z
}
$$
，其中各对象为 $S$ 上的代数空间。设 $X \to Z$ 局部有限型。则存在开子空间
$U(f) \subset X$，使得 $|U(f)| \subset |X|$ 正是 $f$ 非分歧处的点集。
此外，对代数空间的任意态射 $Z' \to Z$，若 $f' : X' \to Y'$ 是 $f$
沿 $Z' \to Z$ 的基变换，则 $U(f')$ 是 $U(f)$ 在投影 $X' \to X$ 下的逆像。
\end{lemma}

\begin{proof}
本引理是《态射》引理
\ref{morphisms-lemma-set-points-where-fibres-unramified} 的类似结果；事实上，
我们将由它推出本引理。由定义 \ref{spaces-morphisms-definition-unramified}，集合
$\{x \in |X| : f \text{ 非分歧于 }x\}$ 在 $X$ 中为开集。因此只需证明最后的
断言。由引理 \ref{spaces-morphisms-lemma-permanence-finite-type}，态射 $X \to Y$ 局部有限型；
由引理 \ref{spaces-morphisms-lemma-base-change-finite-type}，态射 $X' \to Y'$ 局部有限型。

\medskip\noindent
取概形 $W$ 及满 \'etale 态射 $W \to Z$。取概形 $V$ 及满 \'etale 态射
$V \to W \times_Z Y$。取概形 $U$ 及满 \'etale 态射
$U \to V \times_Y X$。最后，取概形 $W'$ 及满 \'etale 态射
$W' \to W \times_Z Z'$。令 $V' = W' \times_W V$、
$U' = W' \times_W U$；于是得到满 \'etale 态射 $V' \to Y'$ 与
$U' \to X'$。以下将不再另行说明地使用如下事实：代数空间的 \'etale 态射在
伴随拓扑空间上诱导开映射（见《空间的性质》引理
\ref{spaces-properties-lemma-etale-open}）。结合引理
\ref{spaces-morphisms-lemma-unramified-local}，可知 $U(f)$ 是 $U$ 中使态射 $U \to V$
非分歧的点集 $T$ 在 $|X|$ 中的像。类似地，$U(f')$ 是 $U'$ 中使态射
$U' \to V'$ 非分歧的点集 $T'$ 在 $|X'|$ 中的像。按构造，图
$$
\xymatrix{
U' \ar[r] \ar[d] & U \ar[d] \\
V' \ar[r] & V
}
$$
在概形范畴中为笛卡尔图。因此，上述《态射》引理
\ref{morphisms-lemma-set-points-where-fibres-unramified} 表明 $T'$ 是 $T$
的逆像。由于 $|U'| \to |X'|$ 满，这就证明了本引理。
\end{proof}

\begin{lemma}
\label{spaces-morphisms-lemma-permanence-unramified}
设 $S$ 为概形，$X \to Y \to Z$ 为 $S$ 上代数空间的态射。
若 $X \to Z$ 非分歧，则 $X \to Y$ 非分歧。
\end{lemma}

\begin{proof}
取交换图
$$
\xymatrix{
U \ar[d] \ar[r] & V \ar[d] \ar[r] & W \ar[d] \\
X \ar[r] & Y \ar[r] & Z
}
$$
，其中竖直箭头均 \'etale 且满（见《空间》引理
\ref{spaces-lemma-lift-morphism-presentations}）。对顶行应用《态射》引理
\ref{morphisms-lemma-unramified-permanence}。
\end{proof}














\section{\'Etale 态射}
\label{spaces-morphisms-section-etale}

\noindent
代数空间的 \'etale 态射这一概念已在《空间的性质》定义
\ref{spaces-properties-definition-etale} 中定义。下面说明态射在一点处
\'etale 的含义。

\begin{definition}
\label{spaces-morphisms-definition-etale}
设 $S$ 为概形，$f : X \to Y$ 为 $S$ 上代数空间的态射，并设
$x \in |X|$。若存在 $x$ 的开邻域 $X' \subset X$，使得
$f|_{X'} : X' \to Y$ 为 \'etale 态射，则称 $f$ {\it 在 $x$ 处 \'etale}。
\end{definition}

\begin{lemma}
\label{spaces-morphisms-lemma-etale-local}
设 $S$ 为概形。
设 $f : X \to Y$ 为 $S$ 上代数空间的态射。
下列条件等价：
\begin{enumerate}
\item $f$ 为 \'etale 态射；
\item 对每个 $x \in |X|$，态射 $f$ 在 $x$ 处 \'etale；
\item 对每个概形 $Z$ 及任意态射 $Z \to Y$，态射
$Z \times_Y X \to Z$ 为 \'etale 态射；
\item 对每个仿射概形 $Z$ 及任意态射 $Z \to Y$，态射
$Z \times_Y X \to Z$ 为 \'etale 态射；
\item 存在概形 $V$ 及满 \'etale 态射 $V \to Y$，使得
$V \times_Y X \to V$ 为 \'etale 态射；
\item 存在概形 $U$ 及满 \'etale 态射 $\varphi : U \to X$，使得复合
$f \circ \varphi$ 为 \'etale 态射；
\item 对每个交换图
$$
\xymatrix{
U \ar[d] \ar[r] & V \ar[d] \\
X \ar[r] & Y
}
$$
其中 $U$、$V$ 为概形，竖直箭头为 \'etale 态射，顶部水平箭头也为
\'etale 态射；
\item 存在交换图
$$
\xymatrix{
U \ar[d] \ar[r] & V \ar[d] \\
X \ar[r] & Y
}
$$
其中 $U$、$V$ 为概形，竖直箭头为 \'etale 态射，且 $U \to X$ 满，
使得顶部水平箭头为 \'etale 态射；
\item 存在 Zariski 覆盖 $Y = \bigcup Y_i$ 以及
$f^{-1}(Y_i) = \bigcup X_{ij}$，使得每个态射
$X_{ij} \to Y_i$ 为 \'etale 态射。
\end{enumerate}
\end{lemma}

\begin{proof}
合用《空间的性质》引理
\ref{spaces-properties-lemma-etale-local}、
\ref{spaces-properties-lemma-base-change-etale} 与
\ref{spaces-properties-lemma-composition-etale}。略去若干细节。
\end{proof}

\begin{lemma}
\label{spaces-morphisms-lemma-composition-etale}
代数空间的两个 \'etale 态射之复合仍为 \'etale 态射。
\end{lemma}

\begin{proof}
这是《空间的性质》引理
\ref{spaces-properties-lemma-composition-etale} 的复述。
\end{proof}

\begin{lemma}
\label{spaces-morphisms-lemma-base-change-etale}
代数空间的 \'etale 态射沿代数空间的任意态射作基变换后仍为 \'etale 态射。
\end{lemma}

\begin{proof}
这是《空间的性质》引理
\ref{spaces-properties-lemma-base-change-etale} 的复述。
\end{proof}

\begin{lemma}
\label{spaces-morphisms-lemma-etale-locally-quasi-finite}
代数空间的 \'etale 态射局部拟有限。
\end{lemma}

\begin{proof}
设 $X \to Y$ 为代数空间的 \'etale 态射，见《空间的性质》定义
\ref{spaces-properties-definition-etale}。由《空间的性质》引理
\ref{spaces-properties-lemma-etale-local}，这意味着存在引理
\ref{spaces-morphisms-lemma-local-source-target} 中的图，其中 $h$ 为 \'etale 态射，且竖直箭头
$a$ 满。由《态射》引理 \ref{morphisms-lemma-etale-locally-quasi-finite}，
$h$ 局部拟有限。因此按定义 $X \to Y$ 局部拟有限。
\end{proof}

\begin{lemma}
\label{spaces-morphisms-lemma-etale-smooth}
代数空间的 \'etale 态射光滑。
\end{lemma}

\begin{proof}
证明与引理 \ref{spaces-morphisms-lemma-etale-locally-quasi-finite} 的证明相同。这里使用如下事实：
概形的 \'etale 态射光滑（由概形的 \'etale 态射之定义）。
\end{proof}

\begin{lemma}
\label{spaces-morphisms-lemma-etale-flat}
代数空间的 \'etale 态射平坦。
\end{lemma}

\begin{proof}
证明与引理 \ref{spaces-morphisms-lemma-etale-locally-quasi-finite} 的证明相同。这里使用
《态射》引理 \ref{morphisms-lemma-etale-flat}。
\end{proof}

\begin{lemma}
\label{spaces-morphisms-lemma-etale-locally-finite-presentation}
\begin{slogan}
\'Etale 蕴含局部有限表示。
\end{slogan}
代数空间的 \'etale 态射局部有限表示。
\end{lemma}

\begin{proof}
证明与引理 \ref{spaces-morphisms-lemma-etale-locally-quasi-finite} 的证明相同。这里使用
《态射》引理 \ref{morphisms-lemma-etale-locally-finite-presentation}。
\end{proof}

\begin{lemma}
\label{spaces-morphisms-lemma-etale-locally-finite-type}
代数空间的 \'etale 态射局部有限型。
\end{lemma}

\begin{proof}
\'Etale 态射局部有限表示，而局部有限表示的态射局部有限型；见引理
\ref{spaces-morphisms-lemma-etale-locally-finite-presentation} 与
\ref{spaces-morphisms-lemma-finite-presentation-finite-type}。
\end{proof}

\begin{lemma}
\label{spaces-morphisms-lemma-etale-unramified}
代数空间的 \'etale 态射非分歧。
\end{lemma}

\begin{proof}
证明与引理 \ref{spaces-morphisms-lemma-etale-locally-quasi-finite} 的证明相同。这里使用
《态射》引理 \ref{morphisms-lemma-etale-smooth-unramified}。
\end{proof}

\begin{lemma}
\label{spaces-morphisms-lemma-etale-permanence}
设 $S$ 为概形。设 $X,Y$ 为在某代数空间 $Z$ 上 \'etale 的代数空间。
则 $Z$ 上的任意态射 $X \to Y$ 都是 \'etale 态射。
\end{lemma}

\begin{proof}
这是《空间的性质》引理
\ref{spaces-properties-lemma-etale-permanence} 的复述。
\end{proof}

\begin{lemma}
\label{spaces-morphisms-lemma-unramified-flat-lfp-etale}
代数空间之间局部有限表示、平坦且非分歧的态射是 \'etale 态射。
\end{lemma}

\begin{proof}
设 $X \to Y$ 为代数空间之间局部有限表示、平坦且非分歧的态射。由
《空间的性质》引理 \ref{spaces-properties-lemma-etale-local}，这意味着存在
引理 \ref{spaces-morphisms-lemma-local-source-target} 中的图，其中 $h$ 局部有限表示、平坦且
非分歧，且竖直箭头 $a$ 满。由《态射》引理
\ref{morphisms-lemma-flat-unramified-etale}，$h$ 为 \'etale 态射。因此按定义，
$X \to Y$ 为 \'etale 态射。
\end{proof}






\section{固有态射}
\label{spaces-morphisms-section-proper}

\noindent
固有态射的概念在代数几何中起着重要作用。下面给出代数空间的固有态射之定义。

\begin{definition}
\label{spaces-morphisms-definition-proper}
设 $S$ 为概形，$f : X \to Y$ 为 $S$ 上代数空间的态射。
若 $f$ 分离、有限型且普遍闭，则称 $f$ {\it 固有}。
\end{definition}

\begin{lemma}
\label{spaces-morphisms-lemma-proper-local}
设 $S$ 为概形，$f : X \to Y$ 为 $S$ 上代数空间的态射。下列条件等价：
\begin{enumerate}
\item $f$ 固有；
\item 对每个概形 $Z$ 及每个态射 $Z \to Y$，投影
$Z \times_Y X \to Z$ 固有；
\item 对每个仿射概形 $Z$ 及每个态射 $Z \to Y$，投影
$Z \times_Y X \to Z$ 固有；
\item 存在概形 $V$ 及满 \'etale 态射 $V \to Y$，使得
$V \times_Y X \to V$ 固有；
\item 存在 Zariski 覆盖 $Y = \bigcup Y_i$，使得每个态射
$f^{-1}(Y_i) \to Y_i$ 固有。
\end{enumerate}
\end{lemma}

\begin{proof}
合用引理 \ref{spaces-morphisms-lemma-separated-local}、
\ref{spaces-morphisms-lemma-finite-type-local},
\ref{spaces-morphisms-lemma-quasi-compact-local} 与
\ref{spaces-morphisms-lemma-universally-closed-local}。
\end{proof}

\begin{lemma}
\label{spaces-morphisms-lemma-base-change-proper}
固有态射的基变换固有。
\end{lemma}

\begin{proof}
见引理 \ref{spaces-morphisms-lemma-base-change-separated}、
\ref{spaces-morphisms-lemma-base-change-finite-type} 与
\ref{spaces-morphisms-lemma-base-change-universally-closed}。
\end{proof}

\begin{lemma}
\label{spaces-morphisms-lemma-composition-proper}
固有态射的复合固有。
\end{lemma}

\begin{proof}
见引理 \ref{spaces-morphisms-lemma-composition-separated}、
\ref{spaces-morphisms-lemma-composition-finite-type} 与
\ref{spaces-morphisms-lemma-composition-universally-closed}。
\end{proof}

\begin{lemma}
\label{spaces-morphisms-lemma-closed-immersion-proper}
代数空间的闭浸入是代数空间的固有态射。
\end{lemma}

\begin{proof}
闭浸入按定义可表，故这由《空间》引理
\ref{spaces-lemma-representable-transformations-property-implication}
及概形态射的相应结果得出；见《态射》引理
\ref{morphisms-lemma-closed-immersion-proper}。
\end{proof}

\begin{lemma}
\label{spaces-morphisms-lemma-universally-closed-permanence}
设 $S$ 为概形。考虑代数空间的交换图
$$
\xymatrix{
X \ar[rr] \ar[rd] & &
Y \ar[ld] \\
& B &
}
$$
，其中各对象在 $S$ 上。
\begin{enumerate}
\item 若 $X \to B$ 普遍闭且 $Y \to B$ 分离，则态射 $X \to Y$ 普遍闭。
特别地，$|X|$ 在 $|Y|$ 中的像为闭集。
\item 若 $X \to B$ 固有且 $Y \to B$ 分离，则态射 $X \to Y$ 固有。
\end{enumerate}
\end{lemma}

\begin{proof}
设 $X \to B$ 普遍闭且 $Y \to B$ 分离。把该态射分解为
$X \to X \times_B Y \to Y$。第一个态射是闭浸入，见引理
\ref{spaces-morphisms-lemma-semi-diagonal}，故普遍闭。投影 $X \times_B Y \to Y$ 是普遍闭态射的
基变换，因而普遍闭，见引理 \ref{spaces-morphisms-lemma-base-change-universally-closed}。
因此，$X \to Y$ 作为普遍闭态射的复合而普遍闭，见引理
\ref{spaces-morphisms-lemma-composition-universally-closed}。这证明了 (1)。为推出 (2)，合用 (1)
与引理 \ref{spaces-morphisms-lemma-compose-after-separated}、
\ref{spaces-morphisms-lemma-quasi-compact-permanence} 及
\ref{spaces-morphisms-lemma-permanence-finite-type}。
\end{proof}

\begin{lemma}
\label{spaces-morphisms-lemma-image-proper-is-proper}
设 $S$ 为概形，$B$ 为 $S$ 上的代数空间，并设
$f : X \to Y$ 为 $B$ 上代数空间的态射。若 $X$ 在 $B$ 上普遍闭且 $f$ 满，
则 $Y$ 在 $B$ 上普遍闭。特别地，若 $Y$ 还在 $B$ 上分离且有限型，
则 $Y$ 在 $B$ 上固有。
\end{lemma}

\begin{proof}
设 $X$ 普遍闭且 $f$ 满。以 $p : X \to B$、$q : Y \to B$ 表示结构态射。
设 $B' \to B$ 为 $S$ 上代数空间的态射。基变换
$f' : X_{B'} \to Y_{B'}$ 满（引理 \ref{spaces-morphisms-lemma-base-change-surjective}），
而基变换 $p' : X_{B'} \to B'$ 闭。若 $T \subset Y_{B'}$ 为闭集，
则 $(f')^{-1}(T) \subset X_{B'}$ 为闭集，故
$p'((f')^{-1}(T)) = q'(T)$ 为闭集。因此 $q'$ 闭。
\end{proof}

\begin{lemma}
\label{spaces-morphisms-lemma-scheme-theoretic-image-is-proper}
设 $S$ 为概形。设
$$
\xymatrix{
X \ar[rr]_h \ar[rd]_f & & Y \ar[ld]^g \\
& B
}
$$
为 $S$ 上代数空间态射的交换图。设
\begin{enumerate}
\item $X \to B$ 为固有态射；
\item $Y \to B$ 分离且局部有限型。
\end{enumerate}
则 $h$ 的概形论像 $Z \subset Y$ 在 $B$ 上固有，且 $X \to Z$ 满。
\end{lemma}

\begin{proof}
$h$ 的概形论像已在第 \ref{spaces-morphisms-section-scheme-theoretic-image} 节中构造。
注意，$h$ 拟紧（引理 \ref{spaces-morphisms-lemma-quasi-compact-quasi-separated-permanence}），
故 $|h|(|X|) \subset |Z|$ 稠密（引理
\ref{spaces-morphisms-lemma-quasi-compact-scheme-theoretic-image}）。另一方面，
$|h|(|X|)$ 在 $|Y|$ 中为闭集（引理
\ref{spaces-morphisms-lemma-universally-closed-permanence}），故 $X \to Z$ 满。
于是 $Z \to B$ 固有（引理 \ref{spaces-morphisms-lemma-image-proper-is-proper}）。
\end{proof}

\begin{lemma}
\label{spaces-morphisms-lemma-separated-diagonal-proper}
设 $S$ 为概形。
设 $f : X \to Y$ 为 $S$ 上代数空间的态射。
下列条件等价：
\begin{enumerate}
\item $f$ 分离；
\item $\Delta_{X/Y} : X \to X \times_Y X$ 普遍闭；
\item $\Delta_{X/Y} : X \to X \times_Y X$ 固有。
\end{enumerate}
\end{lemma}

\begin{proof}
(1) $\Rightarrow$ (3) 由引理 \ref{spaces-morphisms-lemma-closed-immersion-proper} 得出。
在余下的证明中，我们将不再另行说明地使用《空间》引理
\ref{spaces-lemma-representable-transformations-property-implication}。
回顾 $\Delta_{X/Y}$ 是局部有限型的可表单态射，见引理
\ref{spaces-morphisms-lemma-properties-diagonal}。由于概形态射中固有 $\Rightarrow$ 普遍闭，
故 (3) 蕴含 (2)。若 $\Delta_{X/Y}$ 普遍闭，则《\'Etale 态射》引理
\ref{etale-lemma-characterize-closed-immersion} 表明它是闭浸入。
因此 (2) $\Rightarrow$ (1)，证明完毕。
\end{proof}





\section{赋值判据}
\label{spaces-morphisms-section-valuative}

\noindent
本节介绍代数空间态射之赋值判据的基础。进一步结果可参见：
\begin{enumerate}
\item 普遍闭性的赋值判据见第
\ref{spaces-morphisms-section-valuative-criterion-universally-closed} 节；
\item 分离性的赋值判据见第 \ref{spaces-morphisms-section-valuative-separatedness} 节；
\item 固有性的赋值判据见第 \ref{spaces-morphisms-section-valuative-criterion-properness} 节；
\item 其他逆命题见《良好空间》第
\ref{decent-spaces-section-valuative-criterion-universally-closed}
节及《良好空间》引理
\ref{decent-spaces-lemma-re-characterize-universally-closed}；
\item 在 Noether 情形，只需对离散赋值环检查该判据；见《空间的上同调》第
\ref{spaces-cohomology-section-Noetherian-valuative-criterion} 节及
《空间的极限》第 \ref{spaces-limits-section-Noetherian-valuative-criterion} 节；
\item Noether 情形下赋值判据的精细版本见《空间的极限》第
\ref{spaces-limits-section-refined-valuative-criteria}.
\end{enumerate}
我们先正式陈述定义，再讨论它与概形态射情形的差别。

\begin{definition}
\label{spaces-morphisms-definition-valuative-criterion}
设 $S$ 为概形，$f : X \to Y$ 为 $S$ 上代数空间的态射。若对任意由实线箭头
构成的交换图
$$
\xymatrix{
\Spec(K) \ar[r] \ar[d] & X \ar[d] \\
\Spec(A) \ar[r] \ar@{-->}[ru] & Y
}
$$
，其中 $A$ 为分式域是 $K$ 的赋值环，虚线箭头至多存在一个（不要求存在），
则称 $f$ {\it 满足赋值判据的唯一性部分}。若对上述任意实线图，存在域扩张
$K'/K$、支配 $A$ 的赋值环 $A' \subset K'$ 以及态射
$\Spec(A') \to X$，使得下图交换
$$
\xymatrix{
\Spec(K') \ar[r] \ar[d] & \Spec(K) \ar[r] & X \ar[d] \\
\Spec(A') \ar[r] \ar[rru] & \Spec(A) \ar[r] & Y
}
$$
，则称 $f$ {\it 满足赋值判据的存在性部分}。若 $f$ 同时满足存在性部分与
唯一性部分，则称 $f$ {\it 满足赋值判据}。
\end{definition}

\noindent
对代数空间的态射，赋值判据之存在性部分的表述略有不同，因为可能需要扩张
赋值环的分式域。实践中，这一区别几乎从不起作用。
\begin{enumerate}
\item 检查赋值判据的唯一性部分从不涉及分式域扩张，故与概形情形完全相同。
\item 一般而言必须允许域扩张，见例 \ref{spaces-morphisms-example-finite-separable-needed}。
\item 对代数空间的态射，在赋值判据的存在性部分中总可以取有限可分扩张
$K'/K$，见引理 \ref{spaces-morphisms-lemma-finite-separable-enough}。
\item 若 $f : X \to Y$ 是代数空间的分离态射，则检查赋值判据的存在性部分时
总可取 $K = K'$，见引理 \ref{spaces-morphisms-lemma-usual-enough}。
\item 对拟紧拟分离态射 $f : X \to Y$，“$f$ 分离且普遍闭”等价于“$f$
满足通常的赋值判据”，见引理
\ref{spaces-morphisms-lemma-characterize-separated-and-universally-closed}。固有性的赋值判据就是
通常的判据，见引理 \ref{spaces-morphisms-lemma-characterize-proper}。
\end{enumerate}
作为理论的第一步，我们证明：当代数空间的态射可表时，该判据与概形态射所用
的表述相同。

\begin{lemma}
\label{spaces-morphisms-lemma-valuative-criterion-representable}
设 $S$ 为概形，$f : X \to Y$ 为 $S$ 上代数空间的态射，并设 $f$ 可表。
下列条件等价：
\begin{enumerate}
\item $f$ 满足定义 \ref{spaces-morphisms-definition-valuative-criterion} 中赋值判据的存在性部分；
\item 对任意由实线箭头构成的交换图
$$
\xymatrix{
\Spec(K) \ar[r] \ar[d] & X \ar[d] \\
\Spec(A) \ar[r] \ar@{-->}[ru] & Y
}
$$
，其中 $A$ 为分式域是 $K$ 的赋值环，都存在虚线箭头；换言之，$f$ 满足
《概形》定义 \ref{schemes-definition-valuative-criterion} 中赋值判据的存在性部分。
\end{enumerate}
\end{lemma}

\begin{proof}
只需证明：给定定义 \ref{spaces-morphisms-definition-valuative-criterion} 中形如
$$
\xymatrix{
\Spec(K') \ar[r] \ar[d] & \Spec(K) \ar[r] & X \ar[d] \\
\Spec(A') \ar[r] \ar[rru]^\varphi & \Spec(A) \ar[r] & Y
}
$$
的交换图后，也能找到可嵌入此图的态射 $\Spec(A) \to X$。
% P08-E243: frozen source fibre-product formula preserved; see ERRATA_CANDIDATES.jsonl.
令 $X_A = \Spec(A) \times_Y Y$。由于 $f$ 可表，$X_A$ 是概形。态射 $\varphi$
给出态射 $\varphi' : \Spec(A') \to X_A$。设 $x \in X_A$ 为
$\varphi' : \Spec(A') \to X_A$ 的闭点之像。于是有环的交换图
$$
\xymatrix{
K' & K \ar[l] & \mathcal{O}_{X_A, x} \ar[l] \ar[lld] \\
A' \ar[u] & A \ar[l] & A \ar[l] \ar[u]
}
$$
由于 $A$ 为赋值环且 $A'$ 支配 $A$，有 $K \cap A' = A$。因此环映射
$\mathcal{O}_{X_A, x} \to K$ 的像包含在 $A$ 中。于是得到所需态射
$\Spec(A) \to X_A$（见《概形》第 \ref{schemes-section-points} 节）。
\end{proof}

\begin{lemma}
\label{spaces-morphisms-lemma-finite-separable-enough}
设 $S$ 为概形，$f : X \to Y$ 为 $S$ 上代数空间的态射。下列条件等价：
\begin{enumerate}
\item $f$ 满足定义 \ref{spaces-morphisms-definition-valuative-criterion} 中赋值判据的存在性部分；
\item $f$ 满足定义 \ref{spaces-morphisms-definition-valuative-criterion} 中赋值判据的存在性部分，
但将其修改为要求扩张 $K'/K$ 有限可分。
\end{enumerate}
\end{lemma}

\begin{proof}
只需证明 (1) 蕴含 (2)。设给定定义 \ref{spaces-morphisms-definition-valuative-criterion} 中的图
$$
\xymatrix{
\Spec(K') \ar[r] \ar[d] & \Spec(K) \ar[r] & X \ar[d] \\
\Spec(A') \ar[r] \ar[rru] & \Spec(A) \ar[r] & Y
}
$$
，其中 $K \subset K'$ 任意。取概形 $U$ 及满 \'etale 态射 $U \to X$。于是
$$
\Spec(A') \times_X U \longrightarrow \Spec(A')
$$
满且 \'etale。令 $p$ 为 $\Spec(A') \times_X U$ 中映至 $\Spec(A')$ 闭点的点。
令 $p' \leadsto p$ 为 $p$ 的一个泛化，并使其映至 $\Spec(A')$ 的一般点。
这样的泛化存在，因为泛化可沿概形的平坦态射提升；见《态射》引理
\ref{morphisms-lemma-generalizations-lift-flat}。于是 $p'$ 对应于概形
$\Spec(K') \times_X U$ 的一个点。注意
$$
\Spec(K') \times_X U
=
\Spec(K') \times_{\Spec(K)} (\Spec(K) \times_X U)
$$
，故 $p'$ 映至一点 $q' \in \Spec(K) \times_X U$，其剩余域是 $K$ 的有限可分
扩张。最后，$p' \leadsto p$ 映至概形 $U$ 上的特化 $u' \leadsto u$。
用上述记号，得到环的图
$$
\xymatrix{
\kappa(p') & & \kappa(q') \ar[ll] & \kappa(u') \ar[l] \\
& \mathcal{O}_{\Spec(A') \times_X U, p} \ar[lu] & &
\mathcal{O}_{U, u} \ar[ll] \ar[u] \\
K' \ar[uu] & A' \ar[l] \ar[u] & A \ar[l] \ar'[u][uu]
}
$$
这意味着，由 $A$ 与 $\mathcal{O}_{U,u}$ 的像生成的环
$B \subset \kappa(q')$ 映至 $\kappa(p')$ 的一个子环，而该子环包含在
$\mathcal{O}_{\Spec(A') \times_X U,p} \to \kappa(p')$ 的像 $B'$ 中。
注意 $B'$ 是局部环。
% P08-E244: frozen source ideal containment preserved; see ERRATA_CANDIDATES.jsonl.
令 $\mathfrak m \subset B$ 为极大理想。按构造，$A \cap \mathfrak m$
（相应地，$\mathcal{O}_{U,u} \cap \mathfrak m$、$A' \cap \mathfrak m$）
是 $A$（相应地，$\mathcal{O}_{U,u}$、$A'$）的极大理想。令
$\mathfrak q = B \cap \mathfrak m$。这是素理想，且 $A \cap \mathfrak q$
是 $A$ 的极大理想。因此 $B_{\mathfrak q} \subset \kappa(q')$ 是支配 $A$
的局部环。由《代数》引理 \ref{algebra-lemma-dominate}，可找到分式域为
$\kappa(q')$ 且支配 $B_{\mathfrak q}$ 的赋值环
$A_1 \subset \kappa(q')$。局部环映射 $\mathcal{O}_{U,u} \to A_1$ 给出态射
$\Spec(A_1) \to U \to X$
，使得图
$$
\xymatrix{
\Spec(\kappa(q')) \ar[r] \ar[d] & \Spec(K) \ar[r] & X \ar[d] \\
\Spec(A_1) \ar[r] \ar[rru] & \Spec(A) \ar[r] & Y
}
$$
交换。由于 $A_1$ 的分式域为 $\kappa(q')$，且按构造
$\kappa(q')/K$ 有限可分，本引理得证。
\end{proof}

\begin{lemma}
\label{spaces-morphisms-lemma-push-down-solution}
设 $S$ 为概形，$f : X \to Y$ 为 $S$ 上代数空间的分离态射。设给定定义
\ref{spaces-morphisms-definition-valuative-criterion} 中的图
$$
\xymatrix{
\Spec(K') \ar[r] \ar[d] & \Spec(K) \ar[r] & X \ar[d] \\
\Spec(A') \ar[r] \ar[rru] & \Spec(A) \ar[r] \ar@{-->}[ru] & Y
}
$$
，其中 $K \subset K'$ 任意。则存在使该图交换的虚线箭头。
\end{lemma}

\begin{proof}
我们必须证明能找到嵌入此图的态射 $\Spec(A) \to X$。

\medskip\noindent
考虑 $X$ 的基变换 $X_A = \Spec(A) \times_Y X$。则
$X_A \to \Spec(A)$ 是代数空间的分离态射（引理
\ref{spaces-morphisms-lemma-base-change-separated}）。对上图的所有态射作基变换，得到
$$
\xymatrix{
\Spec(K') \ar[r] \ar[d] & \Spec(K) \ar[r] & X_A \ar[d] \\
\Spec(A') \ar[r] \ar[rru] & \Spec(A) \ar@{=}[r] & \Spec(A)
}
$$
因此可以用 $X_A$ 替换 $X$，设 $Y = \Spec(A)$，并设有上图。我们还可以且
确实用一个包含 $|\Spec(A')| \to |X|$ 之像的拟紧开子空间替换 $X$。

\medskip\noindent
由引理 \ref{spaces-morphisms-lemma-quasi-compact-permanence}，态射 $\Spec(A') \to X$ 拟紧。
令 $Z \subset X$ 为 $\Spec(A') \to X$ 的概形论像。则 $Z$ 是 $A$ 上的既约
（引理 \ref{spaces-morphisms-lemma-scheme-theoretic-image-reduced}）、拟紧（因为是 $X$ 的闭
子空间）且分离（同样因为是 $X$ 的闭子空间）的代数空间。考虑态射
$\Spec(A') \to X$ 沿概形的平坦态射 $\Spec(K) \to \Spec(A)$ 的基变换
$$
\Spec(K') = \Spec(A') \times_{\Spec(A)} \Spec(K) \to
X \times_{\Spec(A)} \Spec(K) = X_K
$$
。由引理 \ref{spaces-morphisms-lemma-flat-base-change-scheme-theoretic-image}，此态射的概形论像
是 $Z$ 的基变换 $Z_K$。另一方面，由假设（即上面的交换图），此态射分解经过
态射 $\Spec(K) \to Z_K$，而后者是结构态射 $Z_K \to \Spec(K)$ 的截面。
由于 $Z_K$ 分离，该截面为闭浸入（引理 \ref{spaces-morphisms-lemma-section-immersion}）。
所以 $Z_K = \Spec(K)$。

\medskip\noindent
取满 \'etale 态射 $V \to Z$，其中 $V$ 为仿射概形（《空间的性质》引理
\ref{spaces-properties-lemma-quasi-compact-affine-cover}）。记
$V = \Spec(B)$。由于 $Z$ 分离，$V \times_Z \Spec(A') = \Spec(C)$ 为仿射。
注意 $B \to C$ 单射，因为由引理
\ref{spaces-morphisms-lemma-quasi-compact-scheme-theoretic-image}，$V$ 是
$V \times_Z \Spec(A') \to V$ 的概形论像。另一方面，$A' \to C$ 是
\'etale 的，因为它对应于 $V \to Z$ 的基变换。由于 $A'$ 是无挠
$A$-模，$A' \to C$ 的平坦性蕴含 $C$ 是无挠 $A$-模，进而 $B$ 是无挠
$A$-模。注意，作为 $A$-模无挠等价于平坦（《关于代数的进一步结果》引理
\ref{more-algebra-lemma-valuation-ring-torsion-free-flat}）。接下来写成
$$
V \times_Z V = \Spec(B')
$$
。由于 $V \to Z$ 为 \'etale 态射，两个环映射 $B \to B'$ 都是 \'etale 的。
因 $Z_K = \Spec(K)$，典范满射 $B \otimes_A B \to B'$ 在 $A$ 上与 $K$
张量后成为同构。但由上述讨论，$B \otimes_A B$ 作为 $A$-模无挠。
故 $B' = B \otimes_A B$。由此可知，环映射 $A \to B$ 沿忠实平坦环映射
% P08-E245: frozen source surjectivity justification preserved; see ERRATA_CANDIDATES.jsonl.
$A \to B$ 的基变换是 \'etale 的（注意，因 $X \to \Spec(A)$ 满，
$\Spec(B) \to \Spec(A)$ 满）。因此 $A \to B$ 为 \'etale 映射
（《下降》引理 \ref{descent-lemma-descending-property-etale}）；换言之，
% P08-E246: frozen source codomain preserved; see ERRATA_CANDIDATES.jsonl.
$V \to X$ 为 \'etale 态射。由于
$V \times_Z V = V \times_{\Spec(A)} V$，可知代数空间
$Z = \Spec(A)$（例如由《空间》引理 \ref{spaces-lemma-space-presentation}），
证明完毕。
\end{proof}

\begin{lemma}
\label{spaces-morphisms-lemma-usual-enough}
设 $S$ 为概形，$f : X \to Y$ 为 $S$ 上代数空间的分离态射。下列条件等价：
\begin{enumerate}
\item $f$ 满足定义 \ref{spaces-morphisms-definition-valuative-criterion} 中赋值判据的存在性部分；
\item 对任意由实线箭头构成的交换图
$$
\xymatrix{
\Spec(K) \ar[r] \ar[d] & X \ar[d] \\
\Spec(A) \ar[r] \ar@{-->}[ru] & Y
}
$$
，其中 $A$ 为分式域是 $K$ 的赋值环，都存在虚线箭头；换言之，$f$ 满足
《概形》定义 \ref{schemes-definition-valuative-criterion} 中赋值判据的存在性部分。
\end{enumerate}
\end{lemma}

\begin{proof}
只需证明 (1) 蕴含 (2)。设给定 (2) 中的交换图
$$
\xymatrix{
\Spec(K) \ar[r] \ar[d] & X \ar[d] \\
\Spec(A) \ar[r] & Y
}
$$
。由 (1)，存在定义 \ref{spaces-morphisms-definition-valuative-criterion} 中的交换图
$$
\xymatrix{
\Spec(K') \ar[r] \ar[d] & \Spec(K) \ar[r] & X \ar[d] \\
\Spec(A') \ar[r] \ar[rru] & \Spec(A) \ar[r] & Y
}
$$
，其中 $K \subset K'$ 任意。由引理 \ref{spaces-morphisms-lemma-push-down-solution}，能找到嵌入
此图的态射 $\Spec(A) \to X$，即 (2) 成立。
\end{proof}

\begin{example}
\label{spaces-morphisms-example-finite-separable-needed}
考虑《空间》例 \ref{spaces-example-non-representable-descent} 中构造的代数空间
$X$。回顾它是把原点加倍之仿射直线的 Galois 扭曲；该扭曲对应于二次 Galois
域扩张 $k'/k$。因而它带有态射
$$
\pi : X \longrightarrow S = \mathbf{A}^1_k
$$
，且该态射拟紧。我们断言 $\pi$ 普遍闭。事实上，沿
$\Spec(k') \to \Spec(k)$ 作基变换后，态射 $\pi$ 等同于态射
$$
\text{原点加倍的仿射直线}
\longrightarrow
\text{仿射直线}
$$
，而它普遍闭（略去若干细节）。由于态射 $\Spec(k') \to \Spec(k)$ 普遍闭且满，
追图可见 $\pi$ 普遍闭。另一方面，考虑图
$$
\xymatrix{
\Spec(k((x))) \ar[r] \ar[d] & X \ar[d]^\pi \\
\Spec(k[[x]]) \ar[r] \ar@{..>}[ru] & \mathbf{A}^1_k
}
$$
由于 $X$ 中位于 $0 \in \mathbf{A}^1_k$ 上方的唯一一点对应于单态射
$\Spec(k') \to X$，显然不可能存在虚线箭头！这说明一般而言需要有限可分域扩张。
\end{example}

\begin{lemma}
\label{spaces-morphisms-lemma-base-change-valuative-criteria}
满足赋值判据之存在性部分（相应地，唯一性部分）的代数空间态射，沿代数空间的
任意态射作基变换后仍满足赋值判据的存在性部分（相应地，唯一性部分）。
\end{lemma}

\begin{proof}
设 $f : X \to Y$ 为概形 $S$ 上代数空间的态射，并设 $Z \to Y$ 为 $S$ 上
代数空间的任意态射。考虑下列由实线箭头构成的交换图
$$
\xymatrix{
\Spec(K) \ar[r] \ar[d] & Z \times_Y X \ar[r] \ar[d] & X \ar[d] \\
\Spec(A) \ar[r] \ar@{-->}[ru] \ar@{-->}[rru] & Z \ar[r] & Y
}
$$
。使该图交换的西北方向虚线箭头之集合，与使该图交换的西偏北方向虚线箭头之
集合一一对应。这证明了“唯一性”情形。对存在性部分，设 $f$ 满足赋值判据的
存在性部分。若给定上述实线交换图，则由假设，存在域扩张 $K'/K$、支配 $A$
的赋值环 $A' \subset K'$，以及嵌入下列交换图的态射 $\Spec(A') \to X$
$$
\xymatrix{
\Spec(K') \ar[r] \ar[d] &
\Spec(K) \ar[r] & Z \times_Y X \ar[r] & X \ar[d] \\
\Spec(A') \ar[r] \ar[rrru] & \Spec(A) \ar[r] & Z \ar[r] & Y
}
$$
。由上述说明，斜箭头对应于所需箭头
$\Spec(A') \to Z \times_Y X$。
\end{proof}

\begin{lemma}
\label{spaces-morphisms-lemma-composition-valuative-criteria}
两个满足赋值判据（存在性部分，相应地，唯一性部分）的代数空间态射之复合，
仍满足赋值判据（存在性部分，相应地，唯一性部分）。
\end{lemma}

\begin{proof}
设 $f : X \to Y$、$g : Y \to Z$ 为概形 $S$ 上代数空间的态射。考虑下列
由实线箭头构成的交换图
$$
\xymatrix{
\Spec(K) \ar[dd] \ar[r] & X \ar[d]^f \\
& Y \ar[d]^g \\
\Spec(A) \ar[r] \ar@{-->}[ru] \ar@{-->}[ruu] & Z
}
$$
若 $g$ 满足唯一性部分，则至多存在一个使图交换的西北方向虚线箭头。
若 $f$ 也满足唯一性部分，则至多存在一个使图交换的北偏西方向虚线箭头。
存在性情形的证明来自考察下图
$$
\xymatrix{
\Spec(K'') \ar[r] \ar[dd] &
\Spec(K') \ar[r] &
\Spec(K) \ar[r] &
X \ar[d]^f \\
& & & Y \ar[d]^g \\
\Spec(A'') \ar[r] \ar[rrruu] &
\Spec(A') \ar[r] \ar[rru] &
\Spec(A) \ar[r] &
Z
}
$$
具体而言，$g$ 的存在性部分给出扩张 $K'$、赋值环 $A'$ 及箭头
$\Spec(A') \to Y$；随后，$f$ 的存在性部分给出扩张 $K''$、赋值环 $A''$
及箭头 $\Spec(A'') \to X$。
\end{proof}






\section{普遍闭性的赋值判据}
\label{spaces-morphisms-section-valuative-criterion-universally-closed}

\noindent
对拟紧态射，赋值判据的存在性部分蕴含普遍闭性，见引理
\ref{spaces-morphisms-lemma-quasi-compact-existence-universally-closed}。在概形情形，这是一个
“当且仅当”的命题，但对代数空间的态射则不成立。例
\ref{spaces-morphisms-example-strange-universally-closed} 表明
$\mathbf{A}^1_k/\mathbf{Z} \to \Spec(k)$ 普遍闭，但容易看出赋值判据的存在性
部分不成立。我们将在《良好空间》第
\ref{decent-spaces-section-valuative-criterion-universally-closed} 节中重访这一主题，
并证明当态射的源为良好空间时逆命题成立（相对版本另见《良好空间》引理
\ref{decent-spaces-lemma-re-characterize-universally-closed}）。

\begin{lemma}
\label{spaces-morphisms-lemma-quasi-compact-existence-universally-closed}
设 $S$ 为概形，$f : X \to Y$ 为 $S$ 上代数空间的态射。设
\begin{enumerate}
\item $f$ 拟紧；
\item $f$ 满足赋值判据的存在性部分。
\end{enumerate}
则 $f$ 普遍闭。
\end{lemma}

\begin{proof}
由引理 \ref{spaces-morphisms-lemma-base-change-quasi-compact} 与
\ref{spaces-morphisms-lemma-base-change-valuative-criteria}，性质 (1)、(2) 在任意基变换下保持。
由引理 \ref{spaces-morphisms-lemma-universally-closed-local}，只需证明：每当 $T$ 是映入 $Y$ 的
$S$ 上仿射概形时，$|T \times_Y X| \to |T|$ 闭。因此，只需证明：若 $Y$
是仿射概形，$f : X \to Y$ 拟紧且满足赋值判据的存在性部分，则
$f : |X| \to |Y|$ 闭。此时 $X$ 是拟紧代数空间。由《空间的性质》引理
\ref{spaces-properties-lemma-quasi-compact-affine-cover}，存在仿射概形 $U$
及满 \'etale 态射 $\varphi : U \to X$。设 $T \subset |X|$ 闭。
逆像 $\varphi^{-1}(T) \subset U$ 闭，因而是某仿射闭子概形 $Z \subset U$
的点集。于是由《代数》引理
\ref{algebra-lemma-image-stable-specialization-closed}，只要
$f(T) = f(\varphi(|Z|)) \subset |Y|$ 在特化下封闭，它就是闭集。

\medskip\noindent
设 $y' \leadsto y$ 为 $Y$ 中的特化，且 $y' \in f(T)$。取一点
$x' \in T \subset |X|$，使其在 $f$ 下映至 $y'$。可以用某个域 $K$ 上的态射
$\Spec(K) \to X$ 表示 $x'$。于是有图
$$
\xymatrix{
\Spec(K) \ar[r]_-{x'} \ar[d] & X \ar[d]^f \\
\Spec(\mathcal{O}_{Y, y}) \ar[r] & Y,
}
$$
左侧竖直映射的存在性见《概形》第 \ref{schemes-section-points} 节。取支配环映射
$\mathcal{O}_{Y,y} \to K$ 之像的赋值环 $A \subset K$（这是可能的，因为当
$y' \not = y$ 时，该像是局部环而非域；见《代数》引理
\ref{algebra-lemma-dominate}）。由假设，存在域扩张 $K'/K$、支配 $A$ 的赋值环
$A' \subset K'$，以及嵌入该交换图的态射 $\Spec(A') \to X$。由于 $A'$ 支配
$A$，且 $A$ 支配 $\mathcal{O}_{Y,y}$，可知 $\Spec(A')$ 的闭点映至一点
$x \in X$，满足 $f(x) = y$，且该点是 $x'$ 的特化。由于 $T$ 闭，故
$x \in T$，从而如所需 $y \in f(T)$。
\end{proof}

\noindent
下一个引理将在《良好空间》引理
\ref{decent-spaces-lemma-re-characterize-universally-closed} 中得到推广。

\begin{lemma}
\label{spaces-morphisms-lemma-characterize-universally-closed-quasi-separated}
设 $S$ 为概形，$f : X \to Y$ 为 $S$ 上代数空间的态射。
\begin{enumerate}
\item 若 $f$ 拟分离且普遍闭，则 $f$ 满足赋值判据的存在性部分。
\item 若 $f$ 拟紧且拟分离，则 $f$ 普遍闭，当且仅当赋值判据的存在性部分成立。
\end{enumerate}
\end{lemma}

\begin{proof}
若 (1) 成立，则将其与引理
\ref{spaces-morphisms-lemma-quasi-compact-existence-universally-closed} 合用即得 (2)。
设 $f$ 拟分离且普遍闭。设给定定义
\ref{spaces-morphisms-definition-valuative-criterion} 中的图
$$
\xymatrix{
\Spec(K) \ar[r] \ar[d] & X \ar[d] \\
\Spec(A) \ar[r] & Y
}
$$
。形式论证表明，所需图
$$
\xymatrix{
\Spec(K') \ar[r] \ar[d] & \Spec(K) \ar[r] & X \ar[d] \\
\Spec(A') \ar[r] \ar[rru] & \Spec(A) \ar[r] & Y
}
$$
的存在性由态射 $X_A \to \Spec(A)$ 情形下的存在性得出。由于拟分离和普遍闭
都在基变换下保持，本引理由下一段中的结果得出。

\medskip\noindent
考虑由实线箭头构成的图
$$
\xymatrix{
\Spec(K) \ar[r]_-x \ar[d] & X \ar[d]^f \\
\Spec(A) \ar@{=}[r] \ar@{..>}[ru] & \Spec(A)
}
$$
，其中 $A$ 为分式域是 $K$ 的赋值环。由引理
\ref{spaces-morphisms-lemma-quasi-compact-permanence} 及 $f$ 拟分离这一事实，态射 $x$ 拟紧。
由于 $f$ 普遍闭，特别有 $|f|(\overline{\{x\}})$ 在 $\Spec(A)$ 中为闭集。
由于该像包含 $\Spec(A)$ 的一般点，存在位于 $x$ 的闭包中的点
$x' \in |X|$，映至 $\Spec(A)$ 的闭点。由引理
\ref{spaces-morphisms-lemma-reach-points-scheme-theoretic-image}，可找到交换图
$$
\xymatrix{
\Spec(K') \ar[r] \ar[d] & \Spec(K) \ar[d] \\
\Spec(A') \ar[r] & X
}
$$
，使得 $\Spec(A')$ 的闭点映至 $x' \in |X|$。由此，
$\Spec(A') \to \Spec(A)$ 把闭点映至闭点；换言之，$A'$ 支配 $A$，
证明完毕。
\end{proof}

\begin{lemma}
\label{spaces-morphisms-lemma-characterize-universally-closed-separated}
设 $S$ 为概形，$f : X \to Y$ 为 $S$ 上代数空间的态射。设 $f$ 拟紧且分离。
下列条件等价：
\begin{enumerate}
\item $f$ 普遍闭；
\item 定义 \ref{spaces-morphisms-definition-valuative-criterion} 中赋值判据的存在性部分成立；
\item 对任意由实线箭头构成的交换图
$$
\xymatrix{
\Spec(K) \ar[r] \ar[d] & X \ar[d] \\
\Spec(A) \ar[r] \ar@{-->}[ru] & Y
}
$$
，其中 $A$ 为分式域是 $K$ 的赋值环，都存在虚线箭头；换言之，$f$ 满足
《概形》定义 \ref{schemes-definition-valuative-criterion} 中赋值判据的存在性部分。
\end{enumerate}
\end{lemma}

\begin{proof}
由于 $f$ 分离，由引理 \ref{spaces-morphisms-lemma-usual-enough}，(2) 与 (3) 等价。
(3) 与 (1) 的等价性由引理
\ref{spaces-morphisms-lemma-characterize-universally-closed-quasi-separated} 得出。
\end{proof}

\begin{lemma}
\label{spaces-morphisms-lemma-lift-valuation-ring-through-flat-morphism}
设 $S$ 为概形，$f : X \to Y$ 为 $S$ 上代数空间的平坦态射。设
$\Spec(A) \to Y$ 为态射，其中 $A$ 为赋值环。若 $\Spec(A)$ 的闭点映至
$|Y|$ 中属于 $|X| \to |Y|$ 之像的一点，则存在交换图
$$
\xymatrix{
\Spec(A') \ar[r] \ar[d] & X \ar[d] \\
\Spec(A) \ar[r] & Y
}
$$
，其中 $A \to A'$ 为赋值环的扩张（《关于代数的进一步结果》定义
\ref{more-algebra-definition-extension-valuation-rings}）。
\end{lemma}

\begin{proof}
基变换 $X_A \to \Spec(A)$ 平坦（引理 \ref{spaces-morphisms-lemma-base-change-flat}），且
$\Spec(A)$ 的闭点属于 $|X_A| \to |\Spec(A)|$ 的像（《空间的性质》引理
\ref{spaces-properties-lemma-points-cartesian}）。故可设 $Y = \Spec(A)$。
设 $U \to X$ 为满 \'etale 态射，其中 $U$ 为概形。取 $u \in U$ 映至
$\Spec(A)$ 的闭点。考虑平坦局部环映射
$A \to B = \mathcal{O}_{U,u}$。由《代数》引理
\ref{algebra-lemma-ff-rings}，存在素理想 $\mathfrak q \subset B$，使得
$\mathfrak q$ 位于 $(0) \subset A$ 上方。由《代数》引理 \ref{algebra-lemma-dominate}，可找到
支配 $B/\mathfrak q$ 的赋值环 $A' \subset \kappa(\mathfrak q)$。
所诱导的态射 $\Spec(A') \to U \to X$ 即为本引理问题的解。
\end{proof}

\begin{lemma}
\label{spaces-morphisms-lemma-refined-valuative-criterion-universally-closed}
设 $S$ 为概形，并设 $f : X \to Y$ 与 $h : U \to X$ 为 $S$ 上代数空间的
态射。若
\begin{enumerate}
\item $f$ 与 $h$ 拟紧；
\item $|h|(|U|)$ 在 $|X|$ 中稠密；
\end{enumerate}
并且对任意由实线箭头构成的交换图
$$
\xymatrix{
\Spec(K) \ar[r] \ar[d] & U \ar[r] & X \ar[d] \\
\Spec(A) \ar[rr] \ar@{-->}[rru] & & Y
}
$$
，其中 $A$ 为分式域是 $K$ 的赋值环，都有
\begin{enumerate}
\item[(3)] 至多存在一个使该图交换的虚线箭头；
\item[(4)] 存在域扩张 $K'/K$、支配 $A$ 的赋值环
$A' \subset K'$ 及态射 $\Spec(A') \to X$，使得下图交换
$$
\xymatrix{
\Spec(K') \ar[r] \ar[d] & \Spec(K) \ar[r] & U \ar[r] & X \ar[d] \\
\Spec(A') \ar[r] \ar[rrru] & \Spec(A) \ar[rr] & & Y
}
$$
\end{enumerate}
，则 $f$ 普遍闭。若此外
\begin{enumerate}
\item[(5)] $f$ 拟分离，
\end{enumerate}
则 $f$ 分离且普遍闭。
\end{lemma}

\begin{proof}
设 (1)、(2)、(3)、(4) 成立。我们将验证 $f$ 满足赋值判据的存在性部分；
由引理 \ref{spaces-morphisms-lemma-quasi-compact-existence-universally-closed}，这将蕴含 $f$
普遍闭。为此，考虑交换图
\begin{equation}
\label{spaces-morphisms-equation-start-with}
\vcenter{
\xymatrix{
\Spec(K) \ar[r] \ar[d] & X \ar[d] \\
\Spec(A) \ar[r] & Y
}
}
\end{equation}
，其中 $A$ 为赋值环，$K$ 为 $A$ 的分式域。注意，赋值环与域均既约，故由
《空间的性质》引理 \ref{spaces-properties-lemma-map-into-reduction}，可以分别用
$U$、$X$、$S$ 的既约化替换它们。在此情形，$h(U)$ 稠密这一假设意味着
$h : U \to X$ 的概形论像是 $X$，见引理
\ref{spaces-morphisms-lemma-scheme-theoretic-image-reduced}。

\medskip\noindent
化归至 $Y$ 仿射的情形。取 \'etale 态射 $\Spec(R) \to Y$，使得
$\Spec(A)$ 的闭点映至 $\Im(|\Spec(R)| \to |Y|)$ 中的一个元素。由引理
\ref{spaces-morphisms-lemma-lift-valuation-ring-through-flat-morphism}，可找到赋值环的局部环映射
$A \to A'$ 以及嵌入下列交换图的态射 $\Spec(A') \to \Spec(R)$
$$
\xymatrix{
\Spec(A') \ar[r] \ar[d] & \Spec(R) \ar[d] \\
\Spec(A) \ar[r] & Y
}
$$
。由于定义 \ref{spaces-morphisms-definition-valuative-criterion} 允许赋值环扩张，显然可以用
$A'$ 替换 $A$，用 $\Spec(R)$ 替换 $Y$，用 $X \times_Y \Spec(R)$ 替换 $X$，
并用 $U \times_Y \Spec(R)$ 替换 $U$。

\medskip\noindent
从现在起，设 $Y = \Spec(R)$ 为仿射概形。设 $\Spec(B) \to X$ 为来自仿射
概形的 \'etale 态射，并使态射 $\Spec(K) \to X$ 属于
$|\Spec(B)| \to |X|$ 的像。由于可以用扩张 $K' \supset K$ 替换 $K$，并用
支配 $A$ 的赋值环 $A' \subset K'$ 替换 $A$（它由《代数》引理
\ref{algebra-lemma-dominate} 存在），由 $|X|$ 的定义，可以设态射
$\Spec(K) \to X$ 分解经过 $\Spec(B)$。换言之，可以把 $K$ 视为
$B$-代数。取 $B$ 上的多项式代数 $P$ 及 $B$-代数满射 $P \to K$。
于是 $\Spec(P) \to X$ 作为复合 $\Spec(P) \to \Spec(B) \to X$ 而平坦。
因此，由引理 \ref{spaces-morphisms-lemma-flat-base-change-scheme-theoretic-image}，态射
$U \times_X \Spec(P) \to \Spec(P)$ 的概形论像是 $\Spec(P)$。由引理
\ref{spaces-morphisms-lemma-reach-points-scheme-theoretic-image}，可找到交换图
$$
\xymatrix{
\Spec(K') \ar[r] \ar[d] & U \times_X \Spec(P) \ar[d] \\
\Spec(A') \ar[r] & \Spec(P)
}
$$
，其中 $A'$ 为赋值环，$K'$ 为 $A'$ 的分式域，并且 $\Spec(A')$ 的闭点映至
$\Spec(K) \subset \Spec(P)$。换言之，存在 $B$-代数映射
$\varphi : K \to A'/\mathfrak m_{A'}$。取支配 $\varphi(A)$ 且分式域为
$K'' = A'/\mathfrak m_{A'}$ 的赋值环
$A'' \subset A'/\mathfrak m_{A'}$（《代数》引理
\ref{algebra-lemma-dominate}）。令
$$
C = \{\lambda \in A' \mid \lambda \bmod \mathfrak m_{A'} \in A''\}.
$$
由《代数》引理 \ref{algebra-lemma-stack-valuation-rings}，这是赋值环。
由于 $C$ 是分式域为 $K'$ 的 $R$-代数，得到由实线箭头构成的交换图
$$
\xymatrix{
\Spec(K'_1) \ar@{-->}[r] \ar@{-->}[d] &
\Spec(K') \ar[r] \ar[d] & U \ar[r] & X \ar[d] \\
\Spec(C_1) \ar@{-->}[r] \ar@{-->}[rrru] & \Spec(C) \ar[rr] & & Y
}
$$
，如引理陈述所示。因此假设 (4) 产生 $C \to C_1$ 以及使该图交换的虚线箭头。
令 $A_1' = (C_1)_\mathfrak p$ 为 $C_1$ 在位于
$\mathfrak m_{A'} \subset C$ 上方的素理想 $\mathfrak p \subset C_1$ 处的局部化。
由《关于代数的进一步结果》引理
\ref{more-algebra-lemma-valuation-ring-torsion-free-flat}，$C \to C_1$ 平坦，
故这样的素理想 $\mathfrak p$ 由《代数》引理
\ref{algebra-lemma-local-flat-ff} 与 \ref{algebra-lemma-ff-rings} 存在。
注意，$A'$ 是 $C$ 在 $\mathfrak m_{A'}$ 处的局部化，而 $A'_1$ 是赋值环
（《代数》引理 \ref{algebra-lemma-make-valuation-rings}）。换言之，
$A' \to A'_1$ 是赋值环的局部环映射。假设 (3) 蕴含图
$$
\xymatrix{
\Spec(A'_1) \ar[r] \ar[d] & \Spec(C_1) \ar[r] & X \\
\Spec(A') \ar[r] & \Spec(P) \ar[r] & \Spec(B) \ar[u]
}
$$
交换。因此，态射 $\Spec(C_1) \to X$ 在 $\Spec(C_1/\mathfrak p)$ 上的限制，
在 $\Spec(C_1/\mathfrak p)$ 的一般点处限制为复合
$$
\Spec(\kappa(\mathfrak p)) \to
\Spec(A'/\mathfrak m_{A'}) = \Spec(K'') \to
\Spec(K) \to X
$$
。此外，$C_1/\mathfrak p$ 是赋值环（《代数》引理
\ref{algebra-lemma-make-valuation-rings}），它支配 $A''$，而后者支配 $A$。
因此，态射 $\Spec(C_1/\mathfrak p) \to X$ 如所需见证了图
(\ref{spaces-morphisms-equation-start-with}) 的赋值判据之存在性部分。

\medskip\noindent
接下来，设 (5) 也成立；换言之，态射
% P08-E247: frozen source diagonal base preserved; see ERRATA_CANDIDATES.jsonl.
$\Delta : X \to X \times_S X$ 拟紧。此时，$h$ 与 $\Delta$ 满足假设
(1)--(4)。故证明的第一部分表明 $\Delta$ 普遍闭。由引理
\ref{spaces-morphisms-lemma-separated-diagonal-proper}，可知 $f$ 分离。
\end{proof}

















\section{分离性的赋值判据}
\label{spaces-morphisms-section-valuative-separatedness}

\noindent
我们先证明一个逆命题，然后陈述该判据。

\begin{lemma}
\label{spaces-morphisms-lemma-separated-implies-valuative}
设 $S$ 为概形。设 $f : X \to Y$ 为 $S$ 上代数空间的态射。
若 $f$ 分离，则 $f$ 满足赋值判据的唯一性部分。
\end{lemma}

\begin{proof}
设给定定义 \ref{spaces-morphisms-definition-valuative-criterion} 中的图。假设有两个不同的态射
$a, b : \Spec(A) \to X$ 可嵌入该图。令 $Z \subset \Spec(A)$ 为 $a$ 与 $b$
的等化子。于是
$Z = \Spec(A) \times_{(a, b), X \times_Y X, \Delta} X$。
若 $f$ 分离，则 $\Delta$ 为闭浸入，故这是 $\Spec(A)$ 的闭子概形。
依假设，它包含 $\Spec(A)$ 的一般点。由于 $A$ 是整环，这蕴含
$Z = \Spec(A)$。因此如所需 $a = b$。
\end{proof}

% P08-E248: frozen source lemma-title grammar noted; see ERRATA_CANDIDATES.jsonl.
\begin{lemma}[分离性的赋值判据]
\label{spaces-morphisms-lemma-valuative-criterion-separatedness}
设 $S$ 为概形。设 $f : X \to Y$ 为 $S$ 上代数空间的态射。假设
\begin{enumerate}
\item 态射 $f$ 拟分离；
\item 态射 $f$ 满足赋值判据的唯一性部分。
\end{enumerate}
则 $f$ 分离。
\end{lemma}

\begin{proof}
假设 (1) 意味着 $\Delta_{X/Y}$ 拟紧。我们断言态射
$\Delta_{X/Y} : X \to X \times_Y X$ 满足赋值判据的存在性部分。
设给定由实线箭头构成的交换图
$$
\xymatrix{
\Spec(K) \ar[r] \ar[d] & X \ar[d] \\
\Spec(A) \ar[r] \ar@{-->}[ru] & X \times_Y X
}
$$
。右下方的箭头对应于 $Y$ 上的一对态射
$a, b : \Spec(A) \to X$。由假设 (2)，可知 $a = b$，故以 $a$ 为虚线箭头即可。
于是引理 \ref{spaces-morphisms-lemma-quasi-compact-existence-universally-closed} 适用，从而
$\Delta_{X/Y}$ 普遍闭。由于 $\Delta_{X/Y}$ 总是局部有限型且分离，由
《关于态射的进一步结果》引理 \ref{more-morphisms-lemma-characterize-finite}，
可知 $\Delta_{X/Y}$ 是有限态射（也可使用《空间》引理
\ref{spaces-lemma-representable-transformations-property-implication}
的一般原理）。此时 $\Delta_{X/Y}$ 是可表的有限单态射，因而由《态射》引理
\ref{morphisms-lemma-finite-monomorphism-closed}，它是闭浸入。
\end{proof}

\begin{lemma}
\label{spaces-morphisms-lemma-characterize-separated-and-universally-closed}
设 $S$ 为概形。设 $f : X \to Y$ 为 $S$ 上代数空间的态射。假设 $f$ 拟紧且
拟分离。则下列条件等价：
\begin{enumerate}
\item $f$ 分离且普遍闭；
\item 定义 \ref{spaces-morphisms-definition-valuative-criterion} 中的赋值判据成立；
\item 对任意由实线箭头构成的交换图
$$
\xymatrix{
\Spec(K) \ar[r] \ar[d] & X \ar[d] \\
\Spec(A) \ar[r] \ar@{-->}[ru] & Y
}
$$
，其中 $A$ 为分式域是 $K$ 的赋值环，都存在唯一的虚线箭头；换言之，$f$
满足《概形》定义 \ref{schemes-definition-valuative-criterion} 中的赋值判据。
\end{enumerate}
\end{lemma}

\begin{proof}
由于 $f$ 拟分离，赋值判据的唯一性部分蕴含 $f$ 分离
% P08-E249: frozen source spelling error noted; see ERRATA_CANDIDATES.jsonl.
（引理 \ref{spaces-morphisms-lemma-valuative-criterion-separatedness}）。反之，若 $f$ 分离，
则它满足赋值判据的唯一性部分（引理 \ref{spaces-morphisms-lemma-separated-implies-valuative}）。
综上，在三个情形中，态射 $f$ 均分离且满足赋值判据的唯一性部分。此时本引理是
引理 \ref{spaces-morphisms-lemma-characterize-universally-closed-separated} 的形式推论。
\end{proof}



\section{固有性的赋值判据}
\label{spaces-morphisms-section-valuative-criterion-properness}

\noindent
现给出如下命题。

\begin{lemma}[固有性的赋值判据]
\label{spaces-morphisms-lemma-characterize-proper}
设 $S$ 为概形。设 $f : X \to Y$ 为 $S$ 上代数空间的态射。假设 $f$ 有限型且
拟分离。则下列条件等价：
\begin{enumerate}
\item $f$ 固有；
\item 定义 \ref{spaces-morphisms-definition-valuative-criterion} 中的赋值判据成立；
\item 对任意由实线箭头构成的交换图
$$
\xymatrix{
\Spec(K) \ar[r] \ar[d] & X \ar[d] \\
\Spec(A) \ar[r] \ar@{-->}[ru] & Y
}
$$
，其中 $A$ 为分式域是 $K$ 的赋值环，都存在唯一的虚线箭头；换言之，$f$
满足《概形》定义 \ref{schemes-definition-valuative-criterion} 中的赋值判据。
\end{enumerate}
\end{lemma}

\begin{proof}
这是引理 \ref{spaces-morphisms-lemma-characterize-separated-and-universally-closed} 与各定义的
形式推论。
\end{proof}





\section{整态射与有限态射}
\label{spaces-morphisms-section-integral}

\noindent
我们已经在第 \ref{spaces-morphisms-section-representable} 节中定义了代数空间的可表示态射何谓
整态射（相应地，有限态射）。

\begin{lemma}
\label{spaces-morphisms-lemma-integral-representable}
设 $S$ 为概形。设 $f : X \to Y$ 为 $S$ 上代数空间的可表示态射。则
$f$ 为整态射，相应地为有限态射（依第 \ref{spaces-morphisms-section-representable} 节之意义），
当且仅当对所有仿射概形 $Z$ 及态射 $Z \to Y$，概形 $X \times_Y Z$ 为仿射的，
且在 $Z$ 上为整的，相应地为有限的。
\end{lemma}

\begin{proof}
这直接由概形的整（相应地，有限）态射之定义得出
（《态射》定义 \ref{morphisms-definition-integral}）。
\end{proof}

\noindent
由此可以作如下定义。

\begin{definition}
\label{spaces-morphisms-definition-integral}
设 $S$ 为概形。设 $f : X \to Y$ 为 $S$ 上代数空间的态射。
\begin{enumerate}
\item 若对每个仿射概形 $Z$ 及态射 $Z \to Y$，代数空间 $X \times_Y Z$
都可由一个在 $Z$ 上为整的仿射概形表示，则称 $f$ 为{\it 整态射}。
\item 若对每个仿射概形 $Z$ 及态射 $Z \to Y$，代数空间 $X \times_Y Z$
都可由一个在 $Z$ 上为有限的仿射概形表示，则称 $f$ 为{\it 有限态射}。
\end{enumerate}
\end{definition}

\begin{lemma}
\label{spaces-morphisms-lemma-integral-local}
设 $S$ 为概形。设 $f : X \to Y$ 为 $S$ 上代数空间的态射。下列条件等价：
\begin{enumerate}
\item $f$ 可表且为整态射（相应地，有限态射）；
\item $f$ 为整态射（相应地，有限态射）；
\item 存在概形 $V$ 及满 \'etale 态射 $V \to Y$，使得
$V \times_Y X \to V$ 为整态射（相应地，有限态射）；
\item 存在扎里斯基覆盖 $Y = \bigcup Y_i$，使每个态射
$f^{-1}(Y_i) \to Y_i$ 都为整态射（相应地，有限态射）。
\end{enumerate}
\end{lemma}

\begin{proof}
显然，(1) 蕴含 (2)；取 $V$ 为 $Y$ 上 \'etale 仿射概形的不交并，即知
(2) 蕴含 (3)，见《空间的性质》引理
\ref{spaces-properties-lemma-cover-by-union-affines}。
假设 $V \to Y$ 如 (3) 所示。则对 $V$ 的每个仿射开集 $W$，
$W \times_Y X$ 都是 $V \times_Y X$ 的仿射开集。因此，由
《空间的性质》引理 \ref{spaces-properties-lemma-subscheme}，可知
$V \times_Y X$ 是概形。此外，态射 $V \times_Y X \to V$ 是仿射的。
这意味着可以应用《空间》引理
\ref{spaces-lemma-morphism-sheaves-with-P-effective-descent-etale}，
因为整（相应地，有限）态射类满足所有所需性质（见《态射》引理
\ref{morphisms-lemma-base-change-finite}，以及《下降》引理
\ref{descent-lemma-descending-property-integral}、
\ref{descent-lemma-descending-property-finite} 和 \ref{descent-lemma-affine}）。
应用该引理所得的结论是 $f$ 可表且为整态射（相应地，有限态射），即 (1) 成立。

\medskip\noindent
(1) 与 (4) 的等价性由整态射（相应地，有限态射）这一性质在目标上是扎里斯基
局部的得出（上面的文献事实上表明，整态射或有限态射在目标上是 fpqc 局部的）。
\end{proof}

\begin{lemma}
\label{spaces-morphisms-lemma-composition-integral}
整（相应地，有限）态射的复合仍为整（相应地，有限）态射。
\end{lemma}

\begin{proof}
略。
\end{proof}

\begin{lemma}
\label{spaces-morphisms-lemma-base-change-integral}
整（相应地，有限）态射的基变换仍为整（相应地，有限）态射。
\end{lemma}

\begin{proof}
略。
\end{proof}

\begin{lemma}
\label{spaces-morphisms-lemma-finite-integral}
代数空间的有限态射是整态射。代数空间的局部有限型整态射是有限态射。
\end{lemma}

\begin{proof}
在两种情形中，态射都是可表的，并且可以在用映入 $Y$ 的仿射概形作基变换后
检验该条件，见引理 \ref{spaces-morphisms-lemma-integral-local}。因此，本引理由概形情形的
同一引理得出，见《态射》引理 \ref{morphisms-lemma-finite-integral}。
\end{proof}

\begin{lemma}
\label{spaces-morphisms-lemma-integral-universally-closed}
设 $S$ 为概形。设 $f : X \to Y$ 为 $S$ 上代数空间的态射。下列条件等价：
\begin{enumerate}
\item $f$ 为整态射；
\item $f$ 仿射且普遍闭。
\end{enumerate}
\end{lemma}

\begin{proof}
在两种情形中，态射都是可表的，并且可以在用映入 $Y$ 的仿射概形作基变换后
检验该条件，见引理 \ref{spaces-morphisms-lemma-integral-local}、\ref{spaces-morphisms-lemma-affine-local} 和
\ref{spaces-morphisms-lemma-universally-closed-local}。因此，结论由《态射》引理
\ref{morphisms-lemma-integral-universally-closed} 得出。
\end{proof}

\begin{lemma}
\label{spaces-morphisms-lemma-finite-quasi-finite}
代数空间的有限态射是拟有限态射。
\end{lemma}

\begin{proof}
设 $f : X \to Y$ 为代数空间的态射。由定义 \ref{spaces-morphisms-definition-integral} 及引理
\ref{spaces-morphisms-lemma-quasi-compact-local}、\ref{spaces-morphisms-lemma-quasi-finite-local}，这两个性质
都可在基变换至 $Y$ 上的仿射概形后检验；换言之，可设 $Y$ 仿射。
若 $f$ 有限，则 $X$ 是概形。因此，结论由概形的相应结果得出，见《态射》引理
\ref{morphisms-lemma-finite-quasi-finite}。
\end{proof}

\begin{lemma}
\label{spaces-morphisms-lemma-finite-proper}
设 $S$ 为概形。设 $f : X \to Y$ 为 $S$ 上代数空间的态射。下列条件等价：
\begin{enumerate}
\item $f$ 有限；
\item $f$ 仿射且固有。
\end{enumerate}
\end{lemma}

\begin{proof}
在两种情形中，态射都是可表的，并且可以在用映入 $Y$ 的仿射概形作基变换后
检验该条件，见引理 \ref{spaces-morphisms-lemma-integral-local}、\ref{spaces-morphisms-lemma-affine-local} 和
\ref{spaces-morphisms-lemma-proper-local}。因此，结论由《态射》引理
\ref{morphisms-lemma-finite-proper} 得出。
\end{proof}

\begin{lemma}
\label{spaces-morphisms-lemma-closed-immersion-finite}
闭浸入是有限的（因而更是整的）。
\end{lemma}

\begin{proof}
略。
\end{proof}

\begin{lemma}
\label{spaces-morphisms-lemma-finite-union-finite}
设 $S$ 为概形。设 $X_i \to Y$，$i = 1, \ldots, n$，为 $S$ 上代数空间的
有限态射。则 $X_1 \amalg \ldots \amalg X_n \to Y$ 也有限。
\end{lemma}

\begin{proof}
通过 \'etale 局部化，这由概形情形得出
（《态射》引理 \ref{morphisms-lemma-finite-union-finite}）。
\end{proof}

\begin{lemma}
\label{spaces-morphisms-lemma-finite-permanence}
设 $S$ 为概形。设 $f : X \to Y$ 与 $g : Y \to Z$ 为 $S$ 上代数空间的态射。
\begin{enumerate}
\item 若 $g \circ f$ 有限且 $g$ 分离，则 $f$ 有限。
\item 若 $g \circ f$ 为整态射且 $g$ 分离，则 $f$ 为整态射。
\end{enumerate}
\end{lemma}

\begin{proof}
假设 $g \circ f$ 有限（相应地，为整态射）且 $g$ 分离。由引理
\ref{spaces-morphisms-lemma-base-change-integral}，基变换 $X \times_Z Y \to Y$ 有限
（相应地，为整态射）。由于 $Y \to Z$ 分离，态射
$X \to X \times_Z Y$ 是闭浸入，见引理 \ref{spaces-morphisms-lemma-section-immersion}。
闭浸入是有限态射（相应地，整态射），见引理
\ref{spaces-morphisms-lemma-closed-immersion-finite}。有限态射（相应地，整态射）的复合仍为
有限态射（相应地，整态射），见引理
\ref{spaces-morphisms-lemma-composition-integral}。故结论成立。
\end{proof}





\section{有限局部自由态射}
\label{spaces-morphisms-section-finite-locally-free}

\noindent
我们已经在第 \ref{spaces-morphisms-section-representable} 节中定义了代数空间的可表示态射何谓
有限局部自由态射。

\begin{lemma}
\label{spaces-morphisms-lemma-finite-locally-free-representable}
设 $S$ 为概形。设 $f : X \to Y$ 为 $S$ 上代数空间的可表示态射。则 $f$
为有限局部自由态射（依第 \ref{spaces-morphisms-section-representable} 节之意义），当且仅当
$f$ 仿射且层 $f_*\mathcal{O}_X$ 是有限局部自由 $\mathcal{O}_Y$-模。
\end{lemma}

\begin{proof}
假设 $f$ 为有限局部自由态射（依第 \ref{spaces-morphisms-section-representable} 节之定义）。
这意味着，对每个源为概形的态射 $V \to Y$，基变换
$f' : V \times_Y X \to V$ 都是概形的有限局部自由态射。这又依概形的
有限局部自由态射之定义，这意味着 $f'_*\mathcal{O}_{V \times_Y X}$
是有限局部自由 $\mathcal{O}_V$-模。可以选取满 \'etale 态射 $V \to Y$。
由《空间的性质》引理
\ref{spaces-properties-lemma-pushforward-etale-base-change-modules}，
可知 $f_*\mathcal{O}_X$ 在 $V$ 上的限制有限局部自由。因此，把《站点上的模》
引理 \ref{sites-modules-lemma-local-final-object} 应用于
$Y_{spaces, \etale}$ 上的层 $f_*\mathcal{O}_X$，可知
$f_*\mathcal{O}_X$ 有限局部自由。

\medskip\noindent
反之，假设 $f$ 仿射，且 $f_*\mathcal{O}_X$ 是有限局部自由
$\mathcal{O}_Y$-模。设 $V$ 为概形，并设 $V \to Y$ 为满 \'etale 态射。
再次由《空间的性质》引理
\ref{spaces-properties-lemma-pushforward-etale-base-change-modules}，
可知 $f'_*\mathcal{O}_{V \times_Y X}$ 有限局部自由。因此，
$f' : V \times_Y X \to V$ 为有限局部自由态射（因为它也仿射）。
由《空间》引理
\ref{spaces-lemma-morphism-sheaves-with-P-effective-descent-etale}，
可知 $f$ 为有限局部自由态射（使用《态射》引理
\ref{morphisms-lemma-base-change-finite-locally-free}，
% P08-E251: frozen source citation separators noted; see ERRATA_CANDIDATES.jsonl.
以及《下降》引理 \ref{descent-lemma-descending-property-finite-locally-free}
和 \ref{descent-lemma-affine}）。故结论成立。
\end{proof}

\noindent
由此可以作如下定义。

\begin{definition}
\label{spaces-morphisms-definition-finite-locally-free}
设 $S$ 为概形。设 $f : X \to Y$ 为 $S$ 上代数空间的态射。若 $f$ 仿射且
$f_*\mathcal{O}_X$ 是有限局部自由 $\mathcal{O}_Y$-模，则称 $f$
为{\it 有限局部自由态射}。
% P08-E250: frozen source grammar noted; see ERRATA_CANDIDATES.jsonl.
在此情形，若层 $f_*\mathcal{O}_X$ 有限局部自由且秩为 $d$，则称 $f$
的{\it 秩}或{\it 次数}为 $d$。
\end{definition}

\begin{lemma}
\label{spaces-morphisms-lemma-finite-locally-free-local}
设 $S$ 为概形。设 $f : X \to Y$ 为 $S$ 上代数空间的态射。下列条件等价：
\begin{enumerate}
\item $f$ 可表且为有限局部自由态射；
\item $f$ 为有限局部自由态射；
\item 存在概形 $V$ 及满 \'etale 态射 $V \to Y$，使得
$V \times_Y X \to V$ 为有限局部自由态射；
\item 存在扎里斯基覆盖 $Y = \bigcup Y_i$，使得每个态射
$f^{-1}(Y_i) \to Y_i$ 都为有限局部自由态射。
\end{enumerate}
\end{lemma}

\begin{proof}
显然，(1) 蕴含 (2)；取 $V$ 为 $Y$ 上 \'etale 仿射概形的不交并，即知
(2) 蕴含 (3)，见《空间的性质》引理
\ref{spaces-properties-lemma-cover-by-union-affines}。
假设 $V \to Y$ 如 (3) 所示。则对 $V$ 的每个仿射开集 $W$，
$W \times_Y X$ 都是 $V \times_Y X$ 的仿射开集。因此，由《空间的性质》
引理 \ref{spaces-properties-lemma-subscheme}，可知 $V \times_Y X$ 是概形。
此外，态射 $V \times_Y X \to V$ 是仿射的。这意味着可以应用《空间》引理
\ref{spaces-lemma-morphism-sheaves-with-P-effective-descent-etale}，
因为有限局部自由态射类满足所有所需性质（见《态射》引理
\ref{morphisms-lemma-base-change-finite-locally-free}，
以及《下降》引理 \ref{descent-lemma-descending-property-finite-locally-free}
和 \ref{descent-lemma-affine}）。应用该引理所得的结论是 $f$ 可表且为
有限局部自由态射，即 (1) 成立。

\medskip\noindent
(1) 与 (4) 的等价性由有限局部自由这一性质在目标上是扎里斯基局部的得出
（上面的文献事实上表明，有限局部自由在目标上是 fpqc 局部的）。
\end{proof}

\begin{lemma}
\label{spaces-morphisms-lemma-composition-finite-locally-free}
有限局部自由态射的复合仍为有限局部自由态射。
\end{lemma}

\begin{proof}
略。
\end{proof}

\begin{lemma}
\label{spaces-morphisms-lemma-base-change-finite-locally-free}
有限局部自由态射的基变换仍为有限局部自由态射。
\end{lemma}

\begin{proof}
略。
\end{proof}

\begin{lemma}
\label{spaces-morphisms-lemma-finite-flat}
设 $S$ 为概形。设 $f : X \to Y$ 为 $S$ 上代数空间的态射。下列条件等价：
\begin{enumerate}
\item $f$ 为有限局部自由态射；
\item $f$ 有限、平坦且局部有限表示。
\end{enumerate}
若 $Y$ 局部诺特，则这些条件还等价于
\begin{enumerate}
\item[(3)] $f$ 有限且平坦。
\end{enumerate}
\end{lemma}

\begin{proof}
在三个情形中，态射都是可表的，并且可以在用满 \'etale 态射
$V \to Y$ 作基变换后检验该性质，见引理 \ref{spaces-morphisms-lemma-integral-local}、
\ref{spaces-morphisms-lemma-finite-locally-free-local}、\ref{spaces-morphisms-lemma-flat-local} 和
\ref{spaces-morphisms-lemma-finite-presentation-local}。若 $Y$ 局部诺特，则 $V$ 局部诺特。
因此，结论由概形情形的相应结果得出，见《态射》引理
\ref{morphisms-lemma-finite-flat}。
\end{proof}





\section{有理映射}
\label{spaces-morphisms-section-rational-maps}

\noindent
本节是《态射》第 \ref{morphisms-section-rational-maps} 节的对应版本。
我们将不再特别说明地使用如下事实：拓扑空间中稠密开集的交仍是稠密开集。

\begin{definition}
\label{spaces-morphisms-definition-rational-map}
设 $S$ 为概形。设 $X$、$Y$ 为 $S$ 上的代数空间。
\begin{enumerate}
\item 设 $f : U \to Y$、$g : V \to Y$ 为 $S$ 上代数空间的态射，分别定义在
$X$ 的稠密开子空间 $U$、$V$ 上。若对某个稠密开子空间
$W \subset U \cap V$ 有 $f|_W = g|_W$，则称 $f$ 与 $g$ {\it 等价}。
\item 从 $X$ 到 $Y$ 的{\it 有理映射}是 (1) 中所定义等价关系的一个等价类。
\item 给定 $S$ 上代数空间的态射 $X \to B$ 与 $Y \to B$。若该等价类有一个
代表元 $f : U \to Y$ 是 $B$ 上的态射，则称从 $X$ 到 $Y$ 的有理映射为
{\it 从 $X$ 到 $Y$ 的 $B$-有理映射}。
\end{enumerate}
\end{definition}

\noindent
对定义 (1) 中的两个态射 $f$、$g$，我们说它们定义同一个有理映射，而不说
它们等价。在以后将考虑的许多情形中，代数空间 $X$ 与 $Y$ 分别含有作为概形的
% P08-E252: frozen source agreement error noted; see ERRATA_CANDIDATES.jsonl.
稠密开子空间 $X'$ 与 $Y'$。此时，从 $X$ 到 $Y$ 的有理映射与《态射》定义
% P08-E253: frozen source cross-reference target preserved; see ERRATA_CANDIDATES.jsonl.
\ref{spaces-morphisms-definition-rational-map} 意义下从 $X'$ 到 $Y'$ 的 $S$-有理映射相同。
于是可以运用为概形建立的全部理论。

\begin{definition}
\label{spaces-morphisms-definition-rational-function}
设 $S$ 为概形。设 $X$ 为 $S$ 上的代数空间。$X$ 上的{\it 有理函数}是从
$X$ 到 $\mathbf{A}^1_S$ 的有理映射。
\end{definition}

\noindent
由《态射》定义 \ref{morphisms-definition-rational-function} 后的讨论可知，
当 $X$ 恰为概形时，这与该处定义的概念相同。

\medskip\noindent
回忆，对 $S$ 上任意概形 $T$，都有典范恒等式
$$
\Mor_S(T, \mathbf{A}^1_S) = \Mor(T, \mathbf{A}^1_\mathbf{Z}) =
\Gamma(T, \mathcal{O}_T)
$$
，见《概形》例 \ref{schemes-example-global-sections}。因此，
$\mathbf{A}^1_S$ 是 $S$ 上概形范畴中的环对象。换言之，加法与乘法定义态射
$$
+ : \mathbf{A}^1_S \times_S \mathbf{A}^1_S \to \mathbf{A}^1_S
\quad\text{且}\quad
* : \mathbf{A}^1_S \times_S \mathbf{A}^1_S \to \mathbf{A}^1_S
$$
，并满足环中加法与乘法的公理（照例是带 $1$ 的交换环）。因此，取值于
$\mathbf{A}^1_S$ 的有理映射集也有自然的环结构。

\begin{definition}
\label{spaces-morphisms-definition-ring-of-rational-functions}
设 $S$ 为概形。设 $X$ 为 $S$ 上的代数空间。$X$ 上的{\it 有理函数环}是环
$R(X)$；其元素为有理函数，加法与乘法如上所述。
\end{definition}

\noindent
稍后我们将定义整代数空间的函数域，见《域上的空间》第
\ref{spaces-over-fields-section-integral-spaces} 节。

\begin{definition}
\label{spaces-morphisms-definition-domain-of-definition}
设 $S$ 为概形。设 $\varphi$ 为 $S$ 上两个代数空间 $X$ 与 $Y$ 之间的有理映射。
若 $\varphi$ 有一个代表元 $(U, f)$ 满足 $x \in |U|$，则称 $\varphi$
{\it 在点 $x \in |X|$ 处有定义}。$\varphi$ 的{\it 定义域}是 $\varphi$
有定义的所有点之集。
\end{definition}

\noindent
通过《空间的性质》引理 \ref{spaces-properties-lemma-open-subspaces}，
把定义域视为 $X$ 的开子空间。依此定义，一般而言并非必有 $\varphi$ 的
某个代表元定义在整个定义域上。

\begin{lemma}
\label{spaces-morphisms-lemma-rational-map-from-reduced-to-separated}
设 $S$ 为概形。设 $X$ 与 $Y$ 为 $S$ 上的代数空间。假设 $X$ 既约且 $Y$
在 $S$ 上分离。设 $\varphi$ 为从 $X$ 到 $Y$ 的有理映射，其定义域为
$U \subset X$。则存在唯一的代数空间态射 $f : U \to Y$ 表示 $\varphi$。
\end{lemma}

\begin{proof}
设 $(V, g)$ 与 $(V', g')$ 为 $\varphi$ 的代表元。则 $g, g'$ 在某个稠密开子空间
$W \subset V \cap V'$ 上相等。另一方面，$g|_{V \cap V'}$ 与
$g'|_{V \cap V'}$ 的等化子 $E$ 是 $V \cap V'$ 的闭子空间，因为它是将
$\Delta : Y \to Y \times_S Y$ 沿由 $g|_{V \cap V'}$ 与 $g'|_{V \cap V'}$
给出的态射 $V \cap V' \to Y \times_S Y$ 作基变换所得。现由
$W \subset E$ 可知 $|E| = |V \cap V'|$。由于 $V \cap V'$ 既约，在概形论意义下
$E = V \cap V'$；换言之，$g|_{V \cap V'} = g'|_{V \cap V'}$，见
《空间的性质》引理 \ref{spaces-properties-lemma-map-into-reduction}。
因此，可以把 $\varphi$ 的各代表元 $g : V \to Y$ 粘合为态射
$f : U \to Y$，因为 $\coprod V \to U$ 是 fppf 层的满射，且
$\coprod_{V, V'} V \cap V' = (\coprod V) \times_U (\coprod V)$。
\end{proof}

\noindent
一般而言，复合有理映射没有意义。原因是第一个有理映射的代表元之像可能与
第二个有理映射的定义域没有交。不过，若假设所论空间不可约，并只考察
支配有理映射，则可以复合有理映射。

\begin{definition}
\label{spaces-morphisms-definition-dominant-rational}
设 $S$ 为概形。设 $X$ 与 $Y$ 为 $S$ 上的代数空间。假设 $|X|$ 与 $|Y|$
不可约。若从 $X$ 到 $Y$ 的有理映射的任意代表元 $f : U \to Y$ 都是定义
\ref{spaces-morphisms-definition-dominant} 意义下的支配态射，则称该有理映射是{\it 支配的}。
\end{definition}

\noindent
可以把不可约代数空间 $X$ 与 $Y$ 之间的支配有理映射 $\varphi$，同从 $Y$
到 $Z$ 的任意有理映射 $\psi$ 复合。具体地，选取代表元
$f : U \to Y$，其中 $|U| \subset |X|$ 开且稠密；以及代表元
$g : V \to Z$，其中 $|V| \subset |Y|$ 开且稠密。则
$W = |f|^{-1}(V) \subset |X|$ 开且非空（因为 $|f|$ 的像稠密，故必与非空开集
$V$ 相交）；又因 $|X|$ 不可约，该开集稠密。定义 $\psi \circ \varphi$ 为
$g \circ f|_W : W \to Z$ 的等价类。略去其良定义性的验证。

\medskip\noindent
由此得到一个范畴，其对象为 $S$ 上不可约代数空间，态射为支配有理映射。

\begin{definition}
\label{spaces-morphisms-definition-birational-spaces}
设 $S$ 为概形。设 $X$ 与 $Y$ 为 $S$ 上的代数空间，且 $|X|$ 与 $|Y|$
不可约。若 $X$ 与 $Y$ 在 $S$ 上不可约代数空间及支配有理映射所成的范畴中
同构，则称 $X$ 与 $Y$ {\it 双有理}。
\end{definition}

\noindent
若 $X$ 与 $Y$ 是双有理的不可约代数空间，则对所有代数空间 $Z$，从 $X$
到 $Z$ 的有理映射集与从 $Y$ 到 $Z$ 的有理映射集双射（关于 $Z$ 函子性地）。
% P08-E254: frozen source number agreement noted; see ERRATA_CANDIDATES.jsonl.
对“一般的”不可约代数空间，这只是一个可能的定义。另一种定义是要求
% P08-E255: frozen source missing complementizer noted; see ERRATA_CANDIDATES.jsonl.
$X$ 与 $Y$ 具有同构的有理函数环；有时这两个概念等价
% P08-E256: frozen source editorial placeholder preserved; see ERRATA_CANDIDATES.jsonl.
（此处插入未来文献）。

\begin{lemma}
\label{spaces-morphisms-lemma-birational}
设 $S$ 为概形。设 $X$ 与 $Y$ 为 $S$ 上的代数空间，且 $|X|$ 与 $|Y|$
不可约。
% P08-E257: frozen source number agreement noted; see ERRATA_CANDIDATES.jsonl.
则 $X$ 与 $Y$ 双有理，当且仅当存在非空开子空间 $U \subset X$ 与
$V \subset Y$，它们作为 $S$ 上代数空间同构。
\end{lemma}

\begin{proof}
假设 $X$ 与 $Y$ 双有理。设 $f : U \to Y$ 与 $g : V \to X$ 分别定义从 $X$
到 $Y$ 以及从 $Y$ 到 $X$ 的互逆支配有理映射。缩小 $U$ 后，可以假设
$f : U \to Y$ 分解经过 $V$。由于 $g \circ f$ 作为支配有理映射是恒等映射，
可知复合 $U \to V \to X$ 在 $U$ 的某个稠密开集上为恒等映射。因此，用更小的
开集替换 $U$ 后，可以假设 $U \to V \to X$ 是 $U$ 到 $X$ 的包含映射。
由对称性，存在开子空间 $V' \subset V$，使得
$g|_{V'} : V' \to X$ 分解经过 $U \subset X$，且 $V' \to U \to Y$
为恒等映射。$|U| \to |V|$ 下 $|V'|$ 的逆像是 $|U|$ 的开集，因而等于某个
开子空间 $U' \subset U$ 的 $|U'|$，见《空间的性质》引理
\ref{spaces-properties-lemma-open-subspaces}。于是 $U' \subset U \to V$
分解为 $U' \to V'$。类似地，$V' \to U$ 分解为 $V' \to U'$。
容易验证，$U' \to V'$ 与 $V' \to U'$ 是 $S$ 上代数空间的互逆态射，
证明完毕。
\end{proof}









\section{代数空间的相对正规化}
\label{spaces-morphisms-section-normalization-X-in-Y}

\noindent
本节是《态射》第 \ref{morphisms-section-normalization-X-in-Y} 节的对应版本。

\begin{lemma}
\label{spaces-morphisms-lemma-integral-closure}
设 $S$ 为概形。设 $X$ 为 $S$ 上的代数空间。设 $\mathcal{A}$ 为
$\mathcal{O}_X$-代数的拟凝聚层。存在 $\mathcal{O}_X$-代数的拟凝聚层
$\mathcal{A}' \subset \mathcal{A}$，使得对 $X_\etale$ 的任意仿射对象 $U$，
环 $\mathcal{A}'(U) \subset \mathcal{A}(U)$ 是 $\mathcal{O}_X(U)$
在 $\mathcal{A}(U)$ 中的整闭包。
\end{lemma}

\begin{proof}
设 $U$ 为 $X_\etale$ 的对象。则 $U$ 是概形。用 $\mathcal{A}|_U$ 表示其在
扎里斯基站点上的限制。于是 $\mathcal{A}|_U$ 是 $\mathcal{O}_U$-代数的
拟凝聚层，故可以应用《态射》引理 \ref{morphisms-lemma-integral-closure}，
找到拟凝聚子代数 $\mathcal{A}'_U \subset \mathcal{A}|_U$，使得
$\mathcal{A}'_U$ 在任意仿射开集 $W \subset U$ 上的值如引理陈述所示。
若 $f : U' \to U$ 是 $X_\etale$ 中的态射，则
$\mathcal{A}|_{U'} = f^*(\mathcal{A}|_U)$，其中 $f^*$ 指在扎里斯基拓扑中
沿态射 $f$ 的拉回；这是因为 $\mathcal{A}$ 拟凝聚（见《空间的性质》第
\ref{spaces-properties-section-quasi-coherent} 节的引言，以及该处指向
概形下降一章相关讨论的文献）。由于 $f$ 为 \'etale，《关于态射的进一步结果》
引理 \ref{more-morphisms-lemma-integral-closure-smooth-pullback} 给出典范同构
$f^*(\mathcal{A}'_U) = \mathcal{A}'_{U'}$。这立即说明，我们在 $X_\etale$
上得到 $\mathcal{O}_X$-代数的子预层
$\mathcal{A}' \subset \mathcal{A}$；它对扎里斯基拓扑成层，并在仿射对象上
具有所需的值。又因为每个 $\mathcal{A}'_U$ 在概形 $U$ 上拟凝聚，且对
\'etale 态射 $f : U' \to U$ 有
$\mathcal{A}'_{U'} = f^*(\mathcal{A}'_U)$，所以 $\mathcal{A}'$ 在
$X_\etale$ 上也拟凝聚（这是局部性质，而上面的文献正是以这种方式描述
$U_\etale$ 上的拟凝聚模）。
\end{proof}

\begin{definition}
\label{spaces-morphisms-definition-integral-closure}
设 $S$ 为概形。设 $X$ 为 $S$ 上的代数空间。设 $\mathcal{A}$ 为
$\mathcal{O}_X$-代数的拟凝聚层。$\mathcal{O}_X$
{\it 在 $\mathcal{A}$ 中的整闭包}是上面引理
\ref{spaces-morphisms-lemma-integral-closure} 中构造的拟凝聚 $\mathcal{O}_X$-子代数
$\mathcal{A}' \subset \mathcal{A}$。
\end{definition}

\noindent
特别地，我们将把它应用于
$\mathcal{A} = f_*\mathcal{O}_Y$ 的情形，其中 $f : Y \to X$
为代数空间的拟紧拟分离态射（见引理 \ref{spaces-morphisms-lemma-pushforward}）。
然后可以取该拟凝聚 $\mathcal{O}_X$-代数的相对谱（引理
\ref{spaces-morphisms-lemma-affine-equivalence-algebras}），得到 $X$ 在 $Y$ 中的正规化。

\begin{definition}
\label{spaces-morphisms-definition-normalization-X-in-Y}
设 $S$ 为概形。设 $f : Y \to X$ 为 $S$ 上代数空间的拟紧拟分离态射。设
$\mathcal{O}'$ 为 $\mathcal{O}_X$ 在 $f_*\mathcal{O}_Y$ 中的整闭包。
{\it $X$ 在 $Y$ 中的正规化}是 $S$ 上代数空间的态射
$$
\nu : X' = \underline{\Spec}_X(\mathcal{O}') \to X
$$
。它配有初始态射 $f$ 的自然分解
$$
Y \xrightarrow{f'} X' \xrightarrow{\nu} X
$$
。
\end{definition}

\noindent
为得到这一分解，使用注 \ref{spaces-morphisms-remark-factorization-quasi-compact-quasi-separated}
以及 $\underline{\Spec}$ 构造的函子性。

\begin{lemma}
\label{spaces-morphisms-lemma-properties-normalization}
设 $S$ 为概形。设 $f : Y \to X$ 为 $S$ 上代数空间的拟紧拟分离态射。设
$Y \to X' \to X$ 为 $X$ 在 $Y$ 中的正规化。
\begin{enumerate}
\item 若 $W \to X$ 为 $S$ 上代数空间的 \'etale 态射，则
$W \times_X X'$ 是 $W$ 在 $W \times_X Y$ 中的正规化。
\item 若 $Y$ 与 $X$ 可表，则 $X'$ 可表，并典范同构于《态射》第
\ref{morphisms-section-normalization} 节所构造的概形 $X$ 在概形 $Y$
中的正规化。
\end{enumerate}
\end{lemma}

\begin{proof}
由构造立即可知，形成 $X$ 在 $Y$ 中的正规化与 \'etale 基变换交换，即 (1)
成立。另一方面，若 $X$ 与 $Y$ 为概形，则对仿射开集 $U \subset X$，
$f_*\mathcal{O}_Y(U) = \mathcal{O}_Y(f^{-1}(U))$；因此，
$\nu^{-1}(U)$ 恰为概形相应构造中所得同一环的谱。
\end{proof}

\noindent
下面给出这一构造的刻画。

\begin{lemma}
\label{spaces-morphisms-lemma-characterize-normalization}
设 $S$ 为概形。设 $f : Y \to X$ 为 $S$ 上代数空间的拟紧拟分离态射。
分解 $f = \nu \circ f'$，其中 $\nu : X' \to X$ 是 $X$ 在 $Y$ 中的正规化，
由下列两个性质刻画：
\begin{enumerate}
\item 态射 $\nu$ 为整态射；
\item 对任意分解 $f = \pi \circ g$，其中 $\pi : Z \to X$ 为整态射，都存在交换图
$$
\xymatrix{
Y \ar[d]_{f'} \ar[r]_g & Z \ar[d]^\pi \\
X' \ar[ru]^h \ar[r]^\nu & X
}
$$
，其中态射 $h : X' \to Z$ 唯一。
\end{enumerate}
此外，在 (2) 中，态射 $h : X' \to Z$ 是 $Z$ 在 $Y$ 中的正规化。
\end{lemma}

\begin{proof}
设 $\mathcal{O}' \subset f_*\mathcal{O}_Y$ 为定义
\ref{spaces-morphisms-definition-normalization-X-in-Y} 中 $\mathcal{O}_X$ 的整闭包。
由构造，态射 $\nu$ 为整态射，这证明了 (1)。假设给定 (2) 中的分解
$f = \pi \circ g$，其中 $\pi : Z \to X$ 为整态射。由定义
\ref{spaces-morphisms-definition-integral}，$\pi$ 仿射；因而 $Z$ 是某个拟凝聚
$\mathcal{O}_X$-代数层 $\mathcal{B}$ 的相对谱。
% P08-E258: frozen source domain of g preserved; see ERRATA_CANDIDATES.jsonl.
态射 $g : X \to Z$ 对应于 $\mathcal{O}_X$-代数映射
$\chi : \mathcal{B} \to f_*\mathcal{O}_Y$。对每个 \'etale 映至 $X$ 的
仿射 $U$，$\mathcal{B}(U)$ 在 $\mathcal{O}_X(U)$ 上整
（由定义 \ref{spaces-morphisms-definition-integral}），故由引理
\ref{spaces-morphisms-lemma-integral-closure} 可知
$\chi(\mathcal{B}) \subset \mathcal{O}'$。由相对谱构造的函子性及引理
\ref{spaces-morphisms-lemma-affine-equivalence-algebras}，这给出唯一态射
$h : X' \to Z$。略去该图交换的验证。

\medskip\noindent
显然，(1)、(2) 刻画分解 $f = \nu \circ f'$，因为它们把该分解刻画为某范畴的
始对象。由引理 \ref{spaces-morphisms-lemma-finite-permanence}，(2) 中的态射 $h$ 为整态射。
给定分解 $g = \pi' \circ g'$，其中 $\pi' : Z' \to Z$ 为整态射，可得分解
$f = (\pi \circ \pi') \circ g'$ 以及态射 $h' : X' \to Z'$。唯一性蕴含
$\pi' \circ h' = h$。因此，刻画 (1)、(2) 可应用于态射
$h : X' \to Z$，从而得到本引理的最后一个陈述。
\end{proof}

\begin{lemma}
\label{spaces-morphisms-lemma-normalization-in-reduced}
设 $S$ 为概形。设 $f : Y \to X$ 为 $S$ 上代数空间的拟紧拟分离态射。
设 $X' \to X$ 为 $X$ 在 $Y$ 中的正规化。若 $Y$ 既约，则 $X'$ 也既约。
\end{lemma}

\begin{proof}
这由既约环的子环仍既约这一事实得出。略去一些细节。
\end{proof}

\begin{lemma}
\label{spaces-morphisms-lemma-normalization-generic}
设 $S$ 为概形。设 $f : Y \to X$ 为概形的拟紧拟分离态射。设 $X' \to X$
为 $X$ 在 $Y$ 中的正规化。若 $x' \in |X'|$ 是余维 $0$ 的点
（《空间的性质》定义
\ref{spaces-properties-definition-dimension-local-ring}），则 $x'$
是某个余维 $0$ 的点 $y \in |Y|$ 之像。
\end{lemma}

\begin{proof}
由引理 \ref{spaces-morphisms-lemma-properties-normalization} 及各定义，可设
$X = \Spec(A)$ 仿射。此时 $X' = \Spec(A')$，其中 $A'$ 是 $A$ 在
$\Gamma(Y, \mathcal{O}_Y)$ 中的整闭包，而 $x'$ 对应于 $A'$ 的极小素理想。
选取满 \'etale 态射 $V \to Y$，其中 $V = \Spec(B)$ 仿射。则
$A' \to B$ 为单射，故 $A'$ 的每个极小素理想都是 $B$ 的某个极小素理想之像，
见《代数》引理 \ref{algebra-lemma-injective-minimal-primes-in-image}。引理得证。
\end{proof}

\begin{lemma}
\label{spaces-morphisms-lemma-normalization-in-disjoint-union}
设 $S$ 为概形。设 $f : Y \to X$ 为 $S$ 上代数空间的拟紧拟分离态射。
假设 $Y = Y_1 \amalg Y_2$ 是两个代数空间的不交并。记
$f_i = f|_{Y_i}$。设 $X_i'$ 为 $X$ 在 $Y_i$ 中的正规化。则
$X_1' \amalg X_2'$ 是 $X$ 在 $Y$ 中的正规化。
\end{lemma}

\begin{proof}
略。
\end{proof}

\begin{lemma}
\label{spaces-morphisms-lemma-normalization-in-universally-closed}
设 $S$ 为概形。
% P08-E259: frozen source number agreement noted; see ERRATA_CANDIDATES.jsonl.
设 $f : X \to Y$ 为 $S$ 上代数空间的拟紧、拟分离且普遍闭的态射。
则 $f_*\mathcal{O}_X$ 在 $\mathcal{O}_Y$ 上整。换言之，$Y$ 在 $X$
中的正规化等于注 \ref{spaces-morphisms-remark-factorization-quasi-compact-quasi-separated}
中的分解
$$
X \longrightarrow \underline{\Spec}_Y(f_*\mathcal{O}_X)
\longrightarrow Y
$$
。
\end{lemma}

\begin{proof}
该问题在 $Y$ 上是 \'etale 局部的，故可化归至 $Y = \Spec(R)$ 仿射的情形。
设 $h \in \Gamma(X, \mathcal{O}_X)$。需要证明 $h$ 满足 $R$ 上的首一方程。
把 $h$ 视为下列交换图中的态射
$$
\xymatrix{
X \ar[rr]_h \ar[rd]_f & & \mathbf{A}^1_Y \ar[ld] \\
& Y &
}
$$
。设 $Z \subset \mathbf{A}^1_Y$ 为 $h$ 的概形论像，见定义
\ref{spaces-morphisms-definition-scheme-theoretic-image}。由于 $f$ 拟紧且
$\mathbf{A}^1_Y \to Y$ 分离，态射 $h$ 拟紧，见引理
\ref{spaces-morphisms-lemma-quasi-compact-permanence}。由引理
\ref{spaces-morphisms-lemma-quasi-compact-scheme-theoretic-image}，态射 $X \to Z$
在底层拓扑空间上有稠密像。由引理
\ref{spaces-morphisms-lemma-universally-closed-permanence}，态射 $X \to Z$ 闭。
因此，在集合论意义下 $h(X) = Z$。于是可以应用引理
\ref{spaces-morphisms-lemma-image-proper-is-proper}，得出 $Z \to Y$ 普遍闭（甚至固有）。
由于 $Z \subset \mathbf{A}^1_Y$，可知 $Z \to Y$ 仿射且固有，故由引理
\ref{spaces-morphisms-lemma-integral-universally-closed} 它为整态射。写成
$\mathbf{A}^1_Y = \Spec(R[T])$，可知 $Z$ 的理想 $I \subset R[T]$
含有首一多项式 $P(T) \in R[T]$。因此 $P(h) = 0$，结论成立。
\end{proof}

\begin{lemma}
\label{spaces-morphisms-lemma-normalization-in-integral}
设 $S$ 为概形。设 $f : Y \to X$ 为 $S$ 上代数空间的整态射。
% P08-E260: frozen source space-level terminology preserved; see ERRATA_CANDIDATES.jsonl.
则 $X$ 在 $Y$ 中的整闭包等于 $Y$。
\end{lemma}

\begin{proof}
由引理 \ref{spaces-morphisms-lemma-integral-universally-closed}，这是引理
\ref{spaces-morphisms-lemma-normalization-in-universally-closed} 的特例。
\end{proof}

\begin{lemma}
\label{spaces-morphisms-lemma-nagata-normalization-finite}
设 $S$ 为概形。设 $f : X \to Y$ 为 $S$ 上代数空间的态射。假设：
\begin{enumerate}
\item $Y$ 是纳伽塔的；
\item $f$ 拟分离且有限型；
\item $X$ 既约。
\end{enumerate}
则 $Y$ 在 $X$ 中的正规化 $\nu : Y' \to Y$ 是有限态射。
\end{lemma}

\begin{proof}
该问题在 $Y$ 上是 \'etale 局部的，见引理
\ref{spaces-morphisms-lemma-properties-normalization}。因此可设 $Y = \Spec(R)$ 仿射。
此时 $R$ 是诺特纳伽塔环，而我们需要证明 $R$ 在
$\Gamma(X, \mathcal{O}_X)$ 中的整闭包是有限 $R$-模。由于 $f$ 拟紧，
可知 $X$ 拟紧。选取仿射概形 $U$ 及满 \'etale 态射 $U \to X$
（《空间的性质》引理
\ref{spaces-properties-lemma-quasi-compact-affine-cover}）。
% P08-E261: frozen source sheaf subscript preserved; see ERRATA_CANDIDATES.jsonl.
则 $\Gamma(X, \mathcal{O}_X) \subset \Gamma(U, \mathcal{O}_X)$。
由于 $R$ 诺特，只需证明 $R$ 在 $\Gamma(U, \mathcal{O}_U)$ 中的整闭包是
有限 $R$-模。由于 $U \to Y$ 有限型，这由《态射》引理
\ref{morphisms-lemma-nagata-normalization-finite} 得出。
\end{proof}






\section{正规化}
\label{spaces-morphisms-section-normalization}

\noindent
本节是《态射》第 \ref{morphisms-section-normalization} 节的对应版本。

\begin{lemma}
\label{spaces-morphisms-lemma-prepare-normalization}
设 $S$ 为概形。设 $X$ 为 $S$ 上的代数空间。下列条件等价：
\begin{enumerate}
\item 存在满 \'etale 态射 $U \to X$，其中 $U$ 为概形，且 $U$ 的每个拟紧开集
只有有限多个不可约分支；
\item 对每个概形 $U$ 及每个 \'etale 态射 $U \to X$，$U$ 的每个拟紧开集
只有有限多个不可约分支；
\item 对每个在 $X$ 上 \'etale 的拟紧代数空间 $Y$，$Y$ 的余维 $0$ 点集
（《空间的性质》定义
\ref{spaces-properties-definition-dimension-local-ring}）有限；
\item 对每个在 $X$ 上 \'etale 的拟紧代数空间 $Y$，空间 $|Y|$
只有有限多个不可约分支。
\end{enumerate}
若 $X$ 可表，这意味着 $X$ 的每个拟紧开集只有有限多个不可约分支。
\end{lemma}

\begin{proof}
(1) 与 (2) 的等价性以及最后一个陈述，由《下降》引理
\ref{descent-lemma-locally-finite-nr-irred-local-fppf} 和《空间的性质》引理
\ref{spaces-properties-lemma-type-property} 得出。只考虑作为概形的 $Y$，
即可明显看出 (4) 蕴含 (1)、(2)。类似地，(3) 蕴含 (1)、(2)，因为对概形而言，
余维 $0$ 的点就是其各不可约分支的一般点；例如见《空间的性质》引理
\ref{spaces-properties-lemma-codimension-0-points}。

\medskip\noindent
反之，假设 (2)，并设 $Y \to X$ 为代数空间的 \'etale 态射，其中 $Y$ 拟紧。
可以选取仿射概形 $V$ 及满 \'etale 态射 $V \to Y$
（《空间的性质》引理
\ref{spaces-properties-lemma-quasi-compact-affine-cover}）。
由 (2)，$V$ 只有有限多个不可约分支；又因 $|V| \to |Y|$ 连续且满，由
《拓扑》引理 \ref{topology-lemma-surjective-continuous-irreducible-components}，
可知 $|Y|$ 只有有限多个不可约分支。故 (4) 成立。类似地，由《空间的性质》引理
\ref{spaces-properties-lemma-codimension-0-points}，$V$ 各不可约分支一般点的
像就是 $Y$ 的余维 $0$ 点，从而这些点只有有限多个，即 (3) 成立。
\end{proof}

\begin{lemma}
\label{spaces-morphisms-lemma-locally-noetherian-can-be-normalized}
设 $S$ 为概形。设 $X$ 为 $S$ 上局部诺特的代数空间。则 $X$ 满足引理
\ref{spaces-morphisms-lemma-prepare-normalization} 中的等价条件。
\end{lemma}

\begin{proof}
若 $U \to X$ 为 \'etale 态射且 $U$ 为概形，则 $U$ 是局部诺特概形，见
《空间的性质》第 \ref{spaces-properties-section-types-properties} 节。
局部诺特概形的不可约分支集局部有限（《除子》引理
\ref{divisors-lemma-components-locally-finite}）。因此，$X$ 满足该引理的
条件 (2)。
\end{proof}

\begin{lemma}
\label{spaces-morphisms-lemma-flat-image-point-codimension-0}
设 $S$ 为概形。设 $f : X \to Y$ 为 $S$ 上代数空间的平坦态射。则对
$x \in |X|$ 有：$x$ 余维 $0$ 于
$X \Rightarrow f(x)$ 余维 $0$ 于 $Y$。
\end{lemma}

\begin{proof}
通过《空间的性质》引理
\ref{spaces-properties-lemma-codimension-0-points} 及 \'etale 局部化，
这化为概形态射及不可约分支一般点的情形。此时，结论由一般化可沿概形的
平坦态射提升得出，见《态射》引理
\ref{morphisms-lemma-generalizations-lift-flat}。
\end{proof}

\begin{lemma}
\label{spaces-morphisms-lemma-flat-locally-finite-type-points-codimension-0}
设 $S$ 为概形。设 $f : X \to Y$ 为 $S$ 上代数空间的态射。假设 $f$
平坦且局部有限型，并假设 $Y$ 满足引理
\ref{spaces-morphisms-lemma-prepare-normalization} 中的等价条件。则 $X$ 也满足引理
\ref{spaces-morphisms-lemma-prepare-normalization} 中的等价条件，且对 $x \in |X|$ 有：
$x$ 余维 $0$ 于 $X \Rightarrow f(x)$ 余维 $0$ 于 $Y$。
\end{lemma}

\begin{proof}
最后一个陈述由引理 \ref{spaces-morphisms-lemma-flat-image-point-codimension-0} 得出。
选取满 \'etale 态射 $V \to Y$，其中 $V$ 为概形；再选取满 \'etale 态射
$U \to X \times_Y V$，其中 $U$ 为概形。只需证明 $U$ 的每个拟紧开集只有
有限多个不可约分支。以下不再特别说明地使用《空间的性质》引理
\ref{spaces-properties-lemma-codimension-0-points} 的结果。由已经证明的内容，
$U$ 的余维 $0$ 点位于
% P08-E262: frozen source repeated U preserved; see ERRATA_CANDIDATES.jsonl.
$U$ 中余维 $0$ 的点上方，而后者依假设局部有限。因此，只需证明：对
余维为 $0$ 的 $v \in V$，概形论纤维
$U_v = U \times_V v$ 的余维 $0$ 点局部有限。这成立是因为 $U_v$ 是
$\kappa(v)$ 上局部有限型的概形，因而局部诺特；例如可应用引理
\ref{spaces-morphisms-lemma-locally-noetherian-can-be-normalized}。
\end{proof}

\begin{lemma}
\label{spaces-morphisms-lemma-normalization}
设 $S$ 为概形。对每个满足引理 \ref{spaces-morphisms-lemma-prepare-normalization}
中等价条件的 $S$ 上代数空间 $X$，都存在代数空间的态射
$$
\nu_X : X^\nu \longrightarrow X
$$
，具有下列性质：
\begin{enumerate}
\item 若 $X$ 满足引理 \ref{spaces-morphisms-lemma-prepare-normalization} 中的等价条件，则
$X^\nu$ 正规且 $\nu_X$ 为整态射；
\item 若 $X$ 为概形，且其每个拟紧开集只有有限多个不可约分支，则
$\nu_X : X^\nu \to X$ 是《态射》第
\ref{morphisms-section-normalization} 节所构造的 $X$ 的正规化；
\item 若 $f : X \to Y$ 为 $S$ 上代数空间的态射，且二者都满足引理
\ref{spaces-morphisms-lemma-prepare-normalization} 中的等价条件，且 $X$ 的每个余维 $0$ 点
都被 $f$ 映至 $Y$ 的余维 $0$ 点，则存在唯一的 $S$ 上代数空间态射
$f^\nu : X^\nu \to Y^\nu$，使得
$\nu_Y \circ f^\nu = f \circ \nu_X$；
\item 若 $f : X \to Y$ 为代数空间的 \'etale 或光滑态射，且 $Y$ 满足引理
\ref{spaces-morphisms-lemma-prepare-normalization} 中的等价条件，则 (3) 的假设成立，且态射
$f^\nu$ 诱导同构 $X^\nu  \to X \times_Y Y^\nu$。
\end{enumerate}
\end{lemma}

\begin{proof}
考虑范畴 $\mathcal{C}$：其对象为 $S$ 上的概形 $U$，使得 $U$ 的每个拟紧开集
只有有限多个不可约分支；其态射为 $S$ 上概形的那些态射
$g : U \to V$，它们把 $U$ 每个不可约分支的一般点映至 $V$
某个不可约分支的一般点。我们已经证明：
\begin{enumerate}
\item[(a)] 对 $U \in \Ob(\mathcal{C})$，有《态射》定义
\ref{morphisms-definition-normalization} 中的正规化态射
$\nu_U : U^\nu \to U$；
\item[(b)] 对 $U \in \Ob(\mathcal{C})$，态射 $\nu_U$ 为整态射，且
$U^\nu$ 为正规概形，见《态射》引理
\ref{morphisms-lemma-normalization-normal}；
\item[(c)] 对每个 $g : U \to V \in \text{Arrows}(\mathcal{C})$，存在唯一态射
$g^\nu : U^\nu \to V^\nu$，使得
$\nu_V \circ g^\nu = g \circ \nu_U$；见《态射》引理
\ref{morphisms-lemma-normalization-normal} 的 (4)，把它应用于复合
% P08-E263: frozen source variables in the composition preserved; see ERRATA_CANDIDATES.jsonl.
$X^\nu \to X \to Y$；
\item[(d)] 若 $V \in \Ob(\mathcal{C})$ 且 $g : U \to V$ 为 \'etale 或光滑态射，
则 $U \in \Ob(\mathcal{C})$，$g \in \text{Arrows}(\mathcal{C})$，且态射
$g^\nu$ 诱导同构 $U^\nu \to U \times_V V^\nu$；见引理
\ref{spaces-morphisms-lemma-flat-locally-finite-type-points-codimension-0} 与
《关于态射的进一步结果》引理
\ref{more-morphisms-lemma-normalization-and-smooth}。
\end{enumerate}
我们的任务是把这一构造推广到 $S$ 上代数空间 $X$ 的相应范畴。

\medskip\noindent
设 $X$ 为满足引理 \ref{spaces-morphisms-lemma-prepare-normalization} 中等价条件的 $S$
上代数空间。设 $U \to X$ 为满 \'etale 态射，其中 $U$ 为概形。置
$R = U \times_X U$，其投影为 $s, t : R \to U$，并置
$j = (t, s) : R \to U \times_S U$，于是 $X = U/R$；见《空间》引理
\ref{spaces-lemma-space-presentation}。由对 $X$ 的假设，$U$ 与 $R$
是 $\mathcal{C}$ 的对象，而态射 $s$、$t$ 是 $S$ 上概形的 \'etale 态射。
由 (a)，有正规化态射 $\nu_U : U^\nu \to U$ 与
$\nu_R : R^\nu \to R$；由 (d)，有态射
$s^\nu : R^\nu \to U^\nu$、$t^\nu : R^\nu \to U^\nu$，它们分别定义同构
$R^\nu \to R \times_{s, U} U^\nu$ 与
$R^\nu \to U^\nu \times_{U, t} R$。因此 $s^\nu$ 与 $t^\nu$ 为 \'etale
态射（因为它们同构于 \'etale 态射的基变换）。所诱导的态射
$j^\nu = (t^\nu, s^\nu) : R^\nu \to U^\nu \times_S U^\nu$
是单态射，因为它等于复合
\begin{align*}
R^\nu
& \to
(U^\nu \times_{U, t} R) \times_R (R \times_{s, U} U^\nu) \\
& =
U^\nu \times_{U, t} R \times_{s, U} U^\nu \\
& \xrightarrow{j}
U^\nu \times_U (U \times_S U) \times_U U^\nu \\
& =
U^\nu \times_S U^\nu
\end{align*}
。第一个箭头是 $\nu_R$ 的对角态射。（这说明 $R^\nu$ 是 $R$ 在 $U^\nu$
上限制的子概形。）利用性质 (d) 对纤维积作形式计算，可知
$R^\nu \times_{s^\nu, U^\nu, t^\nu} R^\nu$ 是
$R \times_{s, U, t} R$ 的正规化。因此，由 (d)，\'etale 态射
$c : R \times_{s, U, t} R \to R$ 唯一延拓为 $c^\nu$。态射 $c^\nu$
与投影
$\text{pr}_{13} : U^\nu \times_S U^\nu \times_S U^\nu \to U^\nu \times_S U^\nu$
相容。类似地，存在态射 $i^\nu : R^\nu \to R^\nu$，它与交换因子的态射
$U^\nu \times_S U^\nu \to U^\nu \times_S U^\nu$ 相容；还存在态射
$e^\nu : U^\nu \to R^\nu$，它与对角态射
$U^\nu \to U^\nu \times_S U^\nu$ 相容。综上可知，
$j^\nu : R^\nu \to U^\nu \times_S U^\nu$ 是 \'etale 等价关系。
此时置 $X^\nu = U^\nu/R^\nu$
（《空间》定理 \ref{spaces-theorem-presentation}）。于是由《群胚》引理
\ref{groupoids-lemma-criterion-fibre-product}，有
$U^\nu = X^\nu \times_X U$。

\medskip\noindent
% P08-E264: frozen source grammar preserved semantically; see ERRATA_CANDIDATES.jsonl.
上一段已经证明：对每个满足引理 \ref{spaces-morphisms-lemma-prepare-normalization}
中等价条件的 $S$ 上代数空间 $X$，若选取满 \'etale 态射
$g : U \to X$，其中 $U$ 为概形，则得到代数空间的笛卡尔图
$$
\xymatrix{
X^\nu \ar[d]_{\nu_X} & U^\nu \ar[l]^{g^\nu} \ar[d]^{\nu_U} \\
X & U \ar[l]_g
}
$$
。这立即推出 $X^\nu$ 是正规代数空间，且 $\nu_X$ 是整态射。这就得到引理的
(1)。

\medskip\noindent
下面将证明态射 $\nu_X : X^\nu \to X$ 在唯一同构的意义下与 $g$ 的选择无关。
眼下，若 $X$ 是概形，就选择 $\text{id} : X \to X$，于是引理的 (2) 显然成立。

\medskip\noindent
还须证明 (3) 与 (4)。令 $g : U \to X$、$\nu_X : X^\nu \to X$ 以及
$g^\nu : U^\nu \to X^\nu$ 如上。设 $Z$ 为正规概形，并设
$h : Z \to U$ 与 $a : Z \to X^\nu$ 为 $S$ 上的态射，使得
$g \circ h = \nu_X \circ a$，且 $Z$ 的每个不可约分支都经由 $h$ 支配 $U$
的某个不可约分支。由《态射》引理
\ref{morphisms-lemma-normalization-normal} 的 (4)，存在唯一态射
$h^\nu : Z \to U^\nu$，使得 $h = \nu_U \circ h^\nu$。图示如下：
$$
\xymatrix{
X^\nu \ar[d]_{\nu_X} &
U^\nu \ar[l]^{g^\nu} \ar[d]^{\nu_U} &
Z \ar[l]^{h^\nu} \ar@/_1em/[ll]_a \ar[dl]^h
\\
X & U \ar[l]_g
}
$$
注意 $a = g^\nu \circ h^\nu$。事实上，四个顶点为 $X^\nu$、$X$、$U^\nu$、
$U$ 的方块是笛卡尔的，因此这立即由给定 $h$ 时 $h^\nu$ 的唯一性得出。
换言之，给定上述 $h : Z \to U$（而不指定 $a$），存在唯一态射
$a : Z \to X^\nu$，使得 $\nu_X \circ a = g \circ h$。

\medskip\noindent
令 $f : X \to Y$ 如引理陈述的 (3) 所示。假设已经选取满 \'etale 态射
$U \to X$ 与 $V \to Y$，其中 $U$ 和 $V$ 为概形，并且 $f$ 可提升为态射
$g : U \to V$。于是 $g \in \text{Arrows}(\mathcal{C})$，并得到与
$\nu_U$ 和 $\nu_V$ 相容的唯一态射 $g^\nu : U^\nu \to V^\nu$。然而，此时两个态射
$$
R^\nu = U^\nu \times_{X^\nu} U^\nu
\to U^\nu \to V^\nu \to Y^\nu
$$
% P08-E265: frozen source spelling "in stead" noted; see ERRATA_CANDIDATES.jsonl.
由上一段的论述必然相同（其中用 $Y$ 代替 $X$）。
由于 $X^\nu$ 是取 $U^\nu$ 关于 $R^\nu$ 的商而构造的，遂得到 (3) 所述的
（唯一）态射 $f^\nu : X^\nu \to Y^\nu$。

\medskip\noindent
为说明 $X^\nu$ 的构造与满 \'etale 态射 $g : U \to X$ 的选择无关，把上一段
的构造应用于 $\text{id} : X \to X$ 以及 $X$ 的两个 \'etale 覆盖之间的态射
$U' \to U$。这已经足够，因为给定 $X$ 的任意两个 \'etale 覆盖，总有第三个覆盖
同时支配二者。读者可验证：分别用 $U' \to X$ 与 $U \to X$ 构造的两个正规化之间
的态射，在基变换到 $U'$ 后成为同构，因而原来就是同构。细节从略。

\medskip\noindent
类似地可证 (4)，这里略去证明；提示：使用上面的 (d)。
\end{proof}

\noindent
由此得到下列定义。

\begin{definition}
\label{spaces-morphisms-definition-normalization}
设 $S$ 为概形。设 $X$ 为满足引理 \ref{spaces-morphisms-lemma-prepare-normalization}
中等价条件的 $S$ 上代数空间。定义 $X$ 的{\it 正规化}为态射
$$
\nu_X : X^\nu \longrightarrow X
$$
，它是在引理 \ref{spaces-morphisms-lemma-normalization} 中构造的。
\end{definition}

\noindent
该定义适用于局部诺特代数空间，见引理
\ref{spaces-morphisms-lemma-locally-noetherian-can-be-normalized}。通常只对约化代数空间定义正规化。
按照上述定义，$X$ 的正规化与 $X$ 的约化 $X_{red}$ 的正规化相同。

\begin{lemma}
\label{spaces-morphisms-lemma-normalization-reduced}
设 $S$ 为概形。设 $X$ 为满足引理 \ref{spaces-morphisms-lemma-prepare-normalization}
中等价条件的 $S$ 上代数空间。正规化态射 $\nu$ 经过约化 $X_{red}$ 分解，且
$X^\nu \to X_{red}$ 是 $X_{red}$ 的正规化。
\end{lemma}

\begin{proof}
这可在 $X$ 上 \'etale 局部检验，因而归结为概形的情形，即《态射》引理
\ref{morphisms-lemma-normalization-reduced}。略去若干细节。
\end{proof}

\begin{lemma}
\label{spaces-morphisms-lemma-normalization-normal}
设 $S$ 为概形。设 $X$ 为满足引理 \ref{spaces-morphisms-lemma-prepare-normalization}
中等价条件的 $S$ 上代数空间。
\begin{enumerate}
\item 正规化 $X^\nu$ 是正规的。
\item 态射 $\nu : X^\nu \to X$ 是满整态射。
\item 映射 $|\nu| : |X^\nu| \to |X|$ 在余维 $0$ 点集之间诱导双射
（《空间的性质》定义
\ref{spaces-properties-definition-dimension-local-ring}）。
\item 设 $Z \to X$ 为态射。假设 $Z$ 是正规代数空间，并且对 $z \in |Z|$ 有：
% P08-E266: frozen source uses undefined f(z); preserved; see ERRATA_CANDIDATES.jsonl.
$z$ 余维 $0$ 于 $Z \Rightarrow f(z)$ 余维 $0$ 于 $X$。则存在唯一分解
$Z \to X^\nu \to X$。
\end{enumerate}
\end{lemma}

\begin{proof}
(1)、(2)、(3) 由概形的相应结果（《态射》引理
\ref{morphisms-lemma-normalization-normal}）以及下述事实得出：概形的一点是某个
不可约分支的一般点，当且仅当其局部环的维数为零（《性质》引理
\ref{properties-lemma-generic-point}）。

\medskip\noindent
设 $Z \to X$ 为 (4) 中的态射。取概形 $U$ 以及满 \'etale 态射 $U \to X$。
选取概形 $V$ 及满 \'etale 态射 $V \to U \times_X Z$。关于余维 $0$ 点的条件保证
$V \to U$ 把 $V$ 各不可约分支的一般点映至 $U$ 各不可约分支的一般点。因此由
《态射》引理 \ref{morphisms-lemma-normalization-normal}，得到唯一分解
$V \to U^\nu \to U$。唯一性保证两个映射
$$
V \times_{U \times_X Z} V \to V \to U^\nu
$$
相同，因为它们都是映射 $V \times_{U \times_X Z} V \to V \to U$ 的唯一分解。
由于代数空间 $U \times_X Z$ 等于商 $V/V \times_{U \times_X Z} V$
（见《空间》第 \ref{spaces-section-presentations} 节），便得到典范态射
$U \times_X Z \to U^\nu$。图示如下：
$$
\xymatrix{
U \times_X Z \ar[r] \ar[d] & U^\nu \ar[r] \ar[d] & U \ar[d] \\
Z \ar@/_/[rr] \ar@{..>}[r] & X^\nu \ar[r] & X
}
$$
为得到虚线箭头，注意态射 $U \times_X Z \to U^\nu$ 的构造关于 \'etale 态射
$U \to X$ 是函子性的（略去精确表述与证明）。因此若置
$R = U \times_X U$，其投影为 $s, t : R \to U$，则得到态射
$R \times_X Z \to R^\nu$，它与 $s, t : R \to U$ 及
$s^\nu, t^\nu : R^\nu \to U^\nu$ 交换。回顾
$X^\nu = U^\nu/R^\nu$，见引理 \ref{spaces-morphisms-lemma-normalization} 的证明。由于
$X = U/R$，一个简单的层论论证表明
$Z = (U \times_X Z)/(R \times_X Z)$。于是态射
$U \times_X Z \to U^\nu$ 与 $R \times_X Z \to R^\nu$ 定义所需的态射
$Z \to X^\nu$。
\end{proof}

\begin{lemma}
\label{spaces-morphisms-lemma-nagata-normalization}
设 $S$ 为概形。设 $X$ 为 $S$ 上的 Nagata 代数空间。正规化
$\nu : X^\nu \to X$ 是有限态射。
\end{lemma}

\begin{proof}
由于 Nagata 代数空间 $X$ 是局部诺特的，定义
\ref{spaces-morphisms-definition-normalization} 适用。根据引理 \ref{spaces-morphisms-lemma-normalization}
中 $X^\nu$ 的构造，立即归结为概形的情形，即《态射》引理
\ref{morphisms-lemma-nagata-normalization}。
\end{proof}












\section{分离且局部拟有限的态射}
\label{spaces-morphisms-section-schemehood}

\noindent
本节证明：在某概形上局部拟有限且分离的代数空间是可表的。这推出分离且局部拟有限
的态射可表（见引理 \ref{spaces-morphisms-lemma-locally-quasi-finite-separated-representable}）。
不过先来一个引理……（它将被命题
\ref{spaces-morphisms-proposition-locally-quasi-finite-separated-over-scheme} 取代）。

\begin{lemma}
\label{spaces-morphisms-lemma-neighbourhood-scheme}
设 $S$ 为概形。考虑 $S$ 上代数空间的交换图
$$
\xymatrix{
V' \ar[r] \ar[rd] & T' \times_T X \ar[r] \ar[d] & X \ar[d] \\
& T' \ar[r] & T
}
$$
。假设
\begin{enumerate}
\item $T' \to T$ 是仿射概形的 \'etale 态射；
\item $X \to T$ 是分离且局部拟有限的态射；
\item $V'$ 是 $T' \times_T X$ 的开子空间；
\item $V' \to T'$ 是拟仿射的。
\end{enumerate}
在此情形下，$V'$ 在 $X$ 中的像 $U$ 是 $X$ 的可表拟紧开子空间。
\end{lemma}

\begin{proof}
先作几个简单观察。由引理 \ref{spaces-morphisms-lemma-quasi-affine-local}，$V'$ 可表。
它也是拟紧的（因为它是仿射概形上的拟仿射概形；见《态射》引理
\ref{morphisms-lemma-quasi-affine-separated}）。由于 $T' \times_T X \to X$
是 \'etale 的（《空间的性质》引理
\ref{spaces-properties-lemma-base-change-etale}），映射
$|T' \times_T X| \to |X|$ 是开的；见《空间的性质》引理
\ref{spaces-properties-lemma-etale-open}。令 $U \subset X$ 为与 $|V'|$ 的像
对应的开子空间；见《空间的性质》引理
\ref{spaces-properties-lemma-open-subspaces}。因为 $|V'|$ 拟紧，$|U|$ 也拟紧，
故由《空间的性质》引理 \ref{spaces-properties-lemma-quasi-compact-space}，
$U$ 是拟紧代数空间。

\medskip\noindent
由《态射》引理
\ref{morphisms-lemma-locally-quasi-finite-qc-source-universally-bounded}，态射
$T' \to T$ 是普遍有界的。因此可对限制 $T' \to T$ 各纤维次数的整数 $n$ 作归纳；
关于 \'etale 态射情形下该整数的描述，见《态射》引理
\ref{morphisms-lemma-etale-universally-bounded}。若 $n = 1$，则 $T' \to T$
是开浸入（见《\'Etale 态射》定理
\ref{etale-theorem-etale-radicial-open}），结论显然。假设 $n > 1$。

\medskip\noindent
考虑仿射概形 $T'' = T' \times_T T'$。由于 $T' \to T$ 是 \'etale 的，有分解
（分解成既开又闭的仿射子概形）$T'' = \Delta(T') \amalg T^*$。事实上，由
《态射》引理 \ref{morphisms-lemma-diagonal-unramified-morphism}，
$\Delta = \Delta_{T'/T}$ 是开的；又因为 $T' \to T$ 作为仿射概形间的态射是
分离的，它也是闭的。第二投影
$\text{pr}_1 : T' \times_T T' \to T'$ 是基变换，其各纤维的次数由 $n$ 限制；
见《态射》引理 \ref{morphisms-lemma-base-change-universally-bounded}。另一方面，
$\text{pr}_1|_{\Delta(T')} : \Delta(T') \to T'$ 是同构，每个纤维恰有一点。
因此应用《态射》引理 \ref{morphisms-lemma-etale-universally-bounded} 可知，限制
$\text{pr}_1|_{T^*} : T^* \to T'$ 各纤维的次数由 $n - 1$ 限制。于是把归纳假设
应用于图
$$
\xymatrix{
p_0^{-1}(V') \cap X^* \ar[r] \ar[rd] &
X^* \ar[r]_{p_1|_{X^*}} \ar[d] &
X' \ar[d] \\
& T^* \ar[r]^{\text{pr}_1|_{T^*}} & T'
}
$$
，得到 $p_1(p_0^{-1}(V') \cap X^*)$ 是拟紧概形。这里置
$X'' = T'' \times_T X$、$X^* = T^* \times_T X$、$X' = T' \times_T X$，而
$p_0, p_1 : X'' \to X'$ 是 $\text{pr}_0, \text{pr}_1$ 的基变换。引理的大多数
假设都通过基变换给出上图的相应假设。例如，
$p_0^{-1}(V') = T'' \times_{T'} V'$ 作为基变换，是 $T''$ 上的拟仿射概形。
略去一些验证。

\medskip\noindent
由《空间的性质》引理 \ref{spaces-properties-lemma-subscheme} 可知
$$
p_1(p_0^{-1}(V')) =
V' \cup p_1(p_0^{-1}(V') \cap X^*)
$$
是拟紧概形。此外，$p_1(p_0^{-1}(V'))$ 显然是证明第一段所讨论的拟紧开子空间
$U \subset X$ 的逆像。换言之，$T' \times_T U$ 是概形！注意
$T' \times_T U$ 拟紧，且在 $T'$ 上分离、局部拟有限，因为
$T' \times_T X \to T'$ 是原态射 $X \to T$ 的基变换，故局部拟有限且分离
（见引理 \ref{spaces-morphisms-lemma-base-change-separated} 与
\ref{spaces-morphisms-lemma-base-change-quasi-finite}）。于是由《关于态射的进一步结果》引理
\ref{more-morphisms-lemma-quasi-finite-separated-quasi-affine}，
$T' \times_T U \to T'$ 是拟仿射的。

\medskip\noindent
由《下降》引理 \ref{descent-lemma-descent-data-sheaves}，这给出
$T' \times_T U / T'$ 上相对于 \'etale 覆盖 $\{T' \to W\}$ 的下降资料，其中
$W \subset T$ 是态射 $T' \to T$ 的像。
% P08-E267: frozen source U' preserved; see ERRATA_CANDIDATES.jsonl.
因为 $U'$ 在 $T'$ 上拟仿射，由《下降》引理
\ref{descent-lemma-quasi-affine} 可知该资料有效；再由《下降》引理
\ref{descent-lemma-descent-data-sheaves} 的最后一部分，推出 $U$ 如所需是概形。
略去一些小细节。
\end{proof}

\begin{proposition}
\label{spaces-morphisms-proposition-locally-quasi-finite-separated-over-scheme}
设 $S$ 为概形。设 $f : X \to T$ 为 $S$ 上代数空间的态射。假设
\begin{enumerate}
\item $T$ 可表；
\item $f$ 局部拟有限；
\item $f$ 分离。
\end{enumerate}
则 $X$ 可表。
\end{proposition}

\begin{proof}
令 $T = \bigcup T_i$ 为概形 $T$ 的仿射开覆盖。若能证明开子空间
$X_i = f^{-1}(T_i)$ 可表，则由《空间的性质》引理
\ref{spaces-properties-lemma-subscheme}，$X$ 可表。注意
$X_i = T_i \times_T X$，且局部拟有限与分离在基变换下都稳定；见引理
\ref{spaces-morphisms-lemma-base-change-separated} 与 \ref{spaces-morphisms-lemma-base-change-quasi-finite}。
因此可假设 $T$ 是仿射概形。

\medskip\noindent
由《空间的性质》引理
\ref{spaces-properties-lemma-union-of-quasi-compact}，存在 Zariski 覆盖
$X = \bigcup X_i$，使每个 $X_i$ 都有来自某仿射概形的满 \'etale 覆盖。
再次由《空间的性质》引理 \ref{spaces-properties-lemma-subscheme}，只需对每个
$X_i$ 证明命题。因此可假设存在仿射概形 $U$ 及满 \'etale 态射 $U \to X$。
这样便归结为下一段的情形。

\medskip\noindent
假设有
$$
U \longrightarrow X \longrightarrow T
$$
，其中 $U$ 与 $T$ 为仿射概形，$U \to X$ 为满 \'etale 态射，而
$X \to T$ 分离且局部拟有限。由引理 \ref{spaces-morphisms-lemma-etale-locally-quasi-finite}
与 \ref{spaces-morphisms-lemma-composition-quasi-finite}，态射 $U \to T$ 局部拟有限。由于
$U$ 与 $T$ 都仿射，它是拟有限的。置 $R = U \times_X U$。则
$X = U/R$；见《空间》引理 \ref{spaces-lemma-space-presentation}。因为
$X \to T$ 分离，态射 $R \to U \times_T U$ 为闭浸入；见引理
\ref{spaces-morphisms-lemma-fibre-product-after-map}。特别地，$R$ 也是仿射概形。由于
$U \to X$ 是 \'etale 的，投影态射 $t, s : R \to U$ 也都是 \'etale 的。
特别地，$s$ 与 $t$ 拟有限、平坦且有限表示（见《态射》引理
\ref{morphisms-lemma-etale-locally-quasi-finite}、
\ref{morphisms-lemma-etale-flat} 与
\ref{morphisms-lemma-etale-locally-finite-presentation}）。

\medskip\noindent
令 $(U, R, s, t, c)$ 为 $U$ 上 \'etale 等价关系 $R$ 所对应的群胚。取一点
$u \in U$，并以 $p \in T$ 记其像。把《关于群胚的进一步结果》引理
\ref{more-groupoids-lemma-quasi-finite-over-base-j-proper} 应用于概形 $T$ 上的群胚
$(U, R, s, t, c)$ 以及上述点 $p$、$u$。由上一段的讨论，该引理的假设
(1)--(7) 全部满足。因此得到 \'etale 邻域 $(T', p') \to (T, p)$ 及不交并分解
$$
U_{T'} = U' \amalg W, \quad
R_{T'} = R' \amalg W'
$$
以及满足前述《关于群胚的进一步结果》引理
\ref{more-groupoids-lemma-quasi-finite-over-base-j-proper} 之结论
(a)、(b)、(c)、(d)、(e)、(f)、(g)、(h) 的 $u' \in U'$。必要时缩小 $T'$ 后，
可以并确实假设 $T'$ 仿射。结论 (h) 蕴含
$R' = U' \times_{X_{T'}} U'$，且投影映射与 $s'$、$t'$ 的限制相同。因此，结论
(g) 中的 $(U', R', s'|_{R'}, t'|_{R'}, c'|_{R' \times_{t', U', s'} R'})$
是 \'etale 等价关系。由《空间》引理 \ref{spaces-lemma-finding-opens} 可知，
$U'/R'$ 是 $X_{T'}$ 的开子空间。由结论 (d)，概形 $U'$、$R'$ 仿射，而态射
$s'|_{R'}, t'|_{R'}$ 为有限 \'etale 态射。因此《群胚》命题
\ref{groupoids-proposition-finite-flat-equivalence} 适用，从而 $U'/R'$ 是仿射概形。

\medskip\noindent
由此得出：对上述每一对点 $(u, p)$，都能找到满足
$\kappa(p) = \kappa(p')$ 的 \'etale 邻域 $(T', p') \to (T, p)$，以及映至 $u$
的一点 $u' \in U_{T'}$，使得 $u'$ 在 $|X_{T'}|$ 中的像 $x'$ 有一个开邻域
$V'$，它是 $X_{T'}$ 中的仿射概形。应用引理 \ref{spaces-morphisms-lemma-neighbourhood-scheme}，得到
作为概形且包含 $x$（即 $u$ 在 $|X|$ 中的像）的开子空间 $W \subset X$。
这对每个 $x$ 都成立，故由《空间的性质》引理
\ref{spaces-properties-lemma-subscheme}，$X$ 是概形。证明完毕。
\end{proof}



\section{应用}
\label{spaces-morphisms-section-applications}

\noindent
下列引理的另一证明，是把它看作代数空间（不可表）态射的 Zariski 主定理之推论；
见《关于空间态射的进一步结果》第
\ref{spaces-more-morphisms-section-structure-quasi-finite} 节。事实上，
《关于空间态射的进一步结果》引理
\ref{spaces-more-morphisms-lemma-quasi-finite-separated-quasi-affine} 蕴含：
代数空间的拟有限分离态射是拟仿射的，因而可表。

\begin{lemma}
\label{spaces-morphisms-lemma-locally-quasi-finite-separated-representable}
设 $S$ 为概形。设 $f : X \to Y$ 为 $S$ 上代数空间的态射。若 $f$ 局部拟有限
且分离，则 $f$ 可表。
\end{lemma}

\begin{proof}
这立即由命题 \ref{spaces-morphisms-proposition-locally-quasi-finite-separated-over-scheme}
以及局部拟有限性和分离性在任意基变换下保持这一事实得出；见引理
\ref{spaces-morphisms-lemma-base-change-quasi-finite} 与 \ref{spaces-morphisms-lemma-base-change-separated}。
\end{proof}

\begin{lemma}
\label{spaces-morphisms-lemma-etale-universally-injective-open}
\begin{slogan}
普遍单射的 \'etale 映射是开浸入。
\end{slogan}
设 $S$ 为概形。设 $f : X \to Y$ 为 $S$ 上代数空间的 \'etale 普遍单射态射。
则 $f$ 是开浸入。
\end{lemma}

\begin{proof}
设 $T \to Y$ 为从概形到 $Y$ 的态射。若能证明 $X \times_Y T \to T$ 是开浸入，
便得结论。由于 \'etale 与普遍单射都是在基变换下稳定的态射性质（见引理
\ref{spaces-morphisms-lemma-base-change-etale} 与
\ref{spaces-morphisms-lemma-base-change-universally-injective}），可假设 $Y$ 是概形。注意由引理
\ref{spaces-morphisms-lemma-universally-injective}，对角态射
$\Delta_{X/Y} : X \to X \times_Y X$ 是 \'etale 的满单态射。因此
$\Delta_{X/Y}$ 是同构（见《空间》引理
\ref{spaces-lemma-surjective-flat-locally-finite-presentation}）；特别地，$X$
在 $Y$ 上分离。再由命题
\ref{spaces-morphisms-proposition-locally-quasi-finite-separated-over-scheme}，$X$ 也是概形。
最后，由概形的 \'etale 态射基本定理，$X \to Y$ 是开浸入；见《\'Etale 态射》
定理 \ref{etale-theorem-etale-radicial-open}。
\end{proof}



\section{Zariski 主定理（可表情形）}
\label{spaces-morphisms-section-Zariski}

\noindent
这是可以利用正规化与 \'etale 局部化交换来证明的版本。要证明更强的版本
（关于不可表示态射），还需发展更多工具。见《关于空间态射的进一步结果》第
\ref{spaces-more-morphisms-section-structure-quasi-finite} 节。

\begin{lemma}
\label{spaces-morphisms-lemma-finite-type-separated}
设 $S$ 为概形。设 $f : X \to Y$ 为 $S$ 上代数空间的可表、有限型、分离态射。
令 $Y'$ 为 $Y$ 在 $X$ 中的正规化。图示如下：
$$
\xymatrix{
X \ar[rd]_f \ar[rr]_{f'} & & Y' \ar[ld]^\nu \\
& Y &
}
$$
则存在开子空间 $U' \subset Y'$，使得
\begin{enumerate}
\item $(f')^{-1}(U') \to U'$ 是同构；
\item $(f')^{-1}(U') \subset X$ 是 $f$ 为拟有限的点集。
\end{enumerate}
\end{lemma}

\begin{proof}
设 $W \to Y$ 为满 \'etale 态射，其中 $W$ 为概形。则 $W \times_Y X$ 也是概形。
由引理 \ref{spaces-morphisms-lemma-properties-normalization}，代数空间 $W \times_Y Y'$ 可表，
并且是概形 $W$ 在概形 $W \times_Y X$ 中的正规化。图示如下：
$$
\xymatrix{
W \times_Y X \ar[rd]_{(1, f)} \ar[rr]_{(1, f')} & &
W \times_Y Y' \ar[ld]^{(1, \nu)} \\
& W &
}
$$
由《关于态射的进一步结果》引理
\ref{more-morphisms-lemma-finite-type-separated}，本引理的结论在 $W$ 上成立。
令 $V' \subset W \times_Y Y'$ 为满足下列条件的开子概形：
\begin{enumerate}
\item $(1, f')^{-1}(V') \to V'$ 是同构；
\item $(1, f')^{-1}(V') \subset W \times_Y X$ 是 $(1, f)$ 为拟有限的点集。
\end{enumerate}
由引理 \ref{spaces-morphisms-lemma-locally-finite-type-quasi-finite-part}，存在使 $f$ 拟有限的最大
开点集 $U \subset X$，且 $W \times_Y U = (1, f')^{-1}(V')$。由引理
\ref{spaces-morphisms-lemma-closed-immersion-local}，态射 $f'|_U : U \to Y'$ 是开浸入，因为它
在 $W$ 上的基变换，是同构 $(1, f')^{-1}(V') \to V'$ 后接开浸入
$V' \to W \times_Y Y'$。置 $U' = \Im(U \to Y')$ 即完成证明
（略去对 $(f')^{-1}(U') = U$ 的验证）。
\end{proof}

\noindent
在下列引理中可以去掉可表的假设，因为已经证明局部拟有限的分离态射可表。

\begin{lemma}
\label{spaces-morphisms-lemma-quasi-finite-separated-quasi-affine}
设 $S$ 为概形。设 $f : X \to Y$ 为 $S$ 上代数空间的态射。假设 $f$ 拟有限
且分离。令 $Y'$ 为 $Y$ 在 $X$ 中的正规化。图示如下：
$$
\xymatrix{
X \ar[rd]_f \ar[rr]_{f'} & & Y' \ar[ld]^\nu \\
& Y &
}
$$
则 $f'$ 是拟紧开浸入，而 $\nu$ 是整态射。特别地，$f$ 拟仿射。
\end{lemma}

\begin{proof}
由引理 \ref{spaces-morphisms-lemma-locally-quasi-finite-separated-representable}，态射 $f$ 可表。
故可应用引理 \ref{spaces-morphisms-lemma-finite-type-separated}。于是存在开子空间
$U' \subset Y'$，使得 $(f')^{-1}(U') = X$（！），且 $X \to U'$ 是同构！
换言之，$f'$ 是开浸入。注意 $f'$ 拟紧，因为 $f$ 拟紧，而
$\nu : Y' \to Y$ 分离（引理 \ref{spaces-morphisms-lemma-quasi-compact-permanence}）。因此，对每个
仿射概形 $Z$ 及态射 $Z \to Y$，纤维积 $Z \times_Y X$ 都是仿射概形
$Z \times_Y Y'$ 的拟紧开子概形。故按定义，$f$ 拟仿射。
\end{proof}








\section{泛同胚}
\label{spaces-morphisms-section-universal-homeomorphisms}

\noindent
概形的泛同胚类在复合与任意基变换下封闭，并且在基底上是 fppf 局部的。见
《态射》引理 \ref{morphisms-lemma-composition-universal-homeomorphism}、
\ref{morphisms-lemma-base-change-universal-homeomorphism}，以及《下降》引理
\ref{descent-lemma-descending-property-universal-homeomorphism}。因此，把第
\ref{spaces-morphisms-section-representable} 节的讨论应用于这一概念，就知道代数空间的可表示态射
是泛同胚意味着什么。

\begin{lemma}
\label{spaces-morphisms-lemma-universal-homeomorphism-representable}
设 $S$ 为概形。设 $f : X \to Y$ 为 $S$ 上代数空间的可表示态射。则 $f$ 是泛同胚
（依第 \ref{spaces-morphisms-section-representable} 节之意义），当且仅当对代数空间的每个态射
$Z \to Y$，基变换映射 $Z \times_Y X \to Z$ 都诱导同胚
$|Z \times_Y X| \to |Z|$。
\end{lemma}

\begin{proof}
若对代数空间的每个态射 $Z \to Y$，基变换映射 $Z \times_Y X \to Z$ 都诱导
同胚 $|Z \times_Y X| \to |Z|$，那么当 $Z$ 是概形时也如此；这形式地蕴含
$f$ 是第 \ref{spaces-morphisms-section-representable} 节意义下的泛同胚。反之，若 $f$ 是第
\ref{spaces-morphisms-section-representable} 节意义下的泛同胚，则由《空间》引理
\ref{spaces-lemma-representable-transformations-property-implication} 与《态射》
引理 \ref{morphisms-lemma-universal-homeomorphism}，$X \to Y$ 是整的、普遍单射
且满。故由引理 \ref{spaces-morphisms-lemma-integral-universally-closed}，$f$ 普遍闭；它又普遍
单射且（普遍）满。也就是说，$f$ 是泛同胚。
\end{proof}

\begin{definition}
\label{spaces-morphisms-definition-universal-homeomorphism}
设 $S$ 为概形。称 $S$ 上代数空间的态射 $f : X \to Y$ 为{\it 泛同胚}，
当且仅当对代数空间的每个态射 $Z \to Y$，基变换 $Z \times_Y X \to Z$
都诱导同胚 $|Z \times_Y X| \to |Z|$。
\end{definition}

\noindent
由引理 \ref{spaces-morphisms-lemma-universal-homeomorphism-representable}，该定义与此前对代数空间
可表示态射所作的定义并不冲突。对代数空间的态射而言，泛同胚未必总是整的。

\begin{example}
\label{spaces-morphisms-example-universal-homeomorphism}
这是备注 \ref{spaces-morphisms-remark-universally-injective-not-separated} 的延续。考虑代数空间
$X = \mathbf{A}^1_k/\{x \sim -x \mid x \not = 0\}$.
存在态射
$$
\mathbf{A}^1_k \longrightarrow X \longrightarrow \mathbf{A}^1_k
$$
，使得第一个箭头是满 \'etale 态射，第二个箭头普遍单射，而复合是映射
$x \mapsto x^2$。因此该复合普遍闭。由此可知，映射
$X \to \mathbf{A}^1_k$ 是泛同胚，但 $X \to \mathbf{A}^1_k$ 不分离。
\end{example}

\noindent
设 $S$ 为概形。设 $f : X \to Y$ 为 $S$ 上代数空间的泛同胚。则 $f$ 普遍闭，
因而拟紧；见引理 \ref{spaces-morphisms-lemma-universally-closed-quasi-compact}。但 $f$ 不必分离
（见上例），甚至不必拟分离：一个例子是取无限维仿射空间
$\mathbf{A}^\infty = \Spec(k[x_1, x_2, \ldots])$ 关于下述等价关系的商：翻转
有限多个非零坐标的符号（细节从略）。

\medskip\noindent
先陈述几个必备引理。

\begin{lemma}
\label{spaces-morphisms-lemma-base-change-universal-homeomorphism}
代数空间的泛同胚沿代数空间的任意态射作基变换后仍是泛同胚。
\end{lemma}

\begin{proof}
这立即由定义得出。
\end{proof}

\begin{lemma}
\label{spaces-morphisms-lemma-composition-universal-homeomorphism}
代数空间的两个泛同胚之复合是泛同胚。
\end{lemma}

\begin{proof}
从略。
\end{proof}

\begin{lemma}
\label{spaces-morphisms-lemma-reduction-universal-homeomorphism}
设 $S$ 为概形。设 $X$ 为 $S$ 上的代数空间。典范闭浸入 $X_{red} \to X$
（见《空间的性质》定义
\ref{spaces-properties-definition-reduced-induced-space}）是泛同胚。
\end{lemma}

\begin{proof}
从略。
\end{proof}

\noindent
由于目前没有更合适的位置，暂把下列结果放在这里。

\begin{lemma}
\label{spaces-morphisms-lemma-integral-universally-injective-push-pull}
设 $S$ 为概形。设 $f : Y \to X$ 为 $S$ 上代数空间的普遍单射整态射。
\begin{enumerate}
\item 函子
$$
f_{small, *} : \Sh(Y_\etale) \longrightarrow \Sh(X_\etale)
$$
是全忠实的；其本质像由 $X_\etale$ 上的那些集合层 $\mathcal{F}$ 构成，它们在
$|X| \setminus f(|Y|)$ 上的限制同构于 $*$；
\item 函子
$$
f_{small, *} : \textit{Ab}(Y_\etale) \longrightarrow \textit{Ab}(X_\etale)
$$
% P08-E268: frozen source target-site mismatch preserved; see ERRATA_CANDIDATES.jsonl.
是全忠实的；其本质像由 $Y_\etale$ 上支撑包含于 $f(|Y|)$ 的那些阿贝尔层构成。
\end{enumerate}
在两种情形下，$f_{small}^{-1}$ 都是函子 $f_{small, *}$ 的左逆。
\end{lemma}

\begin{proof}
由于 $f$ 是整的，它普遍闭（引理 \ref{spaces-morphisms-lemma-integral-universally-closed}）。
特别地，$f(|Y|)$ 是 $|X|$ 的闭子集，所以上述陈述有意义。证明的其余部分与引理
\ref{spaces-morphisms-lemma-closed-immersion-push-pull} 的证明相同，只是使用《\'Etale 上同调》
命题 \ref{etale-cohomology-proposition-integral-universally-injective-pushforward}
来代替《\'Etale 上同调》命题
\ref{etale-cohomology-proposition-closed-immersion-pushforward}。
\end{proof}
